% !TEX root = nc-cubic.tex
\documentclass[11pt,a4paper]{ltjsarticle}
\usepackage{amsmath,amssymb,amsthm,mathtools,bm,mathrsfs}
\usepackage{luatexja-fontspec}
\usepackage{enumitem,booktabs,longtable,array,needspace}
\usepackage{xcolor}
\usepackage{etoolbox,xr-hyper}
\usepackage[colorlinks=true,linkcolor=blue!45!black,citecolor=blue!45!black,urlcolor=blue!45!black]{hyperref}
\DeclareMathSizes{9}{9}{6}{5}
\DeclareMathSizes{9.85501}{9}{6}{5}
\IfFileExists{/usr/share/fonts/opentype/noto/NotoSerifCJK-Regular.ttc}{%
  \setmainjfont[Path=/usr/share/fonts/opentype/noto/]{NotoSerifCJK-Regular.ttc}%
}{\setmainjfont{HaranoAjiMincho-Regular}}
\setlength{\textwidth}{15.6cm}
\setlength{\oddsidemargin}{0.25cm}
\setlength{\evensidemargin}{0.25cm}
\setlength{\textheight}{22.6cm}
\setlength{\topmargin}{-0.2cm}
\setlength{\parindent}{1em}
\setlength{\parskip}{1pt}
\numberwithin{equation}{section}
\newtheorem{theorem}{定理}[section]
\newtheorem{maintheorem}{主定理}
\renewcommand{\themaintheorem}{\Alph{maintheorem}}
\newtheorem{lemma}[theorem]{補題}
\newtheorem{proposition}[theorem]{命題}
\newtheorem{corollary}[theorem]{系}
\newtheorem{remark}[theorem]{注意}
\newtheorem{definition}[theorem]{定義}
\newtheorem{example}[theorem]{例}
\newtheorem{conjecture}[theorem]{予想}
\renewcommand{\proofname}{証明}
% 三稿で共通に用いる記号。本文と証明はそれぞれ一つの稿に置く。
\providecommand{\ZZ}{\mathbb Z}
\providecommand{\cK}{\mathcal K}
\providecommand{\cP}{\mathcal P}
\providecommand{\cD}{\mathcal D}
\providecommand{\cC}{\mathcal C}
\providecommand{\cE}{\mathcal E}
\providecommand{\cF}{\mathcal F}
\providecommand{\cM}{\mathcal M}
\providecommand{\fm}{\mathfrak m}
\providecommand{\Perf}{\operatorname{Perf}}
\providecommand{\Ku}{\operatorname{Ku}}
\providecommand{\MF}{\operatorname{MF}}
\providecommand{\HH}{\operatorname{HH}}
\providecommand{\Dsg}{D_{\mathrm{sg}}}
\providecommand{\RHom}{\operatorname{RHom}}
\providecommand{\REnd}{\operatorname{REnd}}
\providecommand{\Ind}{\operatorname{Ind}}
\providecommand{\Def}{\operatorname{Def}}
\providecommand{\thick}{\operatorname{thick}}
\providecommand{\Spf}{\operatorname{Spf}}
\providecommand{\Id}{\operatorname{Id}}
\providecommand{\sym}{\operatorname{sym}}
\providecommand{\Cone}{\operatorname{Cone}}
\providecommand{\Pic}{\operatorname{Pic}}
\providecommand{\otL}{\mathbin{\otimes^{\mathbf L}}}
\newtheorem{problem}[theorem]{問題}

\newcommand{\C}{\mathbb C}
\newcommand{\CC}{\mathbb C}
\newcommand{\PP}{\mathbb P}
\newcommand{\OO}{\mathcal O}
\newcommand{\Sym}{\operatorname{Sym}}
\newcommand{\Spec}{\operatorname{Spec}}
\newcommand{\Proj}{\operatorname{Proj}}
\newcommand{\Hom}{\operatorname{Hom}}
\newcommand{\End}{\operatorname{End}}
\newcommand{\Ext}{\operatorname{Ext}}
\newcommand{\Aut}{\operatorname{Aut}}
\newcommand{\GL}{\operatorname{GL}}
\newcommand{\PGL}{\operatorname{PGL}}
\newcommand{\Gr}{\operatorname{Gr}}
\newcommand{\rank}{\operatorname{rank}}
\newcommand{\qgr}{\operatorname{qgr}}
\newcommand{\QGr}{\operatorname{QGr}}

% 他稿の番号には R (可換復元), M (モジュライ・降下), S (混合族) を添える。
% clean build では scripts/build_papers.py が三稿の aux を順に生成・収束させる。
\newcommand{\SetPaper}[1]{%
  \def\CurrentPaper{#1}%
  \ifstrequal{#1}{main}{}{\externaldocument[main-][nocite]{build/nc-cubic}[nc-cubic.pdf]}%
  \ifstrequal{#1}{moduli}{}{\externaldocument[moduli-][nocite]{build/cubic-moduli-descent}[cubic-moduli-descent.pdf]}%
  \ifstrequal{#1}{mixed}{}{\externaldocument[mixed-][nocite]{build/mixed-sklyanin-cubics}[mixed-sklyanin-cubics.pdf]}%
}
\makeatletter
\def\PaperTag@main{R}
\def\PaperTag@moduli{M}
\def\PaperTag@mixed{S}
\newcommand{\PaperRefText}[2]{\textnormal{[\csname PaperTag@#1\endcsname,~\ref{#1-#2}]}}
\newcommand{\PaperEqRefText}[2]{\textnormal{[\csname PaperTag@#1\endcsname,~\eqref{#1-#2}]}}
\DeclareRobustCommand{\paperref}[2]{\ifmmode\text{\PaperRefText{#1}{#2}}\else\PaperRefText{#1}{#2}\fi}
\DeclareRobustCommand{\papereqref}[2]{\ifmmode\text{\PaperEqRefText{#1}{#2}}\else\PaperEqRefText{#1}{#2}\fi}
\DeclareRobustCommand{\paperpageref}[2]{\textnormal{[\csname PaperTag@#1\endcsname,~p.~\pageref{#1-#2}]}}
\newcommand{\PaperReferenceGuide}{%
  本稿と関連する二稿の参照では、R を可換復元稿\cite{NCCubicRealization}、
  M をモジュライ・降下稿\cite{NCCubicModuliDescent}、
  S を混合族稿\cite{NCCubicMixedSklyanin}の略号とする。
  他稿の結果は、例えば定理 [R, 2.1] のように稿の略号とその稿内の番号を併記する。
  略号のない番号は本稿内の参照である。
}
\makeatother

\SetPaper{main}
\title{非可換三次の可換 Kuznetsov 実現\\
\large 一般点復元・楕円因子の直接比較・境界延長}
\author{}
\date{2026年11月22日・研究稿 v56}
\hypersetup{pdftitle={非可換三次の可換 Kuznetsov 実現},
pdfsubject={一般点復元・楕円因子の直接比較・境界延長}}
\begin{document}
\maketitle
\begin{abstract}
非可換三次の Kuznetsov 成分を同次元の滑らかな可換三次の成分として
実現する問題を調べる。指定された中心三次非可換変形について、
完備曲率移送、有限 window、方程式の降下と coherent 飽和により、
一般係数体の有限拡大上で rotation と指定三次周期を保つ可換復元を与える。
実楕円因子の分離和については、有限行列分解の積と因子比較から、
固定された複素係数で全 coherent Kuznetsov 成分の可換実現を直接構成する。
また、通常 Artin 変形の形式認識、可換実現 locus の開性、
DVR 上の比較核の延長判定を扱う。
全ての内在的な非可換三次の可換実現は未解決である。
点スキームと Poisson 零スキームの一致が一般に成り立つという期待には、
実際の Ku 比較を持つ三様体の非被約局所環が反例を与える。
構造対象付きモジュライと精密な降下は関連するモジュライ稿で扱う。
混合族稿では、固定した二つの Sklyanin 因子の全定数補正について
中心超平面の制限随伴と三項までの余乗法を比較する。
さらに同じ因子の任意個 $k\ge2$ へ拡張し、
$V(G-u^2L+\kappa_kcu^3)$ への普遍比較と、その可換 smooth 開上の相対 perfect 核を得る。
さらに全線形 PBW 補正に対し、非分裂 comonad の比較を三項までに還元し、
三角な多項式係数変更を持つ可換三次への普遍残余線束圏比較を構成する。
可換 smooth 開では全 coherent Ku 同値を得る。
一般の $\kappa_k$ と線形族の比較係数の数値化は未完了である。
指定した五様体族の全 smooth 中心係数で、
$V(\lambda(G-u^2L)+cu^3)$ への相対 perfect 比較核を構成する。
定常 norm の全六十四根にも同じ核が適用され、
実際の可換 Ku 同値を持ちながら被約点スキームが Poisson 零集合と一致しない例となる。
\end{abstract}
\setcounter{tocdepth}{1}
\tableofcontents

% !TEX root = ../nc-cubic.tex
\section{序論}
\label{sec:introduction}

\subsection{可換 Kuznetsov 実現の問題}

本稿の目的は、非可換三次の Kuznetsov 成分を、同じ次元の滑らかな
可換三次超曲面の成分として実現することである。
基礎体を $\C$、次元を $n\ge3$ とし、非可換三次 $n$-fold を
$\mathcal X$ と書く。標準的な三次型の例外部分を除いた dg 残余圏を
$\Ku(\mathcal X)$ とする。中心三次表示 $A=B/(w)$ を用いる場合は、
有限線束から生成する残余部分と全 coherent 残余圏を区別し、
window と飽和を証明した範囲で後者と同定する。

\begin{conjecture}[主予想：任意の非可換三次の可換 Ku 実現]
\label{conj:nc-main}
任意の滑らかな非可換三次 $n$-fold $\mathcal X$ に対して、
滑らかな可換三次超曲面 $X\subset\PP^{n+1}_{\C}$ と
$\C$-線型 dg Morita 同値
\[
 \Ku(\mathcal X)\simeq\Ku(X)
\]
が存在する。
\end{conjecture}

この予想の「任意」は、既知の表示空間の稠密開へ弱めない。
全対象の内在的な定義・認識、全 $\mathcal O$-marked モジュライの構成、
表示による全成分の被覆には未完成部分がある。
以下で証明するのは、各定理に明記する中心三次変形や実楕円因子の
表示を持つ対象についての結果である。
一般体への核の降下や指定 rotation・周期の保持は、有効な強化問題であるが、
複素数体上の全対象についての存在問題とは区別する。

\subsection{主結果とその範囲}

第一の結果は、次数別有限自由な中心三次変形の一般点における可換復元である。
可換な特殊ファイバーと滑らかな剰余三次を持つ、指定された代数的変形に対し、
定理\ref{gdr:thm:generic}は元の関数体上に滑らかな可換三次を構成し、
有限拡大上で全 coherent 残余圏の同値を与える。
この同値は rotation と指定三次周期を保つ。
完備曲率移送、三対象 window、通常の一般点への移行、方程式の降下、
coherent 飽和を別々に証明する。
全閉パラメータでの同値や、形式的な核への $h=1$ の代入は結論しない。

第二の結果は、複素数体上の実楕円因子からなる分離型の直接比較である。
定理\ref{pps:thm:elliptic-existence}と定理\ref{csc:thm:comparison}は、
実際の THCR 商を持つ因子と、有限 Jacobi 環を持つ中心三次式について、
有限行列分解の積表示と周期付き因子比較から可換 Ku 実現を構成する。
存在のために因子の $j$ 不変量が互いに異なることや一般的であることは要求しない。
必要な右 THCR の shift、楕円 window、rotation の補正を
第\ref{cf:sec:elliptic-foundations}節にまとめる。
中心変数のない満支持型、および中心変数が一つで三乗係数が非零の型では、
明記した点対応の仮定の下で点スキームと Poisson 零スキームの同時比較も得る。

第三の結果は、可換実現を族の中で進めるための形式認識と延長判定である。
定理\ref{thm:nf-residual-formal}は $n\ne4$ の滑らかな可換三次について、
通常の可換局所 Artin 基底上の変形を、圏変形の $1$-切断 $\tau_{\le1}$ から認識する。
定理\ref{grl:thm:open-locus}は $n\ne4$ について、実際の smooth proper 圏族における
可換実現 locus の開性と \'etale 局所の比較核を与える。
一般点での比較核と perfect な DVR 模型が与えられた場合には、
命題\ref{grl:prop:trait-kernel-test}の単位・余単位の cone が
特殊点への同値延長を判定する。開性だけで全 smooth locus は覆えない。

可換 Ku 実現から点スキームと Poisson 零スキームの一致まで期待することは、
一般には正しくない。本文の注意\ref{ppz:rem:naive-expectation}に
この素朴な期待を述べ、系\ref{ppz:cor:complex-counterexample}に、
実際の Ku 比較が存在しても非被約局所環が一致しない三様体の例を与える。
混合五様体では被約化後にも一致しない例があり、
混合族稿の系\paperref{mixed}{ppz:cor:reduced-point-poisson-counterexample}で扱う。
これらは点スキームによる可換標的の理解の限界を示す補助例である。

\subsection{先行研究との関係}

三様体の中心一次空間、二次包絡、中心三次の直線と残余圏を結ぶ構成
\cite{Threefold}は出発点の一つである。
一般次元では Serre 関手が回復する twist、Hochschild 群、
Poisson 方程式のスキーム構造が変わるため、次元を置き換えるだけでは足りない。
Poisson 正規形と相対 Poisson 層は未公表研究稿\cite{NCCubicA}、
PBW・AS 正則性、指定 Casimir の中心持上げと量子化の存在は
別稿\cite{CubicQuantization}を用いる。
本稿の一般点復元は、入力となる変形を明記して適用する独立した基礎定理である。

非可換 Proj と導来 section の基礎は Artin--Zhang\cite{LP-AZ}および
相対双対の既存理論\cite{LP-BS}を用いる。
楕円 THCR と Sklyanin 代数の構造\cite{GRSklyaninATV,GRSklyaninST}、
特異圏と行列分解の window\cite{NFMOrlovSingularities}を基に、
本稿では有限模型、内部次数、rotation と実中心元の指定周期を追跡する。
一般点の同値を得た後に全 coherent 圏との一致を証明する順序は、
smoothness を比較の前提に紛れ込ませないために必要である。

三様体の全 Poisson 構造の分類\cite[Theorem~1, Section~2.2]{Mayanskiy}と
categorical Torelli\cite[Theorem~1.1]{NFMThreefoldCategoricalTorelli}は、
点スキーム不一致の例で用いる既知の入力である。
本稿の内容は、孤立零点の局所 Artin 環の障害と、直接比較を持つ非可換族への適用である。
圏変形の形式認識には Lurie\cite{NFDAGX}の tame 圏の理論を用いる。
これらの具体的適用の文献上の新規性について、網羅的な確認はなお未完了である。

\subsection{構成と関連する二稿}

本稿は一般点復元の基礎、実楕円因子の積比較、分離型の点対応、
三様体の不一致例、形式認識と可換実現 locus を扱う。
構造対象付きモジュライ、表示認識、自然な中心直線、
Picard 障害と周期位相の精密な降下はモジュライ・降下稿に置く。
自然な標的を指定する相対比較には、そこでの尺度や scalar 自明化の結果も使う。
固定表示に対する直接の存在定理と、この相対的な追加構造は区別する。
混合族の中心、局所代数、正確な smooth locus、可換候補、Eckardt 模型は混合族稿に置く。
同稿の定理\paperref{mixed}{cmo:thm:root-order-residual}は、零係数の同変残余圏を
分離型四様体の Ku と不変環の残余圏に、実際の核で分解する。
接着は球面導来制限であり、可換側の第二成分は Kn\"orrer 核によって三次三様体の Ku となる。
また定理\paperref{mixed}{mix:thm:h-exact-candidate}は、交換子パラメータの形式近傍と
一般係数体上で、指定された可換候補への実比較を与える。
固定値 $h=1$ への核の延長と、元の中心係数族全体の同値は未解決である。
零係数の対合については、同稿の命題\paperref{mixed}{cmo:prop:equivariant-window-hom}で
三対象 window の全 stable Hom と対合指標を計算する。
この有限な計算は Hochschild homology の指標とは区別し、
比較へ必要な高次積と対角の入力を明記する。
元の中心係数 $c$ を動かす族の指定可換候補への Ku 比較は、一般点でも未解決である。

\PaperReferenceGuide

% !TEX root = ../nc-cubic.tex
\section{楕円因子の window と比較核の基礎}
\label{cf:sec:elliptic-foundations}

分離型の直接比較に必要な楕円因子の規約を固定する。
右 twisted homogeneous coordinate ring の内部 shift、
rotation の共役式、体上の Fourier--Mukai intertwiner の存在判定と、
window1 における剰余体の像を同じ規約で記述する。
ここで扱うのは体上の因子比較であり、族の上の核 gerbe や full 残余核の降下は含めない。
楕円曲線上の twist は評価写像の cone
\[
 T_{\mathcal O_E}(F)=\operatorname{Cone}\bigl(
 R\operatorname{Hom}(\mathcal O_E,F)\otimes\mathcal O_E\longrightarrow F\bigr)
\]
で定める。

\subsection{平行移動・線束と rotation の共役}

楕円曲線の点 $b$ を次数零線束 $P_b=\mathcal O_E(b-0)$ と同一視する。
以下の式はまず代数閉体上で証明し、必要な等式を忠実平坦降下して用いる。
右加群の twisted homogeneous coordinate ring の規約は
\cite[\S3, equation~(3.1)]{NFMStaffordVdB}に合わせる。
すなわち $\sigma$ が点加群の shift を表すとき、
coherent sheaf 上の次数 shift は
\begin{equation}
\label{eq:nfm-right-thcr-shift}
 s=\sigma_*\circ(-\otimes L)
\end{equation}
である。特に $s(k(x))=k(\sigma(x))$ となる。
$\sigma^*$ と取り違えると、以下の平行移動の符号が逆になる。

\begin{lemma}[楕円曲線の rotation の共役式]
\label{lem:nfm-elliptic-rotation-conjugacy}
$E$ を標数零の代数閉体上の楕円曲線、$L$ を次数三の線束とする。
\[
 U_{a,b}=t_a^*\circ(-\otimes P_b),\qquad
 \rho=T_{\mathcal O_E}\circ(-\otimes L)
\]
と置く。関手の同型類について
\begin{equation}
\label{eq:nfm-elliptic-rotation-conjugacy}
 \rho U_{a,b}\rho^{-1}\simeq U_{-2a-b,\,3a+b}
\end{equation}
が成り立つ。
$\sigma=t_p$ とし、右加群の規約に従って
$\rho_p=T_{\mathcal O_E}\circ\sigma_*\circ(-\otimes L)$ と置くと、
\begin{equation}
\label{eq:nfm-elliptic-point-intertwiner}
 U_{a,b}\rho_pU_{a,b}^{-1}\simeq\rho
 \quad\Longleftrightarrow\quad
 3a=0,\qquad b=p.
\end{equation}
特に $\rho$ と可換する $U_{a,b}$ は $U_{a,0}$、$a\in E[3]$ に限られ、
これらは次数零 Picard 類を変えない。
\end{lemma}

\begin{proof}
平行移動は次数零 Picard 類に自明に作用するので、
$U_{a,b}$ の同型類の積は $(a,b)$ の加法である。
$t_a^*L\otimes L^{-1}\simeq P_{-3a}$ より
\[
 \operatorname{Ad}_{-\otimes L}(a,b)=(a,b+3a).
\]
また
\[
 \operatorname{Ad}_{T_{\mathcal O_E}}(a,b)=(a-b,b)
\]
である。後式の graph と線束を確認する。
twist の Fourier--Mukai 核は
$\operatorname{Cone}(\mathcal O_{E\times E}\to\mathcal O_\Delta)
=\mathcal O_{E\times E}(-\Delta)[1]$ であり、
$T_{\mathcal O_E}^{-1}k(x)=\mathcal O_E(x)$ となる。
従って共役 $T_{\mathcal O_E}U_{a,b}T_{\mathcal O_E}^{-1}$ は
$k(x)$ を $k(x-a+b)$ へ送り、$\mathcal O_E$ を $P_b$ へ送る。
点層を shift なしの点層へ送る同値の核は、その台である graph 上の線束である。
この場合の graph は $t_{a-b}$ の pullback、線束は $P_b$ なので、所要の式を得る。
これは上記の twist 核を使った畳込み計算とも一致する。

二つの共役作用の積は
\[
 A=\begin{pmatrix}-2&-1\\3&1\end{pmatrix}
\]
であり、\eqref{eq:nfm-elliptic-rotation-conjugacy}を与える。
平行移動は $\mathcal O_E$ を保つので、その pullback は
$T_{\mathcal O_E}$ と可換する。
従って $\rho_p=U_{-p,0}\rho$ であり、intertwining の条件は
\[
 (I-A)\binom ab=\binom p0,
 \qquad\text{すなわち}\qquad
 3a+b=p,\quad -3a=0
\]
に等しい。これで全ての主張が従う。
\end{proof}

\subsection{全楕円同値の分類と二段階の障害}

次に、平行移動と線束の tensor 以外も含む全ての Fourier--Mukai 同値を扱う。
$K$ を標数零の体、$E,C$ を $K$ 上の滑らかな射影 genus-one 曲線とし、
実際の次数三線束 $L_E,L_C$ を固定する。原点や $K$-有理点は仮定しない。
$p\in J_E(K)=\operatorname{Pic}^0(E)(K)$ に対し
\[
 \rho_{E,p}=T_{\mathcal O_E}\circ(t_p)_*\circ(-\otimes L_E),
 \qquad
 \rho_C=T_{\mathcal O_C}\circ(-\otimes L_C)
\]
と置く。
以下で intertwiner というときは、実際の perfect Fourier--Mukai 核による
$K$-線型同値 $F$ に、$F\rho_{E,p}\simeq\rho_CF$ という
\emph{同型の存在条件}を課す。その同型自体の標識と指定周期との整合は、
最後に別に扱う。

\begin{lemma}[数値 centralizer と graph の認識]
\label{lem:nfm-elliptic-numerical-graph}
幾何的数値格子を $(\operatorname{rank},\deg)$ によって $\mathbf Z^2$ と同定する。
$\rho_{E,p}$ と $\rho_C$ の行列はともに
\[
 R=\begin{pmatrix}-2&-1\\3&1\end{pmatrix},\qquad R^3=I
\]
であり、
\begin{equation}\label{eq:nfm-elliptic-numerical-centralizer}
 \operatorname{Cent}_{\operatorname{SL}_2(\mathbf Z)}(R)
   =\{\pm I,\pm R,\pm R^2\}.
\end{equation}
また、実際の $K$-線型 Fourier--Mukai 同値
$F:D^b(E)\to D^b(C)$ の幾何的数値作用が恒等なら、
$K$-同型 $\phi:E\xrightarrow{\sim}C$、$C$ 上の次数零線束 $M$、
整数 $a$ があって
\begin{equation}\label{eq:nfm-elliptic-numerical-graph}
 F\simeq(-\otimes M)\circ\phi_*[2a]
\end{equation}
となる。これらは同型を除いて一意であり、
$F(\mathcal O_E)$ が shift のない次数零線束なら $a=0$ である。
\end{lemma}

\begin{proof}
次数三線束の tensor と $T_{\mathcal O}$ の数値行列は、それぞれ
\[
 \begin{pmatrix}1&0\\3&1\end{pmatrix},\qquad
 \begin{pmatrix}1&-1\\0&1\end{pmatrix}
\]
であり、平行移動の数値作用は恒等である。従ってその積は $R$ となる。
$B=\left(\begin{smallmatrix}a&b\\c&d\end{smallmatrix}\right)$ とすると、
$BR=RB$ は $c=-3b$、$d=a-3b$ に等しい。
さらに
\[
 1=\det B=a^2-3ab+3b^2
   =\left(a-\frac32b\right)^2+\frac34b^2
\]
なので $|b|\le1$ である。
$b=0$ なら $a=\pm1$、$b=1$ なら $a=1,2$、
$b=-1$ なら $a=-1,-2$ であり、これが表示した六行列を与える。

graph の主張をまず代数閉包上で証明する。
滑らかな曲線では $\operatorname{Ext}^{\ge2}=0$ なので、
有界複体は cohomology 層の shift の直和である。
$F(k(x))$ は indecomposable で数値類 $(0,1)$ を持つため、
ある点 $\phi(x)$ と整数 $a(x)$ に対し
$F(k(x))\simeq k(\phi(x))[2a(x)]$ となる。
kernel の第一射影に関する relative perfect 性と Tor amplitude の開性から
$a(x)$ は局所定数であり、$E$ の連結性によって一定である。
shift を除いた kernel は第一射影について平坦な長さ一の層の族なので、
一つの graph 上の線束となる。
逆 kernel にも同じ議論を適用すると、graph の射は同型である。
これで\eqref{eq:nfm-elliptic-numerical-graph}を得る。
数値作用が恒等なので $\deg M=0$ であり、
$F(\mathcal O_E)=M[2a]$ から最後の主張も従う。
この graph 認識は
\cite[Proposition~4.2, Corollary~4.3 and Theorem~5.3]{NFMHilleVdBFM}
の証明と同じ議論である。

$K$ 上への降下では、kernel が実際の $K$ 上の対象であることを用いる。
忠実平坦な体拡大後に cohomology が一つの次数に集中するので、
$K$ 上でも同じ次数の層となる。
その平坦性、graph 上の可逆性は fpqc 局所的性質である。
scheme 理論的 support の $E$ への射は代数閉包上で同型なので、
$K$ 上でも同型となり、$\phi$ と線束が降下する。
graph とその上の線束の一意性から、表示の一意性も従う。
ここで $K$ 上の数値格子が次数一の対象を含むとは仮定していない。
\end{proof}

\begin{theorem}[構造層の像で規格化した全楕円 intertwiner]
\label{thm:nfm-elliptic-normalized-intertwiners}
$F:D^b(E)\to D^b(C)$ を実際の $K$-線型 Fourier--Mukai 同値とし、
$F(\mathcal O_E)$ は shift のない次数零線束とする。
このとき $F\rho_{E,p}\simeq\rho_CF$ であることは、
\begin{equation}\label{eq:nfm-elliptic-normalized-intertwiners}
 F\simeq(-\otimes M)\circ\phi_*,\qquad
 \phi^*L_C\simeq L_E,\qquad
 [M]=\operatorname{Jac}(\phi)(p)
\end{equation}
を満たす $K$-同型 $\phi:E\to C$ と実線束 $M$ があることと同値である。
ここで $\operatorname{Jac}(\phi)$ は $\phi$ による共変な Picard 同型である。
\end{theorem}

\begin{proof}
$F$ の幾何的数値行列 $B$ は $e_1=(1,0)^{\mathsf t}$ を固定し、
$BR=RB$ を満たす。$e_1$ と $Re_1=(-2,3)^{\mathsf t}$ は
$\mathbf Q$ 上独立なので $B=I$ である。
前補題によって $F=(-\otimes M)\phi_*$ と書ける。
\[
 q=\operatorname{Jac}(\phi)(p),\qquad
 L'=\phi_*L_E,\qquad d=[L'\otimes L_C^{-1}],\qquad m=[M]
\]
と置く。
代数閉包上で原点を取り、補題\ref{lem:nfm-elliptic-rotation-conjugacy}の
$U_{a,b}$ を用いると
\[
 \phi_*\rho_{E,p}\phi_*^{-1}=U_{-q,0}\rho_CU_{0,d}
\]
である。従って同補題の共役式から
\begin{equation}\label{eq:nfm-elliptic-graph-conjugation}
 \begin{split}
 F\rho_{E,p}F^{-1}
 &=U_{0,m}U_{-q,0}\rho_CU_{0,d}U_{0,-m}\\
 &=U_{m-q-d,d}\rho_C
 \end{split}
\end{equation}
となる。
$U_{a,b}\simeq\mathrm{id}$ なら、$\mathcal O$ の像から $b=0$、
点層の像から $a=0$ である。
従って intertwining は $d=0$、$m=q$ に等しい。
特に非偏極的な $\phi$ のずれを線束 $M$ で補償することはできない。

これらは相対 Picard 類の等式なので $K$ 上へ降りる。
$d=0$ は、実線束 $L'\otimes L_C^{-1}$ が
$\operatorname{Spec}K$ からの線束であることを意味し、従って自明である。
逆向きの functor の同型も $K$ 上で存在する。
実際、比較する二つの同値 kernel は代数閉包上で同型であり、
その $\operatorname{Hom}^0$ は $\mathrm{HH}^0(E_{\overline K})=\overline K$
によって一次元である。
proper perfect Hom の体拡大により $K$ 上でも一次元であり、
その非零元は代数閉包上で同型となるため、既に $K$ 上で同型である。
\end{proof}

$p\in J_E(K)$ に対し、$p$ を表す実際の次数零線束の Picard gerbe の類を
$\delta_E(p)\in\operatorname{Br}(K)$ と書く。
その消滅は、$p$ を表す線束が $E$ 上に存在することを意味する。

\begin{corollary}[全数値型を含む active 同値の存在判定]
\label{cor:nfm-elliptic-any-intertwiner-criterion}
$F(\mathcal O_E)$ に条件を課さない場合にも、
実際の $K$-線型 Fourier--Mukai intertwiner
$(D^b(E),\rho_{E,p})\to(D^b(C),\rho_C)$ が存在することは、
次の二条件と同値である：
\begin{enumerate}
 \item $(E,L_E)$ と $(C,L_C)$ の $K$ 上の偏極曲線同型が存在する。
 \item $\delta_E(p)=0$、すなわち $p$ は $E$ 上の実際の次数零線束で表される。
\end{enumerate}
特に同じ $(E,L)$ 上では
\begin{equation}\label{eq:nfm-elliptic-any-intertwiner-obstruction}
 \{F:F\rho_{E,p}\simeq\rho_EF\}\ne\varnothing
 \quad\Longleftrightarrow\quad \delta_E(p)=0.
\end{equation}
この主張は intertwining 同型自体の周期付き rigidification を含まない。
\end{corollary}

\begin{proof}
任意の intertwiner の数値行列は
\eqref{eq:nfm-elliptic-numerical-centralizer}に属する。
$K$ 上で定義された $\rho_C$ の冪と shift を postcompose して、
数値作用を恒等にする。
補題\ref{lem:nfm-elliptic-numerical-graph}によって graph と
次数零線束の偶 shift を得るので、その shift も除けば前定理を適用できる。
従って偏極曲線同型 $\phi$ と
$\operatorname{Jac}(\phi)(p)$ を表す実線束を得る。
実際の $K$-同型 $\phi$ による線束の移送は Picard gerbe の band を保つから
\[
 \delta_C(\operatorname{Jac}(\phi)(p))=\delta_E(p)
\]
であり、必要性を得る。
逆に $p$ の実線束を $\phi$ で移送すれば、前定理から所要の同値を得る。
\end{proof}

\subsection{右座標環の window と内部 shift}

以下の window の議論では基礎体を $k$ と書き、
$\kappa_A$ を増大剰余体 $k=A/A_{\ge1}$ の特異圏における像とする。

\begin{lemma}[右楕円座標環の window1 と内部 shift]
\label{lem:nfm-twisted-elliptic-window}
$A=B(E,L,\sigma)$ を右 twisted homogeneous coordinate ring とする。
$\operatorname{qgr}A\simeq\operatorname{coh}E$ の下で
$\pi A\leftrightarrow\mathcal O_E$、$(1)\leftrightarrow s$ は
\eqref{eq:nfm-right-thcr-shift}に従うものとする。
$A$ に Gorenstein parameter 零の Orlov 同値が適用でき、
window1 を $\Phi_1=q_{\mathrm{sg}}R\Gamma_{\ge1}$ と取る。
このとき
\begin{equation}
\label{eq:nfm-twisted-active-window}
 \kappa_A\longleftrightarrow\mathcal O_E[1],\qquad
 \Phi_1^{-1}(1)\Phi_1=T_{\mathcal O_E}\circ s.
\end{equation}
\end{lemma}

\begin{proof}
この window の存在については
\cite[Theorem~2.5(iii)]{NFMOrlovSingularities}を用いる。
$m>0$ に対して $s^m\mathcal O_E$ は次数 $3m$ の線束なので、
その高次 cohomology は消える。
従って $R\Gamma_{\ge1}\mathcal O_E=A_{\ge1}$ であり、
$0\to A_{\ge1}\to A\to k\to0$ を特異圏へ写すと第一式を得る。
任意の $F$ に対し
\[
 (R\Gamma_{\ge1}F)(1)=R\Gamma_{\ge0}(sF)
\]
であり、次数零の切断は三角形
\[
 R\Gamma_{\ge1}(sF)\longrightarrow R\Gamma_{\ge0}(sF)
 \longrightarrow R\operatorname{Hom}(\mathcal O_E,sF)\otimes k
 \longrightarrow
\]
を与える。
特異圏へ写し、第一式を用いて三角形を回転すると、評価写像の三角形
\[
 R\operatorname{Hom}(\mathcal O_E,sF)\otimes\mathcal O_E
 \longrightarrow sF\longrightarrow\Phi_1^{-1}(1)\Phi_1(F)
 \longrightarrow
\]
を得る。この cone は $T_{\mathcal O_E}(sF)$ である。
従って剰余体と内部 shift の双方について同じ window の規約を得た。
この切断の議論は $A$ の可換性を用いない。
\end{proof}


\section{主定理と証明の構成}\label{gdr:sec:main}

\subsection{入力、残余圏と指定周期の規約}
$n\ge3$、$N=n+2$、$a=N-3$とし、
\[
 R=\mathbb C[h]_{(h)},\qquad \widehat R=\mathbb C[[h]],\qquad
 K=\mathbb C(h),\qquad L=\mathbb C((h))
\]
と置く。非負次数付き中央$R$-代数$A_R$について、
$(A_R)_0=R$、各$(A_R)_j$が有限自由で、
$A_R/hA_R\simeq\mathbb C[x_1,\ldots,x_N]$と仮定する。
$\Omega_R\in Z(A_R)_3$の剰余を$F$とし、$V(F)\subset\mathbb P^{N-1}$は滑らかとする。
$B_R=A_R/(\Omega_R)$と置く。

$B=B_K$について$\mathcal O(j)=\pi B(j)$とし、
\[
 \mathcal C_B=\operatorname{thick}\{\mathcal O(j):j\in\mathbb Z\},\qquad
 \mathcal E_B=\langle\mathcal O,\ldots,\mathcal O(a-1)\rangle
\]
と定める。二つの残余圏
\begin{equation}\label{gdr:eq:residual-defs}
 \mathcal K_B^{\rm lin}=\mathcal C_B\cap\mathcal E_B^\perp,
 \qquad
 \mathcal K_B^{\rm coh}=D^b(\operatorname{qgr}B)\cap\mathcal E_B^\perp
\end{equation}
を区別し、rotationは$\rho_B=L_{\mathcal O}\circ(1)$とする。
滑らかな三次超曲面$Y\subset\mathbb P^{n+1}$についても
$\operatorname{Ku}(Y)=\langle\mathcal O_Y,\ldots,\mathcal O_Y(n-2)\rangle^\perp$、
$\rho_Y=L_{\mathcal O_Y}\circ(1)$を用いる。
内部shiftは$M(j)_q=M_{q+j}$とし、$[r]$はcohomological shiftとする。
全ての圏はdg enhancementと冪等完備化を伴う。
\(\operatorname{MF}^{\mathbb Z}(A,\Omega)\)は有限自由graded MFのdg圏を表し、
Morita同値を述べる際にはその有限coneと冪等完備化も含める。
中心元$\Omega$から定まる三次周期を$\tau^{\rm MF}:(3)\to[2]$と書く。
特異圏では導来自己交叉の境界写像$\eta(3)$がこれに対応し、
残余圏ではwindowを通して$\tau^{\rm rot}:\rho_B^3\to[2]$を定める。
この規格化を図式\eqref{tw:eq:rotation-period-square}で追跡する。
以下の「周期を保つ」は、この指定された自然比較を保つことをいう。

\subsection{一般点の可換模型}

\begin{theorem}[一般次元の中心三次非可換変形の可換復元]
\label{gdr:thm:generic}
以上の仮定の下で、$K$上に幾何学的に滑らかな三次超曲面
$Y\subset\mathbb P^{n+1}_K$が存在し、ある有限拡大$K'/K$上で
\begin{equation}\label{gdr:eq:main}
 (\mathcal K^{\rm coh}_{B_{K'}},\rho_B)
 \simeq(\operatorname{Ku}(Y_{K'}),\rho_Y)
\end{equation}
となる。このMorita同値は、中心元に由来する三次周期比較も保つ。
さらに、$D^b(\operatorname{qgr}B_K)$はsmooth properであり、
$\bigoplus_{j=0}^{N-1}\mathcal O(j)$が生成する。
$Y$の幾何学的な射影同型類は、上のrotation付きpairを実現する
滑らかな三次超曲面の中で一意である。
また、代数閉包上の任意の非零$M\in\mathcal K^{\rm coh}_{B_K}\otimes_K\overline K$について、
$M\oplus\rho_BM\oplus\rho_B^2M$の生成段数は$2(n+2)(n+3)-1$以下である。
\end{theorem}

\begin{proof}
命題\ref{gcd:prop:automatic-as}により一般ファイバーのambientは
noetherian Koszul AS正則であり、$\Omega_R$は正則である。
定理\ref{tw:thm:window}は線束生成残余圏を、原点のstabilizationと
その二つの内部shiftのendomorphism代数で表示する。
この固定三対象に定理\ref{cft:thm:complete}の完備MF比較を適用してから局所化する。
命題\ref{rog:prop:euler}の非零EulerベクトルとBFKの生成定理
（命題\ref{gcd:prop:bound}）により、代数閉拡大体上で像は可換MF圏全体となる。
定理\ref{rog:thm:geometric-generic}の有限降下によって、有限拡大後の
滑らかな可換三次模型と線束生成残余圏の同値を得る。

補題\ref{crs:lem:descent}により残余圏のsmoothnessは$K$へ降下する。
命題\ref{crs:prop:residual-smooth}と定理\ref{crs:thm:coherent}を
$K$と有限拡大上で適用すれば、線束生成部分は全coherent導来圏に一致する。
これにより全coherent残余圏の同値と、全体のsmooth proper性・有限線束生成を得る。

最後に、得られたsmooth properなrotation付き圏へ
定理\ref{gcd:thm:field-model}を適用する。
低rotation積の核を勾配空間として積分することで、可換模型の方程式を$K$上に選べる。
必要なら同値核の有限定義体を拡大すれば式\eqref{gdr:eq:main}を得る。
幾何学的一意性は命題\ref{prop:ej-projective-uniqueness}、
生成段数の上界は命題\ref{gcd:prop:bound}から従う。
\end{proof}

生成段数の上界はBFK\cite[Theorem~5.15]{GCDBFK}に$d=3$を代入した帰結である。
主定理の存在結論と、この既知の定量的帰結を区別する。

\subsection{結論に含まれないことと比較の範囲}\label{gdr:subsec:scope}
主定理、完備移送定理、方程式降下定理の定義体と対象範囲をここに集約する。
\begin{enumerate}[label=(\arabic*)]
\item \textbf{入力と復元する圏。}
入力は指定された次数別有限自由な代数的中心三次対である。
任意の抽象的な非可換三次圏にこの表示があることや、$A_K,B_K$の可換性は結論しない。
復元するのはrotation付き残余pairであり、
$D^b(\operatorname{qgr}B_K)\simeq D^b(Y)$を主張するものでもない。
\item \textbf{方程式と同値核の体。}
方程式は$K$上に選べるが、同値核とその比較図式について保証するのは
ある有限拡大$K'/K$上の存在である。
一意性は共通の代数閉拡大体上の射影同型類について述べる。
$K$上のformは一意とは限らず、例\ref{gcd:ex:forms}がこの違いを示す。
同例は同値核の降下障害が非零であることを証明したものではない。
\item \textbf{線束部分と全coherent圏。}
式\eqref{gdr:eq:residual-defs}の二つの圏は、証明の最初から同一視しない。
局所proper模型は対角のperfect性まで保証するものではなく、
一般点の可換模型からsmoothnessを得た後、
定理\ref{crs:thm:coherent}で全coherent圏との一致を示す。
任意の相対底上でこの一致やenhancedなOrlov同値が成立するとは主張しない。
\item \textbf{完備底と通常の局所化。}
完備移送は各内部次数の有限自由加群内で行い、内部次数についての無限級数は導入しない。
一般点では固定三対象の通常のHom複体を先に構成し、その後$h$を反転する。
任意のLaurent係数MFへ無限Taylor和を直接評価することや、
torsion Morita同値だけから通常の一般点同値を得ることは用いない。
任意の完備noetherian局所底を許すのは定理\ref{cft:thm:complete}であり、
windowや全coherent圏の比較をその一般性へ拡張するものではない。
\item \textbf{特殊化と具体計算。}
$Y$が元の$V(F)$の定数族であること、全閉パラメータでの同値、
特殊点を含む代数的相対同値は結論しない。
基礎体上で積分した方程式が$h=0$で正則であることや、その剰余が$F$であることも
方程式降下定理だけからは従わない。
非可換入力から得る形式的potentialの有限jetと、
$K$上の方程式・rotation積・同値核の全係数の計算は別の問題である。
第\ref{gre:sec:example}節のHesse族では、追加の楕円曲線比較によって
特殊ファイバーを保つ明示的な$K$方程式まで求めるが、
同値核の具体行列と有限拡大の除去は行わない。
\end{enumerate}

% !TEX root = ../../nc-cubic.tex
\section{一般次元の局所代数的模型：仮定と主定理}
\label{lp:sec:main}
本節では、既に構成された次数別flatな量子化から、射影化のproperなdg模型を作る。
Poisson成分のブロック数、純ブロックの楕円性、有限中心性は仮定しない。
ただし、量子化代数そのものの全次数flatnessは入力である。
基底は完備化やArtin切断ではなく、可換点における実在の正則局所環とする。

\begin{definition}[局所入力]\label{lp:def:input}
$(R,\mathfrak m,k)$を有限生成複素代数の局所化である正則局所環とし、
$d=\dim R$、$N\ge5$とする。
$A=\bigoplus_{j\ge0}A_j$は中央$R$-代数、$A_0=R$で、各$A_j$は有限自由$R$-加群、
$A\otimes_R k\simeq k[z_1,\ldots,z_N]$を次数付き代数として満たすとする。
$\Omega\in Z(A)_3$が与えられ、その剰余$F$は滑らかな
$cubic\ X=V(F)\subset\mathbb P^{N-1}_k$を定めるとする。
$B=A/(\Omega)$、$n=N-2$、$a=N-3$と置く。
内部shiftは$M(j)_q=M_{q+j}$、複体はcohomological conventionを用いる。
\end{definition}

\begin{theorem}[局所proper模型と有界Orlov分解]\label{lp:thm:main}
定義\ref{lp:def:input}の下で以下が成立する。
\begin{enumerate}[label=(\arabic*)]
\item $A,B$は左右noetherianである。$A$の$R$上の対角には長さ$N$の有限自由双加群分解があり、
$\operatorname{gldim}A\le N+d$である。$\Omega$は中心正則で、各$B_j$は有限自由、
その形成は任意の通常の底変更と両立する。
\item 左右で
\[
 \operatorname{R\underline{Hom}}_A(R,A)\simeq R(N)[-N],\qquad
 \operatorname{R\underline{Hom}}_B(R,B)\simeq R(a)[-(N-1)].
 \tag{G}\label{lp:eq:G}
\]
$B$は有限左右injective dimensionを持ち、Burke--Stevensonの意味でGorensteinである。
従って通常の有界三角圏について、$D^b(\operatorname{qgr}B)$は
$D_{\rm sg}^b(\operatorname{gr}B)$と$a$個の$\operatorname{Perf}(R)$成分による
Orlov型半直交分解を持つ。特異圏では自然な三次周期性$(3)\simeq[2]$が成立する。
\item $\mathcal O(j)=\pi B(j)$とすると
\[
 \operatorname{Ext}^p_{\operatorname{qgr}B}(\mathcal O(i),\mathcal O(j))\simeq
 \begin{cases}
 B_{j-i}&p=0,\\
 B_{i-j-a}^{\vee}&p=n,\\
 0&\text{その他}
 \end{cases}\qquad (B_q=0\text{ for }q<0).
 \tag{H}\label{lp:eq:H}
\]
ここで$(-)^\vee=\operatorname{Hom}_R(-,R)$である。
この計算の自然な導来Homの底変更写像は、命題\ref{lp:prop:basechange}の通常の底変更で擬同型となる。
\item $\mathcal C_B=\operatorname{thick}_R\{\mathcal O(j):j\in\mathbb Z\}$を
$D^b(\operatorname{qgr}B)$の標準dg enhancement内の冪等完備な部分と定めると、
\[
 E=\bigoplus_{i=0}^{N-1}\mathcal O(i),\qquad
 \mathcal C_B=\operatorname{thick}_R(E)\simeq\operatorname{Perf}(\Lambda_R),
 \quad\Lambda_R=\operatorname{REnd}(E)\in\operatorname{Perf}(R).
 \tag{P}\label{lp:eq:proper-model}
\]
\item $\mathcal K_B^{\rm lin}=\mathcal C_B\cap
\langle\mathcal O,\ldots,\mathcal O(a-1)\rangle^\perp$は
有限個のmutationにより生成対象$G$を持ち、
$\mathcal K_B^{\rm lin}\simeq\operatorname{Perf}(\Gamma_R)$、
$\Gamma_R=\operatorname{REnd}(G)\in\operatorname{Perf}(R)$である。
特殊ファイバーでは$\mathcal C_B\otimes_R^{\mathbf L}k\simeq\operatorname{Perf}(X)$、
$\mathcal K_B^{\rm lin}\otimes_R^{\mathbf L}k\simeq\operatorname{Ku}(X)$を得る。
\end{enumerate}
\end{theorem}

\begin{remark}[主証明における局所模型の役割]\label{lp:rem:limits}
本節はproper模型と有界Orlov分解を具体化する。
$\Gamma_R^e$上の対角のperfect性と全coherent圏への拡張は、
第\ref{gdr:subsec:scope}節(3)の順序に従い、一般点の可換模型を得た後に証明する。
従って、この段階では上付き${\rm lin}$を保って議論する。
\end{remark}

\subsection{先行研究と本節の寄与}
非可換射影圏のsection functorはArtin--Zhangの枠組み\cite{LP-AZ}に属する。
次数零部が体でないGorenstein環の有界Orlov分解には
Burke--Stevenson\cite[Definition 5.1, Lemmas 6.1--6.2, Theorem 6.4]{LP-BS}を用いる。
dg生成対象によるMorita表示はKeller\cite[Theorem 3.8]{LP-Keller}の一般理論である。
これらの定理自体を新しい結果とは扱わない。
以下では、与えられた量子化に対してnoetherian性・Gorenstein型・Homの底変更・有限生成を導き、
相対proper模型の適用条件を、正則局所基底上で具体的に検証する。

\section{局所基底でのnoetherian性と有限Koszul持上げ}
\label{lp:sec:lift}
\begin{lemma}[次数別有限性を用いたnoetherian性]\label{lp:lem:noetherian}
$(R,\mathfrak m)$をnoetherian局所環、$D=\bigoplus_{j\ge0}D_j$を中央$R$-代数とし、
各$D_j$は有限$R$-加群とする。
$D/\mathfrak mD$が右noetherianなら$D$は右noetherianである。左にも同様である。
\end{lemma}
\begin{proof}
$\mathfrak m$の有限生成元を選ぶと、
$(D/\mathfrak mD)[t_1,\ldots,t_e]\twoheadrightarrow\operatorname{gr}_{\mathfrak m}D$
という中央変数の多項式拡大からの全射がある。
従って右辺は右noetherianである。
斉次右イデアル$I\subset D$に誘導filtration $I\cap\mathfrak m^qD$を入れる。
その随伴次数イデアルの有限個の二重斉次生成元を、$f_1,\ldots,f_r\in I$へ持ち上げ、
$J=\sum f_iD$とする。
$x\in I_j$の最低$\mathfrak m$-次数記号を消す操作を繰り返すと、任意の$q$について
$x\in J_j+\mathfrak m^qD_j$を得る。
$D_j/J_j$は有限$R$-加群だからKrull交叉定理より$x\in J_j$である。
従って全ての斉次右イデアルは有限生成である。
任意の右イデアルには内部次数の上昇filtrationを入れ、最高次数項のイデアルの
有限生成元を持ち上げる。内部次数の有限降下により元のイデアルも有限生成である。
この証明は$R$の完備性を使わず、各内部次数における有限性を使う。
\end{proof}

\begin{lemma}[次数別分裂する自由分解の持上げ]\label{lp:lem:resolution}
$D$の各次数が有限自由$R$-加群で、$M$も各次数有限自由とする。
$\overline M$が$\overline D=D/\mathfrak mD$上、有限階数の次数付き自由加群による
有限分解$\overline P_\bullet$を持つなら、同じ項の階数とshiftを持つ
$D$-自由分解$P_\bullet\to M$を持ち上げられる。
分解は各内部次数で$R$-加群の分裂完全列であり、任意の通常の底変更後も完全である。
\end{lemma}
\begin{proof}
$\overline P_0$の生成元の像を$M$の同じ次数へ持ち上げる。
各次数の余核は有限$R$-加群で剰余が零なのでNakayamaの補題で零である。
$P_0\twoheadrightarrow M$は各次数で有限自由加群への全射だから分裂し、
核は各次数有限自由となる。剰余体への底変更後の核も正しく計算される。
核の生成元を$\overline P_1$から同様に持ち上げ、帰納する。
最後の核は各次数有限自由かつ剰余が零であるから零である。
各段階の内部次数別分裂により、底変更後にも完全性が残る。
\end{proof}

\begin{proposition}[全次数の相対Koszul性とambientの対角]\label{lp:prop:koszul}
定義\ref{lp:def:input}の$A$は次数一で生成され、二次関係で定まる。
$R$の左右$A$-加群としての分解と、$A$の$A^e=A\otimes_RA^{\rm op}$-加群としての分解は
それぞれ
\[
 P_j=A(-j)^{\binom Nj},\qquad
 Q_j=A^e(-j)^{\binom Nj}\quad(0\le j\le N)
 \tag{K}\label{lp:eq:koszul}
\]
という項を持ち、全ての通常の底変更と両立する。
特に$A$は$R$上smoothで、$\operatorname{gldim}A\le N+d$である。
\end{proposition}
\begin{proof}
$A_1$の剰余は$N$変数の空間であり、そのliftによる自由代数からの写像は
各次数のNakayamaにより全射となる。
多項式環の$k$のKoszul分解に補題\ref{lp:lem:resolution}を適用すると$P_\bullet$を得る。
第一次関係が次数二で生成されることもこの構成から従う。
同様に$A$を$A^e$-加群として見て、多項式環の対角Koszul分解を持ち上げると$Q_\bullet$を得る。
$A$は$R$上の自由加群（次数別の直和）なので$Q_\bullet$は右$A$-加群としても分裂完全である。
任意の左$A$-加群$M$を右からテンソルすると、項$A\otimes_R M$の有限直和からなる長さ$N$の分解になる。
$\operatorname{gldim}R=d$より各項の$A$-射影次元は高々$d$。
従って$\operatorname{pd}_A M\le N+d$で、右側も同様である。
\end{proof}

\begin{proposition}[相対Gorenstein型と中心商]\label{lp:prop:gorenstein}
式\eqref{lp:eq:G}が成立する。$B$は左右noetherian、各次数有限自由で、
左右injective dimensionは高々$N+d-1$である。
\end{proposition}
\begin{proof}
補題\ref{lp:lem:noetherian}を$A$に適用する。
$P_\bullet$を$A$で双対化した有界複体は、各内部次数で有限自由$R$-複体である。
剰余上ではKoszul双対性により、内部次数$-N$・cohomological degree $N$以外で完全。
有限自由有界複体の剰余が完全なら、単元行列成分の消去とNakayamaにより複体自身も分裂完全である。
内部次数$-N$では最後の項$R$のみ残る。従って$\operatorname{R\underline{Hom}}_A(R,A)=R(N)[-N]$を得る。
左側も同じである。

各次数での$\Omega:A_{j-3}\to A_j$の剰余は非零多項式$F$の乗法で単射。
行列の列数と同じ大きさの小行列式の一つが単元なので、この写像は局所基底変換で分裂単射となる。
よって$\Omega$は正則、$B_j$は自由で、全底変更と両立する。
中央正則元に対するchange-of-ringsは
\[
 \operatorname{R\underline{Hom}}_A(M,A)
 \simeq\operatorname{R\underline{Hom}}_B(M,B)(3)[-1]
 \quad(M\in D(\operatorname{Gr}B))
 \tag{CR}\label{lp:eq:change-rings}
\]
を与える。これは$0\to A(-3)\xrightarrow{\Omega}A\to B\to0$から
$\operatorname{R\underline{Hom}}_A(B,A)=B(3)[-1]$を計算し、導来随伴を使った式である。
$M=R$を代入して$B$の式\eqref{lp:eq:G}を得る。
一般の$M$に適用すると、$A$の有限大域次元から$B$の有限injective dimensionが従う。
商なのでnoetherian性も成り立つ。
\end{proof}

\section{射影化の局所コホモロジーと底変更}
\label{lp:sec:cohomology}
$\tau$を$B_{\ge1}$-torsion functor、$\pi:\operatorname{Gr}B\to\operatorname{QGr}B$を商関手とする。
Burke--Stevenson\cite[Section 4, Lemmas 6.1--6.2]{LP-BS}により、
上記のGorenstein型と$R$-flatnessは$\chi$と必要な有限cohomological dimensionを与える。
従って$D^b(\operatorname{qgr}B)$を$D(\operatorname{QGr}B)$で計算でき、導来sectionの三角
\[
 R\tau(B)\longrightarrow B\longrightarrow R\omega\pi(B)\longrightarrow
 \tag{LC}\label{lp:eq:torsion}
\]
を使用できる。

\begin{lemma}[局所コホモロジーの各次数]\label{lp:lem:local-cohomology}
$d_B=N-1$とすると
\[
 H^q(R\tau B)_t=
 \begin{cases}
  B_{-t-a}^{\vee}&q=d_B,\\
  0&q\ne d_B
 \end{cases}
 \tag{LC'}\label{lp:eq:localcohom}
\]
である。負添字は零とする。
\end{lemma}
\begin{proof}
$M_\ell=B/B_{\ge\ell}$の次数filtrationの商は
$B_j\otimes_R R(-j)$（$0\le j<\ell$）である。
式\eqref{lp:eq:G}と$B_j$の自由性により、各商の$B$への導来Homは次数$d_B$に集中し、
その内部次数は$-a-j$のみである。
長完全列によって$\operatorname{Ext}_B^q(M_\ell,B)=0$（$q\ne d_B$）。
内部次数$t$では$j=-t-a$のただ一層だけが寄与するので、その層が含まれれば群は$B_j^\vee$となる。
$\ell>j$で推移写像は同型である。
$R\tau B=\varinjlim_\ell\operatorname{R\underline{Hom}}_B(M_\ell,B)$のcohomologyを取ればよい。
\end{proof}

\begin{proof}[式\eqref{lp:eq:H}の証明]
$\operatorname{RHom}_{\operatorname{QGr}B}(\pi B,\pi B(t))$
は導来sectionの内部次数$t$に一致する。
式\eqref{lp:eq:torsion}と補題\ref{lp:lem:local-cohomology}を用いると
$H^0=B_t$、$H^{N-2}=B_{-t-a}^\vee$で、それ以外は零となる。
$a>0$なので同じ$t$に対して二群が同時に非零となることはない。
従ってこれはcohomology次元だけでなく、各Hom複体が有限自由加群の一つのshiftであることを示す。
$i$のtwistを外せば式\eqref{lp:eq:H}である。
\end{proof}

\begin{proposition}[線束Homの底変更]\label{lp:prop:basechange}
$R\to R'$を有限大域次元のnoetherian正則環への通常の準同型で、$A'=A\otimes_RR'$が左右noetherianであるものとする。
$B'=B\otimes_RR'$と置けば自然な写像
\[
 \operatorname{RHom}_{\operatorname{qgr}B}(\mathcal O(i),\mathcal O(j))
 \otimes_R^{\mathbf L}R'
 \longrightarrow
 \operatorname{RHom}_{\operatorname{qgr}B'}(\mathcal O'(i),\mathcal O'(j))
 \tag{BC}\label{lp:eq:bc}
\]
は擬同型である。
特に任意の剰余体$R\to\kappa(\mathfrak p)$、および正則noetherian局所環への局所準同型には適用できる。
\end{proposition}
\begin{proof}
$A,B$の有限相対分解と$B$の正則中心商表示は底変更と両立するので、
$A',B'$も対応するGorenstein型、有限injective dimension、$\chi$を持つ。
剰余体については中央の商・局所化によりnoetherian性を得る。
局所準同型の場合は剰余が多項式環なので補題\ref{lp:lem:noetherian}を再適用できる。

まず$R$の有限大域次元により$-\otimes_R^{\mathbf L}R'$は有限Tor振幅を持ち、
有界coherent graded複体を有界coherent複体へ送り、torsionを保つ。
従って射影商の有界導来圏にも底変更が定義される。
自然な底変更写像そのものを確認するため、$M_\ell$を有限階数graded自由$B$-加群の分解$P_\bullet$で計算する。
$M_\ell$は内部次数別$R$-自由なので、そのsyzygyを各次数で分裂させて、分解の全底変更を保てる。
固定した内部次数$t$で$C^\bullet=\operatorname{\underline{Hom}}_B(P_\bullet,B)_t$は
下に有界な有限自由$R$-加群の複体であり、前補題の証明により、
cohomologyは有限自由加群の一つの次数にしか現れない。
$R$の大域次元は有限なので、cohomologyが全て射影なこの複体は、そのcohomology複体にhomotopy同値である。
実際、射影加群からなる完全複体のcycleについてdimension shiftingを$d+1$回適用すると、
全cycleが射影となり複体が分裂する。
従って$C^\bullet$の導来テンソルは通常のテンソルで計算できる。
有限階数自由分解上ではHomと底変更が項ごとに可換であるため、$M_\ell$のHomの底変更を得る。
各$t$で安定するfiltered colimitと式\eqref{lp:eq:torsion}を通せば\eqref{lp:eq:bc}となる。
\end{proof}

\section{有限生成proper部分圏と残余圏}
\label{lp:sec:proper}
\begin{proof}[有界Orlov分解と有限生成の証明]
命題\ref{lp:prop:gorenstein}の結果はBurke--StevensonのDefinition 5.1の仮定を満たす。
$R$は有限大域次元で、$A,B$は中央$R$-flatなので、同論文Lemma 6.2によって$\chi$を満たす。
Theorem 6.4を$A$へ適用すると、$D_{\rm sg}^b(\operatorname{gr}A)=0$であるため
任意の連続した$N$個のtwistが$D^b(\operatorname{qgr}A)$を生成する。
$B$への導来制限は$0\to A(-3)\to A\to B\to0$により有限Tor振幅を持ち、
torsionをtorsionに送るので射影商へ降りる。
この関手は$\mathcal O_A(j)$を$\mathcal O_B(j)$へ送る。
従って全ての$\mathcal O_B(j)$は$E$から有限回のcone・直和因子で作られる。
これが式\eqref{lp:eq:proper-model}の生成性である。

$B$へのTheorem 6.4は、$a>0$より$a$個のtwistとgraded singularity categoryによる半直交分解を与える。
この有界分解の向きは引用定理の規約に従う。以下では線束生成部分における右直交を別に構成する。
この有界三角圏の同値を、任意の導来基底変更と両立するdg同値へ強めたとは主張しない。
\end{proof}

\begin{proposition}[properな残余dg代数]\label{lp:prop:residual}
$\mathcal E=\langle\mathcal O,\ldots,\mathcal O(a-1)\rangle$とする。
各成分は$\operatorname{Perf}(R)$であり、$\mathcal C_B$で
\[
 \mathcal C_B=\langle\mathcal K_B^{\rm lin},\mathcal O,\ldots,\mathcal O(a-1)\rangle
 \tag{SOD}\label{lp:eq:sod}
\]
を持つ。$p=L_{\mathcal O}\circ\cdots\circ L_{\mathcal O(a-1)}$と置くと
$G=p(E)$は$\mathcal K_B^{\rm lin}$を生成する。
$\Gamma_R=\operatorname{REnd}(G)$の底上の複体はperfectであり、
命題\ref{lp:prop:basechange}の底変更でこの構成はMorita同値まで可換である。
\end{proposition}
\begin{proof}
式\eqref{lp:eq:H}から$\operatorname{RHom}(\mathcal O(i),\mathcal O(i))=R$で、
$0\le i<j\le a-1$なら$\operatorname{RHom}(\mathcal O(j),\mathcal O(i))=0$である。
全$\mathcal C_B$の対象間のHomは、$E$間のperfect Homから有限cone・直和因子で作られるので$R$-perfectである。
従って各例外成分の左右随伴はHomとその$R$-双対で定義でき、mutationは有限操作となる。
右端の例外対象から順に左mutationして$p$を作ると、$p(\mathcal C_B)$は正確に$\mathcal E^\perp$である。
$E$の生成性より$G$が生成し、$\operatorname{REnd}(G)$も$R$-perfectとなる。
標準dg enhancementと生成対象のdg Yonedaを用いればMorita表示を得る\cite[Theorem 3.8]{LP-Keller}。

式\eqref{lp:eq:bc}は生成対象間の全導来Homと積に自然なので、有限coneと冪等完備化へ延長すると
線束生成部分圏の底変更で完全忠実性を得る。像は全線束を生成するのでMorita同値である。
評価写像とそのconeも底変更に可換で、$p,G$についても同じである。
特殊ファイバーでは$X$は滑らかで、全twistは$\operatorname{Perf}(X)=D^b(\operatorname{coh}X)$を生成する。
右直交は通常のKuznetsov成分なので定理\ref{lp:thm:main}の第五項が従う。
\end{proof}

\begin{lemma}[局所proper圏の完全忠実性の持上げ]\label{lp:lem:ff-lift}
$\mathcal C,\mathcal D$を$R$-properな冪等完備dg圏とし、対象間のHomが剰余体への導来底変更と可換であるとする。
$\Phi:\mathcal C\to\mathcal D$を底変更と両立するexactなdg関手とする。
$\Phi_k$が完全忠実なら$\Phi$も完全忠実である。
さらに$\Phi_k$が同値で、$\Phi$が右または左随伴を持つなら$\Phi$は同値である。
\end{lemma}
\begin{proof}
任意の対象$M,N$に対し、Hom上の自然な写像のconeは$R$-perfectで、その剰余は零である。
局所環上の導来Nakayamaによりconeは零となるので完全忠実性を得る。
また$Q\in\mathcal D$について$Q_k=0$なら、その自己Homの剰余も零である。
従って$\operatorname{RHom}(Q,Q)=0$、特に恒等射が零なので$Q=0$となる。
$\Phi$の像が片側admissibleであれば、任意の対象は像とその直交成分によるtriangleを持つ。
直交成分の対象$Q$は剰余上で$\Phi_k(\mathcal C_k)$と直交し、$\Phi_k$の同値性から$Q_k=0$。
したがって$Q=0$であり、像は全体である。
\end{proof}

\begin{proposition}[標準rotation候補の完全忠実性]\label{lp:prop:rotation-candidate}
$\widetilde\rho=L_{\mathcal O}\circ(1)$は$\mathcal K_B^{\rm lin}$を保つ完全忠実なexact dg自己関手であり、
特殊ファイバーでは通常のrotationとなる。
上記の底変更と両立する。
本節だけでは全基底上の本質的全射性と$\widetilde\rho^3\simeq[2]$を主張しない。
\end{proposition}
\begin{proof}
$M\in\mathcal E^\perp$なら
$\operatorname{RHom}(\mathcal O(j),M(1))=0$（$1\le j\le a$）。
$L_{\mathcal O}$は$\mathcal O$からのHomを零にする。
$1\le j<a$について$\operatorname{RHom}(\mathcal O(j),\mathcal O)=0$だから、
このmutationは他の必要な直交条件を壊さない。
twistは線束生成部分圏を保ち、mutationも有限操作なので自己関手を得る。
基底変更は前命題の証明と同じである。
特殊ファイバーでは滑らかなcubicの標準rotationであり自己同値なので、
補題\ref{lp:lem:ff-lift}によって$R$上で完全忠実になる。
\end{proof}
\begin{remark}[自己同値化を残す理由]
補題\ref{lp:lem:ff-lift}により、$\widetilde\rho$に有限対象を保つ片側随伴を構成すれば、
本質的全射性も従う。しかしproper性だけから任意のexact関手が随伴を持つとは限らない。
本稿ではこの随伴の存在を仮定なしに付け加えない。
同様に、$\widetilde\rho^3\to[2]$という自然変換を実際に構成し、特殊ファイバーで同型と示せれば、
対象ごとのconeに導来Nakayamaを適用して全体で同型にできる。
自然変換の持上げそのものを、その特殊ファイバーでの存在で代用してはならない。
\end{remark}

\subsection{中心商の導来自己交叉と特異圏の周期性}
\begin{proposition}[三次周期性の代数的な構成]\label{lp:prop:sing-period}
本節の仮定で、graded singularity categoryには自然な同型
\[
 (3)\simeq[2]
 \qquad\text{on }D_{\rm sg}^b(\operatorname{gr}B)
 \tag{PER}\label{lp:eq:period}
\]
が存在する。ここでは射影残余圏のsmoothnessを使わない。
\end{proposition}
\begin{proof}
$Q=B\otimes_A^{\mathbf L}B$をgraded $B$-双加群とする。
中央正則元の長さ一分解から、$H^0(Q)=B$、
$H^{-1}(Q)=(\Omega)/(\Omega)^2\simeq B(-3)$で、他のcohomologyは零である。
この同定は$B$-双加群として成立する。導来双加群の標準truncationは
\[
 B(-3)[1]\longrightarrow Q\longrightarrow B
 \xrightarrow{\eta}B(-3)[2]
\]
を与える。
任意の$M\in D^b(\operatorname{gr}B)$に$M\otimes_B^{\mathbf L}-$を作用させると
\[
 M(-3)[1]\longrightarrow M\otimes_A^{\mathbf L}B\longrightarrow M
 \xrightarrow{\eta_M}M(-3)[2]
\]
となる。$A$はnoetherianで有限大域次元なので$M$は$A$上perfectであり、
中項は$B$上perfectである。従って$\eta_M$はgraded singularity categoryで同型となる。
$\eta$が双加群射から来ることにより$M$に自然である。
内部shiftを三つ戻すと式\eqref{lp:eq:period}を得る。
\end{proof}
\begin{remark}[二種類のrotationを混同しない]
有界Orlov同値で内部shift $(1)$を移せば、全coherent残余圏には
三乗が$[2]$となる自己同値$\rho_{\rm sg}$を定義できる。
他方、命題\ref{lp:prop:rotation-candidate}は線束生成部分に
$\widetilde\rho=L_{\mathcal O}\circ(1)$を具体的に構成する。
この二つの関手を指定したwindowで同定し、全体での底変更とcoherenceを確認する作業は別である。
この同定は定理\ref{tw:thm:window}で、sectionの切断とevaluation三角から構成する。
本節の周期性だけから線束生成残余部分のrotationとの整合が従うわけではない。
また$Q$を$B\oplus B(-3)[1]$と双加群として分裂させてはならない。
その拡張類$\eta$こそが周期性を与えている。
\end{remark}

% Current local-window proof. Compatible with the general comparison in the following section.
\section{三つのstabilizationによる局所残余圏の表示}
\label{tw:sec:window}

本節は、主定理で用いる局所比較(L)を、代数的入力から導く。
形式的MF同値(CF)、相対対角のperfect性、特殊ファイバーからの生成性の持上げは使わない。
完備化はHom複体の有限性を証明するためだけに用い、
関手自体は代数的なsection--windowと余核のdg比較から作る。
したがって、両立するArtin同値の列を逆極限にして関手を得たとは主張しない。

$(R,\mathfrak m,k)$を有限Krull次元の正則noetherian局所$\mathbb C$-代数、
$k=\mathbb C$とし、$\widehat R$をその完備化とする。
$A=\bigoplus_{j\ge0}A_j$は中央$R$-代数、$A_0=R$、
各$A_j$は有限自由で、$A\otimes_R k=k[x_1,\ldots,x_N]$、$N\ge5$とする。
$\Omega\in Z(A)_3$の剰余$F$は滑らかな射影三次超曲面を定めるものとし、
$B=A/(\Omega)$、$a=N-3$、$i=1-a$と置く。
完備化後も$A_{\widehat R}=\bigoplus_j(A_j\otimes_R\widehat R)$という
\emph{次数別}の代数を使い、内部次数についての無限級数は導入しない。

\subsection{完備Hom複体の縮約と代数的基底への降下}

\begin{lemma}[次数別完備な複体の有限縮約]\label{tw:lem:contraction}
$S$を剰余体$\mathbb C$の完備noetherian局所$\mathbb C$-代数とし、
$C^\bullet$を各項が有限自由$S$-加群である複体とする。
$C^\bullet\otimes_S\mathbb C$の全cohomologyの次元の和が有限なら、
$C^\bullet$は有界有限自由複体と$S$-線形にhomotopy同値である。
複体の全長が有限であることは仮定しない。
\end{lemma}
\begin{proof}
$C_0=C\otimes_S\mathbb C$のcohomologyを$H$とする。
体上の分裂から、次数を保つ$i_0:H\to C_0$、$p_0:C_0\to H$と
次数$-1$の$s_0$を
\begin{equation}\label{tw:eq:retract}
 p_0i_0=1,\qquad d_0s_0+s_0d_0=1-i_0p_0,
 \qquad s_0^2=p_0s_0=s_0i_0=0
\end{equation}
となるように選べる。各次数の基底を選び、これらを$S$へ延長する。
元の微分は$d=d_0+\delta$、$\delta\in\mathfrak m_S\operatorname{Hom}^1(C,C)$と書ける。
各固定次数は有限自由なので、次の逆行列はその次数内で収束する：
\begin{equation}\label{tw:eq:perturbation}
 \begin{aligned}
 J&=(1+s_0\delta)^{-1},& L&=(1+\delta s_0)^{-1},\\
 i'&=Ji_0,&p'&=p_0L,\\
 s'&=s_0L,&d_H&=p_0\delta Ji_0.
 \end{aligned}
\end{equation}
これは標準のhomological perturbation公式である
\cite[Theorem~2.3]{TWCrainic}。
同文献のhomotopyの符号は本稿の$-s_0$に対応する。
$d^2=0$と式\eqref{tw:eq:retract}を用いると
$d_H^2=0$、$di'=i'd_H$、$p'd=d_Hp'$、$p'i'=1$、
$ds'+s'd=1-i'p'$となる。
各積は固定次数での有限行列の積であり、逆行列の級数だけが$\mathfrak m_S$進級数である。
$H$は有界有限次元なので$(H\otimes_{\mathbb C}S,d_H)$が所要の有限複体である。
\end{proof}

\begin{lemma}[有限rank MFのproper Homと底変更]\label{tw:lem:proper-hom}
上の$R,A,\Omega$に対し、有限自由graded MFの任意の二対象$P,Q$の
Hom複体は$R$-perfectである。その通常の複体は有界有限自由複体とhomotopy同値であり、
$R\to k,\widehat R,\operatorname{Frac}(R)$について
\begin{equation}\label{tw:eq:hom-bc}
 \operatorname{Hom}_{\mathrm{MF}}(P,Q)\otimes_R^{\mathbf L}S
 \simeq\operatorname{Hom}_{\mathrm{MF}_S}(P_S,Q_S)
\end{equation}
となる。ここで右辺は有限rank行列の通常の底変更から作るHomである。
\end{lemma}
\begin{proof}
周期変数の二重次数を$(2,-3)$とする。
固定cohomological degreeのHom行列成分では周期変数の冪と内部次数が固定される。
従ってHom複体$C$の各項は有限個の$A_j$の直和、すなわち有限自由$R$-加群であり、
各項の形成と通常の底変更は可換である。
剰余のHomは滑らかな可換cubicのgraded MFのHomなので、
古典Orlov比較により有界有限次元cohomologyを持つ。
補題\ref{tw:lem:contraction}を$C\otimes_R\widehat R$へ適用すると、
完備化後の複体は有界有限自由複体へ縮約する。

$R$はnoetherianなので$H^q(C)$は各$q$で有限生成である。
平坦な完備化はcohomologyと可換であり、忠実平坦性から
$H^q(C)=0$は$q$の有限区間の外で成り立つ。
正則$R$の大域次元は有限なので$C$は$R$-perfectである。
さらに有界有限自由複体$D$から$C$への準同型を選ぶと、そのconeは
各項が射影である完全複体になる。
$e=\operatorname{gldim}R$とすれば、完全複体のcycleに対して
$\operatorname{Ext}^1_R(Z^q,M)=\operatorname{Ext}^{e+1}_R(Z^{q+e},M)=0$
というdimension shiftingが適用できる。よって全cycleは射影で、このconeは可縮である。
従って$C$自体が$D$とhomotopy同値、特にK-flatである。
項ごとの底変更の同定が式\eqref{tw:eq:hom-bc}を与える。
有限coneと直和因子にも同じ結論が延長される。
\end{proof}

\subsection{stabilizationとsection--windowの実際のenhancement}

\begin{lemma}[正則局所基底上のdg stabilization]\label{tw:lem:stabilization}
$A,B$は左右noetherian、$A$の大域次元は有限、$\Omega$は正則である。
余核による比較は、内部shiftと中心元に由来する$(3)\simeq[2]$を保つ
\begin{equation}\label{tw:eq:stabilization}
 \mathrm{MF}^{\mathbb Z}(A,\Omega)
 \simeq D^b_{\mathrm{sg,dg}}(\operatorname{gr}B)^{\mathrm{idem}}
\end{equation}
というdg Morita同値を与える。
本節の有限windowに現れる対象では、$k$および$\widehat R$への底変更と整合する。
\end{lemma}
\begin{proof}
noetherian性、相対Koszul分解、$\operatorname{gldim}A\le N+\dim R$、
$\Omega$の正則性、$B_j$の自由性と有限左右injective dimensionは、
補題\ref{lp:lem:noetherian}、\ref{lp:lem:resolution}と
命題\ref{lp:prop:koszul}、\ref{lp:prop:gorenstein}による。
これらの次数別の証明は基底の完備性や有限型性を要求せず、$\widehat R$でも通る。

Chenの余核定理\cite[Theorem~7.2]{AMChen}は三角圏の定理である。
そのままdg同値を引用するのではなく、以下のenhancementを用いる。
有限自由graded MFの成分別split exact圏を$\mathcal F$、
有限生成graded Gorenstein-projective $B$-加群のexact圏を$\mathcal G$とする。
$A$の大域次元が有限なので、$\mathcal G$の全対象は有限$A$-射影次元を持つ。
ChenのCorollary~2.11により、非零$M\in\mathcal G$の$A$-射影次元は一である
（零対象も含めれば高々一）。従って有限生成元を斉次に選ぶと
\begin{equation}\label{tw:eq:graded-presentation}
 0\longrightarrow P_0\xrightarrow{d_0}P_1\longrightarrow M\longrightarrow0
\end{equation}
という有限graded projective $A$-加群による表示を得る。
実際、有限graded freeな$P_1$からの全射の核は$A$上射影であり、
斉次写像の持上げの次数零成分を取ることでgradedにも射影である。
各有限graded projective $A$-加群は、$A_0=R$が局所であることから
有限個の内部shiftを持つ自由加群である。
$\Omega M=0$より、次数零の乗法写像$\Omega_{P_1}:P_1\to P_1(3)$は
$d_0(3)$を経由し、一意な$d_1:P_1\to P_0(3)$を定める。
従って
\begin{equation}\label{tw:eq:graded-factorization}
 d_0(3)d_1=\Omega_{P_1},\qquad d_1d_0=\Omega_{P_0}.
\end{equation}
第二式は第一式と$d_0(3)$の単射性から従う。
これが本稿の内部shift規約における有限自由graded MFで、その余核は$M$である。
この構成はChenのTheorem~7.2とRemark~7.3(2)の本質的全射性を斉次に実現する。
余核はexact関手$\mathcal F\to\mathcal G$であり、projective-injectiveを保つ。
射とhomotopyの持上げも次数零成分を取れるので、同定理の完全忠実性もgraded版で成立する。
右加群については$A^{\mathrm{op}}$へ同じ結果を適用する。
従って、その誘導するstable categoryの関手は同値である。

exact圏の有界導来dg圏をprojective-injectiveの有界複体で割ることで、
実際のdg関手のzigzagを得る。
左側とMFの比較は有限totalizationである。
成分別splitな完全複体に横方向のcontraction $s$を取ると、
縦MF微分$\partial$に対する補正$s(1+\partial s)^{-1}$は横方向の有限性により有限和である。
中心曲率$w$の残差は$sws=ws^2=0$なので、totalizationは可縮である。
右側と特異圏の比較は$\mathcal G$の包含から誘導される。
完全射影分解がstable Homを計算し、有限injective dimensionによる高いsyzygyが
本質的全射性を与える。各zigzagの$H^0$は上記の三角同値である。
全てのshiftにも同じ議論を適用することでHomの全cohomologyの同型となり、
式\eqref{tw:eq:stabilization}を得る。

指定周期の比較は、余核の二拡張で確認できる。
式\eqref{tw:eq:graded-factorization}のMFの余核を$M$とし、
$\overline P_j=P_j\otimes_A B$と書く。その周期分解の連続する二段は
\begin{equation}\label{tw:eq:period-extension}
 0\longrightarrow M(-3)\xrightarrow{\overline d_1(-3)}
 \overline P_0\xrightarrow{\overline d_0}\overline P_1
 \longrightarrow M\longrightarrow0
\end{equation}
である。左端の写像は$d_1(-3)$が余核を経由して誘導する写像を表す。
二つの微分を$A$へ持ち上げた合成は、式\eqref{tw:eq:graded-factorization}により
指定した$\Omega$の乗法そのものである。
従って、$\Omega\bmod(\Omega)^2$を用いて
$(\Omega)/(\Omega)^2=B(-3)$と同定した導来自己交叉の連結射
$\eta_M:M\to M(-3)[2]$は、この二拡張のYoneda類である。
ここでは命題\ref{lp:prop:sing-period}と同じconeの符号規約を用いる。
したがって余核は、MFの二段周期を
\begin{equation}\label{tw:eq:specified-period}
 \tau_M^{\rm sg}:=\eta_M(3):M(3)\xrightarrow{\sim}M[2]
\end{equation}
へ移す。有限totalizationは内部shiftと複体の二重shiftに可換であり、
同じ二拡張をtotalizeするので、この比較を保つ。
$\mathcal G$から特異圏への包含も二拡張をそのまま送る。
これで全てのzigzagに沿う周期の規格化が決まり、任意のスカラー倍を挿入しない。

底変更には余核構成を使う。有限MFの二つの行列の積は$\Omega$であり、
剰余では非零多項式$F$の乗法になるため、各行列は剰余上単射である。
固定内部次数で単元小行列式を選ぶと、余核は有限自由$R$-加群となる。
従って余核の周期分解、有限totalization、windowで使う次数別自由な有限長加群は
通常の剰余底変更と整合する。完備化はさらに平坦である。
Morita逆であるstabilizationの比較は、この余核の自然な比較の逆として選ぶ。
任意のcoherent対象の任意の非正則底変更まで主張してはいない。
\end{proof}

\subsection{比較図式と関手の定義域}\label{tw:subsec:comparison}
上のexact圏$\mathcal F,\mathcal G$のprojective-injective部分圏を
$\mathcal P_{\mathcal F},\mathcal P_{\mathcal G}$とし、
\[
 \mathfrak F=
 \bigl(D^b_{\rm dg}(\mathcal F)/K^b_{\rm dg}(\mathcal P_{\mathcal F})\bigr)^{\rm idem},
 \qquad
 \mathfrak G=
 \bigl(D^b_{\rm dg}(\mathcal G)/K^b_{\rm dg}(\mathcal P_{\mathcal G})\bigr)^{\rm idem}
\]
と置く。ここで$D^b_{\rm dg}$はexact圏の有界導来dg圏、斜線はdg商を表す。
さらに$\mathcal D_B=D^b_{\rm dg}(\operatorname{gr}B)$、
$\mathcal S_B=(\mathcal D_B/\operatorname{Perf}_{\rm dg}(B))^{\rm idem}$、
$\mathcal M_R=\mathrm{MF}^{\mathbb Z}(A,\Omega)$と略記する。
$\mathcal M_R$にも有限coneと冪等完備化を施す。

derived sectionの切断を$\sigma_i=(R\omega(-))_{\ge i}$、
例外成分$\mathcal E_B$を除くmutationを$p$とすると、
stabilizationと次のwindow定理の関手は一つの図式に収まる：
\begin{equation}\label{tw:eq:comparison-diagram}
 \begin{array}{c@{\ }c@{\ }c@{\ }c@{\ }c@{\ }c@{\ }c}
 \mathcal C_B&\xrightarrow{\sigma_i}&\mathcal D_B&\xrightarrow{q}&
 \mathcal S_B&\xleftarrow[\sim]{\overline\iota}&\mathfrak G\\[2pt]
 {\scriptstyle p}\!\downarrow&&&&&&\Vert\\[2pt]
 \mathcal K_B^{\rm lin}&\xrightarrow{\Psi_R}&\mathcal M_R&
 \xleftarrow[\sim]{\operatorname{Tot}}&\mathfrak F&
 \xrightarrow[\sim]{\overline c}&\mathfrak G .
 \end{array}
\end{equation}
ここで$q$はperfect複体によるdg商、$\overline c$は余核exact関手から、
$\overline\iota$は$\mathcal G\hookrightarrow\operatorname{gr}B$から誘導する。
$\operatorname{Tot}$は有界複体の有限totalizationである。
上段の$\sigma_i$の定義域は線束生成部分$\mathcal C_B$であり、
任意のcoherent対象に対する底変更をこの図式で主張してはいない。
記号$\sim$を付けた矢印は補題\ref{tw:lem:stabilization}のdg Morita同値である。
図式の可換性はMorita局所化における自然な比較
\begin{equation}\label{tw:eq:comparison-composite}
 q\sigma_i\simeq
 \overline\iota\,\overline c\,\operatorname{Tot}^{-1}\Psi_Rp
\end{equation}
を意味し、$\Psi_R$の像と完全忠実性は以下で証明する。
式\eqref{tw:eq:period-extension}は
$\operatorname{Tot},\overline c,\overline\iota$の各矢印に沿う指定周期を与える。
$\mathcal C_B$と$\mathcal D_B$自体に三次周期性があるとは述べない。
$q\sigma_i$を残余圏へ制限した後のrotationとの比較は
式\eqref{tw:eq:rotation-period-square}で明示する。

\begin{theorem}[三対象window定理：局所比較(L)の確立]\label{tw:thm:window}
本節の代数的入力だけを仮定する。$K_R=\operatorname{stab}(R)$、
$T_R=K_R\oplus K_R(1)\oplus K_R(2)$と置くと、
\begin{equation}\label{tw:eq:main}
 \mathcal K_{B_R}^{\mathrm{lin}}
 \simeq\operatorname{thick}\{K_R,K_R(1),K_R(2)\}
 \simeq\operatorname{Perf}(\Gamma_R),\qquad
 \Gamma_R=R\operatorname{End}(T_R)\in\operatorname{Perf}(R)
\end{equation}
である。最初の同値は
\begin{equation}\label{tw:eq:window}
 \Psi_R=\operatorname{stab}\circ\bigl((R\omega -)_{\ge i}\bigr),
 \qquad i=1-a
\end{equation}
の制限で与えられる。
これは$k,\widehat R,\operatorname{Frac}(R)$への底変更と整合し、
$L_{\mathcal O}\circ(1)$をMFの内部shiftへ移す。
特にこのrotationは自己同値であり、中心元に由来する三次周期同型を持つ。
\end{theorem}
\begin{proof}
$E_R=\bigoplus_{j=0}^{N-1}\mathcal O(j)$は線束生成部分を生成し、
その全対象対のHomは$R$-perfectである。
例外列$\mathcal E=\langle\mathcal O,\ldots,\mathcal O(a-1)\rangle$を除くmutationを$p$とすると、
$G_R=p(E_R)$が$\mathcal K^{\mathrm{lin}}$を生成する
（命題\ref{lp:prop:residual}）。
必要なsection、有限cohomological dimension、$\chi$の入力は
Burke--Stevenson\cite[Proposition~4.18, Lemmas~6.1--6.2]{LP-BS}による。

\emph{関手の定義。}
線束Homの公式から、$0\le j\le N-1$では
$(R\omega\mathcal O(j))_{\ge i}=B(j)_{\ge i}$である。
$T_j=B(j)_{<i}$とすると、$T_j$の各次数は有限自由で、
$R$のshiftによる有限filtrationを持つ。
$R$の相対Koszul分解と$B$の長さ一の$A$-分解により、
$B(j)_{\ge i}$は$A$-perfectである。
従ってderived section、内部切断、dg商、補題\ref{tw:lem:stabilization}のMorita逆を合成して
式\eqref{tw:eq:window}が実際のdg quasi-functorとして定義される。
$0\le j<a$では$T_j=0$なので$\Psi_R(\mathcal O(j))=0$であり、
$\Psi_Rp\simeq\Psi_R$である。

\emph{完全忠実性。}
この関手の剰余は、古典的なsection--windowを特異圏へ送る関手である。
Burke--Stevenson\cite[Theorem~6.4とその証明]{LP-BS}において$i=1-a$を取ると、
そのkernelは$\mathcal E_k$であり、Verdier商を特異圏と同定する。
本稿の右直交と同文献の半直交分解の向きは異なり得るが、
双方が同じVerdier商を表示するため、$\mathcal E_k^\perp$への制限は同値である。

sourceの線束Homの底変更は命題\ref{lp:prop:basechange}による。
sectionの内部切断は生成線束上で上の加群の表示になり、
stabilizationの底変更は補題\ref{tw:lem:stabilization}による。
有限cone、直和因子、evaluationを通じ、$\Psi_R$のHom上の\emph{自然な写像}は
剰余で古典windowのHom写像へ移る。
targetのHomは補題\ref{tw:lem:proper-hom}により$R$-perfectかつ剰余底変更と可換である。
従って任意の二対象についてHom比較のconeはperfectで剰余が零となり、
局所環上の導来Nakayamaにより零である。これがenhancedな完全忠実性を与える。
単にHomの階数だけを比較しているのではない。

\emph{像の同定。}
残る$j=a,a+1,a+2$について、$B(j)$は特異圏で零なので
\begin{equation}\label{tw:eq:triangular}
 \Psi_R(\mathcal O(j))\simeq\operatorname{stab}(T_j)[-1],\qquad
 \operatorname{gr}(T_j)=\bigoplus_{r=0}^{j-a}B_r\otimes_RR(j-r).
\end{equation}
$B_0=R$である。まず$T_a=R(a)$、次に$T_{a+1}$の高次数部分は
$R(a)$の有限直和、最後に$T_{a+2}$の高次数部分は
$R(a),R(a+1)$の有限extensionである。
その商はそれぞれ$R(a+1),R(a+2)$なので、三角的に
\begin{equation}\label{tw:eq:three-thick}
 \operatorname{thick}(\Psi_R(G_R))
 =\operatorname{thick}\{K_R(a),K_R(a+1),K_R(a+2)\}
 =\operatorname{thick}(T_R).
\end{equation}
最後はMFの$(3)=[2]$による。
完全忠実なexact関手の冪等完備sourceの像はthickなので、これが式\eqref{tw:eq:main}の同値となる。
$\Gamma_R$のproper性は補題\ref{tw:lem:proper-hom}からも従う。
生成性の議論は有限長加群の実際のfiltrationによるものであり、
特殊ファイバーの生成性だけを持ち上げたものではない。

\emph{rotationと三次周期性。}
線束Homの消滅とevaluation三角により、$L_{\mathcal O}(M(1))$は再び残余圏に属する
（命題\ref{lp:prop:rotation-candidate}）。
$M\in\mathcal K^{\mathrm{lin}}$なら
$(R\omega M)_i=R\operatorname{Hom}(\mathcal O(a-1),M)=0$である。
切断の自然な三角により
\begin{equation}\label{tw:eq:rotation-cut}
 ((R\omega M)(1))_{\ge i}
 =(R\omega M)_{\ge i+1}(1)
 \simeq(R\omega M)_{\ge i}(1).
\end{equation}
$\Psi_R(\mathcal O)=0$とevaluationの三角を使うと
$\Psi_R(L_{\mathcal O}M(1))\simeq\Psi_R(M)(1)$を得る。
式\eqref{tw:eq:three-thick}の像は$(1)$と$(-1)$の両方で安定なので、
rotationは自己同値となる。
周期同型はMFで中心元$\Omega$から定まる$(3)=[2]$をこの同値で移したものとする。
より正確には、上で得た自然比較を
$\beta_M:\Psi_R(\rho_BM)\xrightarrow{\sim}\Psi_R(M)(1)$とし、
$\beta_M^{(3)}=\beta_M(2)\circ\beta_{\rho_BM}(1)\circ\beta_{\rho_B^2M}$と置く。
MFの指定周期を$\tau^{\rm MF}$と書くと、$\Psi_R$の完全忠実性により
次の正方形を可換にする$\tau^{\rm rot}$が一意に定まる：
\begin{equation}\label{tw:eq:rotation-period-square}
 \begin{array}{ccc}
 \Psi_R(\rho_B^3M)&\xrightarrow{\Psi_R(\tau_M^{\rm rot})}&\Psi_R(M[2])\\[2pt]
 {\scriptstyle\beta_M^{(3)}}\!\downarrow&&\downarrow\!{\scriptstyle\simeq}\\[2pt]
 \Psi_R(M)(3)&\xrightarrow{\tau_{\Psi_R(M)}^{\rm MF}}&\Psi_R(M)[2].
 \end{array}
\end{equation}
図式\eqref{tw:eq:comparison-diagram}の右側では、下辺の周期は
$\tau^{\rm sg}=\eta(3)$へ移る。
従ってこの正方形は、切断とevaluationが与えるrotation比較と、
中心元の二拡張が与える周期比較とを同時に固定している。
各矢印は余核・切断・evaluationの自然な写像から作ったので、
以下の底変更とも整合する。別に選んだ周期のスカラー倍を保存するとは主張しない。

完備化・分数体への底変更は平坦であり、有限MF行列のHomと積に項ごとに可換する。
sourceも有限線束生成とmutationから作った模型の底変更に一致する。
剰余底変更は上で確認した。これで必要な三対象とrotationの比較をすべて得る。
\end{proof}

% Complete formal comparison and its three-object restriction.
\section{次数別完備な曲率付き移送と形式比較の構成}
\label{cft:sec:transport}

本節では、完備係数上のMF比較を実際の関手と逆関手によって構成し、
その固定三対象への制限を明らかにする。
既知のKontsevich formalityを\emph{曲率を入れる前}に使い、
その後のゲージ経路は有限Artin次数ごとに直接積分する。
曲率の大きさ、全Hochschild cochain、単位、逆写像、有限rank、
二つの次数、指定された周期比較をそれぞれ確認する。
完備化と一般点への局所化の規約は第\ref{gdr:subsec:scope}節に従う。

\subsection{主張と既知結果の使用範囲}
$(R,\mathfrak m)$を剰余体$\mathbb C$の完備noetherian局所$\mathbb C$-代数とする。
ここでは$R$の正則性は不要である。
$S=\mathbb C[x_1,\ldots,x_N]$、$N\ge3$、$\deg x_i=1$とし、
$A^{\rm gr}=\bigoplus_{j\ge0}A_j$をunitalな中央$R$-代数とする。
各$A_j$は有限自由、$A_0=R$、$A^{\rm gr}/\mathfrak m A^{\rm gr}=S$と仮定する。
$\Omega\in Z(A^{\rm gr})_3$の剰余$F$は滑らかな射影三次超曲面を定めるものとする。
次数別に剰余の単項式基底を持ち上げ、単位を保つ加群同定
$A_j\simeq S_j\otimes_{\mathbb C}R$を選ぶ。

\begin{theorem}[完備係数上の実際のgraded MF比較]
\label{cft:thm:complete}
上の入力に対し、$W\in R[x_1,\ldots,x_N]_3$、$W\bmod\mathfrak m=F$と、
有限rank graded MFの実際の$R$-線形unital $A_\infty$同値
\begin{equation}\label{cft:eq:equivalence}
 \operatorname{MF}^{\mathbb Z}(A^{\rm gr},\Omega)
 \simeq\operatorname{MF}^{\mathbb Z}(R[x_1,\ldots,x_N],W)
\end{equation}
が存在する。これは内部shiftと、中心元による$(3)\simeq[2]$を保ち、
剰余では恒等であり、すべての$R/\mathfrak m^{q+1}$への切断と両立する。
有限自由模型上で互いに逆な関手を構成し、その二つのrankと内部shiftを保つ。
構成は$\deg t=(2,-3)$の周期変数についてrelativeであり、$t$を微分しない。
同値は有限coneと冪等完備化へ延長する。

曲率付きbar写像のTaylor成分$F_r$と逆写像の成分$H_r$は
\begin{equation}\label{cft:eq:bound}
 F_0=H_0=0,\quad F_1,H_1=1+O(\mathfrak m),\quad
 F_{2r}=H_{2r}=0,\quad F_{2r+1},H_{2r+1}\in\mathfrak m^r\quad(r\ge1)
\end{equation}
を満たすように取れる。次数付き代数の厳密同型は、この写像の前後に合成する。
この評価はMF行列を挿入する前のbar写像について述べており、
挿入後のMF関手の偶数Taylor成分の消滅を意味しない。
\end{theorem}

Kontsevichの既知結果は、affine formality、HKR、Maurer--Cartan変形問題の比較
\cite[Theorem~4.6, Section~4.6.2, Theorem~6.3]{PKontsevich}である。
曲率付きtwistによるgraded MFの記述には
\cite[Section~2, 特にRemark~2.8とTheorem~2.12]{CT}を使う。
後者を形式的な無限arity関手へ無条件に適用せず、以下で評価式の有限性と収束を証明する。
これらの一般理論の新規性は主張せず、本定理の組合せの初出は未確認である。

\begin{remark}[既知の曲率付きformalityとの比較]\label{cft:rem:prior-formality}
Carqueville--Kay\cite[Proposition~2.3, Remark~2.4, Proposition~3.4]{CK}の
Proposition~3.4は、$S=\mathbb C[x_1,\ldots,x_N]$上の曲率付き代数と
$[-F,-]$を微分とするpolyvector場の変形複体を比較する。
用いられる第二種cochainはarityについての直和であり、各arityでは
全多重線形写像を許すので、変形積がbidifferentialであることを入力条件とはしない。
同論文のTheorem~3.1はさらに固定したMFのendomorphism代数を扱う。
本稿の全graded MF圏の同値や指定周期比較は、同定理の主張そのものではない。

本稿の状況には、$\deg t=(2,-3)$の周期変数を係数として扱い、$t$を微分しない
relativeな比較、二重次数とnormalizedな単位を保つ構成、
$\mathfrak m$進完備化を加える必要がある。
第二種cochainをこの位相で完備化すると、各$R/\mathfrak m^{q+1}$では有限arityとなる。
これは初めからarityの任意の無限直積を許すこととは異なる。
$t$-relativeなpolyvector複体では、勾配Koszul複体により
\begin{equation}\label{cft:eq:prior-jacobian}
 H\bigl(T_{\rm poly}\otimes_{\mathbb C}\mathbb C[t,t^{-1}],[-tF,-]\bigr)
 \simeq\operatorname{Jac}(F)\otimes_{\mathbb C}\mathbb C[t,t^{-1}].
\end{equation}
右辺の関数のLie次数は$-1$であるため、Lie次数一・内部weight零の成分は
$t\operatorname{Jac}(F)_3$である。関数の斉次代表同士のSchouten括弧は零であり、
基点$b_F=\mu_0-tF$の曲率を微分に含めて
$\delta=(\mu_A-\mu_0)-t(\Omega-F)\in\mathfrak m$を扱うことが、
potentialによる形式的存在原理に対応する。

一般の孤立potentialでは、形式係数の極限は多項式次数について有界とは限らず、
例えば$\mathbb C[x][[h]]$と$\mathbb C[[h]][x]$は異なる。
従って、形式的なMC同値だけで本稿の次数別直和上の有限rank同値を結論しない。
以下では補題\ref{cft:lem:model}で元の全cochainから始め、
補題\ref{cft:lem:arity}、\ref{cft:lem:flow}で単位・二重次数と有限評価を確認し、
式\eqref{cft:eq:two-block-matrix}で固定rank・内部shiftのMFに留まることを証明する。
三次特有の定数trivector、一段のゲージ経路、式\eqref{cft:eq:bound}の一様評価は、
この一般原理の本稿での明示的な実現である。
通常のMF圏のHochschild複体と第二種複体を無条件に同一視する必要はない。
\end{remark}

\subsection{通常の結合変形と中心元の有限次元逆像}
\begin{lemma}[次数を保つunitalな星積模型]\label{cft:lem:model}
formal Poisson tensor
$\pi\in\mathfrak m(S_2\otimes\bigwedge^2\langle\partial_i\rangle)\widehat\otimes R$
と、単位・内部次数を保つ厳密な代数同型
\begin{equation}\label{cft:eq:phi}
 \phi:A^{\rm gr}\xrightarrow{\sim}(S\otimes R,*_\pi),\qquad
 \phi\equiv1\pmod{\mathfrak m}
\end{equation}
が存在する。右辺も内部次数ごとの直和とする。
\end{lemma}
\begin{proof}
単位を入力すると零になるnormalized Hochschild cochainを使う。
多項式環の対角のgraded Koszul分解とnormalized bar分解の比較は
内部次数を保つように選べる。従ってpolydifferential cochainから
\emph{全ての}内部次数指定Hochschild cochainへの包含はHKR同型を誘導する。
既存の積がbidifferentialであるという仮定は不要である。
Kontsevichのグラフ写像はaffine線形変換に同変なので、内部weightを保つ。
weight零に制限したformalityと上の包含の合成は準同型であり、
次数別flatな結合積の変形をPoisson変形に対応させる。

この対応を各$R/\mathfrak m^{q+1}$で行う。
一段の核$\mathfrak m^q/\mathfrak m^{q+1}$では、持上げ障害の比較はHKRの
次数二での単射性、積の差の補正は次数一での全射性による。
一入力gaugeを同時に持ち上げるので、すでに選んだ低次jetを保てる。
これは\cite[Theorem~4.6]{PKontsevich}の冪零変形問題の比較の逐次形である。
全Hochschild cochain側のdegree零gaugeは一入力の写像であり、
厳密な代数同型を与える。normalized性から単位を保つ。
極限は各有限自由$A_j$内で存在し、逆写像も同様に存在する。
weight零のbivectorの係数次数は二である。
\end{proof}

\begin{lemma}[中心逆像を有限行列で取り出す]\label{cft:lem:center}
formalityを固定すると、$W\in S_3\otimes R$が一意に存在して
\begin{equation}\label{cft:eq:center}
 C_\pi(W)=\phi(\Omega),\qquad [\pi,W]=0,\qquad W\bmod\mathfrak m=F,
 \quad
 C_\pi(f):=\sum_{r\ge0}\frac1{r!}\mathcal U_{r+1}(f,\pi,\ldots,\pi)
\end{equation}
となる。
\end{lemma}
\begin{proof}
関数一個とbivector $r$個のグラフの出力arityは零であり、内部weightを保つ。
よって$C_\pi$は有限自由加群$S_3\otimes R$の$R$-線形自己写像で、
$\mathcal U_1(f)=f$により$C_\pi=1+O(\mathfrak m)$である。
その逆は各固定次数内で収束する幾何級数であり、$W$を一意に定める。

接写像のchain恒等式は
\begin{equation}\label{cft:eq:chain-center}
 [*_\pi,C_\pi(W)]=\mathcal U^\pi([\pi,W])
\end{equation}
である。左辺は中心性から零。
右辺の入力は係数次数四のvector場であり、接写像の最低次は
vector場を導分として見るHKRの単射である。
$\mathfrak m^q/\mathfrak m^{q+1}$で順にこの単射性を使うと
$[\pi,W]\in\bigcap_q\mathfrak m^q=0$となる。
曲率付きの通常のHochschild cohomologyをJacobian環と同一視する操作は使わない。
\end{proof}

\subsection{三次式の次数制約とゲージ経路}
\begin{lemma}[唯一の定数三重vectorと二回目の作用]\label{cft:lem:primitive}
一意な$\Xi\in\mathfrak m\bigwedge^3\langle\partial_i\rangle\otimes R$が存在し、
\begin{equation}\label{cft:eq:primitive}
 \pi=[\Xi,W],\qquad [\Xi,\pi]=0
\end{equation}
となる。
\end{lemma}
\begin{proof}
滑らかな$F$の全偏微分は多項式環の正則列である。
勾配Koszul contractionは、係数次数とrankの和「係数次数$+2\,\mathrm{rank}$」を保つ。
各斉次部分は有限の有限自由$R$-複体である。
剰余の完全性を使い、$\mathfrak m^q/\mathfrak m^{q+1}$で原像を逐次持ち上げると、
完備化と分離性から正のrankで完全となる。

$[\pi,W]=0$なので、係数次数二・rank二のcycle $\pi$の原像は
係数次数零・rank三に取れる。二つの原像の差の原像は係数次数$-2$を要求するため零。
これで$\Xi$の存在・一意性と$\Xi\bmod\mathfrak m=0$を得る。
$D=[\Xi,\pi]$は係数次数一・rank四であり、graded Jacobiと$[\pi,\pi]=0$から
$[W,D]=-[\pi,\pi]=0$となる。
そのKoszul原像は係数次数$-1$を要求するので$D=0$である。
ここでJacobi条件は必須であり、Casimir条件だけから第二の等式は従わない。
\end{proof}

周期変数$t$の二重次数を$(2,-3)$とし、$t$を微分せず、係数として多重線形に扱う。
$\eta=\Xi/t$はcohomological degree零・内部weight零である。
\begin{equation}\label{cft:eq:path}
 \alpha_s=-tW+(1-s)\pi,\qquad
 \partial_s\alpha_s=[\eta,\alpha_s]=-\pi
\end{equation}
はMaurer--Cartan経路で、始点は非可換模型、終点は可換模型に対応する。
$W$は$\mathfrak m$の倍数でないので、通常の冪零MC定理をこの経路へ直接適用してはいけない。

\begin{lemma}[大きい曲率でも有限なformality評価]\label{cft:lem:arity}
$\alpha_s$の像はcurved associative構造
$b_s=(*_{(1-s)\pi},-tC_{(1-s)\pi}(W))$である。
$\gamma_s=\mathcal U^{\alpha_s}(\eta)$はarity一と三だけを持ち、
$\gamma_s\in\mathfrak m$で、
\begin{equation}\label{cft:eq:lax}
 \partial_sb_s=[\gamma_s,b_s]
\end{equation}
となる。$\gamma_s$の正のarity成分はnormalizedである。
\end{lemma}
\begin{proof}
グラフの辺数と配置空間の次元$2n+m-2$を等しくすると、
関数$k$個とbivector任意個の出力arityは$2-2k$、
さらにtrivector一個を入れた場合は$3-2k$となる。
従って関数入力は高々一個で、残る入力は$\mathfrak m$の倍数である。
mod $\mathfrak m^{q+1}$では必要な全Taylor和が有限となり、
$L_\infty$恒等式も有限和の恒等式として適用できる。
chain恒等式に$\partial_s\alpha_s=-[\alpha_s,\eta]$を代入すると
$\partial_sb_s=-[b_s,\mathcal U^{\alpha_s}(\eta)]$、すなわち式\eqref{cft:eq:lax}である。

正のarityで単位をboundary入力へ置く。
そこへ辺が入れば定数の微分が零であり、入らなければ積分形式はその頂点を忘れた
一つ低次元の配置空間からの引戻しなので最高次積分が零となる。
internal頂点が一個の場合はHKR写像であり、同じ消滅は直接分かる。
この議論は\cite[Sections~6.2--6.4]{PKontsevich}のグラフ表示に適用する。
$t$への多重線形性と内部weightは全ての項で保たれる。
\end{proof}

\subsection{bar写像と逆写像の一様なarity評価}
\begin{lemma}[順序付きflowと逆写像]\label{cft:lem:flow}
$\partial_s\mathcal F_s=\gamma_s\mathcal F_s$、$\mathcal F_0=1$を解くと、
各$R/\mathfrak m^{q+1}$上で実際のbar coalgebra同型を得る。
$\mathcal F_1b_0=b_1\mathcal F_1$であり、写像・逆写像は式\eqref{cft:eq:bound}を満たす。
\end{lemma}
\begin{proof}
mod $\mathfrak m^{q+1}$で
\begin{equation}\label{cft:eq:ordered}
 \mathcal F_s=1+\sum_{j=1}^{q}
 \int_{0\le u_j\le\cdots\le u_1\le s}
 \gamma_{u_1}\cdots\gamma_{u_j}\,du_j\cdots du_1
\end{equation}
とする。各係数は$s$の多項式で、積分は標数零での代数的積分である。
$\gamma$がcoderivationなので、$\Delta\mathcal F_s$と
$(\mathcal F_s\otimes\mathcal F_s)\Delta$は同じ微分方程式・初期値を持つ。
よってcoalgebra写像である。
$\mathcal F_s-1\in\mathfrak m$から逆は
$\sum_{j=0}^{q}(1-\mathcal F_s)^j$で与えられる。
式\eqref{cft:eq:lax}により$\mathcal F_s^{-1}b_s\mathcal F_s$の微分は零である。

零項成分はなく、三入力coderivationを一回使うごとに出力長は二つ減る。
一入力成分は長さを変えない。従ってarity $2r+1$には少なくとも$r$回の
coderivationが必要で、係数は$\mathfrak m^r$に入る。
偶数arityは現れず、逆flowにも同じ評価が成立する。
さらに$t$を含まない入力に制限すれば
\begin{equation}\label{cft:eq:t-weight}
 F_{2r+1}=t^{-r}f_{2r+1},\qquad
 H_{2r+1}=t^{-r}g_{2r+1},
 \qquad \operatorname{wt}(f_{2r+1})=
 \operatorname{wt}(g_{2r+1})=-3r
\end{equation}
である。ここで$\operatorname{wt}$は、出力の多項式次数から入力の次数の和を引いた内部weightである。
実際、$\gamma_s$の三入力成分は関数入力を持たず、$\eta=\Xi/t$から$t^{-1}$を一つ含む。
一入力成分は関数入力$-tW$を一つ持つので$t$の冪は零である。
arityを$2r+1$から一へ減らすflowの項には三入力成分が正確に$r$回現れるため、
$t$の冪は$-r$である。内部weight零から$f_{2r+1},g_{2r+1}$のweightも従う。
出力tensor長$\ell$に寄与する成分の積の位数を足すと、mod $\mathfrak m^{q+1}$では
入力長は高々$\ell+2q$である。従ってtensor長完備barへの延長も各剰余で連続である。
単位はnormalized性からflow中固定され、高arity成分の単位入力は零となる。
同じ形式的係数を切断しているので、商$R/\mathfrak m^{q+1}\to R/\mathfrak m^q$と両立する。
\end{proof}

\subsection{冪零でないMF行列への評価と周期比較}
\begin{proof}[定理\ref{cft:thm:complete}の証明]
上のbar写像を有限個のshiftを持つ自由加群の行列成分に拡張する。
合成可能な添字列に沿う和とsuspensionのKoszul符号を用いるので、
元のbar恒等式は行列上でも成立する。
MF行列$Q$は$\mathfrak m$の倍数でも上三角でもなくてよい。
mod $\mathfrak m^{q+1}$では、像は有限式
\begin{equation}\label{cft:eq:object}
 Q'=\sum_{r=0}^{q}F_{2r+1}(Q,\ldots,Q)
\end{equation}
で与えられる。$k$個の射に対するTaylor成分は、suspended表示で
\begin{equation}\label{cft:eq:insertions}
 \sum_{r_0,\ldots,r_k\ge0}
 F_{k+r_0+\cdots+r_k}
 (Q_k^{\otimes r_k},a_k,Q_{k-1}^{\otimes r_{k-1}},\ldots,a_1,Q_0^{\otimes r_0})
\end{equation}
である。$k+\sum r_i\le2q+1$の項だけが寄与するため有限である。

suspensionを$\varsigma$と書くと、tensor長完備bar上で$\mathcal F(e^{\varsigma Q})=e^{\varsigma Q'}$である。
$\mathcal F b_0=b_1\mathcal F$から対象のMF方程式と全functor恒等式が従う。
同じ構成を逆bar写像に施すと、挿入公式はbar写像の合成と可換するので、
二つのtwisted functorの合成は恒等となる。
したがって各Artin切断で実際の互いに逆な関手がある。

有限rankの通常のgraded MFの形が保存されることも確認する。
基礎自由加群の偶・奇基底のcohomological次数をそれぞれ零・一に取ると、
その行列は、内部shiftを基底に付したまま
\begin{equation}\label{cft:eq:two-block-matrix}
 Q=\begin{pmatrix}0&tq_1\\q_0&0\end{pmatrix},
 \qquad Q^2=tC_\pi(W)\,\mathrm{id}
\end{equation}
と書ける。ここでは$\phi$で星積模型へ移した行列を用い、平方は始点の星積で計算する。
$2r+1$個の$Q$を挿入した合成可能な添字列では、
左下成分に$t$が$r$個、右上成分に$r+1$個現れる。
式\eqref{cft:eq:t-weight}の$t^{-r}$を掛けると、像の$t$の冪はそれぞれ零・一となる。
従って$Q'$は同じ有限個の内部shiftと同じ二つの階数を持つ行列であり、
bar恒等式により$(Q')^2=tW\,\mathrm{id}$を満たす。
逆写像についても同じ計算が成り立つ。

固定した有限自由加群、固定したcohomological次数と内部次数では、
各行列成分の$t$の冪と$x$の次数が固定される。
従って出力は有限個の$S_j\otimes R$の直和に属する。
式\eqref{cft:eq:bound}から、式\eqref{cft:eq:object}と\eqref{cft:eq:insertions}は
この有限自由加群内で$\mathfrak m$進収束する。
すべての恒等式は各Artin商で成立しているため、分離性から$R$上でも成立する。
これは同値の列の形式的な逆極限を宣言するのではなく、同じ基礎行列上の
関手と逆関手の明示的な係数wise極限である。
この極限には有限自由加群の$\mathfrak m$進完備性・分離性と、
各Artin商での有限評価だけを使っており、$R$の正則性や有限大域次元は必要ない。

内部shiftは基礎加群のラベルを変えるだけであり、行列上の公式は不変である。
各$F_r,H_r$は$R[t,t^{-1}]$-多重線形なので、$t$が与える
$(3)\simeq[2]$の自然比較とも可換する。
従って第\ref{gdr:sec:main}節で固定した周期比較がそのまま保存される。
正規化された単位も行列上で保たれる。
補題\ref{cft:lem:model}の厳密同型$\phi$と合成し、有限cone・冪等完備化へ延長して
式\eqref{cft:eq:equivalence}を得る。
\end{proof}

% !TEX root = ../../nc-cubic.tex
\section{一般点への移行と有限降下}\label{rog:sec:generic}

\subsection{残余体stabilizationのEuler形式}

今度は$d=3$とし、$F\in k[x_1,\ldots,x_N]_3$を滑らかな射影cubicを定める三次式とする。
$K_F$を原点の残余体のstabilizationとする。
\begin{proposition}[非零のEulerベクトル]\label{rog:prop:euler}
$r=0,1,2$について
\[
 \chi_r:=\sum_p(-1)^p\dim_k\operatorname{Hom}(K_F,K_F(r)[p])
 =\sum_{\substack{0\le s\le N\\s\equiv-r\pmod3}}(-1)^s\binom Ns.
 \tag{RO3}\label{rog:eq:euler}
\]
この三成分ベクトルは零でない。より正確に
\[
 \chi_0+\chi_1+\chi_2=0,\qquad
 \chi_0^2+\chi_1^2+\chi_2^2=2\cdot3^{N-1}.
 \tag{RO4}\label{rog:eq:norm}
\]
$N=5$では$(\chi_0,\chi_1,\chi_2)=(-9,9,0)$である。
\end{proposition}
\begin{proof}
$F=\sum_i x_if_i$、$f_i=(\partial_iF)/3$、
$f_i=\sum_jx_ja_{ij}$、$a_{ij}=(\partial_jf_i)/2$と置く。
外積作用$e_i\wedge-$の二重次数を$(-1,1)$、
収縮$\iota_i$を$(1,-1)$、周期変数$t$を$(2,-3)$とする。
$K_F$の微分は
\[
 Q=\sum_i x_i\iota_i+t\sum_i f_i(e_i\wedge-),\qquad Q^2=tF.
\]
そのendomorphism代数のsuper交換子微分について
\[
 d(e_i\wedge-)=x_i,\qquad
 \xi_i:=\iota_i-t\sum_j a_{ij}(e_j\wedge-),\qquad d\xi_i=0.
\]
順序を固定した$e_J\xi_I$はendomorphism加群の基底である。
実際、標準の$e_J\iota_I$からの変換は三角的である。
この基底では微分は$x_1,\ldots,x_N$のKoszul微分となる。
また$\xi_i\xi_j+\xi_j\xi_i=-t(a_{ij}+a_{ji})$であり、右辺は原点で零となる一次式なので境界である。
従って、周期変数を残した二重次数付きコホモロジーは
\[
 \bigwedge(\xi_1,\ldots,\xi_N)\otimes_k k[t,t^{-1}],
 \qquad\deg\xi_i=(1,-1)
 \tag{RO5}\label{rog:eq:exterior}
\]
となる。これは\cite[\S5.6]{ROGDyckerhoff}のClifford計算の三次の場合でもある。

$\xi_I t^j$が内部shift $r$のHomに寄与する条件は
$-|I|-3j=r$であり、cohomological次数は$|I|+2j$である。
$0\le|I|\le N$なので各$r$について有限個しか寄与せず、符号は$(-1)^{|I|}$。
これで\eqref{rog:eq:euler}を得る。

$A(z)=(1-z)^N$を$z^3-1$で割った余りの係数が、順序を除いて$\chi_r$である。
原始三乗根$\zeta$において$A(\zeta)=(1-\zeta)^N\ne0$なので三成分は同時には零でない。
$A(1)=0$、$|1-\zeta|^2=3$と三点のFourier公式から\eqref{rog:eq:norm}を得る。
これは整数恒等式なので、基礎体に依らない。$N=5$の値は直接代入する。
\end{proof}

\begin{remark}
一般$N$で自己Euler数$\chi_0$だけを使うのは誤りである。
例えば$N=9$では$\chi_0=0$となるが、三成分ベクトルは非零である。
また上の群は\emph{生成対象間のHom}であり、$\operatorname{Nat}(\mathrm{Id},\rho^r)$ではない。
\end{remark}

\subsection{固定した有限対象上では形式functorを先に局所化できる}

$R=\mathbb C[h]_{(h)}$、$K=\mathbb C(h)$、
$\widehat R=\mathbb C[[h]]$、$L=\mathbb C((h))$とする。
代数的中心三次対$(A_R,\Omega_R)$について、次の既存入力を使用する。
\begin{description}
\item[\textup{(L)}] 次数別flatな局所量子化と有限Koszul分解、
線束生成残余圏のproper模型、regular底のstabilization/window比較がある。
原点の$R$-加群の有限rank stabilizationを$K_R$とすると
\[
 T_R=K_R\oplus K_R(1)\oplus K_R(2),\quad
 \Gamma_R=R\operatorname{End}(T_R),\quad
 \mathcal K^{\rm lin}_{B_R}\simeq\operatorname{Perf}(\Gamma_R),
 \qquad \Gamma_R\in\operatorname{Perf}(R).
 \tag{RO6}\label{rog:eq:source}
\]
この模型とHomは$\widehat R,L$への平坦底変更に両立する。
\item[\textup{(CF)}] 次数別完備底上の有限rank graded MFに、
内部shiftと指定周期性を保つ実際の$\widehat R$-線形$A_\infty$同値
\[
 \Phi:\operatorname{MF}^{\mathbb Z}(A^{\rm gr}_{\widehat R},\Omega)
 \simeq\operatorname{MF}^{\mathbb Z}(\widehat R[x],W_h)
 \tag{RO7}\label{rog:eq:completed}
\]
がある。$W_h\bmod h=F$は滑らかで、$\Phi$は剰余で恒等。
単なる両立したArtin同値の列ではなく、
定理\ref{cft:thm:complete}で構成した固定内部次数内で収束する
実際の関手を用いる。
\end{description}
記号(CF)は定理\ref{cft:thm:complete}で構成した完備底上の関手を指す。
本節では、これらの比較定理を入力として、
固定三対象の局所化から一般点の生成性へ進む部分を独立に記述する。



\begin{lemma}[有限対象への制限とLaurent局所化]\label{rog:lem:localization}
\textup{(L),(CF)}の下で$E=\Phi(K_{\widehat R})$と置くと、
$E$は通常の$\widehat R[x]$上の有限rank graded MFであり、
\[
 \operatorname{Perf}(\Gamma_R\otimes_RL)
 \simeq\operatorname{thick}\{E_L,E_L(1),E_L(2)\}
 \tag{RO8}\label{rog:eq:image}
\]
というrotationと周期性を保つ同値がある。
また$\operatorname{Hom}(E,E(r))$は$r=0,1,2$で$\widehat R$-perfectであり、
その特殊・一般ファイバー間のEuler数は\eqref{rog:eq:euler}に等しい。
\end{lemma}
\begin{proof}
固定した有限rankの基礎加群と固定cohomological次数では、MF行列の各成分の
内部次数と$t$の指数は固定される。従って各成分は有限自由な
$\widehat R$-加群$\widehat R[x]_j$に属する。
(CF)の$F_{2q+1}=O(h^q)$によるtwistingの和は、その有限自由成分内で収束する。
結果は$x$の次数を無制限に増やす級数ではなく、$h$の形式係数を持つ有限次数の多項式行列である。
固定個数の射に対する各Taylor成分も同じ理由で通常の$\widehat R$-多重線形写像となる。
恒等式と逆functorは全Artin剰余で成立し、成分の分離性によって$\widehat R$上で成立する。

従って$T_{\widehat R}$の有限対象部分圏では、(CF)は\emph{通常の}endomorphism複体の
$A_\infty$準同型を与える。$\widehat R\to L$は平坦なので、これをテンソルすると
通常の準同型である。固定cohomological次数のHomの形成はこの平坦底変更と可換する。
$\Phi$を任意のLaurent係数対象へ新たに拡張することは使わず、
この固定部分圏のpretriangulated・冪等包絡を取って\eqref{rog:eq:image}を得る。
有限twisted complexのtwisting行列は厳密上三角に取れるため、包絡への延長では有限個のTaylor成分だけが使われる。
したがって$L$への移行後に新しい無限級数の収束を要求しない。

Homのperfect性は(L)と(CF)から従い、$K_R(r)$は周期性により$T_R$のshiftの直和因子である。
有界有限射影複体のEuler数は特殊化で不変なので、剰余の$K_F$について計算した
\eqref{rog:eq:euler}が一般ファイバーでも保たれる。
内部shiftと周期性の自然比較も、この固定部分圏への制限と平坦底変更で保たれる。
\end{proof}

\subsection{体拡大上の同値から有限代数拡大上の同値へ}

この降下の基礎は、smooth proper dg代数の有限表示・降下
\cite[Th\'eor\`eme~1]{ROGToenDefinition}、perfect加群のfiltered colimitとの両立
\cite[Lemma~2.10]{ROGTV}、および同値の核を単位・余単位によって判定する原理
\cite[Lemma~3.25]{ROGTV}である。
以下では片側のsmoothnessから始め、
rotationと指定した有限個の周期図式を保つ形で、これらの原理を適用する。
そのため、降下核と同値を判定するconeの底上のperfect性を明示する。

\begin{lemma}[固定された可換模型へのgeneric降下]\label{rog:lem:descent}
$K$を標数零の体、$L/K$を任意の体拡大、$\Gamma,\mathsf C$を$K$上properなdg代数とし、
$\mathsf C$をsmoothとする。両側にrotationを表す可逆双加群と、必要なら三次周期同型を指定する。
$L$上でrotation・周期性を保つMorita同値があれば、ある有限代数拡大$K'/K$上でも
同じ二つの代数の間にそれらを保つMorita同値がある。
\end{lemma}
\begin{proof}
$L$上の同値を可換側から非可換側への核$P_L$で表す。
これは右$\Gamma_L$-perfectであり、$\mathsf C_L$のsmoothnessにより
$\mathsf C_L^{\rm op}\otimes_L\Gamma_L$上でもperfectである
\cite[Lemma~2.8]{ROGTV}。
$L$を有限生成$K$-部分代数のfiltered colimitとして表し、
\cite[Lemma~2.10]{ROGTV}により核とrotation比較を一つの有限生成整域$D\subset L$へ降ろす。
$\mathsf C_D$は$D$-properなので、包絡代数上のperfect加群を右$\Gamma_D$へ制限してもperfectである。

$Q_D=R\operatorname{Hom}_{\Gamma_D}(P_D,\Gamma_D)$と置く。
$P_D$が右$\Gamma_D$-perfectで、$\Gamma_D$が$D$-properなので、
$P_D$と$Q_D$はいずれも$D$-perfectである。
$Q_D$を右$\mathsf C_D$-加群と見ると、$\mathsf C_D$のsmoothnessと
\cite[Lemma~2.8(2)]{ROGTV}により右$\mathsf C_D$-perfectでもある。
従って$-\otimes_{\mathsf C_D}^{\mathbf L}P_D$と
$-\otimes_{\Gamma_D}^{\mathbf L}Q_D$はperfect加群の圏の間の随伴対をなし、
その単位・余単位は双加群の射
\[
 \mathsf C_D\longrightarrow P_D\otimes_{\Gamma_D}^{\mathbf L}Q_D,
 \qquad Q_D\otimes_{\mathsf C_D}^{\mathbf L}P_D\longrightarrow\Gamma_D
\]
で表される。第一のテンソル積は$P_D$の右$\Gamma_D$-perfect性と
$Q_D$の$D$-perfect性から、第二のテンソル積は$Q_D$の右$\mathsf C_D$-perfect性と
$P_D$の$D$-perfect性から、それぞれ$D$-perfectである。
両方の対角も$D$-perfectなので、単位・余単位のconeは$D$-perfectとなる。
また$P_D$の右$\Gamma_D$-perfect性により、$Q_D$とこの随伴の形成は
$D\to L$の底変更と可換する。従って両coneは$L$への底変更後に零となる。

$D$はnoetherian整域なので、これらのconeは有限個の有限生成cohomology加群を持つ。
$D\subset L$より各加群は$D$-torsionであり、有限生成元のannihilatorの積を取れば、
ある$0\ne f\in D$の反転後に両coneが零となる。
rotation比較の両辺も、可逆双加群によるテンソルが片側のperfect性を保つことと
両代数のproper性から$D$-perfectであり、そのconeを同じ方法で消せる。
指定周期性の図式はperfect双加群間の有限個の射の等式であるから、
必要なら段階を拡大し、さらに一つの元を反転して等式を保つ。
任意の高次群作用を有限の周期図式と同一視してはいない。

$\operatorname{Spec}D_f$の閉点を一つ取る。その剰余体$K'$はZariskiの補題により$K$の有限拡大であり、
標数零なので分離的である。そこでの底変更が所要の同値を与える。
これは元の$L$上の形式分枝を全次数で再現する主張ではない。
\end{proof}

\begin{theorem}[一変数量子化の幾何学的一般点]\label{rog:thm:geometric-generic}
\textup{(L),(CF)}を満たす一変数中心三次量子化に対し、
ある有限拡大$K'/\mathbb C(h)$と滑らかな三次超曲面$Y/K'$が存在して
\[
 (\mathcal K^{\rm lin}_{B_K}\otimes_KK',\rho)
 \simeq(\operatorname{Ku}(Y),\rho_Y)
 \tag{RO10}\label{rog:eq:generic}
\]
となる。$\mathcal K^{\rm lin}$と全coherent残余圏との同一性は、この定理に含めない。
\end{theorem}
\begin{proof}
補題\ref{rog:lem:localization}により、$L=\mathbb C((h))$上で像は
$E_L,E_L(1),E_L(2)$の厚い閉包になる。
命題\ref{rog:prop:euler}とproper HomのEuler数の底変更により$E_L\ne0$である。
$W_h$の判別式は特殊ファイバーで非零なので$L$上も滑らかである。
$L$の代数閉包$\overline L$へ移し、命題\ref{gcd:prop:bound}を適用すると
像は可換MF圏全体となる。体上のOrlov比較により$\overline L$上で可換
Kuznetsov成分との同値を得る。以下では$L$をこの$\overline L$に置き換えて
有限降下を行う。
$W_h$の有限個の三次係数と判別式の逆を含む有限生成$K$-部分代数$D_0\subset L$を取ると、
$Y_0=V(W_{D_0})$は$D_0$上の滑らかな射影cubic族である。
その線束の有限和を残余成分へ射影した生成子を$H_{D_0}$、
生成dg代数を$\mathsf C_{D_0}=R\operatorname{End}(H_{D_0})$とする。
これはsmooth properで、底変更と両立する。
ここで生成子は相対ampleな全線束が生成することと、座標のKoszul完全列によって
$N$個の連続したtwistへ減らすことから得られる。
射影は有限個のexceptional mutationによる。
この有限線束生成の基底環に関する一般性は\cite[\S36.16]{ROGStacks}を参照する。

$L$上の同値核を、$D_0$を含む有限生成$K$-部分代数へ降ろす。
以後は補題\ref{rog:lem:descent}の証明と同じであり、
properな二つのconeを消して閉点を選ぶ。
その点では$Y_0$も滑らかなcubicへ特殊化し、\eqref{rog:eq:generic}を得る。
この一般存在定理は、$Y$の係数の閉公式や同値核の行列を明示するものではない。
\end{proof}


% !TEX root = ../../nc-cubic.tex
% 2026-11-13. Coherent saturation with explicit admissibility.
\section{線束生成部分の飽和と全coherent残余圏}
\label{crs:sec:saturation}

本節は、一般点復元で用いた線束生成部分と全coherent導来圏の差を解消する。
鍵は、一般点で既に得た線束生成残余圏のsmoothnessを、
全coherent対象に対する有限性と組み合わせることである。
特殊点の生成性を一般点へ持ち上げる議論は使わない。

既知の入力は、Burke--Stevensonの導来section、Gorenstein環の
$\chi$条件と有限cohomological dimension、およびOrlov型分解
\cite[Proposition~4.18, Lemmas~6.1--6.2, Theorem~6.4]{LP-BS}、
smooth dg圏上のpseudo-perfect加群のperfect性
\cite[Lemma~2.8(2)]{ROGTV}、smooth properな充満部分圏のadmissibilityと
半直交分解に沿うsmoothnessの貼合せ
\cite[Lemmas~4.13, 4.15(3)]{CRSPerry}である。
本節ではこれらの仮定を一般点の中心三次商について確認し、
線束生成残余pairの復元を全coherent残余pairへ拡張する。
この組合せの包括的新規性調査・外部査読は未実施である。

\subsection{抽象的な飽和判定}

$k$を体、$\mathcal D$を小さな冪等完備pretriangulated $k$-線形dg圏とし、
$G\in\mathcal D$、$\mathcal C=\operatorname{thick}(G)$、
$\Lambda=\operatorname{RHom}_{\mathcal D}(G,G)$と置く。
以下の直交は全cohomological shiftに関するものである。

\begin{theorem}[有限性とsmoothnessによる全体への飽和]
\label{crs:thm:abstract}
次の三条件を仮定する。
\begin{enumerate}
\item $\Lambda$は$k$上smooth、すなわち
$\Delta_\Lambda\in\operatorname{Perf}(\Lambda^{\mathrm{op}}\otimes_k\Lambda)$である。
\item 任意の$M\in\mathcal D$について
$\operatorname{RHom}(G,M)\in\operatorname{Perf}(k)$である。
\item $\mathcal C^\perp=0$である。
\end{enumerate}
このとき$\mathcal C=\mathcal D$であり、
$\operatorname{Perf}(\Lambda)\to\mathcal D$はdg Morita同値である。
\end{theorem}
\begin{proof}
右$\Lambda$-加群$N_M=\operatorname{RHom}(G,M)$を考える。
smoothな対角の有限cell/retract表示へ
$N_M\otimes_\Lambda^{\mathbf L}(-)$を適用する。
自由な包絡加群の像は$N_M\otimes_k\Lambda$であり、
$N_M\in\operatorname{Perf}(k)$だから右$\Lambda$-perfectである。
従って$N_M\in\operatorname{Perf}(\Lambda)$となる。
これは\cite[Lemma~2.8(2)]{ROGTV}の一対象版の直接証明である。

$G$を用いるMorita関手
$F:\operatorname{Perf}(\Lambda)\xrightarrow{\sim}\mathcal C$
で$N_M$を移し、評価射$F(N_M)\to M$を取る。
ここで$F(N_M)$は有限coneと直和因子だけで作られ、
$\mathcal D$に無限直和があることを要求しない。
$F$と$\operatorname{RHom}(G,-)$の単位は$\Lambda$で同型なので、
有限操作による帰納法で全perfect加群に対して同型である。
従って評価射に$\operatorname{RHom}(G,-)$を施すと同型となる。
評価射のconeは$G$、従って$\mathcal C$と右直交する。
第三条件によりこのconeは零である。
全$M$が$\mathcal C$に属し、enhancedな完全忠実性は
$\operatorname{thick}(G)$におけるMorita表示から従う。
\end{proof}

\subsection{中心三次商のcoherent対象に対する有限性と検出}

$A$を連結局所有限な両側noetherian Koszul AS-regular $k$-代数とし、
global dimensionとGorenstein parameterをともに$N\ge5$とする。
$w\in Z(A)_3$を正則元、$B=A/(w)$、$a=N-3$と置く。
$\operatorname{qgr}B$は有限生成次数付き右加群を有限次元加群で割った圏、
$\mathcal O(j)=\pi B(j)$とする。標準のenhancementを用いる。

\begin{lemma}[必要なGorenstein・section条件]
\label{crs:lem:sections}
$B$は両側noetherian AS-Gorensteinで、parameterは$a$である。
$B$は$\chi$条件と有限cohomological dimensionを満たす。
任意の$M\in D^b(\operatorname{qgr}B)$と任意の$j\in\mathbb Z$について
\[
 \operatorname{RHom}_{\operatorname{qgr}B}(\mathcal O(j),M)
 \in\operatorname{Perf}(k).
 \tag{CS1}\label{crs:eq:finite-Hom}
\]
また、全$j,p$で$\operatorname{Hom}(\mathcal O(j),M[p])=0$なら$M=0$である。
\end{lemma}
\begin{proof}
$0\to A(-3)\xrightarrow{w}A\to B\to0$と導来change-of-ringsにより、
$B$-加群$M$について
\[
 R\underline{\operatorname{Hom}}_B(M,B)
 \simeq R\underline{\operatorname{Hom}}_A(M,A)(-3)[1]
\]
を得る。ここでは$w$の中央性を使い、normal元のtwistを省略したのではない。
左右でinjective dimensionは有限であり、$M=k$を代入すると
$R\underline{\operatorname{Hom}}_B(k,B)\simeq k(N-3)[-(N-1)]$となる。
$B$のnoetherian性は商であることから従う。
従って\cite[Lemmas~6.1--6.2]{LP-BS}を適用できる。

\cite[Proposition~4.18とその直前の導来随伴]{LP-BS}により、任意の整数$b$について
\[
 R\Gamma_{\ge b}:D^b(\operatorname{qgr}B)
 \longrightarrow D^b(\operatorname{gr}_{\ge b}B)
 \tag{CS2}\label{crs:eq:sections}
\]
は完全忠実で、各cohomologyは有限生成次数付き$B$-加群である。
固定内部次数の成分は有限次元であり、cohomologyは有界だから、
$\bigl(R\Gamma_*M\bigr)_{-j}=R\operatorname{Hom}(\mathcal O(j),M)$は
$k$-perfectとなる。任意の$j$に対し$b\le-j$を選べばよい。

全ての線束とのHomが零なら$R\Gamma_{\ge0}M$の各次数のcohomologyが零である。
この複体自体が零となり、完全忠実性から$M=0$を得る。
この検出は高次Extも含み、射の次数零だけを調べているのではない。
\end{proof}

\begin{lemma}[有限線束生成対象]
\label{crs:lem:finite-lines}
\[
 \mathcal C_B=\operatorname{thick}\{\mathcal O(j):j\in\mathbb Z\},
 \qquad G_B=\bigoplus_{j=0}^{N-1}\mathcal O(j)
\]
とすると$\mathcal C_B=\operatorname{thick}(G_B)$であり、$\mathcal C_B$はproperである。
\end{lemma}
\begin{proof}
$k$の$A$上の有限Koszul分解を$B$へ導来テンソルする。
二項の$A$-分解で計算すると$\operatorname{Tor}^A_i(k,B)$は
$i=0,1$にのみあり、それぞれ$k,k(-3)$である。
従ってこのKoszul複体を$qgr$へ移すと完全になる。
両端は一つの線束で、中間の項のshiftは$1,\ldots,N-1$にある。
任意の連続$N$個の線束から、一つ手前と一つ後をそれぞれ生成できる。
前後への帰納により全てのtwistを得る。
proper性は補題\ref{crs:lem:sections}を有限個の$\mathcal O(i),\mathcal O(j)$に適用し、
有限coneと直和因子へ延長して従う。
\end{proof}

\begin{theorem}[線束部分のsmoothnessで全coherent圏を得る]
\label{crs:thm:coherent}
上の中心三次商について、$\mathcal C_B$がsmoothなら
\[
 \boxed{\quad
 \mathcal C_B=D^b(\operatorname{qgr}B).
 \quad}\tag{CS3}\label{crs:eq:full}
\]
従って右辺はsmooth proper dg圏で、$G_B$が古典的生成対象になる。
\end{theorem}
\begin{proof}
任意のcoherent有界対象に対する$\operatorname{RHom}(G_B,-)$の有限性は
補題\ref{crs:lem:sections}による。線束による検出と
補題\ref{crs:lem:finite-lines}から$\mathcal C_B^\perp=0$である。
定理\ref{crs:thm:abstract}を適用する。
ここで全$D^b(\operatorname{qgr}B)$のsmoothnessを仮定していない。
\end{proof}

\subsection{残余部分のsmoothnessと有限拡大からの降下}

\begin{lemma}[有限体拡張からのsmoothness降下]
\label{crs:lem:descent}
有限体拡大$k'/k$と$k$-dg代数$\Gamma$について、
$\Gamma\otimes_k k'$が$k'$上smoothなら$\Gamma$は$k$上smoothである。
\end{lemma}
\begin{proof}
$E=\Gamma^{\mathrm{op}}\otimes_k\Gamma$と置く。
$E\otimes_k k'$は$E$上有限自由なので、そのperfect加群のrestrictionは$E$-perfectである。
従ってsmoothな対角$\Gamma\otimes_k k'$をrestrictionしたものは$E$-perfect。
$\lambda:k'\to k$を$\lambda(1)=1$となる$k$-線形写像とすれば、
$\Gamma\to\Gamma\otimes_k k'\xrightarrow{1\otimes\lambda}\Gamma$は
$E$-線形なretractである。従って$\Gamma$は$E$-perfectである。
\end{proof}

\begin{proposition}[残余のsmoothnessから全線束部分へ]
\label{crs:prop:residual-smooth}
$\mathcal K_B^{\rm lin}=\mathcal C_B\cap
\langle\mathcal O,\ldots,\mathcal O(a-1)\rangle^\perp$とする。
$\mathcal K_B^{\rm lin}$がsmoothなら$\mathcal C_B$もsmoothである。
従ってこの場合
\[
 \mathcal K_B^{\rm lin}=\mathcal K_B^{\rm coh}:=
 D^b(\operatorname{qgr}B)\cap
 \langle\mathcal O,\ldots,\mathcal O(a-1)\rangle^\perp.
 \tag{CS4}\label{crs:eq:residual-full}
\]
\end{proposition}
\begin{proof}
$\mathcal O,\ldots,\mathcal O(a-1)$は例外列である。
これはGorensteinの線束cohomology公式、または
\cite[Theorem~6.4]{LP-BS}のOrlov分解から従う。
$\mathcal C_B$はproperなので、例外列への評価・余評価は有限であり、
その生成部分はadmissibleである。これにより半直交分解
\[
 \mathcal C_B=\langle\mathcal K_B^{\rm lin},
                 \mathcal O,\ldots,\mathcal O(a-1)\rangle
\]
を得る。$\mathcal K_B^{\rm lin}$はproperな$\mathcal C_B$の充満部分圏なのでproperであり、
仮定によりsmoothでもある。従って、その包含
$\mathcal K_B^{\rm lin}\hookrightarrow\mathcal C_B$は
\cite[Lemma~4.13]{CRSPerry}により左右随伴を持つ。
これで残余成分もadmissibleとなる。
全成分のadmissibilityを確認した上で、
smoothnessの貼合せ\cite[Lemma~4.15(3)]{CRSPerry}を適用する。
具体的には各例外成分は$\operatorname{Perf}(k)$でsmooth、
相互Homはproperであり、smoothな成分の間の貼合せ双加群は
\cite[Lemma~2.8(2)]{ROGTV}によりperfectである。
上三角dg圏の対角の有限三角を順に貼り合わせれば全体の対角もperfectとなる。
定理\ref{crs:thm:coherent}から\eqref{crs:eq:residual-full}を得る。
\end{proof}


\section{低rotation積と勾配空間の射影的一意性}
\label{sec:explicit-jacobian-recovery}

ここでは方程式の基礎体降下に必要な二つの入力を示す。
導来自然変換の低次数積の同定はLin--Zhangの計算を用い、
同じ勾配空間から得られる滑らかな三次式の射影的一意性は直接証明する。
復元アルゴリズムの実例や積分不可能な二次関係の例は扱わない。

\begin{proposition}[可換三次超曲面における低次数積]
\label{prop:ej-geometric-low-degree}
$k$を標数零の代数閉体、$F\in k[x_1,\ldots,x_N]_3$を
滑らかな射影三次超曲面$X\subset\mathbb P^{N-1}$の方程式とし、$N\ge5$とする。
$\mathcal C=\operatorname{Ku}(X)$、その対角双加群を$\Delta$、
rotationの双加群を$P$と置く。このとき
\[
 H^0R\operatorname{Hom}_{\mathcal C^e}(\Delta,P)
 \simeq\operatorname{Jac}(F)_1,
 \qquad
 H^0R\operatorname{Hom}_{\mathcal C^e}
 (\Delta,P\otimes_{\mathcal C}^{\mathbf L}P)
 \simeq\operatorname{Jac}(F)_2
\]
であり、合成による積はJacobian環の積である。
特に、この積は対称積を経由し、誘導される全射
\begin{equation}\label{ej:eq:symmetric-product}
 \operatorname{Sym}^2
 H^0R\operatorname{Hom}_{\mathcal C^e}(\Delta,P)
 \longrightarrow
 H^0R\operatorname{Hom}_{\mathcal C^e}
 (\Delta,P\otimes_{\mathcal C}^{\mathbf L}P)
\end{equation}
の核は、上の同定の下で
$\langle\partial_1F,\ldots,\partial_NF\rangle$に一致する。
\end{proposition}
\begin{proof}
Orlovのdg比較をrotationも含めて用いると、$\rho_X$はgraded matrix
factorizationの内部次数シフトとなる。
Lin--Zhang\cite[Theorems~2.2, 3.5, Remark~3.6, Section~5]{LZ}の
対角kernelの計算を適用する。
恒等sectorは$\operatorname{Jac}(F)_t$を与える。
三次式の非恒等sectorは全$N$変数を動かすので、
cohomological degree零の$\operatorname{Hom}(\mathrm{Id},\rho^t)$に寄与できるのは
$N$が偶数で$t=N/2$の場合に限る。
$N$が奇数なら寄与はなく、$N$が偶数でも$N\ge6$なので$t=1,2$では寄与しない。
従って表示した二つの同定を得る。
このsector計算には$\gcd(N,3)=1$を必要としない。
同論文Proposition~3.7の互いに素という条件は、内部シフトを
Serre関手とcohomological shiftで表示してSerre代数へ埋め込むための条件であり、
ここでは$\rho$自体を入力している。

恒等sector上の合成は多項式の積である\cite[Remark~3.6]{LZ}。
Jacobianイデアルは二次式$\partial_iF$で生成され、次数一部分は零である。
滑らかさにより$\partial_iF$は一次独立である。
実際、定数係数の関係があれば線形座標変換後に$F$は一つの変数に依らず、
その変数軸上の射影点が特異点になる。
従って対称積上の次数二の核は勾配空間であり、積は全射である。
$V\otimes V$からの積の核には$\bigwedge^2V$も含まれるので、
勾配空間は必ず式\eqref{ej:eq:symmetric-product}の核として取る。
\end{proof}

rotation付きの滑らかなFano超曲面の幾何学的一意性は
\cite[Theorem~5.1]{LZ}により既知である。
勾配空間が滑らかな斉次式を定数倍まで決定する際のSebastiani--Thom型の例外は
Ueda--Yoshinaga\cite[Lemma~3]{UedaYoshinaga}と
Wang\cite[Theorem~1.1, Corollary~1.3]{WangJacobian}により調べられている。
以下のsymmetrizerによる説明は
Hwang\cite[Propositions~2.2, 2.6, 3.1, 3.2]{HwangSymmetrizer}の
滑らかな三次式の場合の基本的な帰結である。
基礎体降下に用いるため、任意の標数零の体上で必要な内容を証明する。

\begin{lemma}[滑らかな三次式のsymmetrizer代数]
\label{ej:lem:symmetrizer-algebra}
$k$を標数零の体、$U$を有限次元$k$-線形空間とし、
$F\in\operatorname{Sym}^3U^*$は幾何学的に滑らかな射影超曲面を定めるとする。
$F(x)=T(x,x,x)$を満たす対称三重線形形式$T$について
\begin{equation}\label{eq:ej-symmetrizer}
 E_F=\{a\in\operatorname{End}_k(U):
 T(au,v,w)=T(u,av,w)=T(u,v,aw)\text{ for all }u,v,w\}
\end{equation}
と置く。このとき$E_F$は有限\'etaleな可換$k$-代数であり、
その形成は体拡大と可換である。

$s=\dim_kE_F$とすると、代数閉包上で$E_F\otimes_k\overline k\simeq\overline k^s$である。
その原始冪等元による分解
\[
 U_{\overline k}=\bigoplus_{i=1}^sU_i,\qquad
 F_{\overline k}=\sum_{i=1}^sF_i,\qquad
 F_i\in\operatorname{Sym}^3U_i^*
\]
は、各$F_i$がさらに直和分解できない最細の幾何学的Sebastiani--Thom分解である。
この分解は成分の順序を除いて一意である。
特に、$F_{\overline k}$が非自明なSebastiani--Thom分解を持たないことと$E_F=k$は同値である。
\end{lemma}
\begin{proof}
式\eqref{eq:ej-symmetrizer}は有限個の$k$-線形方程式なので、形成は体拡大と可換である。
$a,b\in E_F$について三つの引数の間で写像を移すと
\[
 \begin{aligned}
 T(abu,v,w)&=T(bu,av,w)=T(u,av,bw)\\
 &=T(au,v,bw)=T(au,bv,w)=T(bau,v,w).
 \end{aligned}
\]
滑らかさにより$u\mapsto T(u,-,-)$は単射なので$ab=ba$である。
さらに$T(abu,v,w)=T(u,bav,w)=T(u,abv,w)$であり、第三の引数へも同様に移せる。
従って$ab\in E_F$で、$E_F$は単位を持つ可換部分代数となる。

代数閉包上でこの代数に非零冪零元があれば、その適当な冪を取って
$0\ne b\in E_{F_{\overline k}}$かつ$b^2=0$とできる。
$u=bv\ne0$を選ぶと
\[
 T(u,u,w)=T(bv,bv,w)=T(v,b^2v,w)=0
 \quad\text{for every }w.
\]
これは$[u]$が$V(F_{\overline k})$の特異点であることを意味し、矛盾する。
よって$E_F$は幾何学的に被約な有限次元可換代数、すなわち有限\'etale代数である。

$E_{F_{\overline k}}$の原始冪等元を$e_1,\ldots,e_s$とし、$U_i=e_iU_{\overline k}$と置く。
$e_i$は非零の線形写像なので各$U_i\ne0$である。
$u\in U_i$、$v\in U_j$、$i\ne j$なら
$T(u,v,w)=T(e_iu,v,w)=T(u,e_iv,w)=0$であり、混合項は消える。
従って$F=\sum F_i$となる。
逆に、任意のSebastiani--Thom分解の各成分への射影は$E_{F_{\overline k}}$の冪等元である。
原始冪等元はそれ以上分割できないので各$F_i$は直和分解不能であり、
任意の分解はこの分解をまとめて得られる。
これで最細性と一意性、ならびに最後の同値が従う。
\end{proof}

\begin{proposition}[滑らかな三次式の勾配空間による一意性]
\label{prop:ej-projective-uniqueness}
標数零の代数閉体上で、同じ二次式空間を勾配空間に持つ
二つの滑らかな三次式は線形座標変換で移り合う。
\end{proposition}
\begin{proof}
$F,G$をその二つの三次式とする。
それぞれの勾配は同じ空間の基底なので$\nabla G=A\nabla F$を満たす可逆行列$A$がある。
$F$を、$F(x)=T(x,x,x)$を満たす対称三重線形形式$T$で表す。
$a=A^{\mathsf t}$とすれば、$G$の対称性から$a\in E_F^\times$であり、
$G(x)=T(ax,x,x)$である。
補題\ref{ej:lem:symmetrizer-algebra}の分解で$a=(\lambda_1,\ldots,\lambda_s)$と書くと
\[
 F=\sum_{i=1}^sF_i,\qquad G=\sum_{i=1}^s\lambda_iF_i,
 \qquad\lambda_i\ne0.
\]
各$U_i$上で$\lambda_i^{1/3}$倍する線形変換により$F$を$G$へ移せる。
この線形変換が射影同型を与える。
\end{proof}

% General-dimensional base-field descent and quantitative generation.
\section{一般点復元の基礎体降下と定量的生成}
\label{gcd:sec:descent}

本節では可換模型の\emph{幾何学的一般点での存在}を得た後の議論を整理する。
低rotation積の可積分性により、方程式を元の体上に選ぶ。
方程式と同値核の定義体の区別は、第\ref{gdr:subsec:scope}節(2)の通りである。
必要な低次数積の同定と勾配空間の射影的一意性は
第\ref{sec:explicit-jacobian-recovery}節に証明を収める。

\subsection{次数別flat入力から従う一般ファイバーのAS正則性}

\begin{proposition}[一般点のambientに対する追加仮定の除去]
\label{gcd:prop:automatic-as}
$(R,\mathfrak m,k)$をnoetherian局所整域、$K=\operatorname{Frac}R$とする。
$A_R$は非負次数の中央$R$-代数、$(A_R)_0=R$、各$(A_R)_j$は有限自由で、
$A_R/\mathfrak mA_R\simeq k[x_1,\ldots,x_N]$とする。
このとき$A_K$は左右noetherianな二次Koszul AS正則代数であり、
大域次元とGorenstein parameterは$N$、Hilbert級数は$(1-z)^{-N}$である。
さらに$\Omega_R\in Z(A_R)_d$の剰余が非零なら、$\Omega_R$は左右正則で、
商の各次数は$R$-自由である。
\end{proposition}
\begin{proof}
補題\ref{lp:lem:noetherian}の次数別Krull交叉の議論を$A_R$に適用すると、
左右noetherian性を得る。中央集合$R\setminus\{0\}$による局所化で$A_K$もそうなる。
ここでnoetherian性を有限Hilbert級数だけから推論していない。

補題\ref{lp:lem:resolution}を、augmentation加群$R$と、
$A_R^e=A_R\otimes_RA_R^{\mathrm{op}}$上の対角$A_R$に適用する。
多項式環のKoszul分解と対角分解を持ち上げると、項がそれぞれ
$A_R(-j)^{\binom Nj}$、$A_R^e(-j)^{\binom Nj}$である長さ$N$の分解を得る。
持上げの各段で、剰余上の生成元を持ち上げた写像の全射性を各内部次数の
Nakayamaで示す。有限自由$R$-加群への全射は分裂するので核も各次数で自由であり、
最後の核は剰余が零のため零となる。この構成は任意の通常の底変更に可換である。
したがって$A_K$は二次Koszulで、対角分解から左右大域次元は高々$N$となる。

augmentation分解の$A_R$双対を固定内部次数ごとに取る。
その剰余複体は、内部次数$-N$・コホモロジー次数$N$以外で完全である。
有界有限自由複体の単元行列成分を順に消去すれば、他の内部次数の複体は分裂完全となる。
内部次数$-N$では最後の$R$のみが残る。よって左右で
\begin{equation}\label{gcd:eq:as}
 R\underline{\operatorname{Hom}}_{A_R}(R,A_R)\simeq R(N)[-N]
\end{equation}
となり、底変更後の非零な最上Extから大域次元は正確に$N$である。
次数別自由性はHilbert級数を与え、増大度も多項式的である。
最後に各次数の$\Omega_R$乗法は剰余上で単射である。
列数と同じ大きさの単元小行列式を使えば$R$上でも分裂単射となり、
余核も自由となる。中央性により左右で同じ結論を得る。
\end{proof}

\subsection{smooth properなpairから取り出す低次数積}

$k$を標数零の体、$\mathcal C$をsmooth properな$k$-線形dg圏、
$\rho$をdg Morita自己同値とする。対角を$\Delta$、$\rho$の双加群を$P$と書き、
\begin{equation}\label{gcd:eq:low}
 V=H^0R\operatorname{Hom}_{\mathcal C^e}(\Delta,P),\qquad
 U=H^0R\operatorname{Hom}_{\mathcal C^e}(\Delta,P\otimes_{\mathcal C}^{\mathbf L}P)
\end{equation}
と置く。積$\mu:V\otimes V\to U$は$\rho(\beta)\circ\alpha$である。
通常の三角圏上の自然変換と、この導来自然変換を区別する。

\begin{lemma}[低次数積の体拡大との両立]\label{gcd:lem:basechange}
$V,U$は有限次元であり、任意の体拡大$\ell/k$に対して、その形成と積は底変更と可換である。
\end{lemma}
\begin{proof}
smoothnessにより$\Delta$は$\mathcal C^e$上perfectである。
proper性と$P$の可逆性により、$P$と$P\otimes_{\mathcal C}^{\mathbf L}P$は
基礎体上perfectである。有限cellとretractで$\Delta$を表示すると、
式\eqref{gcd:eq:low}のHom複体も$k$-perfectとなる。
同じ有限表示を使えばHomと平坦な体拡大が可換であり、$H^0$も可換である。
双加群のテンソルと射の合成はこの底変更に自然なので、積も保たれる。
\end{proof}

\begin{theorem}[可換cubic模型の方程式は基礎体上に取れる]
\label{gcd:thm:field-model}
$n\ge3$、$N=n+2$とする。上のpairがある代数閉拡大体$\Omega/k$上で
滑らかなcubic $Y/\Omega$のrotation付きKuznetsov pairとMorita同値であると仮定する。
このとき次を満たす$G\in\operatorname{Sym}^3_k V$が存在する。
\begin{enumerate}
\item $X_G=V(G)\subset\mathbb P(V^*)$は$k$上幾何学的に滑らかで、$X_{G,\Omega}\simeq Y$である。
\item ある有限拡大$k'/k$上で$(\mathcal C_{k'},\rho)\simeq(\operatorname{Ku}(X_{G,k'}),\rho_G)$である。
元の同値に三次周期比較を指定していた場合、それも保存できる。
\item pair $(\mathcal C,\rho)$の任意のrotation付き滑らかなcubic模型は、
共通の代数閉拡大体上で射影同型である。
\end{enumerate}
従って方程式の定義体には有限拡大を必要としない。
一意性と同値核の範囲は第\ref{gdr:subsec:scope}節(2)に従う。
\end{theorem}
\begin{proof}
命題\ref{prop:ej-geometric-low-degree}に用いる
Lin--Zhangのsector計算\cite[Theorem~3.5, Remark~3.6]{LZ}を$\Omega$上で適用する。
三次式の非恒等sectorがコホモロジー次数零へ寄与する内部次数は、
$N$が偶数の場合の$N/2$のみである。$N\ge5$では内部次数$1,2$に干渉しない。
したがって、$\Omega$上で$V,U,\mu$はJacobian環の次数$1,2$と積に一致する。
この計算は標数零の代数閉体上の代数的なfactorization計算である。
補題\ref{gcd:lem:basechange}と忠実平坦性から、$k$上でも
\begin{equation}\label{gcd:eq:gradient-data}
 \dim V=N,\quad \dim U=\binom N2,\quad
 \overline\mu:\operatorname{Sym}^2V\twoheadrightarrow U,\quad
 Q:=\ker\overline\mu,\quad\dim Q=N
\end{equation}
となる。積の対称性も降下する。
$Q$には代数閉包上で共通射影零点がない。これは$\Omega$上で勾配空間となることと
$Y$の滑らかさから従う。

$V$の基底を$x_i$、$Q$の基底を$q_a$とする。可積分性の線形解空間
\begin{equation}\label{gcd:eq:integrability}
 \mathcal L_Q=\left\{A\in\operatorname{Mat}_N(k):
 \sum_a A_{ia}\partial_jq_a=\sum_a A_{ja}\partial_iq_a\quad(i<j)\right\}
\end{equation}
は体拡大と可換である。$Y$の方程式の勾配を$q_a$で展開すると、
$\mathcal L_Q\otimes_k\Omega$に可逆行列が存在する。
$k$-基底$A_1,\ldots,A_r$を選べば
$D(z)=\det(\sum_\nu z_\nu A_\nu)$は零多項式でない。
$k$は無限体なので$D(c)\ne0$となる$c\in k^r$が存在する。
実際、各変数の次数は高々$N$なので$\{0,1,\ldots,N\}^r$のいずれかで非零となる。
$A=\sum c_\nu A_\nu$として
\begin{equation}\label{gcd:eq:model}
 G=\frac13\sum_{i,a}x_iA_{ia}q_a
\end{equation}
と置く。可積分性とEuler等式から$\partial_iG=\sum_aA_{ia}q_a$である。
従って$G$の勾配空間は$Q$、その共通射影零点はなく、$X_G$は幾何学的に滑らかである。

$\Omega$上で$G$と$Y$の方程式は同じ勾配空間を持つ。
命題\ref{prop:ej-projective-uniqueness}の対称三重線形形式の証明は
任意の標数零代数閉体で通るので、両者は線形座標変換で一致する。
この証明はsymmetrizer代数の可換性と非零冪零元の不存在を用い、
固有空間ごとの非零係数を三乗根で吸収するものである。
その座標変換はgraded MFの内部shiftと三次周期性を保つ。
従って$\Omega$上で$(\mathcal C,\rho)$と$X_G$のpairの同値を得る。

$X_G$は既に$k$上にあり、そのKuznetsov成分はsmooth properである。
固定された二つのproper dg模型に補題\ref{rog:lem:descent}を適用する。
perfect同値核、rotation比較、指定周期図式を有限生成$k$-部分代数へ降ろし、
単位・余単位のconeを消す非空開の閉点を選べば、有限拡大$k'/k$上の同値を得る。
第三の主張も、任意のrotation付き模型に式\eqref{gcd:eq:gradient-data}を適用し、
同じ勾配空間の射影的一意性を用いれば従う。
\end{proof}

\begin{remark}[存在証明と実計算の違い]\label{gcd:rem:limits}
この定理は幾何学的な可換模型の存在を\emph{仮定して}、その方程式の定義体を改善する。
任意の非可換圏で次元と積分条件を検査するだけでは圏同値は得られない。
また、$V,U,\mu$の有限性の証明は、その非可換入力での具体行列を計算したという意味ではない。
特殊化の範囲は第\ref{gdr:subsec:scope}節(5)を参照する。
\end{remark}

\subsection{積分空間と基礎体上の方程式の曖昧さ}

積分の線形解空間の次元には、Sebastiani--Thom分解による意味がある。
次の命題は補題\ref{ej:lem:symmetrizer-algebra}と
Hwang\cite[Proposition~2.6, Corollary~2.7]{HwangSymmetrizer}のfiber記述を
基礎体上で用いるために書き直したものである。

\begin{proposition}[symmetrizerによる積分空間の記述]
\label{gcd:prop:integrability-symmetrizer}
$k$を標数零の体、$V$を$N$次元$k$-線形空間とし、
$F\in\operatorname{Sym}^3_kV$は幾何学的に滑らかな射影超曲面を定めるとする。
その勾配空間を$Q\subset\operatorname{Sym}^2V$とし、
\[
 \mathcal I_Q=\{G\in\operatorname{Sym}^3V:\partial_iG\in Q\text{ for all }i\}
\]
と置く。ここで$x_i$は$V$の基底であり、この条件は基底に依らない。
$F(x)=T(x,x,x)$を$V^*$上の対称三重線形形式で表し、
補題\ref{ej:lem:symmetrizer-algebra}の代数を$E_F\subset\operatorname{End}_k(V^*)$とする。
このとき
\begin{equation}\label{gcd:eq:integrability-symmetrizer}
 E_F\xrightarrow{\sim}\mathcal I_Q,\qquad
 a\longmapsto F^a,\qquad F^a(x)=T(ax,x,x)
\end{equation}
は$k$-線形同型である。
式\eqref{gcd:eq:integrability}の$\mathcal L_Q$も$\mathcal I_Q$と線形同型で、
\[
 \dim_k\mathcal L_Q=\dim_kE_F
 =\text{$F_{\overline k}$の最細Sebastiani--Thom分解の成分数}
\]
となる。

勾配空間がちょうど$Q$である三次式全体は$E_F^\times$に対応し、全て幾何学的に滑らかである。
この集合に対し、$b\in E_F^\times$による座標変換$x\mapsto bx$と
$k^\times$による方程式の定数倍だけを許すと、その軌道集合は
\begin{equation}\label{gcd:eq:symmetrizer-form-quotient}
 E_F^\times/\bigl(k^\times(E_F^\times)^3\bigr)
\end{equation}
である。これは全ての射影線形座標変換による同型類の分類ではなく、
圏同値核の降下障害を分類するものでもない。
\end{proposition}
\begin{proof}
$\partial_iF$は$Q$の基底なので、$G\in\mathcal I_Q$に対して
$\nabla G=A\nabla F$となる行列$A$が一意に存在する。
対称三重線形形式へ偏極すると、$G$の対称性は$a=A^{\mathsf t}\in E_F$と同値であり、
$G=F^a$となる。逆に$a\in E_F$なら$F^a$は対称三次式で
$\nabla F^a=a^{\mathsf t}\nabla F$を満たす。これで式\eqref{gcd:eq:integrability-symmetrizer}を得る。
任意の基底$q_a$について、$G\mapsto A$を$\nabla G=A(q_a)_a$で定めると、
可積分性と式\eqref{gcd:eq:model}により$\mathcal I_Q\simeq\mathcal L_Q$である。
次元と分解の関係は補題\ref{ej:lem:symmetrizer-algebra}から従う。

勾配が$Q$の基底となる条件は$a$の可逆性である。
この場合、勾配の共通射影零点は$F$と同じく存在しないので$F^a$も幾何学的に滑らかである。
$a,b\in E_F^\times$に対して、symmetrizer条件と可換性から
\[
 F^a(bx)=T(abx,bx,bx)=T(ab^3x,x,x)=F^{ab^3}(x)
\]
である。従って座標変換と定数倍は$a\mapsto c\,ab^3$、$c\in k^\times$として作用し、
その軌道は式\eqref{gcd:eq:symmetrizer-form-quotient}で与えられる。
\end{proof}

\begin{corollary}[幾何学的に直和分解不能な場合]
\label{gcd:cor:indecomposable-equation}
前命題で$F_{\overline k}$が非自明なSebastiani--Thom分解を持たなければ、
$E_F=k$、$\dim_k\mathcal L_Q=1$である。
従って、同じ勾配空間$Q$を持つ$k$上の三次式は全て$F$の$k^\times$倍である。
\end{corollary}
\begin{proof}
補題\ref{ej:lem:symmetrizer-algebra}と
命題\ref{gcd:prop:integrability-symmetrizer}を適用する。
\end{proof}

この系は、固定した勾配空間から積分した方程式の一意性を述べる。
定理\ref{gcd:thm:field-model}の圏同値が$k$上へ降下することまでは意味しない。

\begin{example}[同じ勾配空間を持つ異なる基礎体上のform]
\label{gcd:ex:forms}
$R=\mathbb C[h]_{(h)}$、$k=\mathbb C(h)$、$L=\mathbb C((h))$、$N\ge5$とし、
\begin{equation}\label{gcd:eq:forms}
 F=\sum_{i=1}^Nx_i^3,\qquad
 G=(1+h)x_1^3+\sum_{i=2}^Nx_i^3
\end{equation}
と置く。両者の勾配空間は$\langle x_1^2,\ldots,x_N^2\rangle$であり、
$R$上で滑らかな射影超曲面族を定め、特殊ファイバーは一致する。
二項級数$(1+h)^{1/3}\in\mathbb C[[h]]^\times$で$x_1$を拡大すれば
$F$を$G$へ移せるので、完備底上、従って$L$上では射影同型である。
しかし、$k$上では射影同型でない。
実際、両者のHessian行列式の零集合は座標超平面の和である。
射影線形同型があればこの配置を保つので、その行列は置換行列と対角行列の積となる。
すると$1+h$と$1$の係数比は$k$内の三乗でなければならないが、
$h+1$に沿う離散付値が$1$であることに矛盾する。
従って勾配空間の復元だけから$k$上の射影同型類の一意性を結論できず、
完備底上の同型から$k$上の同型を結論することもできない。
この例では$E_F\simeq k^N$であり、$G$はその単元$(1+h,1,\ldots,1)$に対応する。
式\eqref{gcd:eq:symmetrizer-form-quotient}での非自明性に加え、上のHessianの議論により
全ての射影線形座標変換を許しても両者は同型でないことを確認している。
これは二つの方程式のformの例であり、圏同値核の降下障害が非零であることを
証明する例ではない。
\end{example}

\subsection{既知の生成定理による最短経路と一様な上界}

\begin{proposition}[三shiftの定量的生成：BFKの帰結]
\label{gcd:prop:bound}
$k$は標数零の代数閉体、$Y\subset\mathbb P^{N-1}_k$は滑らかなcubic、$N\ge5$とする。
任意の非零$E\in\operatorname{Ku}(Y)$について
\begin{equation}\label{gcd:eq:bound}
 \operatorname U\bigl(E\oplus\rho_YE\oplus\rho_Y^2E\bigr)
 \le 2N(N+1)-1
\end{equation}
となる。ここで$\operatorname U$は、有限直和・shift・直和因子を段数零として、
圏の全対象を得るためのcone段数の最小値である。
\end{proposition}
\begin{proof}
Orlov比較により$\rho_Y$はgraded singularity categoryの内部shiftとなる。
Ballard--Favero--Katzarkov\cite[Theorem~5.15]{GCDBFK}は、
$N$変数の次数$d$の孤立斉次超曲面について、任意の非零対象の$d$個のshiftの直和が
生成時間$2N\bigl(N(d-2)+1\bigr)-1$以下のstrong generatorとなることを証明している。
$d=3$を代入すればよい。滑らかな射影cubicのアフィン錐の特異点は原点だけである。
\end{proof}

% !TEX root = ../../nc-cubic.tex
\section{非可換Hesse族の二次jetと全次数の有理復元}
\label{gre:sec:example}

本節では、具体的な非可換中心三次対を入力として、
定理\ref{cft:thm:complete}が与える復元potentialのHesse係数を全次数で決定し、
関数体上の可換三次三様体を明示する。
中心元の公式は古典的Sklyanin理論による。
対称化による積、Kontsevich積への厳密同型、formalityの中心写像を比較し、
中心元の対称化表示から復元potentialへの二次補正を求める。
さらに有限対称群を用いて全形式次数のHesse表示を構成し、
三変数の中心商と楕円曲線の相対比較によって、その係数を有理式と同定する。

\subsection{代数的入力と全次数の有限自由性}
$R=\mathbb C[h]_{(h)}$とし、$D_R$を三生成元$x,y,z$と関係式
\begin{equation}\label{gre:eq:relations}
 \begin{aligned}
 [y,z]&=h x^2+\frac{h^2}{2}(yz+zy),\\
 [z,x]&=h y^2+\frac{h^2}{2}(zx+xz),\\
 [x,y]&=h z^2+\frac{h^2}{2}(xy+yx)
 \end{aligned}
\end{equation}
で定め、$A_R=D_R[u,v]$と置く。$u,v$は中心の次数一元である。
以下の三次元と次数零の係数を用いる：
\begin{equation}\label{gre:eq:central-coefficients}
 \begin{aligned}
 S&=x^3+y^3+z^3,&
 T&=\frac{xyz+yzx+zxy+yxz+xzy+zyx}{6},\\
 u_h&=1+\frac{h^4}{2}+\frac{h^5}{6}-\frac{h^8}{48},&
 v_h&=\frac{(4-h^4)(12h+h^5-8h^2)}{48},\\
 \Omega_h&=\frac{u_h}{3}S+v_hT+\frac{u^3+v^3}{3}.&&
 \end{aligned}
\end{equation}
係数$u_h,v_h$と中心変数$u,v$は別の記号である。

\begin{proposition}[主定理を適用できる非可換な三次三様体の族]
\label{gre:prop:input}
$A_R$の各次数成分は有限自由で、$A_R/hA_R=\mathbb C[x,y,z,u,v]$である。
$\Omega_h$は中心三次元であり、その剰余
\begin{equation}\label{gre:eq:fermat}
 F=\frac{x^3+y^3+z^3+u^3+v^3}{3}
\end{equation}
は滑らかな三次三様体を定める。
$A_K$および$A_K/(\Omega_h)$は非可換である。
\end{proposition}
\begin{proof}
$a=1-h^2/2$、$b=-1-h^2/2$、$c=-h$とすると、
式\eqref{gre:eq:relations}は$ayz+bzy+cx^2=0$とその巡回置換である。
$abc=h(1-h^4/4)$であり、
\begin{equation}\label{gre:eq:sklyanin-discriminant}
 (a^3+b^3+c^3)^3-(3abc)^3
 =-27h^3\left\{
 \left(h+\frac{h^2}{3}\left(1+\frac{h^3}{4}\right)\right)^3
 +\left(1-\frac{h^4}{4}\right)^3\right\}\ne0.
\end{equation}
従って$\overline K$上の一般ファイバーは滑らかな楕円型の三次元Sklyanin代数であり、
そのHilbert級数は$(1-z)^{-3}$である
\cite[Section~1, pp.~19--20]{GRSklyaninST}。
基礎体の忠実平坦拡大は各次数の次元を保つので、$K$上でも同じである。
一方、剰余は三変数多項式環であり、同じHilbert級数を持つ。

ここから\emph{全ての}次数の自由性が従う。
固定した$j$について$(D_R)_j$は有限生成$R$-加群である。
DVR上の構造定理により、その自由部分の階数は一般ファイバーの次元
$\binom{j+2}{2}$、非零torsion巡回成分の個数は
$\dim_{\mathbb C}((D_R)_j/h(D_R)_j)-\binom{j+2}{2}$である。
後者は零なので$(D_R)_j$は自由である。
中心多項式変数$u,v$を加えても各次数の有限自由性は保たれる。
有限個のHilbert関数の検算からPBW性を推論しているのではない。

$C=xyz+yzx+zxy$、$D=yxz+xzy+zyx$と置く。
既知のSklyanin中心三次元は
\begin{equation}\label{gre:eq:known-center}
 G=a(b^3-c^3)C+b(c^3-a^3)D+c(a^3-b^3)S
\end{equation}
である\cite[Section~2, p.~3]{GRSklyaninDLL}。
これは\cite[p.~35, Remark~(1)]{GRSklyaninST}の中心元の巡回和でもある。
関係式の巡回和から$C-D=hS+(h^2/2)(C+D)$となるので、係数を代入すると
\begin{equation}\label{gre:eq:center-identity}
 G=-9h\left(\frac{u_h}{3}S+v_hT\right)
\end{equation}
を得る。この等式は多項式恒等式である。
一般ファイバーでは$h$が可逆であり、$G$の中心性から括弧内は中心となる。
次数四成分の$R$-自由性により交換子の消滅は$R$へ戻る。
$u,v$は中心なので$\Omega_h$も中心である。
剰余は式\eqref{gre:eq:fermat}であり、五つの偏微分の共通零点は原点だけである。

最後に、$h^{-1}[y,z]=x^2+(h/2)(yz+zy)$の剰余は$x^2\ne0$である。
次数二の自由性から$[y,z]$は$A_K$で非零である。
中心三次商は次数二に関係を加えないので、その商でも非零である。
\end{proof}

\subsection{中心写像を経由した二次jetと非零の変形類}
$\widehat R=\mathbb C[[h]]$上で、Weyl対称化による次数別加群同型
\[
 \sigma:\widehat R[x,y,z,u,v]\longrightarrow A_{\widehat R}
\]
を使う。これは各次数の有限自由加群の間の写像で、剰余が恒等なので同型である。

\begin{proposition}[復元potentialの二次jet]
\label{gre:prop:jet}
Kontsevichのharmonic-angle formalityを固定し、
定理\ref{cft:thm:complete}の厳密同型$\phi$が
次数一生成元$x,y,z,u,v$を固定するように選ぶ。
式\eqref{gre:eq:relations}、\eqref{gre:eq:central-coefficients}から得られるpotentialは
\begin{equation}\label{gre:eq:recovered-jet}
 W_{\rm rec}=F+hxyz+\frac{h^2}{3}xyz+O(h^3)
 \quad\text{in }\widehat R[x,y,z,u,v]_3.
\end{equation}
その一次項は$\operatorname{Jac}(F)_3$で非零であり、
恒等を剰余とする射影座標変換と方程式の単元倍では消せない。
\end{proposition}
\begin{proof}
対称化で戻した積を$m_h$とする。
式\eqref{gre:eq:relations}を$h^2$で切断すると、
$h$を$-h$へ送りwordを逆順にする操作は反対代数との同型となる。
Weyl対称化はwordの逆順に不変なので、$m_h=m_0+hm_1+O(h^2)$の
$m_1$は反対称である。結合則の一次項より$m_1$はHochschild二cocycleであり、
可換代数上の反対称二cocycleはbiderivationである。
生成元上の交換子を使えば
\begin{equation}\label{gre:eq:first-product}
 m_h(f,g)=fg+\frac h2\{f,g\}_0+O(h^2),\qquad
 \{y,z\}_0=x^2,\quad\{z,x\}_0=y^2,\quad\{x,y\}_0=z^2,
\end{equation}
となり、$u,v$との括弧は零である。
この括弧は三変数Fermat三次式のJacobian Poisson括弧である。

Kontsevich積の一次項も式\eqref{gre:eq:first-product}に一致する規約を取る。
補題\ref{cft:lem:model}の逐次持上げを、$h^2$を法とする恒等写像から始めれば、
厳密同型$\phi$を$\phi\sigma=1+O(h^2)$となるように選べる。
その次数一成分は$1+O(h^2)$の可逆行列である。
Kontsevich formalityのaffine線形共変性により、標的の座標をこの行列の逆で
取り替えれば、$\phi\sigma=1+O(h^2)$を保って生成元を固定できる。

formal Poisson tensorを$\pi=h\pi_0+h^2\pi_1+O(h^3)$とし、
その行列を$P_h=hP_0+h^2P_1+O(h^3)$と書く。
添字は$x,y,z,u,v$を走り、重複添字について和を取る。
星積の二次項を
\[
 f*g=fg+\frac h2P_0^{ij}f_i g_j
 +h^2\left(\frac12P_1^{ij}f_i g_j+B_2(P_0;f,g)\right)+O(h^3)
\]
と表す。$P_0^{yz}=x^2$、$P_0^{zx}=y^2$、$P_0^{xy}=z^2$であり、
$u,v$を含む成分は零である。
harmonic-angleの全二次項は
\cite[Section~2, equation~(3)]{GRDito}で$\pi=P_0/2$と置いた
\begin{equation}\label{gre:eq:full-second-product}
 \begin{aligned}
 B_2(P_0;f,g)
 ={}&\frac18P_0^{ij}P_0^{kl}f_{ik}g_{jl}\\
 &+\frac1{12}P_0^{ij}(\partial_iP_0^{kl})
       (f_{jl}g_k+f_k g_{jl})
 -\frac1{24}\Theta^{ik}f_i g_k,\\
 \Theta^{ik}={}&(\partial_lP_0^{ij})(\partial_jP_0^{kl})
 \end{aligned}
\end{equation}
である。同文献のこの式は一般Poisson構造についてのものであり、
Lie代数の双対への限定は後のSection~3.1で導入される。
最後の二cycle（eye）項を含む式を用いる。
\cite[Section~1.4.2]{PKontsevich}にあるeye無しの二次で結合的な式は、
この固定したformalityの積とはgaugeが異なる。
$B_2$は対称なので、生成元の関係式を比較すれば
$P_1^{yz}=yz$、$P_1^{zx}=zx$、$P_1^{xy}=xy$となる。

生成元を固定する$\phi$の下では、線形形式$X=a_i x_i$のWeyl対称化は
$\phi\sigma(X^r)=X^{*r}$へ移る。
式\eqref{gre:eq:full-second-product}を$(X^{r-1},X)$に代入すると、
初めの二項は反対称性から消え、eye項だけが残る。
$P_0,P_1$の括弧も零なので、$r$に関する帰納法から
\[
 X^{*r}=X^r-\frac{h^2r(r-1)}{48}
 X^{r-2}\Theta^{ik}a_i a_k+O(h^3)
\]
を得る。polarizationにより
\begin{equation}\label{gre:eq:weyl-second-gauge}
 \phi\sigma=1+h^2\rho_2+O(h^3),\qquad
 \rho_2=-\frac1{48}\Theta^{ik}\partial_i\partial_k.
\end{equation}
本例では
\begin{equation}\label{gre:eq:theta-fermat}
 \begin{aligned}
 (\Theta^{ik})&=\operatorname{diag}(-8yz,-8zx,-8xy,0,0),\\
 \rho_2&=\frac16\left(yz\partial_x^2+zx\partial_y^2+xy\partial_z^2\right),
 &\rho_2F&=xyz.
 \end{aligned}
\end{equation}

次に中心写像$C_\pi$を比較する。
任意の二重vector$q$について、$\mathcal U_2(f,q)$のグラフには
関数頂点一個と二重vector頂点一個があり、出力arityは零である。
二本の辺は両方とも関数頂点へ向かわねばならないが、
同じ向きの多重辺は許されない（許しても反対称tensorと対称二階微分の縮約が零になる）。
従って$\mathcal U_2(f,q)=0$であり、$C_\pi$の一次項と、
Poisson tensorの二次係数を$\mathcal U_2$へ代入する項はともに零となる。
残る二次項$\mathcal U_3(f,\pi_0,\pi_0)/2$では、
各二重vector頂点の二本の辺は他方の二重vector頂点と関数頂点へ一本ずつ向かう。
寄与し得るグラフは、spokeが関数頂点へ向くinner二輪だけである。
そのharmonic-angle重みは零である
\cite[Section~1.3, Theorem]{GRShoikhet}。
従って
\begin{equation}\label{gre:eq:center-second-vanishing}
 C_\pi=1+O(h^3).
\end{equation}
非線形Poisson構造の中心写像の全高次項が零だとは述べていない。

式\eqref{gre:eq:central-coefficients}は
$\sigma^{-1}(\Omega_h)=F+hxyz-2h^2xyz/3+O(h^3)$を与える。
補題\ref{cft:lem:center}の実際の中心逆像
$W_{\rm rec}=C_\pi^{-1}\phi(\Omega_h)$を取ると、
式\eqref{gre:eq:weyl-second-gauge}と\eqref{gre:eq:center-second-vanishing}から
\[
 W_{\rm rec}=F+hxyz+h^2\left(-\frac23xyz+\rho_2F\right)+O(h^3)
\]
となり、式\eqref{gre:eq:recovered-jet}を得る。

最後に、$F$のJacobianイデアルは$(x^2,y^2,z^2,u^2,v^2)$である。
次数三の$xyz$はこのイデアルに入らない。
一次の線形座標変換が加える項は$\sum\ell_i\partial_iF$であり、
方程式の一次の単元倍が加える$F$もEuler恒等式で同じイデアルに属する。
従って上の一次変形類は非零である。
\end{proof}

eye項をgaugeで除くときは、中心写像も同時に変更する。
実際、$\rho=1+h^2\rho_2+O(h^3)$でWeyl積へ戻すと、
中心写像は$\rho^{-1}C_\pi=1-h^2\rho_2+O(h^3)$となる。
eye項と中心写像の補正を両方とも省くと、異なるgaugeを混ぜることになる。
式\eqref{gre:eq:recovered-jet}の二次係数は、指定した生成元・formality・
底のパラメータ$h$に対するものである。

\subsection{有限対称群による全次数のHesse表示}
本小節では$P=xyz$、$S=x^3+y^3+z^3$と置く。
以下の座標変換は内部次数一を保ち、底のパラメータ$h$を固定する。

\begin{proposition}[復元potentialの全次数Hesse表示]
\label{gre:prop:hesse-normalform}
式\eqref{gre:eq:relations}、\eqref{gre:eq:central-coefficients}の入力に対し、
harmonic-angle formalityによる定理\ref{cft:thm:complete}の比較を選び、
さらに剰余が恒等である$\mathbb C[[h]]$上の次数付き線形座標変換を施すことにより、
復元potentialを
\begin{equation}\label{gre:eq:all-order-hesse}
 \begin{aligned}
 W_{\rm H}&=F+\lambda(h)P,\qquad \lambda(h)\in h\mathbb C[[h]],\\
 \lambda(h)&=h+\frac{h^2}{3}+O(h^3)
 \end{aligned}
\end{equation}
と書ける。この座標変換は$u,v$と周期変数$t$を固定する。
従ってMF比較の内部shiftと指定された三次周期を保つ。
固定した有限対称群不変なformality比較の有限jetは、
そのjetを保って高次へ延長できる。
\end{proposition}
\begin{proof}
$\omega$を原始三乗根とし、三変数上の線形変換
\[
 c(x,y,z)=(y,z,x),\qquad
 d(x,y,z)=(x,\omega y,\omega^2z)
\]
が生成する群を$H$とする。
これは位数27の有限群であり、$h$を固定して$D_{\widehat R}$に作用する。
$c$は三つの関係式を巡回させ、$d$は各関係式の両辺を同じ三乗根倍にする。
$S$と$P$、従って指定されたactiveな中心三次元も$H$不変である。
Weyl対称化も$H$同変である。

補題\ref{cft:lem:model}の構成を、まず三変数の$D_{\widehat R}$について行う。
固定したKontsevich formalityは線形変換に同変である。
有限群の不変部分を取る操作は標数零で完全なので、
HKRとformalityの線形項を$H$不変な内部weight零の複体へ制限しても準同型になる。
従って同補題のArtin次数ごとの持上げにおいて、
Poisson tensorと厳密同型$\phi$を$H$同変に選べる。
各平方零の持上げ段階に現れる線形な補正方程式に
Reynolds平均を適用すればよい。
これは完成した非線形解そのものを平均する操作ではなく、
既に選んだ低次jetを保つ補正の選択である。
さらに$\phi$の次数一成分は$H$同変な可逆行列なので、
その逆による線形座標変換も$H$同変である。
formalityの線形共変性により、この座標変換で次数一成分を打ち消せば、
$H$同変性と生成元を固定する規格化を同時に実現できる。
既に生成元を固定する有限jetを延長するときは、
この行列はそのjetで恒等なので、逆座標変換も既存のjetを変更しない。

中心写像$C_\pi$も$H$同変であり、その三次成分上の逆も同変である。
従ってactiveな復元potentialは$H$不変な三次式となる。
$d$不変の三次単項式は$x^3,y^3,z^3,xyz$であり、
さらに$c$不変性を課すと
\[
 \mathbb C[x,y,z]_3^H=\mathbb C S\oplus\mathbb C P.
\]
よってある$A(h),B(h)\in\mathbb C[[h]]$によって
$W_{\rm act}=A(h)S/3+B(h)P$と書ける。

この比較を中心多項式変数$u,v$へ自明に延長する。
全てのpolyvectorと微分作用素は$x,y,z$方向だけを用い、
$\phi$は$\mathbb C[[h]][u,v]$上線形とする。
正規化された中心写像は$\mathbb C[[h]][u,v]$を恒等に保つので、
全potentialは
\[
 W_{\rm rec}=A(h)S/3+B(h)P+(u^3+v^3)/3
\]
となる。中心変数だけの関数に対する非恒等なグラフの項は、
関数入力へのactive方向の微分が零となるため消滅する。
関数頂点が孤立するグラフの重みも正規化により零である。
この$H$同変で生成元を固定する比較に命題\ref{gre:prop:jet}を適用すると、
\[
 A(h)=1+O(h^3),\qquad B(h)=h+\frac{h^2}{3}+O(h^3)
\]
を得る。

$A(0)=1$より、剰余が一である形式三乗根$A(h)^{-1/3}$が一意に存在する。
$x,y,z$だけをこの係数で一様に拡大縮小すると、
上のpotentialは厳密に式\eqref{gre:eq:all-order-hesse}へ移り、
$\lambda(h)=B(h)/A(h)$となる。
方程式の単元倍や$h$の取り替えは用いていない。
この座標変換は$h^3$を法として恒等であり、二次jetを保つ。
potentialを保つ次数付き代数の厳密同型なので、
MFの有限rank、内部shift、$t$に由来する指定周期を保つ。
\end{proof}

\begin{proposition}[底を固定したHesse表示の一意性]
\label{gre:prop:marked-hesse-uniqueness}
$\lambda,\mu\in h\mathbb C[[h]]$とする。
五変数の線形変換$L(h)\in\operatorname{GL}_5(\mathbb C[[h]])$と
単元$c(h)\in\mathbb C[[h]]^\times$が
\[
 L(h)\equiv I\pmod h,\qquad c(h)\equiv1\pmod h,\qquad
 F+\lambda P=c(h)(F+\mu P)\circ L(h)
\]
を満たすなら、$\lambda=\mu$である。
さらに$c(h)^{1/3}L(h)=I$となる。
\end{proposition}
\begin{proof}
$c(0)=1$の形式三乗根を選び、$\widetilde L=c^{1/3}L$と置く。
三次斉次性により
$F+\lambda P=(F+\mu P)\circ\widetilde L$である。
$\widetilde L-I$または$\lambda-\mu$が非零なら、
両者の全係数の$h$進位数の最小値を$q\ge1$とする。
$\widetilde L=I+h^qM+O(h^{q+1})$、
$\lambda-\mu=h^qb+O(h^{q+1})$と書けば、少なくとも一方の$M,b$は非零である。
$\mu\in h\mathbb C[[h]]$なので、$\mu(P\circ\widetilde L-P)$の位数は$q+1$以上である。
従って$h^q$の係数を比較すると、$x_1,x_2,x_3,x_4,x_5=x,y,z,u,v$として
\[
 bxyz=\sum_{i,j=1}^{5}M_{ij}x_i^2x_j
\]
を得る。
右辺の25単項式は全て相異なり、いずれも平方因子を持つ。
従って左辺の$xyz$と線形独立であり、$M=0$かつ$b=0$となって矛盾する。
\end{proof}

\begin{remark}[一意性とパラメータの区別]
前命題は、$h$を固定し、剰余が恒等である射影座標同値についての主張である。
異なるformality写像や比較の全選択が必ずこの座標同値を誘導する、と主張するものではない。
また、$\lambda'(0)=1$なので底の形式パラメータを取り替えることまで許せば
$\widetilde h=\lambda(h)$とできるが、本稿の表示ではこの操作を行わない。
有限群不変性はHesse表示を与えるが、それだけでは係数は決まらない。
二次jetは前小節で星積と中心写像から直接計算し、
以下では点曲線との比較によって全次数の式を得る。
\end{remark}

\subsection{指定された中心商と相対楕円曲線}
activeな中心元とその商を
\[
 \Omega_D=\frac{u_h}{3}S+v_hT,\qquad
 B_R^{\rm act}=D_R/(\Omega_D)
\]
と置く。三次元Sklyanin代数の中心商を点曲線のtwisted homogeneous coordinate ringで
記述する結果は既知である\cite[Theorem~6.8]{GRSklyaninATV}\cite[p.~20]{GRSklyaninST}。
ここでは、指定した$\Omega_D$について、その記述を可換な特殊ファイバーまで延長する。

\begin{proposition}[中心商の相対的な点曲線表示]
\label{gre:prop:elliptic-quotient}
次の有理式は$R=\mathbb C[h]_{(h)}$に属する：
\begin{equation}\label{gre:eq:point-parameter}
 \lambda_{\rm pt}(h)
 =-\frac{a^3+b^3+c^3}{3abc}
 =\frac{h+h^2/3+h^5/12}{1-h^4/4}.
\end{equation}
$E_R=V(S+3\lambda_{\rm pt}xyz)\subset\mathbb P^2_R$は滑らかな楕円曲線族であり、
一般点のSklyanin点曲線の平坦閉包である。
$e=(1:-1:0)$を原点とすると、$p=(a:b:c)$は$E_R$の切断で$p\bmod h=e$となる。
$\tau=t_p$、$\mathcal L=\mathcal O_{E_R}(1)$、
$\mathcal L_n=\bigotimes_{j=0}^{n-1}\tau^{j*}\mathcal L$と置けば、
次数一の座標切断から
\begin{equation}\label{gre:eq:relative-thcr}
 B_R^{\rm act}\simeq
 \mathsf T_R:=\bigoplus_{n\ge0}H^0(E_R,\mathcal L_n),
 \qquad
 \operatorname{qgr}B_R^{\rm act}\simeq\operatorname{coh}E_R
\end{equation}
を得る。右辺の積は$s\cdot t=s\otimes\tau^{m*}t$（$\deg s=m$）である。
これらの構成は$R\to\widehat R$の底変更と両立する。
\end{proposition}
\begin{proof}
$a^3+b^3+c^3=-3h^2-h^3-h^6/4$、$abc=h(1-h^4/4)$なので
式\eqref{gre:eq:point-parameter}を得る。
$1+\lambda_{\rm pt}^3$は単元であり、三つの偏微分を用いれば$E_R$の滑らかさが分かる。
その方程式は$x$についてmonicな三次式で、商の各次数は有限自由であるから、
一般点の点曲線の平坦閉包となる。
$S(a,b,c)+3\lambda_{\rm pt}abc=0$より$p$は切断であり、剰余は$e$である。
滑らかな種数一曲線族の群法により$\tau$は$R$上の自己同型として定まり、
$\tau\bmod h=\mathrm{Id}$となる。
一般点でこれがSklyaninの点自己同型であることは
\cite[Section~1, pp.~37--38]{GRSklyaninATV}による。

$n>0$に対して$\mathcal L_n$は各幾何ファイバー上で次数$3n$を持つ。
従って$H^1$は零、$H^0$の次元は$3n$であり、cohomologyと底変更により
$(\mathsf T_R)_n$は階数$3n$の有限自由$R$-加群である。
次数零は$R$である。
$x,y,z\in H^0(E_R,\mathcal L)$から自由代数$R\langle x,y,z\rangle\to\mathsf T_R$を作る。
一般点のSklyanin理論は$\overline K$上で適用でき、
同じ座標切断から定まる写像について忠実平坦性で$K$へ戻せる。
従って一般点でSklyaninの三つの関係式は零に写り、
$(\mathsf T_R)_2$がtorsionfreeなので、$R$上でも零に写る。
よって$D_R\to\mathsf T_R$が定まる。
式\eqref{gre:eq:center-identity}の$G=-9h\Omega_D$と、
$G$が自然な商写像の核生成元の三倍であることから、
$\Omega_D$の像は一般点で零である。
ここで核生成元は\cite[Sections~1--2, pp.~1, 3]{GRSklyaninDLL}の$c_3$であり、
同文献Section~2の$3c_3$の表示が式\eqref{gre:eq:known-center}である。
これは\cite[pp.~20, 35, Remark~(1)]{GRSklyaninST}の$g$とも一致する。
単に非零の中心三次元であることだけから商核への所属を推論してはいない。
$(\mathsf T_R)_3$のtorsionfree性により、その像は$R$上でも零となる。
従って、規格化された中心元による商$B_R^{\rm act}\to\mathsf T_R$を得る。
ここで$h$の反転は一般点での消滅を調べる際にだけ用いている。

$\Omega_D\bmod h=S/3$の乗法は単射である。
命題\ref{lp:prop:gorenstein}の単元小行列式の議論を各次数に適用すると、
$\Omega_D$は正則で、$(B_R^{\rm act})_n$は有限自由となる。
その階数は$n=0$で一、$n>0$で$3n$である。
一方、$\tau\bmod h=\mathrm{Id}$なので、上の環写像の剰余は
\[
 \mathbb C[x,y,z]/(S)\longrightarrow
 \bigoplus_{n\ge0}H^0(E_0,\mathcal O_{E_0}(n))
\]
という通常の斉次座標環の同型である。
各次数のNakayamaにより$B_R^{\rm act}\to\mathsf T_R$は全射であり、
同じ階数の有限自由加群の間の写像なので同型である。

射影商を同定するため、相対的な$\tau$-amplenessも確認する。
任意のcoherent層$\mathcal F$に対して
$\mathcal G_n=\mathcal F\otimes\mathcal L_n$と置く。
$E_R\to\operatorname{Spec}R$のcohomological dimensionは一なので、
\[
 0\to\mathcal G_n[h]\to\mathcal G_n\xrightarrow{h}h\mathcal G_n\to0,
 \qquad
 0\to h\mathcal G_n\to\mathcal G_n\to\mathcal G_n/h\mathcal G_n\to0
\]
の第一列から$H^1(\mathcal G_n)\to H^1(h\mathcal G_n)$は全射となる。
これと第二列の$H^1(h\mathcal G_n)\to H^1(\mathcal G_n)$の合成は$h$倍なので、
\begin{equation}\label{gre:eq:relative-ampleness}
 H^1(E_R,\mathcal G_n)/hH^1(E_R,\mathcal G_n)
 \simeq H^1\bigl(E_0,(\mathcal F/h\mathcal F)\otimes\mathcal O_{E_0}(n)\bigr).
\end{equation}
右辺は$n\gg0$で零となり、左辺のcohomologyは有限$R$-加群なのでNakayamaにより零となる。
これで右$\tau$-amplenessを得る。
$\tau^{-n*}\mathcal L_n$の剰余も$\mathcal O_{E_0}(n)$なので同じ議論を適用でき、
$H^q(E_R,\mathcal L_n\otimes\tau^{n*}\mathcal F)=0$（$q>0$, $n\gg0$）も従う。
従って左$\tau$-amplenessも成り立つ。
一般noetherian底上のArtin--Zhangのample pair定理
\cite[Theorem~4.5]{LP-AZ}\cite[Proposition~7.2, Lemma~7.6]{GRKeeler}が
式\eqref{gre:eq:relative-thcr}の射影商同値を与える。
左右加群・反対環の規約を取り替えても、$\tau^{-1}$を用いて同じ曲線を得る。
有限自由な各次数のcohomology、群法および環写像は平坦底変更と両立し、
$\widehat R$上でも同じ構成と証明が成立する。
\end{proof}

\subsection{三変数の$a=0$比較と曲線上での底変更}
楕円曲線との比較に用いる部分を、$N=3$に対する補題として述べる。
ここでは全coherent圏を扱い、Fanoの場合の例外列によるmutationは用いない。

\begin{lemma}[正則DVR上の三変数MF比較]
\label{gre:lem:elliptic-mf}
$R_0$を$R$または$\widehat R$とし、$D=\bigoplus_{j\ge0}D_j$を中央$R_0$-代数とする。
$D_0=R_0$、各$D_j$は有限自由、$D/hD=\mathbb C[x,y,z]$と仮定する。
中心三次元$\omega$の剰余が非零なら、$C=D/(\omega)$について
\begin{equation}\label{gre:eq:elliptic-mf}
 \operatorname{MF}^{\mathbb Z}(D,\omega)
 \simeq D^b_{\rm sg,dg}(\operatorname{gr}C)^{\rm idem}
 \simeq D^b_{\rm dg}(\operatorname{qgr}C)^{\rm idem}
\end{equation}
という$R_0$-線形dg Morita同値がある。
さらに$\operatorname{qgr}C\simeq\operatorname{coh}E$が
$R_0$上滑らかで射影的な曲線$E$との$R_0$-線形同値なら、右辺は$\operatorname{Perf}(E)$となる。
\end{lemma}
\begin{proof}
補題\ref{lp:lem:noetherian}は変数の個数に依らず、$D$の左右noetherian性を与える。
補題\ref{lp:lem:resolution}を三変数のKoszul分解と対角のKoszul分解へ適用すると、
長さ三の有限自由分解を得る。
対角分解に任意の加群をテンソルし、$\operatorname{gldim}R_0=1$を使えば
$\operatorname{gldim}D\le4$となる。
各内部次数での双対Koszul計算と、$\omega$の乗法に対する単元小行列式の議論から、
$\omega$は正則、$C_j$は有限自由であり、
\begin{equation}\label{gre:eq:elliptic-gorenstein}
 R\underline{\operatorname{Hom}}_D(R_0,D)\simeq R_0(3)[-3],
 \qquad
 R\underline{\operatorname{Hom}}_C(R_0,C)\simeq R_0[-2]
\end{equation}
を得る。第二式は$0\to D(-3)\xrightarrow{\omega}D\to C\to0$の
change of ringsによる。
同じ公式を一般の$C$-加群へ適用すれば、$C$の左右injective dimensionは高々三である。
これらは命題\ref{lp:prop:koszul}、\ref{lp:prop:gorenstein}の証明の
次元に依らない部分であり、$N\ge5$の仮定を必要としない。

従って$C$はBurke--Stevensonの意味で$a=0$のGorenstein環であり、
中央$R_0$-flatnessと\cite[Lemma~6.2]{LP-BS}から$\chi$を満たす。
\cite[Theorem~6.4(3)]{LP-BS}により、
$q(R\omega_{\rm sec}(-))_{\ge0}$は
$D^b(\operatorname{qgr}C)$と$D^b_{\rm sg}(\operatorname{gr}C)$の同値を与える。
ここで$\omega_{\rm sec}$はsection functor、$q$はperfect複体による商であり、
中心元$\omega$とは別の記号である。
derived section、次数切断およびdg商は標準enhancementの間の
$R_0$-線形dg quasi-functorとして実現できる。
引用定理の完全忠実性を全てのshiftへ適用すると、Homの全cohomologyで同型となる。
本質的全射性と合わせて、第二のdg Morita同値を得る。

第一の同値には、補題\ref{tw:lem:stabilization}の
余核・Gorenstein-projective加群・有限totalizationのdg比較を用いる。
必要な仮定は、ambientの有限大域次元、中心元の正則性および商の有限injective dimensionであり、
いずれも上で確認した。
Chen\cite[Theorem~7.2]{AMChen}の三角同値をそのままdg同値と読むのではなく、
同補題のexact圏のdg商と余核関手のzigzagを用いる。
この構成も変数の個数に依存しない。
最後に$E$は正則DVR上滑らかなので正則であり、
$D^b(\operatorname{coh}E)=\operatorname{Perf}(E)$となる。
\end{proof}

\begin{theorem}[復元potentialの全次数有理式]
\label{gre:thm:exact-potential}
命題\ref{gre:prop:hesse-normalform}のHesse係数は
\begin{equation}\label{gre:eq:exact-potential}
 \lambda(h)=\lambda_{\rm pt}(h)
 =\frac{h+h^2/3+h^5/12}{1-h^4/4}.
\end{equation}
従って$h,u,v,t$を固定し、剰余が恒等である次数付き線形座標変換の後、
完備復元potentialは$W_{\rm H}=F+\lambda_{\rm pt}(h)xyz$と全次数で書ける。
\end{theorem}
\begin{proof}
命題\ref{gre:prop:hesse-normalform}では、formality比較をactiveな三変数で構成した。
従って定理\ref{cft:thm:complete}の$N=3$の場合と、
同命題の$A(h)^{-1/3}$による正規化を用いて
\[
 \operatorname{MF}^{\mathbb Z}(D_{\widehat R},\Omega_D)
 \simeq
 \operatorname{MF}^{\mathbb Z}\bigl(\widehat R[x,y,z],S/3+\lambda(h)xyz\bigr)
\]
を得る。
$E_{\lambda,\widehat R}=V(S+3\lambda(h)xyz)$は$\operatorname{Spec}\widehat R$上の
滑らかな射影曲線である。
命題\ref{gre:prop:elliptic-quotient}と補題\ref{gre:lem:elliptic-mf}を両側へ適用すると、
\begin{equation}\label{gre:eq:relative-elliptic-equivalence}
 \operatorname{Perf}(E_{\widehat R})
 \simeq\operatorname{Perf}(E_{\lambda,\widehat R})
\end{equation}
という$\widehat R$-線形dg Morita同値となる。

この同値の底変更は曲線側で行う。
$X$が$\widehat R$上滑らかで射影的な曲線なら、$L=\mathbb C((h))$について
\begin{equation}\label{gre:eq:curve-basechange}
 \operatorname{Perf}(X)\otimes_{\widehat R}^{\mathbf L}L
 \simeq\operatorname{Perf}(X_L)
\end{equation}
がMorita同値として成り立つ\cite[Section~3.4]{GRAKW}。
本件では直接にも確認できる。
perfect複体間の導来Homは平坦局所化と可換なので、制限関手は底変更後に完全忠実である。
また任意の$X_L$上のcoherent層は、noetherianスキーム$X$へcoherentに延長できる。
$X$は有限Krull次元の正則スキームなのでこの延長はperfectであり、
その制限は$\operatorname{Perf}(X_L)=D^b(\operatorname{coh}X_L)$を厚生成する。
これで式\eqref{gre:eq:curve-basechange}を得る。
従って式\eqref{gre:eq:relative-elliptic-equivalence}から$L$上の曲線の導来同値を得られる。
任意のLaurent係数MFへ完備関手の無限Taylor和を新たに評価する操作は必要ない。

両曲線は$e=(1:-1:0)$という$L$-有理点を持つ楕円曲線である。
Antieau--Krashen--Ward\cite[Theorem~2.2]{GRAKW}により$L$上同型となり、
その$j$不変量は一致する。
Hesse曲線$E_\nu=V(S+3\nu xyz)$について
\begin{equation}\label{gre:eq:hesse-j}
 j(E_\nu)=\frac{27\nu^3(8-\nu^3)^3}{(1+\nu^3)^3}
\end{equation}
である。実際、$X=z$、$Y=x-y$、$Z=x+y-\nu z$とすると方程式は
\[
 3ZY^2+4(1+\nu^3)X^3+9\nu^2X^2Z+6\nu XZ^2+Z^3=0
\]
となる。$Z=1$で$X,Y$をともに$-4(1+\nu^3)/3$倍したWeierstrass表示では
$c_4=16\nu(\nu^3-8)$、$\Delta=-4096(1+\nu^3)^3/27$となり、
$j=c_4^3/\Delta$から式\eqref{gre:eq:hesse-j}を得る。

$g(s)=27s(8-s)^3/(1+s)^3$と置けば$g(0)=0$、$g'(0)=13824\ne0$なので、
$g$は原点で形式的な合成逆を持つ。
$j(E_{\lambda,L})=j(E_{\lambda_{\rm pt},L})$は$\widehat R$内の等式でもあるから、
この逆を適用して$\lambda^3=\lambda_{\rm pt}^3$を得る。
両者は$h+O(h^2)$で始まり、その比は$1+h\widehat R$に属する。
三乗が一であるこの比は一なので、式\eqref{gre:eq:exact-potential}が従う。
この分枝の固定には一次項だけで足り、
命題\ref{gre:prop:jet}の二次係数$1/3$は全次数式に対する独立の照合にもなっている。
\end{proof}

\subsection{明示した可換三次三様体への有限降下}
\begin{corollary}[固定された有理方程式への一般点復元]
\label{gre:cor:rational-model}
式\eqref{gre:eq:relations}、\eqref{gre:eq:central-coefficients}の
$B_R=A_R/(\Omega_h)$に対して、
\begin{equation}\label{gre:eq:rational-model}
 \begin{aligned}
 Q_h={}&(4-h^4)(x^3+y^3+z^3+u^3+v^3)
       +(12h+4h^2+h^5)xyz,\\
 Y_R={}&V(Q_h)\subset\mathbb P^4_R
 \end{aligned}
\end{equation}
は滑らかな可換三次三様体の族である。
可換側の周期は、次の規格化した方程式から定める：
\begin{equation}\label{gre:eq:normalized-rational-model}
 W_h^{\rm alg}=\frac{Q_h}{3(4-h^4)}=F+\lambda_{\rm pt}(h)xyz.
\end{equation}
ある有限拡大$K'/K$上で
\begin{equation}\label{gre:eq:explicit-generic-equivalence}
 (\mathcal K^{\rm coh}_{B_{K'}},\rho_B)
 \simeq(\operatorname{Ku}(Y_{K'}),\rho_Y)
\end{equation}
となり、同値は$\Omega_h$と$W_h^{\rm alg}$に由来する指定三次周期を保つ。
\end{corollary}
\begin{proof}
$W_h^{\rm alg}$の係数の分母は$R$の単元である。
その$u,v$についての偏微分は$u^2,v^2$であり、
残る三つの偏微分の射影共通零点は$1+\lambda_{\rm pt}^3$が単元なので存在しない。
従って$Y_R$は滑らかである。
$Q_h$は$W_h^{\rm alg}$の分母を払った同じ零点を持つ方程式である。
周期の比較には式\eqref{gre:eq:normalized-rational-model}の規格化を保つ。

定理\ref{gre:thm:exact-potential}により、五変数の完備比較の標的は
$W_{\rm H}=(W_h^{\rm alg})_{\widehat R}$である。
補題\ref{rog:lem:localization}を固定した三対象へ適用し、
命題\ref{rog:prop:euler}の非零Eulerベクトルと
命題\ref{gcd:prop:bound}の生成性を使えば、$\overline L$上で
\[
 (\operatorname{Perf}(\Gamma_K\otimes_K\overline L),\rho_B)
 \simeq(\operatorname{Ku}(Y_{\overline L}),\rho_Y)
\]
を得る。内部shiftと指定周期は完備比較、正規化の座標同型およびwindowを通じて保たれる。
$\Gamma_K$はproperであり、$\operatorname{Ku}(Y_K)$には
$K$上smooth properな生成dg代数がある。
従って標的$Y_K$を固定したまま補題\ref{rog:lem:descent}を適用でき、
ある有限$K'/K$上にrotation・指定周期付きの同値が降下する。
この同値により$\Gamma_{K'}$はsmoothとなり、
補題\ref{crs:lem:descent}、命題\ref{crs:prop:residual-smooth}および
定理\ref{crs:thm:coherent}から
線束生成残余圏は全coherent残余圏に一致する。
これが式\eqref{gre:eq:explicit-generic-equivalence}である。
\end{proof}

以上により、指定した非可換入力から、原点で正則な有理係数を持つ方程式$Q_h$を得た。
中心スカラー$A(h)$やformality比較の全成分を算出することは必要とせず、
$A(h)^{-1/3}$のactive座標変換と楕円曲線の比較によってHesse比を決定した。
有限拡大$K'/K$の次数、同値核の具体的な行列、
非可換側から直接抽出する低rotation積の全行列は明示していない。
また、特殊点を含む$R$上の代数的相対圏同値は本系の結論に含めない。

付属の\texttt{computations/verify\_recovery\_example.py}は、
式\eqref{gre:eq:sklyanin-discriminant}と\eqref{gre:eq:center-identity}を
$\mathbb Q[h]$上で確認し、三つの中心交換子を$h^4$を法として、
また次数三以下の単項式間の積を$h^2$を法として厳密計算する。
さらに全二次式\eqref{gre:eq:full-second-product}を評価し、
$\Theta$、$\rho_2F$、任意の線形形式の二乗・三乗に対するWeyl補正式、
$\phi(\Omega_h)$の二次像、およびJacobian環での係数を確認する。
出力の\texttt{rational\_hesse}は、点曲線からの$\lambda_{\rm pt}$の有理式、
Weierstrass座標変換と$c_4,\Delta,j$、
$Q_h=3(4-h^4)(F+\lambda_{\rm pt}xyz)$の恒等式、
独立に計算した二次jetとの一致、および特殊ファイバーの滑らかさを厳密に確認する。
これらは有限な多項式・有理式の恒等式の検算であり、
Shoikhetの配置空間積分、全次数の自由性、formality、THCR表示、Orlov比較、
全次数のHesse係数同定、完備MF同値や一般点降下の計算機証明ではない。

% !TEX root = ../nc-cubic.tex
\section{満支持族の相対 Gorenstein 双対と反標準切断環}
\label{mbsd:sec:serre}

本節では、中心一次変数を持たない満支持表示に対し、
全 $\OO$-保存同値を調べるための相対双対と切断飽和を構成する。
中心一次変数を二つ使う補題\ref{lem:nfm-central-saturation}は適用しない。
代わりに有限な切断加群の Ext を計算する。
また、族の中では体上の Serre 式に現れない基底の直線束を保持する。
この直線束を入れた反標準切断環は、共通の grading scalar を除いた
表示スタックそのものへ降下する。

\subsection{仮定と向きの直線束}

基礎体は $\C$、$k\geq2$ とし、
\[
 N=3k,\qquad r=N-3,\qquad n=N-2
\]
と置く。$R$ は任意の可換 $\C$-代数である。
第\ref{mbs:sec:presentations}節の普遍表示を
$\Spec R$ へ引き戻したもの、または fpqc 局所的にそのようになる
実際の次数付き代数表示を考える。
各 $W_i=(S_i)_1$ は階数三、
\[
 S_i=T_R(W_i)/(R_i),\qquad
 B=S_1\otimes_R\cdots\otimes_RS_k,\qquad
 A=B/(\ell)
\]
とする。ここで $R_i$ は二次関係束、
$\ell\subset\bigoplus_i Z(S_i)_3\subset B_3$ は各中心直線への射が
同型となる直線部分束である。
従って局所的には $\ell=Rw$、
$w=\sum_i a_i\Omega_i$, $a_i\in R^\times$ と書ける。
因子の幾何学的仮定は、滑らかな楕円曲線と平行移動 $t_{p_i}$ に対する
$p_i,2p_i,3p_i\ne0$ である。
異なる $j$ 不変量という仮定は表示スタックで用いるが、本節の双対計算には不要である。

補題\ref{mbs:lem:factor-family}と定理\ref{mbs:thm:presentation-stack}により、
全ての次数成分は有限射影であり、列
\begin{equation}\label{mbsd:eq:regular-line}
 0\longrightarrow B_{a-3}\otimes_R\ell
 \longrightarrow B_a\longrightarrow A_a\longrightarrow0
\end{equation}
は任意の底変更後にも完全である。
この列は係数加群として分裂する。
本節で非被約な $R$ を許すのは、この普遍族からの引戻しとしてである。
任意の非被約底上の代数について、幾何ファイバーの次元だけから
これらの結論を推論するものではない。

次の二つの直線束を区別する。
\begin{equation}\label{mbsd:eq:orientation}
 V=\bigoplus_{i=1}^kW_i,\qquad
 \mathscr V=\det V,\qquad
 \mathscr D=\mathscr V\otimes_R\ell^\vee.
\end{equation}
$\ell$ は内部次数三の関係を表すが、式中の直線束そのものは係数加群として扱う。
以下、$M^\vee$ は内部次数を反転する係数双対
$(M^\vee)_a=\Hom_R(M_{-a},R)$ を表す。
有限射影な次数成分に対して用いるので、通常の双対と導来双対は一致する。

\begin{lemma}[一因子の potential 直線と行列式]
\label{mbsd:lem:potential-line}
各因子の三次 syzygy
\[
 P_i=(W_i\otimes_RR_i)\cap(R_i\otimes_RW_i)
       \subset W_i^{\otimes3}
\]
は直和因子となる直線束で、形成は任意の底変更と両立する。
外積への自然な射は同型
\begin{equation}\label{mbsd:eq:potential-det}
 P_i\xrightarrow{\sim}\det W_i
\end{equation}
を与える。$P_i$ は cyclic tensor からなる。
\end{lemma}
\begin{proof}
まず滑らかで被約な普遍因子チャート $\mathcal T$ 上で考える。
三次元二次 Sklyanin 代数の Koszul 解消では、三次 syzygy の次元は一である
\cite[Theorems~6.6, 6.8]{GRSklyaninATV}。
$P_i$ を核として表示する有限行列の階数は全幾何点で一定なので、
補題\ref{mbs:lem:factor-family}と同じ小行列式の議論で、
$P_i$ は直線の直和因子となり、その形成は任意底変更に両立する。

各幾何点では Hesse 座標を選び、関係を
\[
 a yz+b zy+c x^2,\quad
 a zx+b xz+c y^2,\quad
 a xy+b yx+c z^2
\]
と書ける\cite[\S1, pp.~20--21]{GRSklyaninST}。
このとき $P_i$ の生成元は
\[
 \Psi=a(xyz+yzx+zxy)+b(xzy+zyx+yxz)+c(x^3+y^3+z^3)
\]
であり、cyclic で、その外積像は $3(a-b)x\wedge y\wedge z$ である。
もし $a=b$ なら全ての二次関係の多重線型化は二つの点の交換に不変となる。
点対応は $t_{p_i}$ のグラフなので、グラフの対称性は
$t_{p_i}=t_{p_i}^{-1}$、従って $2p_i=0$ を意味する。
これは仮定に反する。よって外積像は全幾何点で非零であり、
直線束間の射 \eqref{mbsd:eq:potential-det} は同型である。
cyclic 性も、被約な普遍底上の有限局所自由束の射が全幾何点で零になることから従う。
これらの同型と恒等式を引き戻すので、非平坦・非被約な底変更にも成立する。
\end{proof}

\begin{theorem}[tensor 周囲代数の普遍有限解消]
\label{mbsd:thm:ambient-duality}
$B$ は長さ $N$ の有限射影な相対二次 Koszul 双加群解消 $K_B$ を持つ。
その第 $-j$ 項は $B\otimes_RU_j\otimes_RB$ の形で、
$U_j$ の内部次数は $j$、階数は $\binom Nj$ であり、
$U_0=R$, $U_N\simeq\mathscr V$ である。
解消とその両側の片側 augmentation は任意の底変更後にも完全である。
真正な次数付き双加群として
\begin{equation}\label{mbsd:eq:ambient-duality}
 \Theta_{B/R}:=
 \mathbf R\underline{\Hom}_{B\otimes_RB^{\mathrm{op}}}
       (B,B\otimes_RB^{\mathrm{op}})
 \simeq B(N)\otimes_R\mathscr V^\vee[-N]
\end{equation}
があり、この同型も任意底変更およびブロックの置換に両立する。
\end{theorem}
\begin{proof}
各因子には potential 複体
\[
 0\longrightarrow S_i\otimes_RP_i\otimes_RS_i
 \longrightarrow S_i\otimes_RR_i\otimes_RS_i
 \longrightarrow S_i\otimes_RW_i\otimes_RS_i
 \longrightarrow S_i\otimes_RS_i\longrightarrow S_i\longrightarrow0
\]
がある。局所的な cyclic tensor $\Psi$ に対し、第一の微分は変数の左右差、
第二の微分は cyclic 微分の非可換 Hessian、第三の微分は
$\Psi\in(W_i\otimes R_i)\cap(R_i\otimes W_i)$ の二つの表示の左右差である。
cyclic 恒等式によりこれは複体となり、幾何ファイバーでは Sklyanin の
Koszul 双加群解消である。
この複体とその自己双対性は potential 解消の標準的な構成であり
\cite[Theorem~3.6 の証明、Lemma~3.7]{NFMBeTa}、
本稿の定理\ref{thm:nfm-relative-ambient-cy}の長さ三の計算と同じである。
ただしここでは $P_i$ の生成元を固定せず、その直線を保持する。

普遍チャート上で全内部次数を固定すると、augmentation を含む各複体は
有限局所自由係数の有界複体である。
各ファイバーで完全なので、右端から順に全射を分裂させれば
全ての kernel が局所自由となり、複体全体が係数加群として局所分裂する。
従って任意の底変更後にも完全である。
片側 $S_i$-加群としては有限射影解消で、終端の $S_i$ は射影だから分裂し、
augmentation $S_i\to R$ との tensor 後も完全となる。

cyclic tensor による contraction は $R_i\simeq W_i^\vee\otimes_RP_i$ を与える。
この対応で複体を双対化すると、Hessian の転置は左右反転した Hessian となり、
左右差は符号を変える。
従って定理\ref{thm:nfm-relative-ambient-cy}の符号付き鎖写像により
\[
 K_i^\vee\simeq K_i(3)\otimes_RP_i^\vee[-3]
\]
を得る。potential の生成元の変更は両辺で $P_i^\vee$ の変換として一致する。
ここで現れる双加群の左右作用は通常の作用であり、
Nakayama 自己同型による twist はない。

$k$ 個の解消を $R$ 上で tensor して全複体を取る。
各内部次数が係数射影で普遍的に完全であることから、
これは $B$ の普遍的に完全な有限双加群解消である。
係数束の階数は $(1+t)^{3k}$ の係数となる。
各因子の二次 Koszul 複体の tensor は、異なる因子間の交換関係を含む
$B$ の二次 Koszul 複体である。
項別双対と補題\ref{mbsd:lem:potential-line}により
\eqref{mbsd:eq:ambient-duality}を得る。
ブロックを交換すると、長さ三の複体の最上項は $(-1)^{3\cdot3}=-1$ を受ける。
これは $\det(\bigoplus_iW_i)$ の二つの三次元 summand を交換する符号と一致する。
従って順序付きの積 $\bigotimes_i\det W_i$ だけで符号を忘れることなく、
$\mathscr V=\det V$ を用いた同型がブロック置換にも両立する。
\end{proof}

\subsection{有限切断の Ext と中心一次変数を使わない飽和}

\begin{theorem}[相対 Gorenstein 恒等式]
\label{mbsd:thm:gorenstein}
$M$ を、係数加群として有限射影な次数付き右 $A$-加群とする。
通常の左作用を持つ次数付き加群として、自然な同型
\begin{equation}\label{mbsd:eq:relative-gorenstein}
 \mathbf R\underline{\Hom}_A(M,A)
 \simeq M^\vee(r)\otimes_R\mathscr D^\vee[-(N-1)]
\end{equation}
がある。特に
\[
 \mathbf R\underline{\Hom}_A(R,A)
 \simeq R(r)\otimes_R\mathscr D^\vee[-(N-1)].
\]
同じ結論は左右を交換しても成立し、全ての同型は任意底変更に両立する。
体上では $A$ は左右の単射次元 $N-1$、Gorenstein parameter $r$ の
AS--Gorenstein 代数である。
ここでは AS--Gorenstein 条件を、この有限単射次元と左右の自明加群の Ext 条件の意味で用いる。
任意の $R$ 上で主張するのは \eqref{mbsd:eq:relative-gorenstein}という相対恒等式であり、
$R$ 自身の正則性・Gorenstein 性は仮定しない。
\end{theorem}
\begin{proof}
有限射影な係数を持つ $M$ に対して、有限双加群解消の項別 tensor--Hom 随伴は
\[
 \mathbf R\underline{\Hom}_B(M,B)
 \simeq\Theta_{B/R}\otimes_B^{\mathbf L}M^\vee
 \simeq M^\vee(N)\otimes_R\mathscr V^\vee[-N]
\]
を与える。$M$ に追加の可換な作用があれば、その作用も保持する。
一方、中心関係直線の二項解消
\[
 0\longrightarrow B(-3)\otimes_R\ell\longrightarrow B\longrightarrow A\longrightarrow0
\]
の双対は、通常の両側作用を持つ $A$-双加群として
\[
 \mathbf R\underline{\Hom}_B(A,B)
 \simeq A(3)\otimes_R\ell^\vee[-1]
\]
である。中心関係なので normalizing automorphism は現れない。
change of rings の随伴から
\[
 \mathbf R\underline{\Hom}_B(M,B)
 \simeq\mathbf R\underline{\Hom}_A(M,A)(3)\otimes_R\ell^\vee[-1]
\]
となる。この式を解けば \eqref{mbsd:eq:relative-gorenstein}を得る。
有限射影解消と \eqref{mbsd:eq:regular-line} の分裂により、
以上の計算は任意の係数変更と両立する。

体上では $K_B$ を任意の片側 $B$-加群と tensor すると、長さ $N$ の射影解消を得る。
従って $B$ の左右の大域次元は高々 $N$ である。
任意の $A$-加群についての同じ change of rings により、
$A$ の左右の単射次元は高々 $N-1$ となる。
$M=R$ の非零な第 $N-1$ Ext が等号を与え、AS--Gorenstein 条件も従う。
一般の係数環について、この絶対的な単射次元の主張はしない。
\end{proof}

\begin{corollary}[全整数の普遍切断飽和]
\label{mbsd:cor:saturation}
$\OO=\pi A\in\QGr A$ とする。
全ての整数 $j$ に対し、左乗法から来る標準写像
\begin{equation}\label{mbsd:eq:saturation}
 A_j\xrightarrow{\sim}\Hom_{\QGr A}(\OO,\OO(j))
\end{equation}
は同型である。負の次数では両辺とも零である。
この同型は任意の非平坦・非被約な底変更と両立する。
\end{corollary}
\begin{proof}
$M_q=A/A_{\geq q}$ は係数加群として有限射影である。
定理\ref{mbsd:thm:gorenstein}により
$\underline{\Hom}_A(M_q,A)=\underline{\Ext}^1_A(M_q,A)=0$ となる。
従って $0\to A_{\geq q}\to A\to M_q\to0$ の Hom 列は、全整数 $j$ に対し
\[
 A_j=\Hom_{\Gr A}(A,A(j))
 \xrightarrow{\sim}\Hom_{\Gr A}(A_{\geq q},A(j))
\]
を与える。
また $A$ の torsion 元は、ある $M_q$ の次数 shift からの射の像に入るので、
同じ Hom 消滅により $A$ は torsion-free である。
$A$ は次数一で有限生成だから torsion 部分圏は局所化部分圏であり、
Gabriel 商の Hom 公式は
\[
 \Hom_{\QGr A}(\OO,\OO(j))
 =\varinjlim_q\Hom_{\Gr A}(A_{\geq q},A(j))
\]
となる\cite[Chapitre~III, \S1]{NFMGabrielQuotient}。
実際、$A$ の商が torsion となる部分加群は、$1$ の像が torsion であるため、
ある $A_{\geq q}$ を含み、右辺の切断系は cofinal である。
よって標準写像が同型となる。
この証明を任意の変更後の係数環に適用でき、写像自体も底変更に自然である。
\end{proof}

\subsection{直線束を保持した相対 Serre 双対}

\begin{theorem}[普遍線コホモロジーと Serre 関手]
\label{mbsd:thm:serre}
任意の係数変更後に、全ての $\OO(i)$ は $D(\QGr A)$ の compact 対象であり、
$\bigoplus_{i=0}^{N-1}\OO(i)$ は compact generator である。
$a=j-i$ と置くと、
\begin{equation}\label{mbsd:eq:line-cohomology}
 \mathbf R\Hom_{\QGr A}(\OO(i),\OO(j))\simeq
 \begin{cases}
 A_a[0],&a\geq0,\\
 0,&-r<a<0,\\
 A_{-a-r}^\vee\otimes_R\mathscr D^\vee[-n],&a\leq-r.
 \end{cases}
\end{equation}
compact 対象間の RHom は係数 perfect であり、任意の導来底変更と両立する。
compact 部分の相対 Serre 関手は
\begin{equation}\label{mbsd:eq:serre}
 \mathsf S_A(K)=K(-r)\otimes_R\mathscr D[n]
\end{equation}
である。従って compact $K,L$ に対し、合成 pairing と両変数に自然な同型
\[
 \mathbf R\Hom_R\bigl(\mathbf R\Hom(K,L),R\bigr)
 \simeq\mathbf R\Hom\bigl(L,K(-r)\otimes_R\mathscr D[n]\bigr)
\]
がある。この同型と trace は任意底変更に両立する。
\end{theorem}
\begin{proof}
定理\ref{mbsd:thm:ambient-duality}と列\eqref{mbsd:eq:regular-line}は、
定理\ref{thm:nfm-relative-line-cohomology}および
定理\ref{thm:nfm-relative-serre}の一般部分の入力を与える。
これらの証明は、周囲代数が一ブロックであることを使わず、
普遍有限 Koszul 解消と中心正則三次の二項解消を用いている。
以下ではその構成における係数直線を明記する。

$M_q=A/A_{\geq q}$ に対する \eqref{mbsd:eq:relative-gorenstein}の colimit は、
真正な $A$-双加群として
\begin{equation}\label{mbsd:eq:local-duality}
 \Gamma:=\mathbf R\Gamma_A(A)
 \simeq A^\vee(r)\otimes_R\mathscr D^\vee[-(N-1)]
\end{equation}
を与える。左右の作用は有限切断の計算の段階から保持されている。
周囲代数上の有限幅 torsion 模型と中心商への restriction により、
定理\ref{thm:nfm-relative-line-cohomology}の証明と同様に
非有界 Gabriel 局所化三角
\[
 \mathbf R\Gamma_A(M)\longrightarrow M\longrightarrow
 \mathbf R\omega\pi M\longrightarrow
\]
が構成される。
その $M=A$、内部次数 $a$ の部分に \eqref{mbsd:eq:local-duality}を代入すれば
\eqref{mbsd:eq:line-cohomology}を得る。
周囲代数上では有限 Koszul 解消が連続する $N$ 個の twist を compact generators とし、
中心商への induction はこれらを $\OO_A(i)$ に送る。
restriction は保守的かつ直和保存なので、この像も compact generator である。
線対象間の Hom の式は係数 perfect かつ任意底変更に両立する。
compact 部分は generator の thick 閉包なので、同じ結論が全 compact 対象に延びる。

Serre pairing の作用を確認するため、線対象の dg 切断代数を $Q$ と書く。
定理\ref{thm:nfm-relative-serre}の証明で用いた導来 Gabriel 局所化により
$Q\otimes_A^{\mathbf L}Q\simeq Q$ である。
\eqref{mbsd:eq:local-duality}の双対は
\[
 \Gamma^\vee\simeq A(-r)\otimes_R\mathscr D[N-1]
\]
となる。
局所化三角 $\Gamma\to A\to Q\to$ を係数双対化して両側局所化を施すと、
$A^\vee$ は torsion なので消え、
\[
 Q^\vee\simeq Q(-r)\otimes_R\mathscr D[N-2]
\]
を真正な $Q$-双加群として得る。
perfect な次数付き $Q$-加群の tensor--Hom 評価により
\eqref{mbsd:eq:serre}と両変数に自然な双対性が従う。
trace は恒等射での評価から得られる。
有限射影模型、局所化三角、係数双対、および評価は全て任意底変更と両立するので、
pairing と trace も同じ性質を持つ。
\end{proof}

\begin{remark}[主張する範囲]
体上で volume と中心方程式の生成元を選べば $\mathscr D$ は自明化され、
Serre 式は $(-r)[n]$ となる。
族ではこの自明化を暗黙に固定しない。
また上の定理は $D(\QGr A)^c$ の properness と Serre 双対を与えるが、
$\QGr A$ の有限大域次元、tails の dg smoothness、
あるいは $D(\QGr A)^c=D^b(\qgr A)$ は結論しない。
coherent 対象と compact 対象の一致を使う場合は、
第\ref{crs:sec:saturation}節の独立な仮定を確認する必要がある。
\end{remark}

\subsection{構造対象からの反標準環と scalar の消去}

\begin{corollary}[全同値の自然自己同型]
\label{mbsd:cor:natural-scalars}
任意底変更後に $\End(\operatorname{id}_{\QGr A})=R$ である。
従って $R$-線型完全同値 $F:\QGr A\to\QGr A'$ の自然自己同型群は $R^\times$ であり、
$F(\OO)\simeq\OO'$ の同定を保つ自然自己同型は恒等だけである。
これは全 $\OO$-保存自己同値の\emph{同型類}の群が自明という意味ではない。
\end{corollary}
\begin{proof}
切断飽和により $\End(\OO(i))=R$ である。
恒等関手の自然自己準同型の $\OO(i)$ 上の成分を $b_i$ とすると、
$\Hom(\OO(i),\OO(i+1))=A_1$ への自然性から
$(b_{i+1}-b_i)A_1=0$ となる。
$A_1$ は正の一定階数の有限射影加群で忠実なので、全ての $b_i$ は一つの $b\in R$ に等しい。
任意の次数付き加群は自由加群の直和からの商であり、$\pi$ は完全だから、
全ての $\OO(i)$ は Abel 圏 $\QGr A$ の生成族である。
従って自然自己準同型は全対象上で $b$ 倍である。
同値に対する結論は逆同値を合成すれば従う。
\end{proof}

\begin{corollary}[反標準切断環の内在性と表示スタックへの降下]
\label{mbsd:cor:anticanonical}
compact 部分上の $\mathsf H_A=\mathsf S_A^{-1}[n]$ は
$(r)\otimes_R\mathscr D^\vee$ であり、これは Abel 圏全体の完全自己同値でもある。
構造対象の反標準 orbit algebra は
\begin{equation}\label{mbsd:eq:anticanonical-ring}
 \mathcal R_{\mathrm{ac}}(A)
 :=\bigoplus_{q\geq0}\Hom(\OO,\mathsf H_A^q\OO)
 \simeq\bigoplus_{q\geq0}A_{rq}\otimes_R(\mathscr D^\vee)^{\otimes q}.
\end{equation}
右加群の左乗法に合わせ、次数 $p,q$ の射 $x,y$ の積を
$\mathsf H_A^q(x)\circ y$ とする。
この規約の下で右辺の積は $A$ の通常の積である。

任意の $R$-線型完全同値 $F:\QGr A\to\QGr A'$ と
$F(\OO)\simeq\OO'$ の同定は、Serre 比較を一貫して反復することにより、
$\mathcal R_{\mathrm{ac}}(A)$ と $\mathcal R_{\mathrm{ac}}(A')$ の次数付き同型を与える。
Serre 比較の scalar 変更は、この同型の grading scalar 変更に対応する。
従って全 $\OO$-保存同値から得る次数付き同型は、grading scalar を除けば内在的である。
この構成は任意底変更に両立する。

普遍表示チャート $\mathcal X_k$ 上の右辺の代数は、
次数付き表示スタック $\mathfrak P_k$ 上へ降下する。
これは各 $A_j$ 自身の降下を要求しない。
\end{corollary}
\begin{proof}
最初の式は定理\ref{mbsd:thm:serre}と系\ref{mbsd:cor:saturation}による。
完全同値は非有界導来圏、その compact 部分、および係数線型 RHom を同定する。
Serre 双対の表現対象の一意性により
$F\mathsf S_A\simeq\mathsf S_{A'}F$ を得る。
その逆を正負に反復し、構造対象の同定を使うと、
反標準 orbit の全 Hom と合成が同定される。
二つの Serre 比較の比を各 $\OO(i)$ 上で見ると、
同値による End の同定と切断飽和から、その成分は単元 $b_i\in R^\times$ である。
隣接する twist 間の Hom への自然性と $A_1$ の忠実性により、
全ての $b_i$ は一つの単元に等しい。
係数の可逆加群との tensor にも両立するので、
反標準 orbit 上の反復は次数 $q$ にその単元の $q$ 乗を掛ける。
従って grading scalar 以外の曖昧さを生まない。

最後に降下を調べる。
$\GL_3^k\rtimes\mathfrak S_k$ への共変性は、各直線と積の構成から従う。
共通 grading scalar $t$ は $A_{rq}$ に重み $rq$、
$\mathscr V=\det V$ に重み $N$、$\ell$ に重み三で作用する。
従って $\mathscr D$ の重みは $r$ であり、
$A_{rq}\otimes(\mathscr D^\vee)^{\otimes q}$ の重みは零となる。
作用は $G_k=(\GL_3^k/\mathbf G_m^{\mathrm{diag}})\rtimes\mathfrak S_k$ を通り、
積もこの作用に両立するから、代数全体が $[\mathcal X_k/G_k]$ へ降下する。
\end{proof}

\begin{remark}[次数一認識との区別]
局所的に $\mathscr D$ を自明化すると、反標準環は $A^{(3k-3)}$ である。
従って本節は、全 $\OO$-保存同値から Veronese 環へ移るための
相対的・非被約底上の入力を供給する。
この Veronese から $A$ の次数一 shift を回復することは別問題であり、
系\ref{tbp:cor:all-blocks}の非自明な block $\mu_3$ twist は、
本節の結論と両立してなお残る。
また、本節は一般の Abel 圏族がこの表示族に fpqc 局所的に属することを主張しない。
\end{remark}

% !TEX root = ../nc-cubic.tex
\section{中心三乗根群の作用と rotation の指標}
\label{mbza:sec:central-action}

系\ref{mbrr:cor:central-band}が認識した中心群を、関手の自然同型と合成則も含めて扱う。
標準ブロック作用は、基準ブロックを選ばずに coherent な作用として降下する。
その同変化も同じ意味で基準に依存しない。
一方、次数一 shift と rotation の標準同変構造は、基準の変更で指標だけ変わる。
この差を記述する有限 torsor は、ブロック数が三の倍数でないときに限り
置換不変な切断を持つ。
共通 $\mathbf G_m$ を除いた表示 gerbe による次数一関手の降下問題は、
この有限群の線型化問題とは区別する。

作用・同変対象・指標による twist の一般論は既知である
\cite[Definitions~3.1, 3.5, \S4]{MEDElagin}。
随伴と半直交分解の同変性、作用の strictification には
\cite[Proposition~3.13, Theorems~5.4, 6.2]{MBZAShinder}を参照できる。
以下では本稿に必要な比較写像を実際に書き、その全ての合成則を確認する。

\subsection{中心拡大の全ての分裂}

基礎は $\C$ 上とし、まず順序付き満支持表示
$A=(\bigotimes_{i=1}^kS_i)/(\sum_i a_i\Omega_i)$、$k\ge2$ を固定する。
\[
 \widetilde Z=\mu_3^k,\qquad
 \Delta=\mu_3^{\mathrm{diag}},\qquad Z_0=\widetilde Z/\Delta
\]
と置く。$h=(h_i)\in\widetilde Z$ は第 $i$ ブロックの次数一空間を $h_i$ 倍する。
各 $\Omega_i$ とその係数を固定するので、これは $A$ の次数付き自己同型である。
共通スカラー $\lambda_t$ は次数 $n$ に $t^n$ を掛ける自己同型を表す。

\begin{proposition}[分裂の有限 torsor]
\label{mbza:prop:splittings}
中心拡大
\begin{equation}\label{mbza:eq:central-extension}
 1\longrightarrow\Delta\longrightarrow\widetilde Z
   \longrightarrow Z_0\longrightarrow1
\end{equation}
の群スキームとしての分裂は、有限集合
\begin{equation}\label{mbza:eq:splitting-space}
 \mathscr S_k=
 \{\boldsymbol a=(a_1,\ldots,a_k)\in\mathbf F_3^k:
                   \textstyle\sum_i a_i=1\}
\end{equation}
によって分類される。対応する分裂は
\begin{equation}\label{mbza:eq:splitting-formula}
 s_{\boldsymbol a}([h])_i
   =h_i\left(\prod_jh_j^{a_j}\right)^{-1}.
\end{equation}
$\widehat Z_0=\operatorname{Hom}(Z_0,\mu_3)$ とすると
\[
 \widehat Z_0=\{\boldsymbol d\in\mathbf F_3^k:\textstyle\sum_i d_i=0\},
 \qquad |\mathscr S_k|=3^{k-1},
\]
であり、$\mathscr S_k$ は $\widehat Z_0$ の torsor である。
二つの分裂について
\begin{equation}\label{mbza:eq:change-character}
 s_{\boldsymbol a'}(z)
   =\chi_{\boldsymbol a\boldsymbol a'}(z)s_{\boldsymbol a}(z),
 \qquad
 \chi_{\boldsymbol a\boldsymbol a'}([h])
      =\prod_i h_i^{a_i-a_i'}
\end{equation}
となる。基準ブロック $b$ を固定する分裂は $\boldsymbol a=\boldsymbol e_b$ であり、
基準を $b$ から $c$ へ変えると $\chi_{bc}([h])=h_b/h_c$ である。
\end{proposition}
\begin{proof}
三乗根の原始元を選んで群を $\mathbf F_3$-ベクトル空間と見れば、
分裂は対角直線への retraction と同値である。
その retraction は線型形式 $v\mapsto\sum_i a_iv_i$ であり、
対角上で恒等である条件が $\sum_i a_i=1$ となる。
分裂は $v\mapsto v-(\sum_i a_iv_i)(1,\ldots,1)$ であるから、
乗法表示へ戻すと式\eqref{mbza:eq:splitting-formula}を得る。
この式は $h$ を共通三乗根倍しても変わらず、原始元の選択にも依存しない。
二つの retraction の差は対角を消す線型形式であり、
これが $\widehat Z_0$ と式\eqref{mbza:eq:change-character}を与える。
\end{proof}

\subsection{共通スカラーを消す自然同型}

本稿の構造移送の規約を固定する。
次数付き自己同型 $h$ に対して $\Phi_h(M)$ の右作用は
\begin{equation}\label{mbza:eq:transport}
 m*a=m\,h^{-1}(a)
\end{equation}
であり、$\Phi_h\Phi_g=\Phi_{hg}$ である。
内部 shift は $M(j)_n=M_{n+j}$ とする。

\begin{lemma}[scalar 自明化と shift の重みの再掲]
\label{mbza:lem:scalar}
共通 grading scalar $\lambda_t$ には自然同型
\begin{equation}\label{mbza:eq:scalar-trivialization}
 \tau_t:\Phi_{\lambda_t}\xrightarrow{\sim}\operatorname{id},
 \qquad \tau_{t,M}(m)=t^n m\quad(m\in M_n)
\end{equation}
がある。この同型は $t$ の積と厳密に両立し、全ての次数付き構造移送と可換する。
torsion 商と複体にも同じ式で誘導される。一方、全整数 $j$ について
\begin{equation}\label{mbza:eq:shift-weight}
 \tau_{t,M(j)}=t^{-j}(\tau_{t,M})(j)
\end{equation}
である。
\end{lemma}
\begin{proof}
命題\ref{mbsd2:prop:scalar-weight}の自然同型と、その shift 比較を再掲したものである。
特に比較の乗法性は、関手の同型類だけでなく自然同型自体について成立する。
\end{proof}

\begin{theorem}[中心作用と同変化の coherent な降下]
\label{mbza:thm:central-action}
各 $\boldsymbol a\in\mathscr S_k$ は
\[
 \alpha^{\boldsymbol a}_z=\Phi_{s_{\boldsymbol a}(z)}
 \qquad(z\in Z_0)
\]
という厳密な $Z_0$-作用を $\mathcal C_A=\QGr A$ に与える。
異なる分裂による作用の間には、恒等関手を underlying functor とする
標準の coherent な同値があり、三つ以上の分裂の比較は cocycle 条件を満たす。
これは構造対象の標準線型化も保つ。

系\ref{mbrr:cor:stack-isomorphism}の一般開では、この作用は
内在的な有限 \'etale 群 $Z_0$ の作用として全 $\OO$-局所型へ降下する。
順序を忘れた底では $Z_0$ 自体を因子置換で降下した群スキームと解釈する。
同変対象の圏も基準ブロックによらず降下する。

標準導来 compact 部分と、定理\ref{mbsd2:thm:projector-descent}の
compact 残余圏にも、任意底変更に両立して同じ結論が成り立つ。
全 coherent 残余や MF 模型へ移す場合は、既存の window・飽和・比較の仮定を別に課す。
\end{theorem}
\begin{proof}
$\chi=\chi_{\boldsymbol a\boldsymbol a'}$ と置く。
式\eqref{mbza:eq:change-character}と補題\ref{mbza:lem:scalar}から
\begin{equation}\label{mbza:eq:action-comparison}
 \epsilon_{\boldsymbol a\boldsymbol a',z}:
 \alpha_z^{\boldsymbol a'}\xrightarrow{\sim}\alpha_z^{\boldsymbol a},
 \qquad m\in M_n\longmapsto\chi(z)^n m
\end{equation}
を得る。$\chi(zz')=\chi(z)\chi(z')$ と $\tau$ の乗法性は、
これらが群の合成と両立することを与える。
さらに
\[
 \chi_{\boldsymbol a\boldsymbol a'}\chi_{\boldsymbol a'\boldsymbol a''}
   =\chi_{\boldsymbol a\boldsymbol a''}
\]
なので、三つの分裂を比較する二つの合成は厳密に等しい。

構造対象の線型化は $\Phi_h(A)\to A$, $x\mapsto h(x)$ である。
$s_{\boldsymbol a'}=\lambda_\chi s_{\boldsymbol a}$ を用いると、
この写像は式\eqref{mbza:eq:action-comparison}と両立する。
系\ref{mbsd:cor:natural-scalars}により構造対象の同定を保つ自然自己同型はないので、
得られた比較には構造対象付き作用として余分なスカラー二コサイクルはない。

実際の次数付き同型 $f:A\to A'$ はブロックを置換し、
$f\,s_{\boldsymbol a}(z)f^{-1}=s_{\pi\boldsymbol a}(\pi z)$ を満たす。
従って対応する移送関手は上の作用と比較を保つ。
$f$ の持上げを共通 grading scalar だけ変えた場合も、
比較は $\tau_t$ によって一致し、その合成則は補題\ref{mbza:lem:scalar}で確認済みである。
相対認識により、全 $\OO$-同値は fpqc 局所でこの形である。
従ってこれらの比較は全 $\OO$-局所型に対する descent datum を与える。

具体的に、同変対象を $\theta_z:\alpha_zM\to M$ で書くと、
二つの同変圏の間の同値は
\begin{equation}\label{mbza:eq:equivariant-change}
 J_{\boldsymbol a\boldsymbol a'}:
 (\mathcal C_A,\alpha^{\boldsymbol a})^{Z_0}
   \xrightarrow{\sim}(\mathcal C_A,\alpha^{\boldsymbol a'})^{Z_0},
 \quad
 (M,\theta_z)\longmapsto
 (M,\theta_z\epsilon_{\boldsymbol a\boldsymbol a',z})
\end{equation}
である。上の計算はこの線型化の群法則と $J$ の全 cocycle を同時に証明する。

複体上では式\eqref{mbza:eq:action-comparison}は各項に作用し、微分と可換する。
従って標準 dg 局所化、torsion 商および冪等完備化へ同じ比較が移る。
各 $\OO(j)$ は作用によって対象として保たれるので、
定理\ref{mbsd2:thm:projector-descent}の admissible 部分圏と
その compact 直交残余圏も保たれる。
生成対象の線型化を変えても、その Hom と対象を tensor した evaluation では
二つの指標が相殺される。従って残余への随伴射影にも比較が誘導される。
これは随伴の同変性の一般論とも一致する
\cite[Proposition~3.13, Theorem~6.2]{MBZAShinder}。
同変 dg 圏を用いる際には、既存の enhancement と有限群作用からの構成
\cite[Theorems~8.7, 8.9]{MEDElagin}、または
補題\ref{med:lem:crossed-perfect}の模型を使う。
compact 残余作用には追加の window 仮定は要らない。
全 coherent 残余や MF 圏との同定、および smooth proper 性は、既存の入力の範囲で用いる。
\end{proof}

\subsection{同変 shift と rotation の変換則}

指標 $\eta\in\widehat Z_0$ に対し、同変化上の自己同値を
\begin{equation}\label{mbza:eq:character-twist}
 T_\eta(M,\theta_z)=(M,\eta(z)\theta_z)
\end{equation}
とする。これは既知の双対群による作用である
\cite[\S4, before Theorem~4.2]{MEDElagin}。

\begin{proposition}[指標による rotation の差]
\label{mbza:prop:shift-character}
$\mathsf s=(1)$ とし、分裂 $\boldsymbol a$ による標準の同変 shift を
$\mathsf s_{\boldsymbol a}$ と書く。
$\chi=\chi_{\boldsymbol a\boldsymbol a'}$ に対し、全整数 $j$ について
\begin{equation}\label{mbza:eq:shift-character}
 J_{\boldsymbol a\boldsymbol a'}\mathsf s_{\boldsymbol a}^{j}
 \simeq T_{\chi^{-j}}\mathsf s_{\boldsymbol a'}^{j}
                     J_{\boldsymbol a\boldsymbol a'}.
\end{equation}
命題\ref{mbsd2:prop:rotation-serre}の compact 残余上の rotation
$\rho=L_{\OO}\circ(1)$ についても
\begin{equation}\label{mbza:eq:rotation-character}
 J_{\boldsymbol a\boldsymbol a'}\rho_{\boldsymbol a}
 \simeq T_{\chi^{-1}}\rho_{\boldsymbol a'}
                      J_{\boldsymbol a\boldsymbol a'}
\end{equation}
である。ここで mutation は underlying 残余模型の evaluation から誘導した関手である。
特に第三冪については指標が消え、
\begin{equation}\label{mbza:eq:third-period}
 J_{\boldsymbol a\boldsymbol a'}\rho_{\boldsymbol a}^{3}
   \simeq\rho_{\boldsymbol a'}^{3}J_{\boldsymbol a\boldsymbol a'}.
\end{equation}
MF/window 比較を使って命題\ref{mbsd2:prop:line-period}の三次周期が得られる範囲では、
その係数直線付きの周期比較も、この同定の下で保たれる。
\end{proposition}
\begin{proof}
$J\mathsf s^j$ の線型化は
$(\mathsf s^j\theta_z)\epsilon_{z,M(j)}$、
$\mathsf s^jJ$ の線型化は
$(\mathsf s^j\theta_z)(\epsilon_{z,M})(j)$ である。
補題\ref{mbza:lem:scalar}より前者は後者の $\chi(z)^{-j}$ 倍なので、
式\eqref{mbza:eq:shift-character}となる。
この符号は規約\eqref{mbza:eq:transport}による。

$J$ は構造対象の標準線型化を保つため、
$\operatorname{RHom}(\OO,-)\otimes\OO\to\operatorname{id}$ とその cone に両立する。
指標 twist もこの評価写像と可換する。
従って式\eqref{mbza:eq:shift-character}を $L_{\OO}$ と合成して
式\eqref{mbza:eq:rotation-character}を得る。
ここで同変化の中の一つの線型化付き $\OO$ だけに対する mutation を使ってはいない。
必要な評価は $Z_0$-表現全体を含むものであり、
underlying の mutation を同変化へ移した関手である。

$T_\chi$ は rotation と可換し、$T_\chi^3=\operatorname{id}$ なので第三冪の式が従う。
周期が与えられた場合、その内部次数は三で、分裂変更が与える係数は $\chi^{-3}=1$ である。
中心三次直線も $\widetilde Z$ に固定される。
従って、指定された周期写像自体も比較と可換する。
これは周期直線の大域的自明化を与える主張ではない。
\end{proof}

\begin{proposition}[線型化の選択の完全性]
\label{mbza:prop:linearization-torsor}
代数閉体 $K$ 上で $\mathcal C_A=\QGr A$ の一つの中心作用を固定する。
内部 shift $(1)$ の $Z_0$-同変構造の同型類は $\widehat Z_0$ の torsor であり、
分裂の標準構造を式\eqref{mbza:eq:action-comparison}で移したものが全てを尽くす。

残余 rotation についても、その dg 模型 $\mathcal K$ が
\[
 HH^0(\mathcal K)=K,\qquad HH^q(\mathcal K)=0\quad(q<0)
\]
を満たす場合には同じ結論が成り立つ。
第\ref{med:sec:descent}節の smooth proper 一般点では、有限体拡大後の
三次 Kuznetsov 圏との比較と命題\ref{prop:nf-hh}からこの条件が従う。
\end{proposition}
\begin{proof}
二つの同変構造の各 $z$ における比は、同値関手の自然自己同型である。
系\ref{mbsd:cor:natural-scalars}により、それは $K^\times$ の元である。
同変構造の群法則は、この比が $Z_0$ の指標であることと同値になる。
関手の自然自己同型もスカラーなので、非自明な指標差を自然同型で消すことはできない。
命題\ref{mbza:prop:splittings}と式\eqref{mbza:eq:shift-character}から、
全ての指標差が分裂の変更によって得られる。

残余の場合、可逆 Morita 関手の自己準同型は恒等双加群の自己準同型に帰着する。
その次数零は $HH^0(\mathcal K)=K$、高い homotopy の曖昧さは
負次数の $HH$ によって消えるため、同じ指標の議論を適用できる。
smooth proper 性の下で Hochschild コホモロジーは体拡大と両立し、
忠実平坦な体拡大後に上の等式が成立すれば元の体でも成立する。
\end{proof}

\subsection{順序を忘れる降下と三の倍数の障害}

\begin{corollary}[置換不変な分裂の必要十分条件]
\label{mbza:cor:permutation-obstruction}
$\mathfrak S_k$ は $\mathscr S_k$ の座標を置換する。
$\mathfrak S_k$-同変な分裂 $Z_0\to\widetilde Z$ が存在する必要十分条件は
$3\nmid k$ である。その場合には一意であり、
\begin{equation}\label{mbza:eq:symmetric-splitting}
 a_i=k^{-1}\in\mathbf F_3\quad(1\le i\le k)
\end{equation}
で与えられる。

$3\mid k$ の場合、分裂の torsor の障害類は
\begin{equation}\label{mbza:eq:splitting-obstruction}
 0\ne\delta(1)\in H^1(\mathfrak S_k,\widehat Z_0),
 \qquad
 0\longrightarrow\widehat Z_0\longrightarrow\mathbf F_3^k
     \xrightarrow{\sum}\mathbf F_3\longrightarrow0
\end{equation}
から来て、位数三である。
これは中心作用そのものや同変化の降下障害ではない。
標準の次数一同変構造を、ブロックの順序によらず一つ選ぶことの障害である。
命題\ref{mbza:prop:linearization-torsor}の仮定の下では、
rotation の同変構造についても同じ意味で完全な障害となる。
\end{corollary}
\begin{proof}
置換不変な $\boldsymbol a$ は全座標が一つの $c\in\mathbf F_3$ に等しい。
条件 $\sum_i a_i=1$ は $kc=1$ となり、主張の必要十分条件と一意性を得る。
完全列の不変部分を取った接続写像は、$1$ の逆像である torsor
$\mathscr S_k$ の類を与える。$3\mid k$ なら不変部分上の $\sum$ は零なので
$\delta(1)\ne0$ であり、係数が $\mathbf F_3$ なのでその位数は三である。
各 $\boldsymbol a$ による作用は既に定理\ref{mbza:thm:central-action}の比較で降下している。
残る指標の変換は命題\ref{mbza:prop:shift-character}に正確に記述され、
線型化の選択を全て尽くすことは命題\ref{mbza:prop:linearization-torsor}による。
\end{proof}

順序付き表示被覆を $\widetilde{\mathfrak M}^{\circ}$ とすると、
\begin{equation}\label{mbza:eq:splitting-cover}
 \widetilde{\mathfrak M}^{\circ}
       \times^{\mathfrak S_k}\mathscr S_k
       \longrightarrow\mathfrak M^{\circ}
\end{equation}
は階数 $3^{k-1}$ の有限 \'etale な分裂選択被覆である。
その上では中心拡大\eqref{mbza:eq:central-extension}の分裂を選べる。
rotation の有限群線型化は、この被覆上の実表示 gerbe へ引き戻した
rotation について選ぶものであり、共通 $\mathbf G_m$-gerbe を中立化する操作ではない。
順序付き表示チャートは既約なので、この被覆の切断は
$\mathscr S_k$ の $\mathfrak S_k$-固定点と同値である。
従って $3\mid k$ なら、非空な Zariski 開へ縮めるだけでこの障害を除くことはできない。
基準ブロック一つを選ぶ $k$ 重の有限 \'etale 被覆でも分裂を選べる。

より具体的に、$\mathscr S_k$ の置換軌道は
\[
 n_0+n_1+n_2=k,\qquad n_1+2n_2\equiv1\pmod3
\]
を満たす三つの非負整数で分類され、その軌道の大きさは
$k!/(n_0!n_1!n_2!)$ である。
例えば $k=3$ では、九つの分裂は三つの三元軌道からなり、固定点はない。

\begin{remark}[半直積としての中心拡大]
同じ障害は
\[
 1\longrightarrow\Delta\longrightarrow
 \widetilde Z\rtimes\mathfrak S_k
 \longrightarrow Z_0\rtimes\mathfrak S_k\longrightarrow1
\]
の群準同型による分裂にも現れる。
分裂を $Z_0$ に制限すれば式\eqref{mbza:eq:central-extension}の分裂となる。
$\widetilde Z$ が可換なので、$\sigma\in\mathfrak S_k$ の持上げを
$\widetilde Z$ の元だけ変えても、その共役作用は変わらない。
従って制限した分裂は必ず置換同変である。
逆向きは明らかなので、この半直積の拡大も $3\nmid k$ のときに限り分裂する。
しかし共通 scalar には式\eqref{mbza:eq:scalar-trivialization}の自然自明化があるため、
分裂しない場合にも coherent な圏作用の降下は妨げられない。
\end{remark}

\subsection{同変 MF との一致と二ブロック}

\begin{proposition}[拡大 grading 群による同じ中心作用]
\label{mbza:prop:gamma-action}
\[
 \Gamma=\{(\lambda_1,\ldots,\lambda_k)\in\mathbf G_m^k:
                         \lambda_i^3=\lambda_j^3\text{ for all }i,j\}
\]
とし、共通三乗を $\upsilon:\Gamma\to\mathbf G_m$ と書く。
自然に
\[
 \Gamma=(\mathbf G_m\times\widetilde Z)/
   \{(t,(t^{-1},\ldots,t^{-1})):t\in\mu_3\},
 \qquad
 1\longrightarrow\mathbf G_m^{\mathrm{diag}}
       \longrightarrow\Gamma\longrightarrow Z_0\longrightarrow1
\]
である。
potential が $\upsilon$ による半不変量である同変 MF の模型では、
対角 $\mathbf G_m$ による通常の内部 grading の後に現れる $Z_0$ の弱作用は、
定理\ref{mbza:thm:central-action}の作用と一致する。
従ってこの $Z_0$-同変化は、$\Gamma$ による同変 MF と同じ圏を与える。
\end{proposition}
\begin{proof}
fpqc 局所で共通三乗の三乗根 $t$ を選ぶと、各 $\lambda_i/t$ は三乗根となる。
その選択の差が対角 $\mu_3$ であり、最初の表示を得る。
分裂 $s_{\boldsymbol a}$ を選べば
\[
 \mathbf G_m\times Z_0\xrightarrow{\sim}\Gamma,
 \qquad(t,z)\longmapsto(t,\ldots,t)s_{\boldsymbol a}(z)
\]
で、$\upsilon(t,z)=t^3$ である。
従って $\Gamma$-同変な factorization は、内部 grading と、
potential を固定する $s_{\boldsymbol a}(Z_0)$ の線型化を持つ factorization である。
これは標準の $Z_0$-同変 graded MF の定義そのものである。
分裂変更は各内部次数 $n$ で $\chi(z)^n$ を掛ける写像を与える。
二つの factorization 写像のうち shift を含む方でも、内部 shift は三であり、
$\chi^3=1$ なのでこの写像は可換する。
従って比較は式\eqref{mbza:eq:action-comparison}と同じであり、全 cocycle も一致する。
\end{proof}

\begin{corollary}[二ブロックの対称な選択]
\label{mbza:cor:two-block}
$k=2$ で $u=h_1/h_2$ により $Z_0\simeq\mu_3$ と書くと、
三つの分裂は
\[
 s_2(u)=(u,1),\qquad
 s_1(u)=(1,u^{-1}),\qquad
 s_{\mathrm{sym}}(u)=(u^2,u)
\]
である。最後の分裂だけが因子交換と両立する。
因子交換は $u\mapsto u^{-1}$ を誘導する。
\end{corollary}
\begin{proof}
式\eqref{mbza:eq:splitting-space}の三点は $(0,1),(1,0),(2,2)$ である。
式\eqref{mbza:eq:splitting-formula}を代表 $(u,1)$ に適用すれば表示した式を得る。
例えば $s_{\mathrm{sym}}(u^{-1})=(u,u^2)$ は $s_{\mathrm{sym}}(u)$ の因子交換である。
\end{proof}

第\ref{med:sec:descent}節の、一方のブロックだけを三乗根倍する作用は $s_2$ に対応する。
他方を基準に取る場合、同じ内在的群 $Z_0$ の座標を使えば作用は
$(1,u^{-1})$ であり、$(1,u)$ と書く際には群の inversion も含める必要がある。
三つの選択の同変残余圏は定理\ref{mbza:thm:central-action}によって同定されるので、
Hu の比較を適用した同変化自体はこの基準選択に依存しない。
ただし標準 rotation の比較には式\eqref{mbza:eq:rotation-character}の指標を含める。

\begin{remark}[有限の指標障害と表示 gerbe の区別]
\label{mbza:rem:scope}
本節が降下したのは中心群の coherent な作用と同変化である。
置換不変な分裂があっても、共通 $\mathbf G_m$ を除いた表示 gerbe が
自明になったことにはならない。
式\eqref{mbza:eq:shift-weight}は任意の $t\in\mathbf G_m$ について
次数一 shift に $t^{-1}$ の差が残ることを示す。
従って、全 $\OO$-底上の次数一関手・rotation・周期直線を
追加の線束や gerbe のデータなしに一括して自明化したとは主張しない。
また、二つの同変圏の間に比較核が存在することと、
その核に全 $\OO$-降下の cocycle が選べることは別問題である。
第\ref{med:sec:descent}節の指定した一般点比較や有限拡大上の核を、
本節だけから全モジュライへ降ろすことはしていない。
本節の scalar 比較と有限 torsor は直接証明したものであり、
既知の同変圏の一般論との組合せの新規性の網羅的確認は未完了である。
\end{remark}

% !TEX root = ../nc-cubic.tex
\section{多ブロック族の同変一般点復元}
\label{med:sec:descent}

本節では、有限群作用を保つ完備MF比較から、一般点上の同変Morita同値を取り出す。
中心となるのは、同変核をcrossed product上のperfect加群として降下させることである。
群作用のcoherenceを有限個の通常の射だけで代用することはしない。
さらに低rotation積の不変な積分空間を使うと、ブロック分解を持つ三次方程式を
元の関数体上に選べる。同値核には有限代数拡大を許す。
これは定理\ref{cft:thm:complete}と定理\ref{gdr:thm:generic}を入力とする帰結であり、
これらの基礎定理の独立な再証明ではない。
一般的な核の有限表示と同変圏の構成には既知の結果を用いる。
この組合せによる多ブロック入力への適用の新規性は、現段階では独立には確認していない。

\subsection{有限群のcoherenceを含むperfect核の降下}

有限群$G$の作用はdg模型上で与える。
生成対象$H$が群で固定されていない場合は、有限和
$\bigoplus_{g\in G}g^*H$とその置換による線形化を用いる。
これにより、以下では厳密な$G$作用を持つdg代数で記述できる。
同変圏には有限coneと冪等完備化を含める。
この完備化を伴うdg同変圏と三角圏の同変化の比較は
Elagin\cite[Lemma~8.6, Theorems~8.7, 8.9]{MEDElagin}に従う。
同変核は、厳密な半線形群作用を持つdg双加群の導来圏で扱う。
同変quasi-functorからはequivariant bar模型によってこのような核を得る。
bar模型の項が無限個あっても、以下のperfect性は有限cellのretractとしての性質である。
従って、coherenceの無限列を項ごとに有限型環へ降ろす必要はない。

\begin{lemma}[crossed productとperfect性]\label{med:lem:crossed-perfect}
$D$を可換環、$|G|\in D^\times$とし、$\mathsf E$を厳密な$D$線形$G$作用を持つ
$D$上のdg代数とする。導来テンソルを計算できるflat模型を使い、
\[
 \mathsf E\rtimes G=\bigoplus_{g\in G}\mathsf E u_g,
 \qquad (a u_g)(b u_h)=a\,g(b)u_{gh}
\]
と置く。$\mathsf E\rtimes G$上のdg加群$M$について、
$M$がperfectであることと、忘却した$\mathsf E$加群がperfectであることは同値である。
また$\mathsf E$が$D$上smooth properなら、$\mathsf E\rtimes G$もsmooth properである。
\end{lemma}
\begin{proof}
証明では左加群の規約を使う。
crossed productは左右の$\mathsf E$加群として有限自由なので、
制限はperfect性を保ち、誘導はperfect加群をperfect加群へ送る。
誘導・制限の余単位
$(\mathsf E\rtimes G)\otimes_{\mathsf E}M\to M$には、
\[
 m\longmapsto |G|^{-1}\sum_{g\in G}u_g\otimes u_{g^{-1}}m
\]
という$\mathsf E\rtimes G$線形な分裂がある。
これは微分と可換し、導来圏でも分裂を与える。
従って忘却後perfectな$M$は、誘導されたperfect加群のretractである。

proper性は$\mathsf E\rtimes G$が$D$複体として$|G|$個の$\mathsf E$の直和であることから従う。
smooth性については、$\mathsf E$のperfectな対角を誘導して得る
\[
 (\mathsf E\rtimes G)\otimes_{\mathsf E}^{\mathbf L}(\mathsf E\rtimes G)
\]
が$(\mathsf E\rtimes G)^e$上perfectである。
その乗法写像は、separability element
$|G|^{-1}\sum_g u_g\otimes u_{g^{-1}}$によって双加群として分裂する。
従って$(\mathsf E\rtimes G)$の対角もperfectである。
\end{proof}

\begin{lemma}[同変同値核の有限代数拡大上の存在]
\label{med:lem:equivariant-descent}
$\kappa$を標数零の体、$\Omega/\kappa$を任意の体拡大とする。
$\mathsf C,\Gamma$を厳密な$G$作用を持つsmooth properな$\kappa$-dg代数とする。
両側に$G$同変な可逆rotation双加群と、必要なら指定した三次周期同型を与える。
$\Omega$上に$G$同変でrotation・指定周期を保つMorita同値があれば、
ある有限代数拡大$\kappa'/\kappa$上にもそのような同値が存在する。
ここで同変性の入力はenhancement上のcoherentな同変核であり、
核の同型類が$G$不変であるという条件だけではない。
\end{lemma}
\begin{proof}
同値を表す$(\mathsf C_\Omega,\Gamma_\Omega)$双加群を$P_\Omega$とする。
右$\Gamma_\Omega$-perfect性と$\mathsf C_\Omega$のsmooth性により、これは
$\mathsf E_\Omega=\mathsf C_\Omega^{\rm op}\otimes_\Omega\Gamma_\Omega$上perfectである
\cite[Lemma~2.8(2)]{ROGTV}。
同変核の群作用とcoherenceは、対角作用に関する
$\mathsf E_\Omega\rtimes G$上の一つのdg加群として表す。
必要なdg加群の置換を行った後、補題\ref{med:lem:crossed-perfect}により
これはcrossed product上でもperfectである。

$\Omega$を有限生成$\kappa$-部分代数のfiltered colimitとして表す。
perfect加群のfiltered colimitとの両立
\cite[Lemma~2.10]{ROGTV}、\cite[F4]{ROGToenDefinition}により、
$P_\Omega$はある有限生成整域$D\subset\Omega$上の
$(\mathsf C_D^{\rm op}\otimes_D^{\mathbf L}\Gamma_D)\rtimes G$-perfect加群$P_D$から来る。
従って群作用の乗法則を含むcoherenceは、この段階ですでに降下している。
忘却後の$P_D$も包絡代数上perfectであり、$\mathsf C_D$がproperなので右$\Gamma_D$-perfectである。

右随伴の核
$Q_D=R\operatorname{Hom}_{\Gamma_D}(P_D,\Gamma_D)$には自然な$G$作用がある。
補題\ref{rog:lem:descent}の証明と同じperfect性の判定により、
随伴の単位・余単位のconeは$D$-perfectとなる。
これらの写像は同変で、$\Omega$への底変更後は同型である。
有限個のcohomology加群のannihilatorを取れば、ある$0\ne f\in D$の反転後に
両coneを零にできる。従って$P_{D_f}$は同変Morita同値を与える。

rotation比較も、対角$G$作用を持つperfect双加群の圏における射として降下させる。
その両辺は忘却後perfectであり、補題\ref{med:lem:crossed-perfect}により
同じcrossed product上でもperfectである。
上記のfiltered colimitの完全忠実性によって比較射、その逆、および
指定周期図式の等式を有限段階へ降ろせる。
必要なら$D$を拡大し、さらに非零元を反転すればよい。
ここで有限個の周期図式を扱う議論と、群coherenceをcrossed productで扱う議論は別である。

非空な有限型$\kappa$-scheme $\operatorname{Spec}D_f$の閉点を取る。
その剰余体$\kappa'$はZariskiの補題により有限代数拡大であり、
そこでの底変更が主張を与える。
\end{proof}

\begin{remark}[存在の降下と指定形式核の降下]
\label{med:rem:formal-versus-existence}
この証明は、元の$\Omega$上の核を有限型$\kappa$-代数上へ広げた後、
その別の閉点で特殊化する。
従って両体を含む共通拡大体へ戻しても、得た核が元の指定核と一致するとは限らない。
形式係数が$\kappa$上代数的であることや、全jetを保存する代数化も従わない。
有限拡大上の存在にはこの議論で十分であり、Artin近似による全形式分枝の再現を仮定していない。
\end{remark}

\begin{example}[指定された完備核は有限拡大へ降りるとは限らない]
\label{med:ex:formal-kernel}
$K=\mathbb C(h)$、$\widehat R=\mathbb C[[h]]$とし、$a=\exp(h)-1\in\widehat R$と置く。
$\mathbb P^1_{\widehat R}$の自己同型
$\sigma_a([X:Y])=[X+aY:Y]$のgraphの構造層は、自己Morita同値を与えるperfect核である。
この核の一般ファイバーが、ある有限拡大$K'/K$上の核から来て、
共通拡大体上で指定核と同型になったと仮定する。
忠実平坦降下によって、そのsupportは$K'$上のgraphとなり、
$[0:1]$の像から$a\in K'$を復元できる。
しかし$a$は$K$上超越的である。
実際、代数関係があれば分母を払って
$\sum_{j=0}^m p_j(h)\exp(jh)=0$という非自明な関係を得るが、
実軸で$h\to+\infty$とし、最大の指数項で割ると非零多項式$p_m(h)$が零に近づくことになり矛盾する。
従ってこの指定核は有限代数拡大上へ降りない。
一方、同じ二つの圏の間には恒等核が$K$上にある。
これは補題\ref{med:lem:equivariant-descent}の存在結論と、指定核の降下との違いを示す。
\end{example}

\subsection{群作用を保つ方程式の基礎体降下}

\begin{lemma}[同じ勾配空間を持つ不変三次式の同変比較]
\label{med:lem:invariant-symmetrizer}
$\Omega$を標数零の代数閉体、$V$を有限次元の$G$表現とする。
二つの$G$不変な滑らかな三次式$F,F'\in\operatorname{Sym}^3V$が
同じ勾配空間を持つなら、$G$と可換する線形座標変換で移り合う。
\end{lemma}
\begin{proof}
$x\in V^*$について$F(x)=T(x,x,x)$とする。
以下の$E_F$と座標変換は$V^*$に作用する。
命題\ref{prop:ej-projective-uniqueness}の証明により、
一意な$a\in E_F^\times$について$F'(x)=T(ax,x,x)$となる。
$F,F'$の不変性から$gag^{-1}=a$であり、$a\in(E_F^G)^\times$である。
補題\ref{ej:lem:symmetrizer-algebra}より$E_F$は有限\'etaleな可換代数である。
代数閉体上では$E_F=\Omega^s$であり、$G$はその因子を置換するので、
$E_F^G$も有限個の$\Omega$の積となる。
各因子で三乗根を取り、$b\in(E_F^G)^\times$かつ$b^3=a$とできる。
symmetrizerの定義から
$F(bx)=T(bx,bx,bx)=T(b^3x,x,x)=F'(x)$であり、$b$は$G$と可換する。
\end{proof}

\begin{proposition}[同変cubic模型の方程式降下]
\label{med:prop:equivariant-field-model}
$\kappa$を標数零の体とし、$(\mathcal C,\rho)$をsmooth properなdg pairとする。
有限群$G$がenhancement上で作用し、$\rho$と指定周期は$G$同変とする。
ある代数閉拡大体$\Omega/\kappa$上で、このpairが
$G$不変な滑らかな三次式$F_\Omega$のKuznetsov pairと$G$同変にMorita同値であると仮定する。
変数の数は$N\ge5$とし、可換側の$G$作用は変数上の線形作用から来るものとする。
このとき、式\eqref{gcd:eq:low}の$\kappa$-空間$V$上に
\[
 F_\kappa\in(\operatorname{Sym}^3_\kappa V)^G
\]
を選べる。その零点集合は幾何学的に滑らかであり、
ある有限拡大$\kappa'/\kappa$上で元のpairと$G$同変にMorita同値である。
同値はrotationと指定周期も保つ。
\end{proposition}
\begin{proof}
$V,U$と積$\mu$には、$\rho$の同変構造から自然な$G$作用が入る。
その形成は体拡大と可換する（補題\ref{gcd:lem:basechange}）。
可換側では、一次式の乗法による自然変換
$\mathrm{Id}\to(1)$と二次式の乗法による自然変換$\mathrm{Id}\to(2)$は
線形座標変換に自然である。
命題\ref{prop:ej-geometric-low-degree}の恒等sector計算はこれらを使うので、
低次数のJacobian環との同定および積は$G$同変である
（\cite[Theorem~3.5, Remark~3.6]{LZ}の多項式乗法による表示）。
従って
$Q=\ker(\operatorname{Sym}^2V\to U)$は$\kappa$上の$G$不変な$N$次元空間であり、
$\Omega$上では$F_\Omega$の勾配空間となる。

積分空間を
\[
 \mathcal I_Q=
 \{H\in\operatorname{Sym}^3V:\partial_\lambda H\in Q\text{ for every }\lambda\in V^*\}
\]
と定める。これは$G$不変な$\kappa$-線形空間であり、その不変部分の形成も体拡大と可換する。
各$H\in\mathcal I_Q$は線形写像
$dH:V^*\to Q$を定める。
基底を選んで$\det(dH)$を$\mathcal I_Q^G$上の多項式と見ると、
$H=F_\Omega$で非零なので零多項式ではない。
$\kappa$は無限体であるから、$\det(dF_\kappa)\ne0$となる
$F_\kappa\in\mathcal I_Q^G$を選べる。
この三次式の勾配空間はちょうど$Q$である。
$Q_\Omega$には共通射影零点がないので、$F_\kappa$は幾何学的に滑らかである。

補題\ref{med:lem:invariant-symmetrizer}により、$\Omega$上で
$F_\kappa$と$F_\Omega$を移す座標変換を$G$と可換するように取れる。
この座標変換は方程式を正確に移すため、MFの内部shiftと指定三次周期を保つ。
したがって$\Omega$上で$G$同変なrotation付きMorita同値を得る。
両側の$\kappa$-dg模型に補題\ref{med:lem:equivariant-descent}を適用すればよい。
\end{proof}

\subsection{多ブロック一変数族への適用}

$k\ge2$、$R=\mathbb C[h]_{(h)}$、$K=\mathbb C(h)$とする。
各$i=1,\ldots,k$について、非負次数付き中央$R$-代数$D_i$と
中心三次元$\Omega_i\in Z(D_i)_3$を与え、次を仮定する。
\begin{enumerate}
\item $(D_i)_0=R$で、各$(D_i)_d$は有限自由である。
\item $D_i/hD_i\simeq\mathbb C[x_{i1},x_{i2},x_{i3}]$である。
\item $\Omega_i\bmod h=f_i$は滑らかな平面三次曲線を定める。
\end{enumerate}
$c_i\in R^\times$を取り、通常のテンソル積と全次数により
\begin{equation}\label{med:eq:multiblock-input}
 B_R=\bigotimes_{i=1}^kD_i,\qquad
 w=\sum_{i=1}^k c_i\Omega_i,\qquad A_R=B_R/(w)
\end{equation}
と置く。各$\Omega_i$は対応するテンソル因子内の元と解釈する。
本節の周囲代数$B_R$、商$A_R$は第\ref{gdr:sec:main}節の記号と逆である。
$G=(\mu_3)^{k-1}$は、$i$番目の因子を$i<k$で独立に次数ごと三乗根倍し、
第$k$因子を固定する。この作用は$w$を保つ。

\begin{theorem}[多ブロック一般点の同変復元]
\label{med:thm:multiblock-generic}
式\eqref{med:eq:multiblock-input}の入力に対して、$K$上の三次式$F_i$を用いた
滑らかな三次超曲面
\begin{equation}\label{med:eq:split-model}
 Y=V\bigl(F_1(z_{11},z_{12},z_{13})+\cdots+
 F_k(z_{k1},z_{k2},z_{k3})\bigr)\subset\mathbb P^{3k-1}_K
\end{equation}
が存在する。各$C_i=V(F_i)\subset\mathbb P^2_K$も滑らかである。
ある有限拡大$K'/K$上で
\begin{equation}\label{med:eq:equivariant-pair}
 \bigl(\mathcal K^{\rm coh}_{A_{K'}},\rho_A,\tau_A\bigr)
 \simeq_G
 \bigl(\operatorname{Ku}(Y_{K'}),\rho_Y,\tau_Y\bigr)
\end{equation}
という$G$同変Morita同値がある。
ここで$Y$上の作用は第$i<k$ブロックを対応する三乗根倍する標準作用であり、
$\tau$は指定した三次周期比較である。
\end{theorem}
\begin{proof}
補題\ref{tbr:lem:input}の次数別分裂単射の証明は$k$因子でも同じである。
$B_R$のHilbert級数は$(1-t)^{-3k}$であり、$w$は左右正則、商の各次数は有限自由となる。
特殊ファイバーの勾配は$k$個の平面三次の勾配に分かれるため、その共通アフィン零点は原点だけである。
従って$\sum_i\overline c_i f_i$は滑らかな射影三次式である。
定理\ref{gdr:thm:generic}を$N=3k$に適用すると、
$\mathcal K^{\rm coh}_{A_K}$はsmooth properであり、線束生成残余圏と一致する。
この段階では非同変の一般点定理だけを用いる。

次に$\widehat R=\mathbb C[[h]]$上で、命題\ref{tbr:prop:complete-split}の構成を$G$に適用する。
同命題の証明は有限群の不変部分の完全性と線形共変性を用いるので、
巡回群からこの直積群への変更でも成立する。
各平方零段階の線形補正を不変部分で選び、非線形な解そのものは平均しない。
中心写像、唯一の定数trivector、周期変数を固定したbar flowをすべて同変に取れるので、
実際の完備MF同値とその逆の各Taylor成分は$G$同変となる。
不変三次空間は
\begin{equation}\label{med:eq:invariant-cubics}
 \bigl(\operatorname{Sym}^3\textstyle\bigoplus_i\widehat R^3_i\bigr)^G
 =\bigoplus_{i=1}^k\operatorname{Sym}^3\widehat R^3_i
\end{equation}
である。実際、各$i<k$ブロックの次数は三の倍数でなければならず、
総次数が三なので一つのブロックに全次数が集中する。
したがって完備potentialは$W=\sum_iW_i$と分裂し、$W_i\bmod h=\overline c_if_i$である。

固定三対象とその有限個の$G$軌道へこの実際の同値を制限し、
補題\ref{rog:lem:localization}の順序で通常のHom複体を構成してから$h$を反転する。
群作用はこの有限対象模型に作用し、同変なTaylor成分と可換する。
stabilization、内部shift、windowで用いる切断とevaluationは係数の群作用に自然であり、
rotationおよび指定周期の比較も同変である。
従ってここで得るのは、単なる核の不変な同型類ではなく、coherentな同変quasi-functorである。
具体的にはそのbar双加群の各tensor因子へ$G$を対角的に作用させると、
微分が$G$同変となり、群の乗法則は厳密に成立する。
$L=\mathbb C((h))$の代数閉包$\overline L$へ移すと、
命題\ref{rog:prop:euler}と命題\ref{gcd:prop:bound}により像は可換残余圏全体となる。
これで$\overline L$上の$G$同変Morita同値を得る。
任意のLaurent係数MFへ無限Taylor和を直接評価する必要はない。

命題\ref{med:prop:equivariant-field-model}を$\kappa=K$、$\Omega=\overline L$に適用する。
低rotation空間$V$の$G$表現は、$\overline L$上で各三次元のブロック指標空間の直和となる。
$K$は三乗根を含むため、指標冪等元によって
$V=V_1\oplus\cdots\oplus V_k$、$\dim_KV_i=3$と$K$上でも分解する。
式\eqref{med:eq:invariant-cubics}を$K$上で用いると、復元した不変三次式は
$F_K=\sum_iF_i$となる。各$V_i$の基底を選べば式\eqref{med:eq:split-model}を得る。
全体の滑らかさは同命題から従う。
ある$F_i$の平面三次が特異なら、その特異点と他の零ブロックを合わせた点が
$Y$の特異点になるため、各$C_i$も滑らかである。
同命題の最後の同変核降下が式\eqref{med:eq:equivariant-pair}を与える。
\end{proof}

\subsection{二ブロックの場合の楕円曲線の積}

Huの定理\cite[Theorems~0.1, 5.8]{TBHu}は複素数体上で述べられている。
以下では、その本文の体の仮定を無断で一般化せず、有限定義体とperfect核の降下を使う。

\begin{lemma}[Huの比較の有限拡大上への移行]
\label{med:lem:hu-field}
$\kappa$を三乗根を含む標数零の体とし、
$Y=V(F_1(x)+F_2(y))\subset\mathbb P^5_\kappa$を滑らかな三次四様体とする。
$G=\mu_3$は一方の三変数ブロックを三乗根倍し、$C_i=V(F_i)$とする。
ある有限拡大$\kappa'/\kappa$上で
\[
 \operatorname{Ku}(Y_{\kappa'})^G
 \simeq D^b(C_{1,\kappa'}\times C_{2,\kappa'})
\]
というdg Morita同値がある。
\end{lemma}
\begin{proof}
三乗根と$F_1,F_2$の有限個の係数を含む、$\mathbb Q$上有限生成の部分体
$\kappa_0\subset\kappa$を選ぶ。
$\kappa_0$上の三次四様体および二曲線は幾何学的に滑らかであり、
その同変Kuznetsov成分と曲線の積の導来圏はsmooth properなdg模型を持つ。
前者については、残余圏のsmooth proper模型と
補題\ref{med:lem:crossed-perfect}を用いてもよい。

$\kappa_0\hookrightarrow\mathbb C$という体の埋込みを選ぶ。
その複素底変更へHuの定理を適用する。
同論文\S5.1のderived McKay関手はFourier--Mukai核を持ち、
\S5.4、Theorem~5.8の比較はこれとmutationおよびblow-upの関手の合成である。
従って得られる同値はdg enhancement上のperfect核で表される。
固定された$\kappa_0$-dg模型に補題\ref{rog:lem:descent}を、rotationの指定なしで適用すると、
ある有限拡大$\kappa_1/\kappa_0$上に同値核が存在する。
$\kappa_1$を$\kappa$の代数閉包へ$\kappa_0$上で埋め込み、合成体を$\kappa'$とすれば、
$\kappa'/\kappa$は有限であり、所要の同値が得られる。
\end{proof}

\begin{corollary}[二ブロック残余圏の同変化の幾何模型]
\label{med:cor:elliptic-product}
定理\ref{med:thm:multiblock-generic}で$k=2$とする。
そこで得た$K$上の二つの滑らかな平面三次$C_1,C_2$に対し、
ある有限拡大$K''/K$上で
\begin{equation}\label{med:eq:elliptic-product}
 \bigl(\mathcal K^{\rm coh}_{A_{K''}}\bigr)^{\mu_3}
 \simeq D^b(C_{1,K''}\times C_{2,K''})
\end{equation}
となる。必要なら$K''$をさらに有限拡大し、各$C_i$に原点を選ぶことで、
右辺を楕円曲線の積の導来圏と呼べる。
\end{corollary}
\begin{proof}
同変Morita同値は同変対象の圏のMorita同値を誘導する。
この形成はcrossed productの表示により体拡大と両立する。
式\eqref{med:eq:equivariant-pair}と補題\ref{med:lem:hu-field}を合成すればよい。
各滑らかな平面三次は種数一であり、有限拡大後には閉点を有理点にできる。
二つの点の定義体の合成を取れば、両曲線に原点を選べる。
\end{proof}

式\eqref{med:eq:elliptic-product}は同変対象の圏についての主張である。
Huの可換模型における既知の同値を、新たな同変一般点比較へ適用したものであり、
その可換側の同値自体を本研究の成果とは数えない。
通常の残余圏を楕円曲線の積の導来圏と同定する主張や、
本節のHu比較だけから積側のrotationを指定する主張は含まない。
全ブロック数の積表示と、BFK外積による因子rotationの記述は
定理\ref{mbrd:thm:split-product}、命題\ref{mbrd:prop:factor-rotations}で与える。

\begin{remark}[得られた降下の範囲]\label{med:rem:scope}
式\eqref{med:eq:split-model}の方程式は$K$上、同変比較核はある有限拡大上に存在する。
不変積分空間から選んだ方程式が$h=0$で正則であることや、
その剰余が指定された$\sum_i\overline c_if_i$であることは結論していない。
各$C_i$を入力Sklyanin因子の点曲線と同定する結果も用いていない。
明示的 Jacobian tensor 族という追加仮定の下では、
定理\ref{mpc:thm:point-model}が点三次からこの特殊ファイバーも保つ模型を構成し、
系\ref{mpc:cor:block-j}が各点曲線との同型を与える。
指定した完備核と全jetを保つ代数化、特殊点を含む相対同値、
任意の固定非零パラメータでの比較、および$\mathcal O$-保存同値のモジュライへの降下は、
本節の一般点における存在定理からは従わない。
\end{remark}

% !TEX root = ../nc-cubic.tex
\section{Jacobian tensor族の復元曲線と点曲線}
\label{mpc:sec:point-comparison}

第\ref{med:sec:descent}節の分裂三次模型について、各平面三次を入力因子の点曲線と比較する。
ここでは因子間に混合関係のないJacobian型の一変数族を扱う。
この追加仮定の下では、生成元固定のharmonic-angle formalityがブロックの直和を保ち、
中心写像も各因子への制限と可換する。
既存の一ブロックの全次数Hessian--Casimir恒等式を使うと、
復元potentialの各ブロックは点曲線の方程式の形式単元倍になる。
その単元をブロックごとの座標変換で吸収することで、点曲線から直接書ける
代数的な分裂方程式へ、有限拡大上の同変比較を取れる。

この節は定理\ref{cft:thm:complete}、定理\ref{gdr:thm:generic}、
補題\ref{lem:nfm-active-restriction-all-orders}を入力とする。
点曲線のHessian公式そのものはMori--Smithの既知結果
\cite[Theorems~1.2, 7.2]{MoriSmith}である。
ここでの多ブロック比較への組合せの新規性は、独立には確認していない。

\subsection{入力の範囲とprimitiveな中心三次}

$R=\mathbb C[h]_{(h)}$、$\widehat R=\mathbb C[[h]]$、$K=\mathbb C(h)$とし、$k\ge2$とする。
各$i=1,\ldots,k$について三変数$x_i=(x_{i1},x_{i2},x_{i3})$、
$g_i\in R[x_i]_3$および$0\ne\varepsilon_i\in hR$を固定する。
$\overline g_i=g_i\bmod h$は滑らかな平面三次を定めるものとする。
次数付き代数$D_i$を、巡回的に三つの関係
\begin{equation}\label{mpc:eq:factor-relations}
 [x_{i2},x_{i3}]
 =\varepsilon_i\operatorname{sym}(\partial_{x_{i1}}g_i)
\end{equation}
で定める。$\operatorname{sym}$は二次式のWeyl対称化である。
各次数成分が有限自由であるというPBW条件を仮定する。
これはこの節で量子化の存在を証明するという意味ではない。

中心三次$\Omega_i\in Z(D_i)_3$は$R$上primitive、すなわち
$\Omega_i\notin h(D_i)_3$とし、$c_i\in R^\times$を取る。
\begin{equation}\label{mpc:eq:input}
 B_R=\bigotimes_{i=1}^kD_i,
 \qquad w=\sum_{i=1}^k c_i\Omega_i,
 \qquad A_R=B_R/(w)
\end{equation}
と置く。テンソル積は通常のもので、異なるブロックの生成元は可換である。
以下のHessianは各ブロックの固定座標で計算し、
\begin{equation}\label{mpc:eq:point-polynomials}
 J_i=g_i+\frac{\varepsilon_i^2}{24}\det\operatorname{Hess}(g_i),
 \qquad E_{i,R}=V(J_i)\subset\mathbb P^2_R
\end{equation}
と書く。

\begin{lemma}[中心の先頭項と点曲線]\label{mpc:lem:primitive-point}
一意な$\lambda_i\in\mathbb C^\times$が存在して
\begin{equation}\label{mpc:eq:central-leading-term}
 \Omega_i\bmod h=\lambda_i\overline g_i
\end{equation}
となる。各$E_{i,R}$は滑らかで、その一般ファイバー$E_i$は$D_{i,K}$の点曲線である。
$E_{i,R}$はその平坦閉包であり、$h=0$の多項式環の点スキーム全体ではない。
\end{lemma}
\begin{proof}
$\varepsilon_i=h^{\nu_i}u_i$、$u_i\in R^\times$と書く。
式\eqref{mpc:eq:factor-relations}の最初の非可換項は、
$u_i(0)\iota_{d\overline g_i}
(\partial_{x_{i1}}\wedge\partial_{x_{i2}}\wedge\partial_{x_{i3}})$である。
中心性の最初の係数を取ると、$\Omega_i\bmod h$はこのJacobian Poisson構造の三次Casimirとなる。
滑らかな$\overline g_i$の三つの偏微分は正則列なので、
補題\ref{lem:nfm-relative-casimir-tensor}と同じKoszul完全性により、その微分は
$d\overline g_i$の定数倍である。Euler恒等式から式\eqref{mpc:eq:central-leading-term}を得る。
primitive性によりこの定数は非零であり、一意性も明らかである。

$p=(x_{i1},x_{i2},x_{i3})^{\mathsf t}$とし、$C(p)q=p\times q$と書く。
二次関係の双線型化は
\[
 M_i(p)q=0,
 \qquad M_i(p)=C(p)-\frac{\varepsilon_i}{2}\operatorname{Hess}(g_i)(p)
\]
である。命題\ref{prop:nfm-correction-point-cubic}の行列式計算により
\begin{equation}\label{mpc:eq:point-determinant}
 \det M_i=-3\varepsilon_i J_i.
\end{equation}
これは\cite[Theorem~7.2]{MoriSmith}のHessian公式を、本稿の交換子規約で書いたものである。
$J_i\bmod h=\overline g_i$より判別式は$R$の単元であり、$E_{i,R}$は滑らかである。
$K$上では$\varepsilon_i$を逆にできる。
右射影の行列式は$\varepsilon_i$を$-\varepsilon_i$へ替えた同じ計算で得られ、
$J_i$は$\varepsilon_i$について偶関数なので、点対応の左右の射影像はともに$E_i$となる。
曲線上で$M_i$の階数が一以下なら行列式の全偏微分が消え、滑らかさに反する。
従って階数は二であり、可逆な二次小行列式を用いて核直線を取ると、
点対応は両射影が同型であるグラフとなる。
これが点曲線を与える。
最後に$J_i$の乗法は各次数で剰余上単射なので、
単元小行列式の議論により商の各次数は$R$-自由である。
従って$V(J_i)$は一般点曲線の平坦閉包である。
\end{proof}

\subsection{生成元固定formalityの直和性}

以下では、補題\ref{lem:nfm-normalized-star-model}と同じ
harmonic-angleのKontsevich写像を全ての変数の個数で固定する。
曲率付きMF圏のtensor積公式を仮定するのではなく、
曲率を移す前の星積と中心写像を直接比較する。

\begin{lemma}[独立ブロックに関するグラフの消滅]
\label{mpc:lem:graph-direct-sum}
$\pi_i$を$x_i$だけに係数を持ち、$x_i$方向だけを微分する形式bivectorとし、
$\pi_i\equiv0\pmod h$とする。$\pi=\sum_i\pi_i$と置く。
$f_i,g_i'$を$i$番目の変数だけの多項式とすると
\begin{equation}\label{mpc:eq:factor-star-center}
 \begin{aligned}
 f_i*_{\pi}g_i'&=f_i*_{\pi_i}g_i',\\
 f_i*_{\pi}f_j&=f_if_j=f_j*_{\pi}f_i\quad(i\ne j),\\
 C_{\pi}(f_i)&=C_{\pi_i}(f_i)
 \end{aligned}
\end{equation}
が全形式次数で成立する。最後の$C_\pi$は関数上のformality接写像である。
\end{lemma}
\begin{proof}
Kontsevichのグラフ表示\cite[Sections~5.1, 6.1--6.3]{NFMKontsevich}で、
各bivector頂点をそのブロック番号で色付けする。
異なる色の内部頂点へ向かう辺は、係数を別ブロックの変数で微分するので零になる。
従って非零の項では、bivector頂点同士の辺だけからなるグラフの各連結成分は単色である。

$r>0$個のbivector頂点からなる成分には、出る辺がちょうど$2r$本ある。
この成分が接する境界頂点が$b\le1$個なら、その辺の角一形式の積は
配置空間$C_{r,b}$からの引戻しである。
しかし$\dim C_{r,b}=2r+b-2<2r$なので、この$2r$形式は零である。
境界が零個の場合は成分全体の実affine変換、境界が一個の場合は
その実点を中心とする拡大縮小が、同じ消滅を直接与える。
これは重みの和における相殺ではなく、該当グラフの積分形式自体の消滅である。

$f_i*_{\pi}g_i'$では、他の色の成分は境界に接することができないため消える。
残るグラフは三変数だけの式そのものであり、第一式を得る。
$f_i*_{\pi}f_j$、$i\ne j$では、各単色成分は境界の高々一つにしか接することができず、
正の頂点数を持つグラフは全て消える。これが第二式である。

$C_\pi(f_i)$では$f_i$は出辺が零本の指定内部頂点である。
他の色のbivector成分はこの頂点に接することもできないので、
その$2r$形式は$C_{r,0}$からの引戻しとなって消える。
従って残るグラフは全て$i$番目のブロックに属し、第三式を得る。
各固定$h$次数では有限個のグラフだけを使うので、等式は形式極限でも成立する。
\end{proof}

\begin{proposition}[tensor族の正規化模型と中心の直和]
\label{mpc:prop:normalized-direct-sum}
各$D_i\otimes_R\widehat R$の生成元固定模型を$(\pi_i,\phi_i)$とする。
式\eqref{mpc:eq:input}の周囲代数に対する生成元固定模型は
\[
 \pi=\sum_i\pi_i,
 \qquad \phi|_{D_i}=\phi_i
\]
である。$W_i=C_{\pi_i}^{-1}\phi_i(\Omega_i)$と書けば、その中心逆像は
\begin{equation}\label{mpc:eq:direct-potential}
 W=C_\pi^{-1}\phi(w)=\sum_i c_iW_i
\end{equation}
である。
\end{proposition}
\begin{proof}
補題\ref{cft:lem:model}により各因子の模型が存在し、
補題\ref{lem:nfm-normalized-star-model}により生成元を固定できる。
異なるブロックのSchouten括弧は零なので、$\sum_i\pi_i$はPoisson tensorである。
補題\ref{mpc:lem:graph-direct-sum}の最初の二式により、
各$\phi_i$の像は互いに可換な部分代数となり、tensor積から星積代数への準同型$\phi$を定める。
各内部次数でその剰余は恒等であり、有限自由加群間の完備Nakayamaにより同型である。
これは全生成元を固定するので、正規化模型の一意性により表示した$(\pi,\phi)$が得られる。

各$W_i$は$i$番目の変数だけの三次式であり、
同補題の第三式と中心写像の線型性から
\[
 C_\pi\Bigl(\sum_i c_iW_i\Bigr)
 =\sum_i c_i C_{\pi_i}(W_i)
 =\phi(w)
\]
となる。三次成分上で$C_\pi$は可逆なので式\eqref{mpc:eq:direct-potential}を得る。
\end{proof}

\subsection{Hessian三次から書ける指定代数的模型}

\begin{proposition}[各復元ブロックと点三次の全次数比較]
\label{mpc:prop:exact-point-potential}
一意な形式単元$b_i\in\widehat R^\times$、$b_i(0)=\lambda_i$が存在して
\begin{equation}\label{mpc:eq:exact-potential}
 W_i=b_iJ_i,
 \qquad W=\sum_i c_ib_iJ_i
\end{equation}
となる。従って、$h$とMFの周期変数を固定し、剰余が恒等である
ブロック対角の線型座標変換の後、完備比較のpotentialを
\begin{equation}\label{mpc:eq:algebraic-potential}
 H_{\mathrm{pt}}=\sum_i c_i\lambda_iJ_i\in R[x_1,\ldots,x_k]_3
\end{equation}
に取れる。この変更は内部shift、標準の$G=(\mu_3)^{k-1}$作用および指定三次周期を保つ。
\end{proposition}
\begin{proof}
補題\ref{lem:nfm-active-restriction-all-orders}の三変数の普遍恒等式を各因子へ用いる。
同補題の各形式係数は三次式の係数の多項式なので、係数へ$g_i(h)$を代入し、
その形式パラメータへ$\varepsilon_i(h)$を代入してよい。
$\varepsilon_i\in hR$であるため、この代入は各$h$次数で有限和となる。
生成元固定再帰の一意性により、得られるtensorは上の$\pi_i$に一致する。
従って
\[
 \iota_{dJ_i}\pi_i=0,
 \qquad \pi_i/\varepsilon_i\equiv
 \iota_{d\overline g_i}
 (\partial_{x_{i1}}\wedge\partial_{x_{i2}}\wedge\partial_{x_{i3}})
 \pmod h
\]
である。$\pi_i$の$\varepsilon_i$による可除性も、同じ普遍再帰の代入から従う。
$J_i$の三つの偏微分は、正則な特殊ファイバーの勾配の完備変形である。
補題\ref{lem:nfm-relative-casimir-tensor}により
\[
 \pi_i=\varepsilon_i u_i(h)\,\iota_{dJ_i}
 (\partial_{x_{i1}}\wedge\partial_{x_{i2}}\wedge\partial_{x_{i3}}),
 \qquad u_i(h)\in1+h\widehat R
\]
と書ける。

$W_i$は$\pi_i$のCasimirである。$\widehat R$は整域で
$\varepsilon_i u_i\ne0$なので、$dJ_i\wedge dW_i=0$を得る。
勾配Koszul複体の完全性と内部次数から$dW_i=b_i\,dJ_i$、$b_i\in\widehat R$となり、
Euler恒等式により$W_i=b_iJ_i$である。
補題\ref{mpc:lem:primitive-point}と$C_{\pi_i},\phi_i\equiv1\pmod h$から
$b_i(0)=\lambda_i\ne0$となる。これで単元性と一意性を得る。
命題\ref{mpc:prop:normalized-direct-sum}が全potentialの式を与える。

各$i$について$s_i=(\lambda_i/b_i)^{1/3}\in1+h\widehat R$を取る。
剰余が一という条件で、この形式三乗根は一意である。
三次斉次性により
\[
 W(s_1x_1,\ldots,s_kx_k)=\sum_i c_i\lambda_iJ_i=H_{\mathrm{pt}}.
\]
この座標変換はpotentialを正確に移し、ブロックごとの三乗根作用と可換する。
周期変数には作用しないので、指定周期を含む主張が従う。
\end{proof}

\begin{remark}[中心生成元の自由度]
三次Casimirは上のKoszul計算から$\widehat R J_i$に一致する。
従って$C_{\pi_i}^{-1}\phi_i$は
$Z(D_i\otimes_R\widehat R)_3$をこの直線へ同定し、primitiveな$\Omega_i$はその生成元となる。
中心性は三つの生成元との交換子からなる有限行列の核で判定できるので、
完備化の忠実平坦性により$Z(D_i)_3=R\Omega_i$でもある。
別のprimitive生成元との差は$R^\times$であり、$b_i$へ吸収される。
従って特定の高次Hesse中心公式を選ぶことは、この平面曲線の同定には必要ない。
ただし中心元の指定を保持した圏比較では、上のpotentialを正確に移す座標変換を使う。
\end{remark}

\begin{theorem}[点三次の和への同変一般点比較]
\label{mpc:thm:point-model}
式\eqref{mpc:eq:input}の仮定の下で
\[
 Y_{\mathrm{pt},R}=V(H_{\mathrm{pt}})\subset\mathbb P^{3k-1}_R
\]
は滑らかな三次超曲面族であり、その特殊方程式は$w\bmod h$に等しい。
ある有限拡大$K'/K$上で
\begin{equation}\label{mpc:eq:point-model-equivalence}
 (\mathcal K^{\rm coh}_{A_{K'}},\rho_A,\tau_A)
 \simeq_G
 (\operatorname{Ku}(Y_{\mathrm{pt},K'}),\rho_{\mathrm{pt}},\tau_{H_{\mathrm{pt}}})
\end{equation}
となる。ここで各点曲線$E_i$の方程式$J_i$は、式\eqref{mpc:eq:algebraic-potential}の
対応する三変数ブロックそのものである。
\end{theorem}
\begin{proof}
$c_i\lambda_i$は単元であり、各$J_i$は滑らかな平面三次を定める。
勾配がブロックごとに分かれるので、全勾配の共通アフィン零点は原点だけである。
従って$Y_{\mathrm{pt},R}$は滑らかであり、特殊方程式の主張は
式\eqref{mpc:eq:central-leading-term}から従う。

定理\ref{gdr:thm:generic}を適用できることは
補題\ref{mpc:lem:primitive-point}およびテンソル積の次数別自由性から分かる。
これで一般点の全coherent残余圏のsmooth proper性を確保する。
一方、命題\ref{mpc:prop:exact-point-potential}は、標的を
既に$R$上で定義された$H_{\mathrm{pt}}$に固定した完備比較を与える。
生成元固定模型の共変性と第\ref{med:sec:descent}節の同変bar flowにより、これは$G$同変である。
固定した有限個の対象とその$G$軌道上のHomを先に構成してから局所化し、
非零EulerベクトルとBFKの生成性を用いると、$\overline{\mathbb C((h))}$上で
全残余圏の同変Morita同値を得る。

ここで両方のdg模型は$K$上にある。
補題\ref{med:lem:equivariant-descent}を、標的$Y_{\mathrm{pt},K}$を固定して適用すればよい。
有限型模型の閉点で特殊化するのは比較核と比較図式であり、
点曲線の係数$J_i$を別の係数へ特殊化する操作は行わない。
従って式\eqref{mpc:eq:point-model-equivalence}を得る。
\end{proof}

\subsection{復元方程式の選択に依存しない各ブロックの同型類}

\begin{corollary}[任意の同変復元ブロックの$j$]
\label{mpc:cor:block-j}
上の入力について、第\ref{med:sec:descent}節の方法により
$Y_F=V(F_1+\cdots+F_k)$を$K$上の分裂模型として選び、$C_i=V(F_i)$とする。
同じ標準$G$作用とrotationの同変構造を保持することを要求する。
このとき各$i$について
\begin{equation}\label{mpc:eq:block-j}
 (\mathbb P^2_{\overline K},C_{i,\overline K})
 \simeq(\mathbb P^2_{\overline K},E_{i,\overline K}),
 \qquad j(C_i)=j(E_i)\in K
\end{equation}
である。同じ結論は、有限拡大上でそのような同変rotation付き比較を持つ任意の分裂模型に成り立つ。
\end{corollary}
\begin{proof}
定理\ref{mpc:thm:point-model}と$Y_F$への比較を、共通の代数閉拡大体上で合成する。
これにより二つの可換Kuznetsov pairの間の同変rotation付き同値を得る。
低rotation積の同定により、次数一空間の$G$同変線型同型は勾配空間を移す。
両方の三次式を同じ次数一空間へ移し、
補題\ref{med:lem:invariant-symmetrizer}を適用すると、
さらに$G$と可換する座標変換で両方の三次式が一致する。

標準$G$作用の$k$個のブロック指標は相異なるので、
これらの$G$同変線型同型はブロック対角である。
従って$i$番目の平面三次は$i$番目の点曲線へ移る。
同型を与える有限個の行列係数はある有限型$K$-代数上にあり、
その同型schemeの幾何点の存在から$\overline K$上でも同型となる。
両$j$は$K$に属し、体拡大後に等しいので$K$上で等しい。
\end{proof}

\begin{corollary}[順序付き平面同型の有限被覆]
\label{mpc:cor:plane-isom-cover}
前系でさらに$j(E_i)\notin\{0,1728\}$を仮定する。
順序付きの平面同型のscheme
\[
 I=\prod_{i=1}^k
 \underline{\operatorname{Isom}}_{\rm plane}
 \bigl((\mathbb P^2_K,E_i),(\mathbb P^2_K,C_i)\bigr)
\]
は$K$上有限\'etaleかつ階数$18^k$であり、その上で各平面同型を選べる。
因子の$j$が相異なるという条件はこの階数には不要である。
\end{corollary}
\begin{proof}
代数閉包上で同じ$j$を持つ種数一曲線は同型である。
両方の次数三線束の差は、平行移動によるPicard類の変化が三倍写像であることから吸収できる。
従って平面同型が存在する。
一つの滑らかな平面三次の射影自己同型群は、原点をflexに選ぶと
$E[3]\rtimes\operatorname{Aut}(E,0)$である。
$j\ne0,1728$では後者は$\{\pm1\}$なので、群schemeは有限\'etale階数十八である。
平面同型のschemeはこの群のtorsorであり、
各因子が階数十八となる。積を取れば主張を得る。
この議論は命題\ref{prop:nfm-point-residual-isom-cover}の曲線ごとの計算と同じである。
\end{proof}

\begin{corollary}[二ブロック同変残余圏の点曲線積]
\label{mpc:cor:point-product}
$k=2$の場合、ある有限拡大$K''/K$上で
\[
 (\mathcal K^{\rm coh}_{A_{K''}})^{\mu_3}
 \simeq D^b(E_{1,K''}\times E_{2,K''})
\]
となる。必要ならさらに有限拡大して各$E_i$に原点を選べる。
\end{corollary}
\begin{proof}
定理\ref{mpc:thm:point-model}に補題\ref{med:lem:hu-field}を適用する。
比較先の分裂三次の二曲線は、ここでは点曲線$E_1,E_2$そのものである。
\end{proof}

\subsection{Hesse表示と比較の限界}

例えば$g_i=p_i(x_{i1}^3+x_{i2}^3+x_{i3}^3)/3+q_ix_{i1}x_{i2}x_{i3}$なら、
式\eqref{mpc:eq:point-polynomials}は
\[
 \begin{aligned}
 J_i&=A_i(x_{i1}^3+x_{i2}^3+x_{i3}^3)/3+B_ix_{i1}x_{i2}x_{i3},\\
 A_i&=p_i\left(1-\frac{\varepsilon_i^2q_i^2}{4}\right),
 &B_i&=q_i+\frac{\varepsilon_i^2(4p_i^3+q_i^3)}{12}
 \end{aligned}
\]
と書ける。$A_i\ne0$の範囲で$\mu_i=B_i/A_i$と置くと、
式\eqref{gre:eq:hesse-j}より
\[
 j(C_i)=j(E_i)
 =\frac{27\mu_i^3(8-\mu_i^3)^3}{(1+\mu_i^3)^3}
\]
である。この式は個々の中心元の高次係数を計算せずに得られる。

本節で代数的に固定したのは、式\eqref{mpc:eq:algebraic-potential}の方程式と、
各点曲線の平面同型類である。中心写像が与える形式単元$b_i$、
その三乗根、および指定完備比較核そのものの代数性は主張しない。
有限被覆$I$は各平面同型を選ぶ被覆であり、比較核の指定された同型類を選ぶ被覆ではない。
点曲線上のshiftに対応する平行移動切断と比較核との整合も、別の問題である。

また、入力のブロック分解と標準$G$作用を用いている。
任意の全$\mathcal O$-保存同値がこの$G$作用やrotationの同変構造を保つとは証明していないため、
この比較をそのまま全$\mathcal O$-markedモジュライ上の射とは呼ばない。
一般の混合多ブロック関係や、式\eqref{mpc:eq:factor-relations}で表示されていない変形への拡張も含めない。
滑らかな方程式族$Y_{\mathrm{pt},R}$の存在と、特殊点を含む$R$上の相対同値の存在は区別する。

% !TEX root = ../nc-cubic.tex
\section{残余圏の降下と分離和の同変積表示}
\label{mbrd:sec:residual-descent}

本節では、残余圏と比較核の降下を区別する。
系\ref{mbrr:cor:stack-isomorphism}の全O開では、次数付き表示の射影商から
compact残余圏とその内在的中心作用を降ろせる。
二つの残余dg族が与えられれば、それらを比較する同変perfect核の
モジュライも降下する。しかし、このモジュライに大域切断が存在すること、
または指定された形式核がその切断から来ることは別の問題である。

具体的な可換標的については、Ballard--Favero--Katzarkovの
Morita積公式\cite[Definition~3.50, Corollary~5.18]{NFBFKKernels}と
Orlov比較のdg版\cite[Theorem~6.13]{NFBFKKernels}を反復する。
これにより、$k$個の三変数三次式の分離和に対し、
中心群による同変Kuznetsov圏を$k$個の種数一曲線の積の導来圏として記述する。
二ブロックのcategorical coveringはすでに
\cite[Example~4.7]{MBRDBFKFurther}に現れ、
Hu\cite[Theorems~0.1, 5.8]{TBHu}もこの型の同変Kuznetsov圏を扱う。
以下の可換積表示は既知の一般定理の本設定への適用である。
本節の非可換側への帰結では、定理\ref{med:thm:multiblock-generic}、
定理\ref{mpc:thm:point-model}の仮定をそのまま保持する。
一変数の仮定を外す体上の直接比較は
第\ref{mbnp:sec:product}節、第\ref{mbrc:sec:rotation-comparison}節で別に証明する。

\subsection{compact残余圏と比較核のモジュライ}

$\mathcal M_k^\circ\simeq\mathcal P_k^\circ$を
系\ref{mbrr:cor:stack-isomorphism}の非空開とし、
\[
 Z_0=Z(K_k)=(\mu_3^k)/\Delta\mu_3
\]
をその中心慣性とする。順序を忘れた底では、右辺は因子置換で降下する
有限\'etale群の局所系と解釈する。

\begin{corollary}[全O開上の同変compact残余圏]
\label{mbrd:cor:residual-family}
表示チャート上で
\[
 \mathcal C_A=\operatorname{thick}\{\mathcal O_A(j):j\in\mathbb Z\},
 \qquad
 \mathcal K_A^{\mathrm{line}}
 =\langle\mathcal O_A,\ldots,\mathcal O_A(3k-4)\rangle^\perp
    \cap\mathcal C_A
\]
と置く。$\mathcal C_A$、$\mathcal K_A^{\mathrm{line}}$、
残余projector、および$\mathcal K_A^{\mathrm{line}}$のcoherentな$Z_0$作用と
その同変化は、$\mathcal M_k^\circ$上に降下する。
この主張だけから全coherent圏との一致やsmooth性は結論しない。
\end{corollary}
\begin{proof}
定理\ref{mbsd2:thm:projector-descent}により、相対evaluation
\[
 R\operatorname{Hom}(\mathcal O_A(j),-)\otimes\mathcal O_A(j)
       \longrightarrow\operatorname{Id}
\]
とそのdg cofiberを用いる残余projectorは共通scalarの重み零を持つ。
実際、$\mathcal O_A(j)$を底の直線束の$j$乗で変えると、Hom側の逆重みと相殺する。
従って射影的次数付き表示の重なりでcoherentに貼り合わさる。
全O射との一致は系\ref{mbrr:cor:stack-isomorphism}による。
中心作用と基準ブロックを変えたときの比較のcoherenceは
定理\ref{mbza:thm:central-action}で得られ、projectorを保つ。
同変化はこの作用図式のhomotopy fixed pointとして取る。
局所的にはcrossed productのdg加群模型であり、
降下のhomotopy limitと群作用のhomotopy limitは可換する。
\end{proof}

次の命題では、降下する対象を明示するため、比較の両側をあらかじめ固定する。
これは$\mathcal M_k^\circ$上の可換復元先を自動的に構成する命題ではない。

\begin{proposition}[同変Morita核のstackと有限被覆]
\label{mbrd:prop:kernel-stack}
$S$を標数零の体上のschemeまたは代数stackとする。
$\mathcal C,\mathcal D$を$S$上のdg圏の族で、fpqc局所に
smooth properなdg代数によって表されるものとする。
$G/S$を有限\'etale群とし、両側にcoherentな$G$作用を与える。
各$T\to S$に対し、$G_T$同変なMorita同値
$\mathcal C_T\simeq\mathcal D_T$を表すperfect双加群の最大$\infty$-groupoidを
$\mathscr E_G(T)$と置く。このとき次が成り立つ。
\begin{enumerate}
\item $\mathscr E_G$はfpqc降下を満たす。
従って、$\mathcal C,\mathcal D$と群作用が射影的表示から降下していれば、
比較核のstack自体も同じ商へ降下する。
\item 両側に降下済みのrotation双加群と指定周期図式を与えた場合、
これらと両立する比較同型およびcoherenceを核のデータに含めても、
同じ降下の主張が成り立つ。
\item $S$が標数零の体上有限型な整なschemeとする。
この項ではさらに、両一般ファイバーがある有限代数拡大後に
実際のsmooth properなdg代数模型を持つと仮定する。
幾何学的一般点で$\mathscr E_G$が非空なら、
ある非空開$S^\circ\subset S$と有限\'etale全射
$T^\circ\to S^\circ$上に実際の同変Morita核が存在する。
同じ結論は、前項の有限個のrotation・周期比較を指定した場合にも成り立つ。
\end{enumerate}
\end{proposition}
\begin{proof}
局所的に$G$を定数群とし、群軌道の直和を生成対象に用いて作用を厳密化する。
そこで
\[
 \mathsf E=\mathsf C^{\mathrm{op}}\otimes_S^{\mathbf L}\mathsf D
\]
と置けば、群作用の全coherenceを持つ双加群は
$\mathsf E\rtimes G$上のdg加群として表せる。
補題\ref{med:lem:crossed-perfect}により、そのperfect性は忘却後のperfect性と同値である。
従って群の乗法則を核の同型類の不変性で代用してはいない。

dg加群のfpqc降下は、複体の忠実平坦降下を作用写像とともに適用して得る。
導来mapping spaceを含む形では
\cite[Lemma~3.1とその証明]{ROGTV}の構成を用いる。
そこで使われる忠実平坦なcohomological descentはfpqc被覆にも適用でき、
perfect性はその底変更に関して局所的である。
smooth proper性の下ではperfectとpseudo-perfectの判定は
\cite[Lemma~2.8]{ROGTV}により一致する。
Morita同値であることも局所的である。
実際、核$P$の右随伴核$Q$を取り、単位と余単位
\[
 \mathsf C\longrightarrow P\otimes_{\mathsf D}^{\mathbf L}Q,\qquad
 Q\otimes_{\mathsf C}^{\mathbf L}P\longrightarrow\mathsf D
\]
のconeの消滅を検査すればよく、忠実平坦底変更はこれを検出する。
これで第1項を得る。第2項の比較射とcoherentな図式も
同じmapping space内のデータであり、そのhomotopy limitとして降下する。
三角圏の非関手的なconeの選択を貼り合わせる議論ではない。

第3項では$K=k(S)$とし、仮定した両側のdg代数模型を得る有限拡大を先に取る。
必要ならさらに有限拡大によって$G_K$を定数群にし、
生成対象の有限軌道和によって作用を厳密化する。
幾何学的一般点上の核に
補題\ref{med:lem:equivariant-descent}を適用して、
ある有限拡大$K'/K$上に同変核を取る。
perfect核の有限表示性は\cite[Lemma~2.10]{ROGTV}による。
$S$を非空開に縮めて正規とし、$K'$における正規化$T\to S$を取る。
これは有限であり、標数零なので非空開上\'etaleとなる。
核、その逆、単位・余単位と、指定した有限個の比較射は$T$の非空開へ広がる。
群coherenceは最初からcrossed product加群の構造に含まれている。
比較射が同型であることと、有限個の周期図式の等式は、さらに縮めれば保たれる。
最後に除いた閉集合の$S$への像を取り除く。
有限射は閉写像なので、得られた非空開の全逆像上で核が定義され、
その射は有限\'etale全射である。
本稿の定理\ref{med:thm:multiblock-generic}と
定理\ref{mpc:thm:point-model}では、両側の実際のdg模型が
もとの関数体上で与えられているので、この追加条件を満たす。
\end{proof}

\begin{remark}[核の型、gerbe、切断]
\label{mbrd:rem:kernel-section}
命題\ref{mbrd:prop:kernel-stack}は核の集合が一つのgerbeになるとは述べていない。
異なる同値核には、自己同値を合成した互いに同型でない核も含まれる。
局所的に同型な核の型を一つ固定し、その型の核$P$について
任意底変更後も
$\operatorname{Ext}^{<0}(P,P)=0$、
$\operatorname{End}(P)=\mathcal O$
を仮定した場合に限り、その型のformsは$\mathbb G_m$-gerbeとなる。
このgerbeの大域切断は、その中立化を意味し、降下するstackの存在とは別である。
例えば体$\kappa$上の非自明な中心単純代数$D$について、
$\operatorname{Perf}(\kappa)$と$\operatorname{Perf}(D)$を比較する核は
分裂体上で存在するが、$\kappa$上のMorita核が存在するのは$[D]=0$の場合だけである。
従って局所的な核の存在だけで大域切断を結論することはできない。
本研究で得る有限拡大上の存在も、
注意\ref{med:rem:formal-versus-existence}のとおり指定形式核の全jetの代数化とは異なる。
\end{remark}

\begin{remark}[rotationを付加する前提]
\label{mbrd:rem:rotation-framing}
命題\ref{mbrd:prop:kernel-stack}(2)では、rotation自体が両側で降下していること、
または同じ追加の標識空間上にあることを仮定した。
共通scalarに対するrotationの重みは一であり、
その三乗は関係直線の双対を伴う周期
$\rho^3\otimes\ell^\vee\simeq[2]$を持つ
（命題\ref{mbsd2:prop:scalar-weight}、命題\ref{mbsd2:prop:line-period}）。
また中心群のsplitを変えると、rotationの線型化には
命題\ref{mbza:prop:shift-character}のcharacter twistが生じる。
従って、標識なしの核stackの降下から、
選択済みのrotation比較や周期比較の降下を追加の確認なしに推論しない。
命題\ref{mbza:prop:splittings}の有限\'etaleなsplit選択被覆は
この線型化の選択を扱うが、核のgerbeを自動的に中立化するものではない。
さらにsourceとtargetの表示gerbeを独立に取ると、
$P\rho_{\mathrm s}\to\rho_{\mathrm t}P$のHomはscalar
$(\lambda_{\mathrm s},\lambda_{\mathrm t})$に対して
$\lambda_{\mathrm t}/\lambda_{\mathrm s}$の重みを持つ
（注意\ref{mbsd2:rem:kernel-synchronization}）。
この相対重みは共通対角scalarへの同期により消えるが、
同期の選択自体から比較射や核の大域切断が存在するわけではない。
\end{remark}

\subsection{分離和の同変Kuznetsov圏}

以下の可換比較では、まず$\kappa$を代数閉な標数零の体とする。
$k\ge2$とし、互いに独立な三変数組$x_i=(x_{i1},x_{i2},x_{i3})$に対し
$f_i(x_i)$を滑らかな平面三次を定める斉次式とする。
\[
 E_i=V(f_i)\subset\mathbb P^2_\kappa,\qquad
 F=\sum_{i=1}^kf_i,\qquad
 X=V(F)\subset\mathbb P^{3k-1}_\kappa
\]
と置く。各$f_i$の勾配の共通零点は原点だけなので、$X$は滑らかである。
$Z_0=(\mu_3^k)/\Delta\mu_3$は各ブロックの三乗根倍によって$X$に作用する。

重みを共有する対角化可能群を
\begin{equation}\label{mbrd:eq:gamma}
 \Gamma_k=
 \{(\lambda_1,\ldots,\lambda_k)\in\mathbb G_m^k:
      \lambda_1^3=\cdots=\lambda_k^3\},\qquad
 \chi(\lambda)=\lambda_i^3
\end{equation}
と定める。これは三乗characterによる$k$個の$\mathbb G_m$のfiber productであり、
\[
 1\longrightarrow\mathbb G_{m,\mathrm{diag}}
  \longrightarrow\Gamma_k\longrightarrow Z_0\longrightarrow1
\]
を持つ。基準ブロック$b$を選ぶと、
$\lambda\mapsto(\lambda_b,(\lambda_i/\lambda_b)_i)$により
$\Gamma_k\simeq\mathbb G_m\times Z_0$となる。
第2因子の代表では$b$番目の座標を一とする。
この分裂で得るweak actionは
命題\ref{mbza:prop:gamma-action}の内在的中心作用と一致し、
分裂を変えた比較はscalarの自然自明化によってcoherentである。

\begin{theorem}[分離和の中心同変圏の積表示]
\label{mbrd:thm:split-product}
上の仮定の下で、dg Morita同値
\begin{equation}\label{mbrd:eq:split-product}
 \operatorname{Ku}(X)^{Z_0}
 \simeq
 D^b(E_1\times\cdots\times E_k)
\end{equation}
が存在する。左辺の作用は上で指定したものとし、
同変化にはenhancementと冪等完備化を含める。
この同変圏は基準ブロック$b$に依存しない。
\end{theorem}
\begin{proof}
$\operatorname{MF}_{\mathrm{gr}}(f_i)$を標準整数次数による
有限rank MFのdg圏のMorita完備化とし、
$\operatorname{MF}_{\Gamma_k}(F)$を
式\eqref{mbrd:eq:gamma}の重みによるMF圏とする。
Ballard--Favero--Katzarkovの外積は
\begin{equation}\label{mbrd:eq:thom-sebastiani}
 \operatorname{MF}_{\mathrm{gr}}(f_1)
     \mathbin{\circledast}\cdots\mathbin{\circledast}
 \operatorname{MF}_{\mathrm{gr}}(f_k)
 \simeq \operatorname{MF}_{\Gamma_k}(F)
\end{equation}
を与える\cite[Definition~3.50, Corollary~5.18]{NFBFKKernels}。
ここで$\circledast$は導来tensor積のperfect加群を取るMorita積である。
適用条件を確認すると、周囲は滑らかなアフィン空間、作用群は対角化可能でreductive、
共通の三乗characterは非torsionである。
同系のsupport付き圏は各段階でsupportを外せる。
実際、部分和の勾配も対応するブロックの原点でのみ消えるため、
和の零超曲面の特異点は各項の零集合の積に含まれる。
その外ではMFはacyclicである。

各三変数三次ではGorenstein parameterが零なので、
Orlov比較により
$\operatorname{MF}_{\mathrm{gr}}(f_i)\simeq D^b(E_i)$となる。
全体ではparameterが$3k-3$であり、
$\operatorname{MF}_{\mathrm{gr}}(F)\simeq\operatorname{Ku}(X)$である。
これらはdg比較として取れる
\cite[Theorem~6.13]{NFBFKKernels}。
標準windowはブロックの線型作用と可換するので、この全体の比較も$Z_0$同変である。
基準$b$による$\Gamma_k$の分裂を用いれば
\[
 \operatorname{MF}_{\Gamma_k}(F)
 \simeq \operatorname{MF}_{\mathrm{gr}}(F)^{Z_0}
 \simeq \operatorname{Ku}(X)^{Z_0}.
\]
群coherenceはequivariant factorizationそのものに含まれる。
最後にsmooth properなschemeのperfect圏のMorita積は積schemeのperfect圏であり
\cite[式~(1.1)]{MBRDBFKFurther}、
式\eqref{mbrd:eq:thom-sebastiani}から主張を得る。
基準変更への独立性は、命題\ref{mbza:prop:gamma-action}のcoherentな比較による。
\end{proof}

\subsection{積上のcharacter作用と双対同変化}

符号を固定するため、次数付き加群では$M(1)_n=M_{n+1}$、
$\Gamma_k$の加群上の重みを次数とする。
$e_i:\Gamma_k\to\mathbb G_m$を$e_i(\lambda)=\lambda_i$と書き、
$\mathsf T_\eta$をcharacter $\eta$をtensorする関手とする。
この規約で$i$番目の内部shift $(1)$は$\mathsf T_{e_i^{-1}}$に対応する。
$e_i^3=\chi$なので、共通周期は
$\mathsf T_{\chi^{-1}}\simeq[2]$である。

各楕円因子上で
\[
 R_{E_i}=T_{\mathcal O_{E_i}}\circ
       (-\otimes\mathcal O_{E_i}(1)),\qquad
 T_{\mathcal O_{E_i}}(M)=
 \operatorname{Cone}\bigl(
 R\operatorname{Hom}(\mathcal O_{E_i},M)\otimes\mathcal O_{E_i}\to M
 \bigr)
\]
と置き、積の第$i$因子にこの関手を作用させるものを$R_i$と書く。
$\mathcal O_{E_i}$はsphericalであり、ここでの$T_{\mathcal O_{E_i}}$は
exceptional対象のmutationではなくspherical twistである。

\begin{proposition}[characterと因子rotation]
\label{mbrd:prop:factor-rotations}
式\eqref{mbrd:eq:split-product}の同値を、
外積と上のOrlov windowから構成する。
このとき$\mathsf T_{e_i^{-1}}$は$R_i$に対応し、
\begin{equation}\label{mbrd:eq:character-rotation}
 \mathsf T_{e_i/e_b}\quad\longleftrightarrow\quad R_i^{-1}R_b,
 \qquad R_i^3\simeq[2].
\end{equation}
左辺の$e_i/e_b$は$\Gamma_k/\mathbb G_{m,\mathrm{diag}}=Z_0$のcharacterである。
従って$\widehat Z_0=\operatorname{Hom}(Z_0,\mathbb G_m)$は、
これらの比の関手を通じて積の導来圏にcoherently作用する。
基準$b$による標準rotationの線型化は積側の$R_b$に対応する。
\end{proposition}
\begin{proof}
外積の定義\cite[Definition~3.51]{NFBFKKernels}は、
各因子のcharacter twistをそのまま$\Gamma_k$への制限characterのtwistへ送る。
各因子のOrlov比較では内部$(1)$が
$T_{\mathcal O_{E_i}}\circ(-\otimes\mathcal O_{E_i}(1))$に対応する
\cite[Theorem~6.13]{NFBFKKernels}。
二因子でのこのrotation比は\cite[Example~4.12]{MBRDBFKFurther}でも明示されている。
従って最初の主張と式\eqref{mbrd:eq:character-rotation}を得る。
periodと群法則のcoherenceはcharacterによる厳密な作用から移送する。
各$R_i^3\simeq[2]$という同型類の等式だけを使って群作用を作るのではない。

基準$b$の分裂では$e_b$は整数次数のcharacterである。
従ってその逆characterによるtwistは線型化された内部$(1)$となる。
基準$c$へ移すと
$e_b^{-1}/e_c^{-1}=e_c/e_b$であり、
これは$\chi_{bc}=z_b/z_c$に対して$\chi_{bc}^{-1}$である。
従って命題\ref{mbza:prop:shift-character}の
$J_{bc}\rho_b\simeq\mathsf T_{\chi_{bc}^{-1}}\rho_cJ_{bc}$と一致する。
\end{proof}

\begin{corollary}[対称分裂に対応する積側のrotation]
\label{mbrd:cor:symmetric-rotation}
固定した上の可換分離和について$3\nmid k$とする。
$c\in\{1,2\}$を$ck\equiv1\pmod3$で定め、$m=(ck-1)/3$と置くと、
\[
 t:\Gamma_k\longrightarrow\mathbb G_m,\qquad
 t(\lambda)=\left(\prod_{i=1}^k\lambda_i\right)^c\chi(\lambda)^{-m}
\]
は対角$\mathbb G_m$への因子置換不変なretractionである。
この対称分裂で線型化された標準rotationは、積側では
\begin{equation}\label{mbrd:eq:symmetric-rotation}
 R_{\mathrm{sym}}=(R_1\cdots R_k)^c[-2m],
 \qquad R_{\mathrm{sym}}^3\simeq[2]
\end{equation}
となる。特に$k=2$では
$R_{\mathrm{sym}}\simeq(R_1R_2)^{-1}[2]$である。
これは固定した可換表示の有限中心群の対称分裂に関する記述であり、
実表示の$\mathbb G_m$-gerbeの中立化を主張しない。
\end{corollary}
\begin{proof}
$ck-3m=1$なので、対角上で$t(\lambda,\ldots,\lambda)=\lambda$、
また$\Gamma_k$上で$t^3=\chi$である。
$\mu_3^k$上では$t(z)=(\prod_i z_i)^c$となるので、
$z_i/t(z)$による分裂は、命題\ref{mbza:prop:splittings}の
定数ベクトル$a_i=c$による対称分裂に一致する。
内部$(1)$は$\mathsf T_{t^{-1}}$に対応し、
\[
 t^{-1}=(e_1^{-1}\cdots e_k^{-1})^c\chi^m
\]
である。命題\ref{mbrd:prop:factor-rotations}と
$\mathsf T_{\chi^m}\simeq[-2m]$から式\eqref{mbrd:eq:symmetric-rotation}の最初の式を得る。
その三乗は$[2ck-6m]=[2]$であり、周期のcoherenceは
characterの恒等式$t^3=\chi$から移送される。
$k=2$では$c=2,m=1$、$(R_1R_2)^3\simeq[4]$なので
$(R_1R_2)^2[-2]\simeq(R_1R_2)^{-1}[2]$となる。
\end{proof}

\begin{corollary}[双対作用による通常のKuznetsov圏]
\label{mbrd:cor:dual-action}
命題\ref{mbrd:prop:factor-rotations}のcoherentな$\widehat Z_0$作用を用いると、
\begin{equation}\label{mbrd:eq:dual-action}
 \operatorname{Ku}(X)\simeq
 D^b(E_1\times\cdots\times E_k)^{\widehat Z_0}
\end{equation}
となる。この作用は一般には曲線積のscheme自己同型による作用ではなく、
character twistを移送した導来自己同値の作用である。
\end{corollary}
\begin{proof}
有限abelian群に対する二重同変化
$(\mathcal C^{Z_0})^{\widehat Z_0}\simeq\mathcal C$
\cite[Theorem~4.2]{MEDElagin}を
$\mathcal C=\operatorname{Ku}(X)$へ適用する。
標数零、冪等完備性、additive性の仮定を満たす。
character作用と比較はdg模型で構成され、
\cite[Lemma~8.6, Theorems~8.7, 8.9]{MEDElagin}により三角構造も保持される。
式\eqref{mbrd:eq:split-product}でこの作用を移送すればよい。
\end{proof}

\subsection{非可換一般点への適用と降下の範囲}

\begin{lemma}[任意の定義体上の有限拡大後の積表示]
\label{mbrd:lem:field-product}
$\kappa$を標数零の体とし、三乗根を含むとする。
$f_i$を$\kappa$上の幾何学的に滑らかな三変数三次式とする。
ある有限拡大$\kappa'/\kappa$上で
\[
 \operatorname{Ku}(V(\textstyle\sum_i f_i)_{\kappa'})^{Z_0}
 \simeq D^b\bigl(\textstyle\prod_i V(f_i)_{\kappa'}\bigr)
\]
となる。曲線の係数はこの操作で変更しない。
\end{lemma}
\begin{proof}
代数閉包$\overline\kappa$上に定理\ref{mbrd:thm:split-product}を適用する。
両側の$\kappa$-dg圏はsmooth properである。
左辺は滑らかな射影超曲面のadmissible残余圏の有限群同変化なので、
補題\ref{med:lem:crossed-perfect}を使える。
右辺は滑らかな射影schemeのperfect圏である。
得られた比較核へ、群を自明とした
補題\ref{med:lem:equivariant-descent}を適用すればよい。
降下するのは核であり、方程式$f_i$は元の$\kappa$上に固定したままである。
ここではBFKの代数閉体という仮定を無断で取り除かず、
有限拡大上での存在だけを結論した。
\end{proof}

\begin{corollary}[全ブロック数の非可換同変残余圏]
\label{mbrd:cor:nc-product}
定理\ref{med:thm:multiblock-generic}の入力
式\eqref{med:eq:multiblock-input}と、そこで用いる
定理\ref{cft:thm:complete}、定理\ref{gdr:thm:generic}の下で、
ある有限拡大$K''/K$上に
\begin{equation}\label{mbrd:eq:nc-product}
 \bigl(\mathcal K^{\mathrm{coh}}_{A_{K''}}\bigr)^{Z_0}
 \simeq D^b(C_{1,K''}\times\cdots\times C_{k,K''})
\end{equation}
が存在する。$C_i=V(F_i)$は同定理の復元曲線である。
さらに定理\ref{mpc:thm:point-model}の明示的Jacobian tensor/PBW入力を仮定すると、
右辺を実際のSklyanin点曲線の積
$D^b(E_{1,K''}\times\cdots\times E_{k,K''})$と取れる。
\end{corollary}
\begin{proof}
定理\ref{med:thm:multiblock-generic}の$G=(\mu_3)^{k-1}$は、
最後のブロックを基準にした$Z_0$の分裂表示である。
その同変比較を同変化し、補題\ref{mbrd:lem:field-product}と合成する。
二つの有限拡大を含む有限拡大を$K''$と取れば式\eqref{mbrd:eq:nc-product}を得る。
点曲線の場合は、定理\ref{mpc:thm:point-model}の標的を用いる。
その式は$\sum_i c_i\lambda_iJ_i$であり、
$c_i\lambda_i\ne0$なので各平面曲線は$V(J_i)=E_i$そのものである。
\end{proof}

系\ref{mbrd:cor:nc-product}は二ブロックの
系\ref{med:cor:elliptic-product}と系\ref{mpc:cor:point-product}を、
同じ入力の範囲で全$k\ge2$へ拡張する。
全O開上では左辺のcompact版は
系\ref{mbrd:cor:residual-family}により内在的に定まる。
しかし、このことと一般点上の積表示を合わせただけでは、
全O開上の比較核の大域切断、
任意の固定非零パラメータでの比較、または特殊点を含む相対比較は得られない。
点曲線積の族は命題\ref{mbfm:prop:product-family}で、
自然な可換三次への射と表示 gerbe・関係直線の同期は
定理\ref{mbpl:thm:commutative-family}の非空開で別に構成する。
可換側の実表示族では定理\ref{mbrp:thm:relative-product}が相対比較核を与える。
これらを非可換側と結ぶ相対核の存在は、一般点比較から自動的には従わない。
そのような族の両側が構成され、幾何学的一般点の比較が確認された場合に限り、
命題\ref{mbrd:prop:kernel-stack}(3)を有限\'etale被覆上の実際の核へ進むために使える。

% !TEX root = ../nc-cubic.tex
\section{実表示上の相対曲線積比較}
\label{mbrp:sec:relative-product}

定理\ref{mbrd:thm:split-product}の体上の比較を、
順序付きのベクトル束と実際の方程式を持つ族へ相対化する。
相対 Orlov 比較には Ballard--Deliu--Favero--Isik--Katzarkov の
具体的な hypersurface fibration の定理を使い、
Thom--Sebastiani の外積については基底上の Hom 計算と生成性を確認する。
体上で各ファイバーに同値が存在することだけから相対同値を推論しない。
また、以下で構成する核は可換な分離和の実表示上の核である。
非可換族との比較核や、表示を忘れた底上の核の切断は別の問題である。

\subsection{共通の係数直線を持つ三次族}

$S$ を滑らかで連結な準射影 $\mathbb C$-多様体、$k\geq2$ とする。
$W_i$ を $S$ 上の階数三のベクトル束、$L$ を直線束とし、
\begin{equation}\label{mbrp:eq:data}
 f_i\in H^0(S,\Sym^3W_i\otimes L),\qquad
 E_i=V(f_i)\subset\mathbb P_S(W_i),\qquad 1\leq i\leq k
\end{equation}
を取る。ここでは $\mathbb P_S(W_i)=\Proj_S\Sym W_i$ の規約を用いる。
各 $E_i\to S$ が滑らかな平面三次曲線族であると仮定する。
\[
 W=\bigoplus_iW_i,\quad F=\sum_i f_i,\quad
 X=V(F)\subset\mathbb P_S(W),\quad
 Y=E_1\times_S\cdots\times_SE_k,\quad r=3k-3
\]
と置く。$X\to S$ も滑らかである。
実際、各幾何学的ファイバーで $d f_i=0$ は原点のみを含むから、
$dF=0$ の射影的な解は存在しない。

群 $\Gamma_k$ と $Z_0$ は式\eqref{mbrd:eq:gamma}のものを用いる。
その指標群を加法的に書くと
\begin{equation}\label{mbrp:eq:lattice}
 \Lambda_k=\mathbb Z^k/
    \langle3e_i-3e_j\rangle_{i,j},\qquad
 \chi=3e_i,\qquad Z_0=(\mu_3^k)/\Delta\mu_3.
\end{equation}
$\Gamma_k$ は $W_i$ の変数に指標 $e_i$ で作用する。
共通対角 $\mathbb G_m$ を除くと $Z_0$ の射影的な $X$ への作用を得る。

有限局所自由な twisted graded MF の dg 圏を Morita 完備化し、
$S$ 上の線型圏の stack として扱う。
その一因子版を $\mathcal M_i=\MF_{\mathrm{gr},S}(f_i;L)$、
$\Lambda_k$ 次数版を $\mathcal M_F=\MF_{\Gamma_k,S}(F;L)$ と書く。
これらの記号は、アフィン開での有限 MF とその perfect 加群の圏を
局所自由束・微分・加群作用の降下で貼り合わせたものを表す。
内部次数の規約は $M(1)_n=M_{n+1}$ である。
一因子では
\begin{equation}\label{mbrp:eq:line-period}
 [2]=(3)\otimes L,
\end{equation}
多因子では $(3)$ を共通の指標 twist
$\mathsf T_{\chi^{-1}}$ と読む。
関係直線を使う記法では $\ell=L^\vee$ であり、
式\eqref{mbrp:eq:line-period}は
命題\ref{mbsd2:prop:line-period}の $[2]=(3)\otimes\ell^\vee$ と一致する。

\begin{lemma}[相対 Orlov window]
\label{mbrp:lem:relative-orlov}
式\eqref{mbrp:eq:data}の仮定の下で、$S$-線型 dg Morita 同値
\begin{equation}\label{mbrp:eq:relative-orlov}
 \mathcal M_i\simeq\Perf(E_i/S),\qquad
 \MF_{\mathrm{gr},S}(F;L)\simeq\Ku_{\Perf}(X/S)
\end{equation}
を取れる。ここで $\Perf(E_i/S)$ は $E_i$ の perfect 圏の相対 stack、
右辺の残余圏は
\[
 \Ku_{\Perf}(X/S)=
 \bigl\langle p^*\Perf(S),p^*\Perf(S)\otimes\OO_X(1),
                  \ldots,p^*\Perf(S)\otimes\OO_X(r-1)\bigr\rangle^\perp
\]
である。全体の比較は $Z_0$ 作用と両立し、従って
\begin{equation}\label{mbrp:eq:gamma-window}
 \mathcal M_F\simeq\Ku_{\Perf}(X/S)^{Z_0}
\end{equation}
を与える。右辺は homotopy-coherent な同変化と冪等完備化を含む。
\end{lemma}
\begin{proof}
\cite[Section~3, Corollary~3.4]{MBRPBDFIK}は、滑らかな基底、
ベクトル束 $W$、係数直線束 $L$、
$\Sym^dW\otimes L$ の切断に対する相対 Orlov 定理を与える。
ここでは零超曲面も滑らかなので、同文献の $D^b(\operatorname{coh})$ は
perfect 圏に等しい。階数三、次数三の場合は parameter が零であり、
階数 $3k$ の場合は parameter が $r$ である。

使用する比較は抽象的な三角同値だけではない。
同文献の構成は、係数直線の全空間を加えた二つの GIT phase の window と、
正則切断に対する Koszul 比較から成る。
window の制限、導来引戻し・押出し、Koszul factorization は dg 関手である。
引用した完全忠実性を全 shift に適用すると Hom 複体の準同型となるので、
これらの関手は dg Morita 比較を与える。
window の制限同値の逆は Morita $\infty$-圏で取る。
同値の逆を選ぶ空間は可縮なので、局所比較の合成にもこの coherence を用いる。

引用の半直交分解は例外成分を左に置く規約である。
残余像を $\OO_X(-r)$ で tensor すれば、ここで採用した右直交補になる。
これは $\omega_{X/S}=\OO_X(-r)\otimes p^*D$ という相対随伴公式と
相対 Serre 双対による。$D$ は基底上の直線束であり、残余部分圏を変えない。
一因子では $r=0$ なのでこの調整は不要である。

window は共通対角 $\mathbb G_m$ の整数 charge で定義し、
各 charge の全 $Z_0$ 指標を保持する。
ブロックの線型作用はこの window と Koszul 比較を保つ。
また $3\mid r$ なので、上の $\OO_X(-r)$ による調整も $Z_0$ と両立する。
$\Gamma_k$ の分裂を一つ使ってこの記述を通常の同変化に直せば
式\eqref{mbrp:eq:gamma-window}を得る。
分裂を変えた同定は命題\ref{mbza:prop:gamma-action}の scalar 自明化による。
従って作用そのものと同変圏には基準ブロックの選択を残さない。
\end{proof}

\subsection{相対外積の直接構成}

\begin{lemma}[共通係数直線を持つ相対外積]
\label{mbrp:lem:relative-ts}
上の仮定の下で、実際の外積 dg 関手が $S$-線型 Morita 同値
\begin{equation}\label{mbrp:eq:relative-ts}
 \mathcal M_1\otimes^{\mathbf L}_{\Perf(S)}\cdots
       \otimes^{\mathbf L}_{\Perf(S)}\mathcal M_k
       \xrightarrow{\ \sim\ }\mathcal M_F
\end{equation}
を与える。相対 tensor 積は線型圏の stack の Morita tensor 積とする。
\end{lemma}
\begin{proof}
まず $S=\Spec R$ とし、束と係数直線の frame を取る。
二因子の外積は、\cite[Definition~3.51]{NFBFKKernels}の行列によって作る。
$P^{-1}\to P^0\to P^{-1}(3)\otimes L$ という MF の規約では、
共通の周期 functor を $T=\mathsf T_{\chi^{-1}}\otimes L$ として
\begin{align}
 (P\boxtimes Q)^{-1}
   &=P^{-1}\boxtimes Q^0\oplus P^0\boxtimes Q^{-1},\\
 (P\boxtimes Q)^0
   &=P^0\boxtimes Q^0\oplus T(P^{-1}\boxtimes Q^{-1}).
 \label{mbrp:eq:exterior-components}
\end{align}
微分は Koszul 符号付きの和であり、その平方は $f_P+f_Q$ である。
二つの周期を $\Lambda_k$ の共通次数 $\chi$ へ写してからこの式を使う。
特に奇成分どうしの積には $L$ が一つ付く。
これは通常の二周期圏を、各因子に独立の周期変数を残したまま tensor する操作ではない。
反復外積では $\epsilon\in\{0,1\}^k$ に対応する項は
$T^{\lfloor|\epsilon|/2\rfloor}(\boxtimes_iP_i^{-\epsilon_i})$ であり、
$|\epsilon|$ の偶奇に従う成分に置かれる。
この表示は括弧の付け方によらず、共通 $L$ の冪も一意に指定する。

Hom の比較を $R$ 上で確認する。
有限次数付き射影 MF $P_i,Q_i$ に対して、外積は自然な写像
\begin{equation}\label{mbrp:eq:relative-hom}
 \bigotimes_{i=1}^k
    \Hom^\bullet_{\MF_{\mathrm{gr},R}(f_i)}(P_i,Q_i)
 \longrightarrow
 \Hom^\bullet_{\MF_{\Gamma_k,R}(F)}
       (\boxtimes_iP_i,\boxtimes_iQ_i)
\end{equation}
を与える。これは複体の同型である。
実際、各項を自由次数付き加群の場合に分解すると、
右辺の多重次数は $\mathbb Z^k\to\Lambda_k$ の核
$\langle3e_i-3e_j\rangle$ に沿った和である。
これを各因子の偶数 cohomological shift に配分すると、
全 cohomological degree が指定された左辺の tensor 複体の項になる。
奇数 shift の組合せでは式\eqref{mbrp:eq:exterior-components}の周期補正が同じ配分を与える。
微分と合成の符号は通常の複体の tensor の符号に一致する。
固定した全次数では多項式の正次数性により和は有限である。
直和因子を取れば有限射影加群にも同じ同型を得る。
これは \cite[Lemma~3.52の証明]{NFBFKKernels}の指標分解を、
体上の cohomology の K\"unneth 公式を使わず複体の段階で行ったものである。

この通常の tensor が導来 tensor を計算することも確認する。
各 Hom 複体の項は有限射影 $R$-加群であり、下に有界である。
補題\ref{mbrp:lem:relative-orlov}と $E_i\to\Spec R$ の properness により、
その導来同型類は $R$-perfect である。
$R$ は有限大域次元を持つ正則環なので、
有界な有限射影模型からの準同型の cone は可縮である。
ここでは射影加群からなる非有界 acyclic 複体でも、
短完全列の dimension shifting を大域次元回だけ繰り返すと全 cycle が射影になることを使う。
従って元の Hom 複体は有界な有限射影複体と homotopy 同値であり、$K$-flat である。
式\eqref{mbrp:eq:relative-hom}は Morita 外積の完全忠実性を与える。

生成性には零切断の Koszul stabilization を用いる。
各 $f_i$ について Euler 恒等式
\[
 f_i=\sum_{a=1}^3 x_{ia}\,\frac{\partial f_i}{3\partial x_{ia}}
\]
から有限 Koszul MF $K_i=K(x_i,\tfrac13d f_i)$ を作る。
この式は垂直微分と tautological section の pairing なので座標に依存しない。
零切断の全 character shift と $\Perf(S)$ の作用による thick 閉包が
MF 圏を生成することは
\cite[Propositions~4.1, 4.2, 4.8]{MBRPBDFIK}の実 dg Koszul 模型から従う。
適用条件は、全空間の零超曲面の特異 locus が零切断であることであり、
各平面三次およびその分離和の滑らかさがこれを保証する。

同じ生成主張を $\Gamma_k$ へ移すには、証明のために一つ分裂を選べば十分である。
有限群 $Z_0$ への induction を施した生成対象は、
標数零の平均化によって全同変対象を生成する。
その induction は零切断の全 $Z_0$ 指標の直和である。
従って $\Gamma_k$ では零切断の全 $\Lambda_k$ shift が生成する。
一方、Koszul 外積には自然な同型
\[
 \boxtimes_i K_i\simeq K_{F}
\]
があり、$\mathbb Z^k\to\Lambda_k$ は全射なので、
$K_F$ の全 shift が外積の像に入る。
これで本質的全射性を得る。
非アフィン基底では $\OO_S$ 一対象での生成を主張せず、
$\Perf(S)$ の作用を含む相対 thick 閉包を用いている。

最後に frame の変更を確認する。
式\eqref{mbrp:eq:exterior-components}は共通の直線束 $L$ を含む式であり、
$L$ の frame の変更は奇成分の積に付けた $L$ の遷移関数で補正される。
ベクトル束の変更にも垂直微分と Koszul 外積は自然である。
従って局所関手と上の Hom 同型は重なりで一致し、三重重なりの条件を満たす。
有限局所自由な因子化と perfect 加群の降下により、
式\eqref{mbrp:eq:relative-ts}は $S$ 上に貼り合う。
この議論では、\cite[Conjecture~5.1.2]{MBRPBDFIK}に関係する
最小 $A_\infty$ 模型への置換を一切使わない。
\end{proof}

\subsection{実核、底変更、置換}

\begin{theorem}[相対曲線積の perfect 核]
\label{mbrp:thm:relative-product}
式\eqref{mbrp:eq:data}の仮定の下で、$S$-線型 dg Morita 同値
\begin{equation}\label{mbrp:eq:relative-product}
 \Phi_S:\Perf(Y)\xrightarrow{\ \sim\ }
                 \Ku_{\Perf}(X/S)^{Z_0}
\end{equation}
が存在する。これは相対 Orlov window と補題\ref{mbrp:lem:relative-ts}の外積の合成であり、
その同変残余圏への包含を含めた関手は実際の核
\begin{equation}\label{mbrp:eq:kernel}
 \mathcal P\in\Perf^{Z_0}(Y\times_S X)
\end{equation}
による積分変換として表される。
$Z_0$ は $Y$ に自明に、$X$ にブロック倍によって作用する。
同値と核の構成には、基底の有限拡大や $L$ の三乗根の選択を要しない。
\end{theorem}
\begin{proof}
補題\ref{mbrp:lem:relative-orlov}と\ref{mbrp:lem:relative-ts}を合成する。
滑らかで proper な $S$-scheme の perfect 圏の相対 Morita 積は
fiber product の perfect 圏である。
この主張はアフィン $S=\Spec R$ 上で
\cite[Lemma~8.11とTheorem~8.15]{MBRPToenMorita}の積公式による。
同文献の基礎環は一般の可換環であり、ここでは $R$ を用いている。
積・導来外積の構成は基底上で自然なので、perfect 複体の降下によって $S$ 上へ貼り合う。
これで式\eqref{mbrp:eq:relative-product}を得る。

核の存在を単に三角同値の表現可能性から推論する必要はない。
構成した dg 関手を $\Perf(X)$ へ合成し、
各アフィン開 $U\subset S$ 上で
\cite[Theorem~8.15]{MBRPToenMorita}を $Y_U,X_U$ に適用する。
両者は $U$ 上 smooth proper なので、関手は
$\Perf(Y_U\times_U X_U)$ の実際の perfect 核に対応する。
核と dg 関手の対応は morphism 複体も含む同値であるから、
関手の $Z_0$ 同変構造は核の coherent な $Z_0$ 線型化へ移る。
同じ理由で、重なり上の関手の比較とその coherence は核へ移る。
perfect 複体の降下を適用して式\eqref{mbrp:eq:kernel}を得る。
\end{proof}

\begin{corollary}[非平坦底変更後の比較]
\label{mbrp:cor:base-change}
定理\ref{mbrp:thm:relative-product}の核を固定する。
任意の準コンパクト・分離的 noetherian $S$-scheme $T$ に対して、
その導来引戻し $\mathcal P_T$ は
\begin{equation}\label{mbrp:eq:base-change}
 \Perf(Y_T)\simeq\Ku_{\Perf}(X_T/T)^{Z_0}
\end{equation}
を与える。$T\to S$ の flatness や $T$ の被約性は仮定しない。
引戻しと比較は、合成した底変更とも coherent に両立する。
$T$ も滑らかな準射影基底なら、これは引戻したデータからの上記構成に一致する。
\end{corollary}
\begin{proof}
$X,Y$ は $S$ 上 smooth proper なので、
関手の随伴と残余 projector も perfect 核で表される。
相対例外成分の projector の核は
$\OO_X(-j)\boxtimes_S\OO_X(j)$ と評価写像から構成される。
この核の $Z_0$ 線型化は二つの重みが相殺するため定まる。
従って $\OO_X(j)$ の個別の線型化を選ぶ必要はない。
有限個の mutation の cone によって得る残余 projector を $\Pi$ とする。
これは $X\times_SX$ 上の diagonal $Z_0$-同変核である。
同変圏上の projector の核は
\[
 \widehat\Pi=
 \operatorname{Ind}_{\Delta Z_0}^{Z_0\times Z_0}\Pi
 \quad\text{on }[X/Z_0]\times_S[X/Z_0]
\]
と読む。従って全 $\operatorname{Rep}(Z_0)$ 成分を保持する underlying mutation の同変化であり、
一つの線型化を持つ $\OO_X(j)$ への invariant Hom だけを取る投影ではない。
同値の逆を残余 projector と合成すれば、その核 $\mathcal Q$ とともに
\[
 \mathcal Q\circ\mathcal P\simeq\OO_{\Delta_Y},\qquad
 \mathcal P\circ\mathcal Q\simeq\widehat\Pi
\]
を得る。第2式は同変 perfect 圏上の projector の恒等式である。
$\mathcal P$ と $\mathcal Q$ の核はそれぞれ
$Y\times_S[X/Z_0]$ と $[X/Z_0]\times_SY$ 上にあり、
convolution では中間の $Z_0$ の不変部分を取る。

これらの核・評価・unit と counit を $T$ へ導来引戻しする。
各 fiber product の射は $S$ 上 flat で、積分変換の射影は proper perfect であるから、
核の convolution とこの引戻しは両立する。
有限群不変部分も平均化 idempotent により底変更と両立する。
引戻した $\Pi_T$ は同じ相対例外成分の残余 projector であり、
その induction は $\widehat\Pi_T$ である。
これは射影束のコホモロジーと超曲面の短完全列による計算で、
$T$ の正則性を使わず確認できる。
上の二つの恒等式が $T$ 上にも成立するので、式\eqref{mbrp:eq:base-change}を得る。
導来引戻しの合成同型が coherence を与える。

滑らかな $T$ では有限 MF の引戻し、Koszul 外積、
window の制限も同じ pullback と両立するため、直接の構成に一致する。
一般の $T$ についてここで主張するのは perfect 圏の比較であり、
$D^b(\operatorname{coh}Y_T)=\Perf(Y_T)$ や
特異基底上の絶対的な coherent MF 圏への拡張ではない。
\end{proof}

\begin{corollary}[既存の比較体上での曲線積表示]
\label{mbrp:cor:comparison-field}
$K$ を $\mathbb C$ 上有限生成な体とし、
$f_i\in K[x_{i1},x_{i2},x_{i3}]_3$ が滑らかな平面三次を定めるとする。
このとき追加の体拡大なしに
\[
 \Ku(V(\textstyle\sum_i f_i))^{Z_0}
       \simeq D^b\bigl(\textstyle\prod_iV(f_i)\bigr)
\]
が $K$ 上の実 perfect 核で得られる。
従って定理\ref{med:thm:multiblock-generic}の仮定の下で
同変可換比較を得た有限比較体上では、
系\ref{mbrd:cor:nc-product}の可換積比較のための追加有限拡大を省ける。
定理\ref{mpc:thm:point-model}の明示的 Jacobian tensor 族でも、
同じ有限比較体上で入力の点曲線の積との比較を得る。
\end{corollary}
\begin{proof}
有限個の係数を含む $\mathbb C$ 上有限生成な整域を選び、
その分数体を $K$ に取る。
その非空開に縮めて基底を滑らかにし、
さらに平面三次の判別式を可逆にすれば、$f_i$ は
定理\ref{mbrp:thm:relative-product}の族へ延びる。
系\ref{mbrp:cor:base-change}を一般点 $\Spec K$ に適用すればよい。
係数は元の $f_i$ のままであり、比較体を変更していない。
既存の二つの非可換適用で得る有限比較体は $\mathbb C(h)$ の有限拡大なので、
この場合に含まれる。
その体上で既に存在する同変可換比較と、今構成した積比較を合成する。
これは指定された形式比較核の全 jet の代数化を主張するものではない。
\end{proof}

\begin{proposition}[因子置換との整合性]
\label{mbrp:prop:permutations}
定理\ref{mbrp:thm:relative-product}の相対比較は、
データ $(W_i,f_i;L)$ の同型と因子置換に対して coherent に取れる。
置換は $\Gamma_k$ と $Z_0$ の指標も同時に置換する。
\end{proposition}
\begin{proof}
外積は複体の標準的な Koszul braiding を用いる。
従って隣接交換の二乗および braid 関係を、その符号を含めて満たす。
Koszul 対象の
$\bigwedge(\bigoplus_iW_i)\simeq\bigotimes_i\bigwedge W_i$
もこの braiding に関して自然である。
相対 Orlov window は共通対角 charge によって定義され、
全 $Z_0$ 指標を保持するから因子の順序を用いない。
これらの実 dg 関手の coherence は、核と関手の対応によって
式\eqref{mbrp:eq:kernel}の核の coherence に移る。
\end{proof}

\begin{remark}[相対可換比較と未解決の降下]
\label{mbrp:rem:scope}
本節は、共通係数直線まで含む可換分離和の実表示上で比較核を構成したものである。
相対外積の構成を仮定に置いた条件付き主張でも、
幾何点での同値からの存在降下だけでもない。
一方、非可換残余圏からこの可換分離和へ至る相対核は
定理\ref{med:thm:multiblock-generic}の一般点の有限降下とは別に必要である。
実表示を持たない全 $\OO$-モジュライ上へ選んだ核が降りることも、
本節の基底変更・置換整合性のみからは従わない。
その際は表示 gerbe の重みと source/target の同期を確認する必要がある。
さらに標準 rotation の線型化を保持するなら、
命題\ref{mbza:prop:shift-character}の character 変換を含めなければならない。
ここで得た置換整合性は、線型化された rotation の分裂選択の存在を主張しない。
\end{remark}

% !TEX root = ../nc-cubic.tex
\section{Sklyanin 中心和の直接積比較}
\label{mbnp:sec:product}

本節では、Sklyanin 因子の中心商を直接用いて、中心群で同変化した
非可換残余圏を点曲線の積と比較する。
一変数の可換特殊化、完備化、復元三次の選択は仮定しない。
入力は、実際の Sklyanin 代数と中心元を持つ体上の順序付き表示である。
証明では、Burke--Stevenson の非可換 Orlov 定理、
有限行列分解の余核比較、Shamash 型の augmentation 分解を組み合わせる。
可換環に対する Thom--Sebastiani 定理を非可換環へそのまま引用せず、
必要な外積の Hom と生成性を以下で証明する。
この適用の包括的な新規性調査は未完了である。

\subsection{実表示と noetherian 性}

$K\supset\mathbb C$ を体、$k\geq2$ とする。
各 $i$ について、$K$ 上の滑らかな平面 genus-one 曲線 $E_i$、
実際の次数三線束 $L_i$、平行移動 $\sigma_i=t_{p_i}$ を取り、
三次元 quadratic Sklyanin 代数 $S_i$ とその正則中心三次 $\Omega_i$ が
\begin{equation}\label{mbnp:eq:factor-data}
 T_i:=S_i/(\Omega_i)\simeq B(E_i,L_i,\sigma_i)
\end{equation}
を満たすとする。
ここで $S_i$ は連結局所有限、Koszul、global dimension と
Gorenstein parameter がともに三の AS-regular 代数である。
本稿の基本楕円開に属する実表示はこれらの条件を満たす。
点曲線を三次元空間内に固定するための追加条件や $j_i$ の相異性は、
本節の積比較自体には不要である。
以下の全ての加群は右加群とする。

$a_i\in K^\times$ を取り、
\begin{equation}\label{mbnp:eq:sum-data}
 w_i=a_i\Omega_i,\quad B=\bigotimes_{i=1}^kS_i,\quad
 w=\sum_{i=1}^kw_i,\quad A=B/(w),\quad N=3k,\quad r=N-3
\end{equation}
と置く。$w_i$ は方程式の frame も含む指定元である。
各因子の有限行列分解の周期は、この $w_i$ による余核二拡張から定める。
$w_i$ をスカラー倍しても中心商 $T_i$ は変わらないが、
指定周期の同型まで同じと見なすことはしない。

\begin{lemma}[独立な楕円因子の tensor 環]
\label{mbnp:lem:noetherian}
$\bigotimes_iT_i$、$B$、$A$ は両側 noetherian である。
$B$ は Koszul AS-regular で、global dimension と parameter は $N$、
$w$ は正則である。
$A$ は両側有限 injective dimension を持つ AS-Gorenstein 環であり、
parameter は $r$ である。従って $A$ に
\cite[Proposition~4.18, Lemmas~6.1--6.2, Theorem~6.4]{LP-BS}
の導来 section と Orlov 分解を適用できる。
\end{lemma}
\begin{proof}
まず $T=\bigotimes_iT_i$ の noetherian 性を証明する。
二因子の場合は \cite[Corollary~3.5]{MBNPKeelerMultiple}にも含まれるが、
三因子以上の中間 tensor 環は通常の一変数 THCR ではないので、
同系の単純な反復とはしない。
体を代数閉包へ拡大しても切断と $T$ の tensor 表示は変わらず、
noetherian 性は忠実平坦な体拡大から降下する。
従ってまず $K$ を代数閉体としてよい。

$X=\prod_iE_i$ とし、$\boldsymbol n\in\mathbb N^k$ に対して
\[
 L_{\boldsymbol n}=\boxtimes_i
    (L_i\otimes\sigma_i^*L_i\otimes\cdots
                  \otimes(\sigma_i^{n_i-1})^*L_i),\qquad
 T_{\boldsymbol n}=H^0(X,L_{\boldsymbol n})
\]
と書く。$n_i=0$ の因子は $\OO_{E_i}$ とする。
因子ごとの平行移動は $X$ の数値的 Picard 群に恒等に作用し、
全ての $n_i>0$ なら $L_{\boldsymbol n}$ は ample である。
従って \cite[Theorem~2.7]{MBNPKeelerMultiple}により、
この可換な可逆双加群の組は NC-ample である。
また全ての $L_{\boldsymbol n}$ は大域生成される。

多重斉次右イデアル $J\subset T$ に対し
\[
 I_{\boldsymbol n}
  =\operatorname{im}\bigl(J_{\boldsymbol n}\otimes
                    L_{\boldsymbol n}^{-1}\longrightarrow\OO_X\bigr)
\]
と置く。積の定義と大域生成性から
$I_{\boldsymbol n}\subset I_{\boldsymbol n+\boldsymbol m}$ である。
$X$ の noetherian 性を対角の増大列に適用すると、
ある $\boldsymbol n_0$ とイデアル層 $I$ があって
$I_{\boldsymbol n}=I$ が全ての $\boldsymbol n\geq\boldsymbol n_0$ で成り立つ。
評価全射
$J_{\boldsymbol n_0}\otimes\OO_X\to I L_{\boldsymbol n_0}$
の核を $Q$ とする。
これを $\sigma_{\boldsymbol n_0}^*L_{\boldsymbol m}$ で tensor し、
NC-ampleness による $H^1$ の消滅を使うと、全ての十分大きい各成分を持つ
$\boldsymbol m$ に対して
\[
 J_{\boldsymbol n_0}T_{\boldsymbol m}
    =H^0(X,I L_{\boldsymbol n_0+\boldsymbol m})
    =J_{\boldsymbol n_0+\boldsymbol m}
\]
となる。最後の等号では、左辺が右イデアル $J$ に含まれることも使った。

残る多重次数は、有限個の条件 $n_i<c_i$ の和集合に含まれる。
$i$ と $d<c_i$ を固定した部分 $J_{n_i=d}$ は
$(T_i)_d\otimes\bigotimes_{j\ne i}T_j$ の右
$\bigotimes_{j\ne i}T_j$-部分加群である。
因子数に関する帰納法により、各部分は有限生成である。
帰納法の始点は因子零の環 $K$ とする。
これら有限個の面の生成元と $J_{\boldsymbol n_0}$ から $J$ 全体が生成される。
\cite[Lemma~4.4(1)]{MBNPKeelerMultiple}により、
多重斉次右イデアルの ACC は通常の右 noetherian 性を与える。
opposite は $(E_i,(\sigma_i^{-1})^*L_i,\sigma_i^{-1})$ の同種の環なので
\cite[Lemma~2.4]{MBNPKeelerMultiple}、同じ証明で左 noetherian 性も従う。

各 $S_i\to T_i$ の斉次 $K$-線型切断を選ぶと、
$w_i$ の正則性と正の次数による帰納法から、
$S_i$ は $K[w_i]$ 上自由である。
従って $B$ は $K[w_1,\ldots,w_k]$ 上自由で、
$w_1,\ldots,w_k$ は中心正則列を成す。
$I=(w_1,\ldots,w_k)$ に関する随伴次数環は
\[
 \operatorname{gr}_I B\simeq T[t_1,\ldots,t_k]
\]
という中心多項式環であり、両側 noetherian である。
斉次イデアルの初項の有限生成元を $B$ へ持ち上げる。
内部次数 $d$ を固定すれば $I$ の次数は高々 $\lfloor d/3\rfloor$ なので、
初項を順に消去する操作は有限回で終わる。
従って全ての斉次左右イデアルは有限生成であり、
通常の内部次数の初項を取れば $B$ 自身も両側 noetherian となる。
また多項式環上の自由性から $\sum_iw_i$ は正則である。

各因子の有限 Koszul 分解および有限双加群分解を tensor すると、
$B$ の global dimension は $N$、AS 双対の parameter も $N$ となる。
これは定理\ref{mbsd:thm:ambient-duality}の体上の構成である。
中心元の二項分解による change of rings は
\[
 R\underline{\operatorname{Hom}}_A(K,A)\simeq K(r)[-N+1]
\]
と左右有限 injective dimension を与える。
商の noetherian 性は既に従っているので、
\cite[Definition~5.1, Lemmas~6.1--6.2]{LP-BS}の仮定が満たされる。
ここでは環 $A_0=K$ の有限大域次元も用いた。
\end{proof}

\subsection{augmentation が生成する window}

有限次数付き射影行列分解の dg 圏に有限 cone と冪等完備化を施したものを
$\MF^{\mathbb Z}(B,w)$ と書く。
$K$ の stabilization を $K_w$ とし、
\[
 \mathcal C_A=\operatorname{thick}\{\OO_A(j):j\in\mathbb Z\},\qquad
 \mathcal K_A^{\mathrm{line}}
   =\langle\OO_A,\ldots,\OO_A(r-1)\rangle^\perp\cap\mathcal C_A
\]
と置く。$r=0$ の一因子の場合は例外部分を空と読む。

\begin{lemma}[体上の三対象 window]
\label{mbnp:lem:window}
上の環論的仮定の下で、実際の dg quasi-functor
\begin{equation}\label{mbnp:eq:window}
 \Psi_A=\operatorname{stab}\circ q_{\mathrm{sg}}
                    \circ(R\omega(-))_{\geq1-r}
\end{equation}
は同値
\[
 \mathcal K_A^{\mathrm{line}}\simeq
  \operatorname{thick}\{K_w,K_w(1),K_w(2)\}
    \subset\MF^{\mathbb Z}(B,w)
\]
を与える。$k\geq2$ では residual rotation
$\rho_A=L_{\OO_A}\circ(1)$ を内部 shift $(1)$ へ移し、
$w$ から定まる指定周期 $(3)\simeq[2]$ と両立する。
一因子では、同じ window1 が
\begin{equation}\label{mbnp:eq:factor-window}
 \MF^{\mathbb Z}(S_i,w_i)\simeq\Perf(E_i),\qquad
 K_{w_i}\longleftrightarrow\OO_{E_i}[1]
\end{equation}
を与え、$K_{w_i},K_{w_i}(1),K_{w_i}(2)$ が左辺全体を生成する。
\end{lemma}
\begin{proof}
まず有限 MF と特異 dg 圏との比較を確認する。
ambient の大域次元が有限なので、商の有限生成 Gorenstein-projective 加群は
ambient 上の射影次元が高々一であり、斉次な長さ一の射影表示から MF を作れる。
余核の完全忠実性と本質的全射性は
\cite[Corollary~2.11, Theorem~7.2, Remark~7.3(2)]{AMChen}による。
右加群については opposite 環へ適用し、斉次な射影表示と射の持上げは
次数零成分を取って行う。従って引用の左加群・非次数付きの表示から
本節の有限 graded 版を得る。
dg 化は補題\ref{tw:lem:stabilization}の証明と同じく、
成分別 split exact な有限 MF の圏と Gorenstein-projective 加群の exact 圏を使う。
有界導来 dg 圏を projective-injective 複体で割り、有限 totalization を取る。
全 shift の stable Hom 比較が準同型を与えるので、これは dg Morita 同値である。
この構成に特殊化や完備局所環は必要ない。
二段の余核分解を ambient へ持ち上げた合成は指定した $w$ の乗法であり、
式\eqref{tw:eq:period-extension}と同じ規約で指定周期も保持される。

補題\ref{mbnp:lem:noetherian}により、
\cite[Theorem~6.4とその証明]{LP-BS}を $i=1-r$ で適用できる。
$q_{\mathrm{sg}}(R\omega(-))_{\geq i}$ の核は
$\langle\OO_A,\ldots,\OO_A(r-1)\rangle$ であり、
その Verdier 商を特異圏と同定する。
半直交分解の表示の向きによらず、同じ商を表す右直交補への制限は完全忠実である。
section、内部切断、dg 商は enhanced な関手であり、
この完全忠実性を全 shift に適用して dg の完全忠実性を得る。
従って体上の本補題では、特殊ファイバーからの Nakayama による持上げを使わない。

$\OO_A(0),\ldots,\OO_A(N-1)$ は $\mathcal C_A$ を生成する。
これは ambient の augmentation Koszul 分解を $A$ へ制限し、
有限長の homology を tails で消すことによる。
線コホモロジーの公式から、この範囲の $j$ について
$(R\omega\OO_A(j))_{\geq i}=A(j)_{\geq i}$ である。
$U_j=A(j)_{<i}$ と置くと、$j<r$ では $U_j=0$、
$j=r,r+1,r+2$ では
\begin{equation}\label{mbnp:eq:triangular}
 \Psi_A(\OO_A(j))\simeq\operatorname{stab}(U_j)[-1],\qquad
 \operatorname{gr}U_j=
   \bigoplus_{u=0}^{j-r}A_u\otimes_KK(j-u)
\end{equation}
となる。$A_0=K$ なので、三つの有限 filtration を三角的に解けば、像は
$K_w(r),K_w(r+1),K_w(r+2)$ の thick 閉包である。
$(3)=[2]$ と $3\mid r$ により所要の三対象を得る。
$r>0$ における rotation の比較は
式\eqref{tw:eq:rotation-cut}の切断三角と evaluation による。
その計算も剰余体や変形を使わず、指定周期の比較正方形
\eqref{tw:eq:rotation-period-square}をそのまま与える。

一因子では $T_i=B(E_i,L_i,\sigma_i)$ かつ parameter が零なので、
\cite[Theorem~6.4(3)]{LP-BS}は特異圏と
$D^b(\operatorname{coh}E_i)=\Perf(E_i)$ 全体を比較する。
この場合も三つの線対象が生成することを直接確認できる。
ambient の三段 Koszul 分解は tails 上で
\[
 0\to\OO(j-3)\to\OO(j-2)^{\oplus3}
          \to\OO(j-1)^{\oplus3}\to\OO(j)\to0
\]
を与える。従って三つの連続する twist が全 twist を生成する。
THCR の ample pair によるこれらの twist は
$D(\operatorname{Qcoh}E_i)$ を生成し、各 twist は compact なので、
三つの線対象が $\Perf(E_i)$ を厚生成する。
最後に $R\Gamma_{\geq1}\OO_{E_i}=(T_i)_{\geq1}$ と
$0\to(T_i)_{\geq1}\to T_i\to K\to0$ から
$K_{w_i}\leftrightarrow\OO_{E_i}[1]$ を得る。
\end{proof}

\subsection{有限 Shamash 模型と外積}

\begin{lemma}[augmentation stabilization の加法性]
\label{mbnp:lem:shamash}
各 $S_i$ の augmentation に有限斉次射影分解を選ぶ。
その higher homotopies による有限 folding で $K_{w_i}$ を表すことができる。
独立因子の外積について、共通周期を同定した dg MF 圏の中で
\begin{equation}\label{mbnp:eq:augmentation-product}
 K_{w_1}\boxtimes\cdots\boxtimes K_{w_k}\simeq K_w
\end{equation}
である。この同定は各 $w_i$ と $w=\sum_iw_i$ による指定周期を使う。
一次 nullhomotopy の平方が零であることは仮定しない。
\end{lemma}
\begin{proof}
まず ambient $D$、中心正則三次 $v$、augmentation 分解
$P_\bullet\to K$ を一つ固定する。
$P_j$ は homological degree $j\geq0$ に置き、$s_0=d$ と書く。
$v$ は $K$ に零として作用するので、その乗法は nullhomotopic である。
負の Ext の消滅を順に使うと、斉次写像
\begin{equation}\label{mbnp:eq:higher-homotopies}
 s_a:P_j\longrightarrow P_{j+2a-1}(3a),\qquad
 \sum_{a+b=q}s_as_b=
 \begin{cases}v\operatorname{id},&q=1,\\0,&q\ne1\end{cases}
\end{equation}
を取れる。$q=0$ は $d^2=0$ である。
$q\geq2$ の障害は、既に満たされた等式から閉じた負 cohomological degree の
$\operatorname{End}(P_\bullet)$ の元となる。
$P_\bullet\simeq K$ なのでその cohomology は零であり、次の $s_q$ を選べる。
分解は有限なので写像の可能な範囲を越える $s_q$ は零である。
特に一因子の三段分解では $s_1,s_2$ だけを加えればよい。
全ての写像は右 $D$-線型であり、中心性は $v$ の乗法をこの計算で可換にする。

\begin{equation}\label{mbnp:eq:fold}
 F^0=\bigoplus_{j\ \mathrm{even}}P_j(3j/2),\qquad
 F^{-1}=\bigoplus_{j\ \mathrm{odd}}P_j(3(j-1)/2)
\end{equation}
と置く。$\sum_as_a$ は
$F^{-1}\to F^0\to F^{-1}(3)$ の微分を与え、その平方は $v$ である。
これは有限階数の行列分解である。

この folding が実際に augmentation を安定化することを確認する。
$\overline P=P\otimes_DD/(v)$ として、非負次数の複体
\begin{equation}\label{mbnp:eq:shamash-resolution}
 C_n=\bigoplus_{q\geq0}\overline P_{n-2q}(-3q)
\end{equation}
を取る。微分の $s_a$ 成分は添字 $q$ を $q-a$ へ送り、
$q-a<0$ の成分は零とする。
式\eqref{mbnp:eq:higher-homotopies}から $C_\bullet$ は複体となる。
$q$ で濾過すると、第一頁は
\[
 \operatorname{Tor}^{D}_j(K,D/(v))=
 \begin{cases}K,&j=0,\\K(-3),&j=1,\\0,&\text{その他}\end{cases}
\]
である。$s_1$ は、$q\geq1$ の $j=0$ の項から $q-1$ の $j=1$ の項へ
$\pm\operatorname{id}$ を誘導する。
実際、$D(-3)\xrightarrow{v}D\to D/(v)$ の連結射で
$[s_1x]$ の像を計算すると、$ds_1+s_1d=v$ により $[x]$ になる。
従って第二頁では $(q,j)=(0,0)$ の $K$ だけが残る。
各全 homological degree の濾過は有限なので収束し、$C_\bullet$ は $K$ の分解である。
十分大きい $m$ では
$C_{2m}=F^0(-3m)$、$C_{2m+1}=F^{-1}(-3m)$ である。
従って $\operatorname{coker}(F^{-1}\to F^0)$ は
$\operatorname{Syz}^{2m}(K)(3m)$ であり、特異圏では
$K[-2m](3m)\simeq K$ となる。
この最後の同型も $v$ の二段分解による指定周期である。

各因子で上の系を取り、$\bigotimes_iP_{i,\bullet}$ 上で
$s_a=\sum_i s_{i,a}$ とする。各 $s_{i,a}$ を他の因子へ延長する際は
通常の homological Koszul 符号を入れる。
これらの写像の homological degree は全て奇数なので、異なる因子の項は反交換する。
従って cross term は消え、曲率は $\sum_iw_i$ となる。
内部多重次数では $s_{i,a}$ の次数は $3ae_i$ である。
指標群で $3e_i=3e_j$ を課した後、全項は同じ周期を持つ。
また
\[
 \left\lfloor\frac j2\right\rfloor+
 \left\lfloor\frac l2\right\rfloor+
 \begin{cases}1,&j,l\text{ がともに奇数},\\0,&\text{その他}\end{cases}
 =\left\lfloor\frac{j+l}{2}\right\rfloor
\]
により、tensor 分解の folding は MF の外積の奇成分どうしの周期補正と一致する。
従って左辺は総 augmentation 分解の folding そのものであり、
式\eqref{mbnp:eq:augmentation-product}を得る。
\end{proof}

\subsection{中心同変化した残余圏の積表示}

\[
 \Gamma_k=\{(\lambda_i)\in\mathbb G_m^k:
                 \lambda_1^3=\cdots=\lambda_k^3\},\quad
 \Lambda_k=\mathbb Z^k/\langle3e_i-3e_j\rangle,
 \quad Z_0=(\mu_3^k)/\Delta\mu_3
\]
と置く。$\Gamma_k$ の共通三次指標を $\chi$ と書く。
$S_i$ の次数一生成元の重みは $e_i$、$w_i$ の重みは $\chi$ である。
以下の $Z_0$ 作用は、定理\ref{mbza:thm:central-action}の
scalar 自明化による coherent な中心作用である。
一つの基準因子による $\Gamma_k\simeq\mathbb G_m\times Z_0$ の分裂を用いれば
通常の同変加群として表せるが、作用と同変圏の同型類はその分裂に依存しない。

\begin{theorem}[Sklyanin 中心和の直接積比較]
\label{mbnp:thm:product}
式\eqref{mbnp:eq:factor-data}--\eqref{mbnp:eq:sum-data}の仮定の下で、
体 $K$ 上の実際の dg Morita 同値
\begin{equation}\label{mbnp:eq:product}
 \bigl(\mathcal K_A^{\mathrm{line}}\bigr)^{Z_0}
     \simeq \Perf(E_1\times_K\cdots\times_KE_k)
\end{equation}
が存在する。これは各因子の window1 と有限 MF の外積によって与えられる。
体の有限拡大、可換特殊化、復元定理の一般性仮定を必要としない。
両辺は smooth proper であり、$\mathcal K_A^{\mathrm{line}}$ も smooth proper である。
全 coherent 残余圏との一致は系\ref{mbnp:cor:coherent}で従う。
\end{theorem}
\begin{proof}
$\MF^{\Lambda_k}(B,w)$ を共通曲率次数 $\chi$ の有限 MF 圏とする。
二因子の外積は有限射影右加群の外積と Koszul 符号付きの和微分で定義する。
その成分は、共通内部 shift を $T=(\chi)$ と書けば
\[
 (P\boxtimes Q)^{-1}=P^{-1}\boxtimes Q^0\oplus P^0\boxtimes Q^{-1},\qquad
 (P\boxtimes Q)^0=P^0\boxtimes Q^0\oplus T(P^{-1}\boxtimes Q^{-1})
\]
である。ここで $(\chi)$ は内部次数を下げる規約なので、
群表現の tensor としては指標 $\chi^{-1}$ を表す。
この式は \cite[Definition~3.51]{NFBFKKernels}と同じ有限構成であり、
独立な二つの代数の元が tensor 環で可換することだけを使う。

有限自由次数付き加群で始めると、各 Hom の係数について
\[
 \operatorname{Hom}_{S_i\otimes S_j}
     (P_i\otimes P_j,Q_i\otimes Q_j)
 =\operatorname{Hom}_{S_i}(P_i,Q_i)
      \otimes_K\operatorname{Hom}_{S_j}(P_j,Q_j)
\]
という有限射影加群の同型を使える。
指標を分解し、$\mathbb Z^k\to\Lambda_k$ の核
$\langle3e_i-3e_j\rangle$ に沿う和を各因子の偶数 cohomological shift に配分すると、
外積は複体の同型
\begin{equation}\label{mbnp:eq:hom}
 \bigotimes_i\operatorname{Hom}^\bullet_{\MF^{\mathbb Z}(S_i,w_i)}(P_i,Q_i)
 \xrightarrow{\sim}
 \operatorname{Hom}^\bullet_{\MF^{\Lambda_k}(B,w)}
       (\boxtimes_iP_i,\boxtimes_iQ_i)
\end{equation}
を与える。奇数 shift の組合せは上の奇成分どうしの補正と一致する。
各因子の graded Hom 複体は、有限個の自由 shift と代数の正次数性から下に有界である。
従って固定した全 cohomological degree でこの配分の和は有限である。
微分と合成の符号は複体の tensor の Koszul 符号に一致する。
直和因子を取れば有限射影加群にも同じ結論が成り立つ。
体上ではこれらの tensor は導来 tensor を計算する。
従って外積は Morita tensor 積からの完全忠実な dg 関手を与える。
これは \cite[Lemma~3.52の証明]{NFBFKKernels}の Hom 分解の、
本節の独立な非可換因子への直接の検証である。

各因子は補題\ref{mbnp:lem:window}の三つの augmentation shift で生成される。
補題\ref{mbnp:lem:shamash}により、その外積像は
$K_w(\lambda)$、$\lambda\in\Lambda_k$ の thick 閉包である。
実際、$\mathbb Z^k\to\Lambda_k$ は全射であり、
全ての因子内部 shift は各因子の三対象の thick 閉包に含まれる。

一方、window はブロックの次数付き作用と可換する。
基準因子を一つ取ると、$\Lambda_k$ 次数付き MF は
$\mathbb Z$ 次数付き MF の $Z_0$ 同変化と同定される。
$K_w$ とその整数 shift はこの作用の線型化を持つ。
有限群の誘導と平均化により、その thick 閉包の同変化は、
これらの対象を全 $Z_0$ 指標で twist した対象で生成される。
実際、任意の同変対象は忘却後の対象を誘導したものの retract であり、
誘導された各線型化対象は全指標 twist の直和に分解する。
この生成部分はちょうど $\operatorname{thick}\{K_w(\lambda)\}_{\lambda\in\Lambda_k}$ である。
分裂の変更は命題\ref{mbza:prop:gamma-action}の自然な比較による。
これで外積の本質的全射性を、未確認の全 MF の生成性を使わずに得た。

因子の window と smooth projective scheme の外積比較から
\[
 \bigotimes_i^{\mathbf L}\MF^{\mathbb Z}(S_i,w_i)
  \simeq\bigotimes_i^{\mathbf L}\Perf(E_i)
  \simeq\Perf\left(\prod_iE_i\right)
\]
となる。最後は perfect 圏の Morita tensor と積の標準比較であり、
\cite[Section~8.3]{MBRPToenMorita}を用いられる。
外積 generator とその Hom の K\"unneth 同型によっても確認できる。

$Y=\prod_iE_i$ は smooth proper なので、得られた同変圏は smooth proper である。
$K\supset\mathbb C$ は $Z_0$ の全指標を含むため、
Elagin の二重同変化 \cite[Theorem~4.2, Section~8]{MEDElagin}が
\[
 \mathcal K_A^{\mathrm{line}}
   \simeq\left((\mathcal K_A^{\mathrm{line}})^{Z_0}\right)^{\widehat Z_0}
\]
を与える。ここでの証明は有限 Fourier 分解と平均化であり、
代数閉体という表記を外しても、$|Z_0|$ が可逆で全指標を含む本節の体 $K$ 上で通る。
双対群の作用は dg 同変圏の指標 twist から移す。
有限軌道の直和を生成対象に取り、同変 dg 模型の endomorphism 代数に移すと、
有限群の作用を厳密な dg 代数作用として表せる。
補題\ref{med:lem:crossed-perfect}をこの有限作用へ適用すると、
$\mathcal K_A^{\mathrm{line}}$ 自身も smooth proper となる。
最後に両辺の有限 generator を取れば、上で構成した dg 関手は
それらの dg endomorphism 代数の間の実際の双加群で表される。
smoothness と properness によりこの双加群は perfect であり、
Morita 逆と単位・余単位も $K$ 上に存在する。
従って式\eqref{mbnp:eq:product}は幾何学的点での抽象的な三角同値だけではない。
\end{proof}

\begin{proposition}[因子 rotation と双対指標作用]
\label{mbnp:prop:rotations}
定理\ref{mbnp:thm:product}で使う一因子の比較の下で、内部 shift $(1)$ は
\begin{equation}\label{mbnp:eq:factor-rotation}
 R_{i,p_i}=T_{\OO_{E_i}}\circ(\sigma_i)_*\circ(-\otimes L_i)
\end{equation}
に対応する。ここで
$T_{\OO_E}(F)=\operatorname{Cone}(R\operatorname{Hom}(\OO_E,F)\otimes\OO_E\to F)$
は spherical twist である。
$R_{i,p_i}^3\simeq[2]$ の同型は、指定した $w_i=a_i\Omega_i$ の MF 周期を
この window で移したものとする。

$\mathsf T_\eta$ を $\Gamma_k$ の指標 $\eta$ の tensor と書けば、
\begin{equation}\label{mbnp:eq:character-rotation}
 \mathsf T_{e_i^{-1}}\longleftrightarrow R_{i,p_i},\qquad
 \mathsf T_{e_i/e_b}\longleftrightarrow R_{i,p_i}^{-1}R_{b,p_b}
\end{equation}
である。後式は $\widehat Z_0$ の作用を与え、
その群法則と三乗・交換の coherence は $\Gamma_k$ 次数付き MF から移す。
基準 $b$ に対応する線型化付き residual rotation は $R_{b,p_b}$ へ移る。
\end{proposition}
\begin{proof}
右 THCR の shift は
式\eqref{eq:nfm-right-thcr-shift}の
$s_i=(\sigma_i)_*\circ(-\otimes L_i)$ である。
補題\ref{lem:nfm-twisted-elliptic-window}の切断三角を適用できる。
その Orlov 仮定は補題\ref{mbnp:lem:window}で既に確認した。
具体的には任意の $F\in\Perf(E_i)$ に対し
\[
 R\Gamma_{\geq1}(s_iF)\longrightarrow
 R\Gamma_{\geq0}(s_iF)\longrightarrow
 R\operatorname{Hom}(\OO_{E_i},s_iF)\otimes K\longrightarrow
\]
を特異圏へ写す。
$K\leftrightarrow\OO_{E_i}[1]$ を用いて回転すれば、
evaluation の cone が内部 $(1)$ に対応することが分かる。
この議論は可換特殊化や DVR を使わない。
$(\sigma_i)_*$ と $-\otimes L_i$ の順序は右加群の規約によって固定される。

内部 shift の規約は $M(1)_n=M_{n+1}$ なので、
因子 $i$ の $(1)$ は表現の指標 $e_i^{-1}$ の tensor に等しい。
外積はこの指標作用とその自然な合成を保つため、
式\eqref{mbnp:eq:character-rotation}を得る。
各因子の三次周期は共通の $\chi^{-1}$ の twist に対応する。
従って比の三乗を自明化する同型も共通の MF 模型から得られ、
自己関手の抽象的な三乗同型を独立に選んで coherence を仮定しているのではない。
基準変更に関する指標の逆向きも
命題\ref{mbza:prop:shift-character}と一致する。
\end{proof}

\begin{corollary}[全 coherent 残余圏への飽和]
\label{mbnp:cor:coherent}
式\eqref{mbnp:eq:factor-data}--\eqref{mbnp:eq:sum-data}の体上の実表示について
\begin{equation}\label{mbnp:eq:coherent}
 \mathcal C_A=D^b_{\mathrm{dg}}(\operatorname{qgr}A),\qquad
 \mathcal K_A^{\mathrm{line}}=\mathcal K_A^{\mathrm{coh}},\qquad
 (\mathcal K_A^{\mathrm{coh}})^{Z_0}
       \simeq\Perf\left(\prod_iE_i\right)
\end{equation}
である。必要な圏は全て冪等完備化して読む。
全 coherent tails 圏と残余圏は smooth proper である。
従ってこの場合は、全有限 MF に対しても
\[
 \bigotimes_i^{\mathbf L}\MF^{\mathbb Z}(S_i,w_i)
    \simeq\MF^{\Lambda_k}(B,w)
\]
という Morita 同値が成り立つ。
\end{corollary}
\begin{proof}
補題\ref{mbnp:lem:noetherian}は第\ref{crs:sec:saturation}節の全環論的仮定を満たす。
具体的には、ambient は連結局所有限、両側 noetherian、Koszul AS-regular で、
dimension と parameter が $N=3k\geq6$、商の中心元は正則三次である。
従って補題\ref{crs:lem:sections}により、全 coherent 対象 $M$ に対する
$R\operatorname{Hom}(\OO_A(j),M)$ は体上 perfect で、全 twist は対象を検出する。
補題\ref{crs:lem:finite-lines}はその twist 全体を有限個の線対象で生成する。

定理\ref{mbnp:thm:product}では、この全 coherent 圏との一致をまだ使わずに
$\mathcal K_A^{\mathrm{line}}$ の smoothness を得た。
有限例外部分との半直交分解に沿って smoothness を貼り合わせ、
命題\ref{crs:prop:residual-smooth}および定理\ref{crs:thm:coherent}を適用すれば
$\mathcal C_A=D^b(\operatorname{qgr}A)$ と残余圏の一致が従う。
properness は線対象間の有限 Hom とこの有限生成性による。
最後に全 coherent 圏への Orlov 分解を適用すると、
補題\ref{mbnp:lem:window}の augmentation 部分が全特異圏に一致する。
従って前の外積比較の像は全 $\Lambda_k$ 次数付き MF である。
これは生成性を最初から仮定した議論ではない。
\end{proof}

\begin{remark}[体上の比較、translation、相対核の区別]
\label{mbnp:rem:scope}
本節の直接比較は、基本楕円開の任意の体上の実表示に適用される。
従って定理\ref{med:thm:multiblock-generic}と
定理\ref{mpc:thm:point-model}における明示的一変数の一般点という仮定を、
\emph{中心同変化した残余圏と点曲線積との比較については}取り除く。
これらの以前の定理が保持する、指定された完備復元核との関係まで置き換えるものではない。

通常の残余圏は、積側の
$R_{i,p_i}^{-1}R_{b,p_b}$ による双対同変化として得られる。
可換分離和の $R_{i,0}^{-1}R_{b,0}$ との同定には、
各因子の Picard 補正と指定周期の同期を用いる。
体上のこの補正と有限拡大後の通常残余圏比較は
定理\ref{mbrc:thm:finite-point-comparison}で得られる。
順序付き実表示上の相対積核は定理\ref{mbrn:thm:relative-product}で構成し、
任意の noetherian 底変更には系\ref{mbrn:cor:base-change}の
compact/perfect 比較を用いる。
さらに、自然な可換標的への Picard 補正型の相対核と指定周期の障害は
定理\paperref{moduli}{mbpd:thm:coherent-corrections}および
系\paperref{moduli}{mbpd:cor:descent-criterion}で分類する。
表示を忘れた全O stack 上の大域切断、全ての比較核の型、
非正則な底変更後の全 coherent 圏の比較へは、これらの結論を拡大しない。
核の存在 stack の降下と選んだ核の大域切断は、
命題\ref{mbrd:prop:kernel-stack}の場合と同様に区別される。
\end{remark}

% !TEX root = ../nc-cubic.tex
\section{実 Sklyanin 表示上の相対非可換積比較}
\label{mbrn:sec:relative-nc}

第\ref{mbnp:sec:product}節の直接比較を、滑らかな基底上の実表示へ相対化する。
相対 tensor 環の noetherian 性には、訂正を含む Keeler の
相対 Fujita 消滅\cite[Theorem~1.5, Definition~2.9, Theorem~E2.5]{NFMKeelerAmple}
を使う。これによって Burke--Stevenson の window の仮定を基底上で確認し、
有限 Koszul--Shamash 模型と外積を実際の dg 関手として構成できる。
各ファイバーでの同値の存在だけから相対同値を推論するものではない。
比較は関係直線を保持し、非正則な底変更後には compact 圏と perfect 圏について述べる。

\subsection{順序付き実表示と共通の関係直線}

$S$ を滑らかな準射影 $\mathbb C$-scheme、$k\geq2$、
$N=3k$、$r=N-3$ とする。
第\ref{mbs:sec:presentations}節の基本楕円開から来る、
順序付きの実際の相対代数
\[
 S_i=T_{\OO_S}(W_i)/(R_i),\qquad
 B=\bigotimes_{\OO_S,i}S_i,\qquad A=B/(\ell)
\]
を固定する。各 $W_i$ は階数三、点曲線
$\pi_i:E_i\to S$ は滑らかな平面 genus-one 曲線族、
$L_i=\OO_{E_i}(1)$、$\sigma_i=t_{p_i}$ とする。
$p_i,2p_i,3p_i\ne0$ は全ての幾何ファイバーで満たされると仮定する。
$j_i$ の相異性は本節の比較には必要ない。
中心直線を $C_i\subset(S_i)_3$ とし、
\begin{equation}\label{mbrn:eq:relation-line}
 \ell\subset\bigoplus_i C_i,\qquad
 \ell\xrightarrow{\sim}C_i,\qquad
 \mathscr L=\ell^\vee,\qquad Y=\prod_SE_i
\end{equation}
とする。局所 frame $e$ に対する中心元を $w_i$ と書けば、
$w=\sum_iw_i$ は同じ $e$ に対する総中心元である。
曲率を座標に依存せず書くと
$w_i\in(S_i)_3\otimes\mathscr L$、
$w\in B_3\otimes\mathscr L$ となる。

以下の MF は内部 shift $M(1)_d=M_{d+1}$ の規約で
\begin{equation}\label{mbrn:eq:mf-period}
 P^{-1}\longrightarrow P^0\longrightarrow
           P^{-1}(3)\otimes\mathscr L,\qquad
 [2]=(3)\otimes\mathscr L
\end{equation}
と書く。係数直線 $\mathscr L$ は内部次数零として数え、
曲率の内部次数三は別に保持する。
$\Gamma_k,\Lambda_k,Z_0$ は第\ref{mbnp:sec:product}節と同じものとし、
$\chi=3e_i\in\Lambda_k$ と置く。
多重次数では共通周期は $(\chi)\otimes\mathscr L$ である。

相対 compact 圏を $\mathscr C_A=D(\QGr A)^c$、
その有限例外部分の右直交を
\[
 \mathscr K_A=
 \langle\OO_A,\OO_A(1),\ldots,\OO_A(r-1)\rangle^\perp
       \subset\mathscr C_A
\]
と書く。非アフィン基底では $\Perf(S)$ の作用と Zariski 降下を含む
相対線型圏の stack として読む。
定理\ref{mbsd:thm:serre}により、この圏は局所的に
有限個の線対象と mutation によって生成され、
その Hom は係数 perfect で任意の導来底変更と両立する。

\subsection{任意の noetherian 基底における環論}

\begin{lemma}[中心商、相対消滅、tensor 環]
\label{mbrn:lem:noetherian}
上の実表示を任意の noetherian $\mathbb C$-代数 $R$ へ引き戻す。
各因子の中心商には、任意の底変更に両立する同型
\begin{equation}\label{mbrn:eq:relative-thcr}
 T_i:=S_i/(C_i)\simeq
  \bigoplus_{n\geq0}H^0(E_i,L_{i,n}),\qquad
 L_{i,n}=L_i\otimes\sigma_i^*L_i\otimes\cdots
                        \otimes(\sigma_i^{n-1})^*L_i
\end{equation}
がある。$T_i$、$\bigotimes_RT_i$、$B$、$A$ は両側 noetherian である。
これらの noetherian 性は、さらに任意の noetherian 係数環への変更後にも成り立つ。
また $\operatorname{qgr}T_i\simeq\operatorname{coh}E_i$ であり、
右加群の shift は $(\sigma_i)_*\circ(-\otimes L_i)$ に対応する。
\end{lemma}
\begin{proof}
まず式\eqref{mbrn:eq:relative-thcr}を普遍チャート上で作る。
式\eqref{mbs:eq:relative-relations}による生成元の評価は、
$S_i$ から右 THCR への自然な次数付き準同型を与える。
普遍底は被約であり、$C_i$ の像は各幾何点で零なので、束の写像として零である。
各次数で $C_i\otimes(S_i)_{n-3}\to(S_i)_n$ は
全ファイバーで単射である。
従ってその余核は有限局所自由で、底変更に両立する。
他方、$n>0$ の $L_{i,n}$ は相対次数 $3n$ であり、
その $H^1$ は零、$H^0$ は階数 $3n$ の有限局所自由加群である。
$n=0$ では $(\pi_i)_*\OO_{E_i}=\OO_S$ である。
各次数の自然な比較は幾何ファイバーで既知の中心商同型なので、束の同型となる。
これを引き戻せば非被約な試験基底でも式\eqref{mbrn:eq:relative-thcr}を得る。

次に $X=\prod_RE_i$ とし、
$L_{\boldsymbol n}=\boxtimes_iL_{i,n_i}$、
$H=\boxtimes_iL_i$ と置く。
$H$ は $\Spec R$ 上 ample である。
任意の coherent 層 $F$ に対し、
\cite[Theorem~1.5]{NFMKeelerAmple}による整数 $m_0$ を取る。
全ての $n_i\geq m_0$ について
\[
 N_{\boldsymbol n}=L_{\boldsymbol n}\otimes H^{-m_0}
\]
は相対 nef である。実際、各幾何ファイバーでこれは
次数 $3(n_i-m_0)\geq0$ の因子線束の外積である。
引用の nef は \cite[Definition~2.9]{NFMKeelerAmple}の相対 nef を意味する。
従って
\begin{equation}\label{mbrn:eq:uniform-vanishing}
 H^q(X,F\otimes L_{\boldsymbol n})=0
 \quad(q>0,\ n_i\geq m_0)
\end{equation}
を得る。同定理の証明の訂正は同文献 Theorem~E2.5 に含まれる。
この適用は $F$ の $R$-flatness を要求しない。
同じ議論を $\sigma_{\boldsymbol a}^*L_{\boldsymbol n}$ に適用できる。

各 $L_{\boldsymbol m}$ は相対大域生成され、
$T_{\boldsymbol n}:=H^0(X,L_{\boldsymbol n})$ は
$\bigotimes_i(T_i)_{n_i}$ に等しい有限射影 $R$-加群である。
多重斉次右イデアル $J\subset T=\bigotimes_iT_i$ に対し、
補題\ref{mbnp:lem:noetherian}と同じ評価イデアル層
\[
 I_{\boldsymbol n}=
  \operatorname{im}(J_{\boldsymbol n}\otimes L_{\boldsymbol n}^{-1}
                      \longrightarrow\OO_X)
\]
を取る。大域生成性により
$I_{\boldsymbol n}\subset I_{\boldsymbol n+\boldsymbol m}$ であり、
noetherian な $X$ 上で対角に沿って安定する。
従ってある $\boldsymbol n_0$ 以降では $I_{\boldsymbol n}=I$ となる。
評価全射
$J_{\boldsymbol n_0}\otimes_R\OO_X\to I L_{\boldsymbol n_0}$
の核に式\eqref{mbrn:eq:uniform-vanishing}を適用すると、
全ての十分大きい各成分を持つ $\boldsymbol m$ について
\[
 J_{\boldsymbol n_0}T_{\boldsymbol m}
   =H^0(X,I L_{\boldsymbol n_0+\boldsymbol m})
   =J_{\boldsymbol n_0+\boldsymbol m}
\]
を得る。
ここで $J_{\boldsymbol n_0}$ は有限 $R$-加群であり、射影とは限らない。
使用する等式
\[
 H^0(X,J_{\boldsymbol n_0}\otimes_R
                    \sigma_{\boldsymbol n_0}^*L_{\boldsymbol m})
 =J_{\boldsymbol n_0}\otimes_R
                    H^0(X,\sigma_{\boldsymbol n_0}^*L_{\boldsymbol m})
\]
は、$m_i>0$ のとき右側の切断複体が degree zero の有限射影加群であることと、
$X/R$ の flatness および射影公式による。
従って非平坦な評価イデアルにも同じ議論を適用できる。

残る有限幅の面 $n_i=d<c_i$ は、
有限射影な $(T_i)_d\otimes_R\bigotimes_{j\ne i}T_j$ の
右部分加群である。因子数に関する帰納法によりこれらは有限生成となり、
$J$ 全体も有限生成となる。
多重次数を全次数、次に辞書式で並べて初項を取ると、
同じ有限生成性は非斉次イデアルにも移る。
初項を消去するたびにこの順序は下がり、固定全次数以下の多重次数は有限なので、
消去は有限回で終わる。
opposite の THCR に同じ証明を適用すれば左右の noetherian 性を得る。

局所的に中心直線の frame を取り、$w_i$ をその生成元とする。
各次数の射影な商 $S_i/(w_i)$ の $R$-線型切断を取ると、
$S_i$ は $R[w_i]$ 上、有限射影 $R$-加群から誘導した加群の次数付き直和である。
体上の自由性を、そのまま $R$ 上で仮定しているのではない。
従って中心列の随伴次数環は
\[
 \operatorname{gr}_{(w_1,\ldots,w_k)}B
       \simeq T[t_1,\ldots,t_k]
\]
であり、両側 noetherian である。
固定内部次数では中心イデアルによる濾過が有限なので、
補題\ref{mbnp:lem:noetherian}の初項持上げがそのまま通る。
これで $B$、その商 $A$ の両側 noetherian 性を得る。
中心直線を自明化する有限アフィン開被覆からの降下も、
中心係数の局所化によるイデアルの有限生成性の判定で行える。

一因子について、式\eqref{mbrn:eq:uniform-vanishing}および同じ議論による大域生成性は
右 THCR の ample pair 条件を与える。
\cite[Proposition~7.2, Lemma~7.6]{NFMKeelerAmple}により、
$\operatorname{qgr}T_i\simeq\operatorname{coh}E_i$ となる。
右加群の shift の向きは式\eqref{eq:nfm-right-thcr-shift}の規約である。
全議論は任意の noetherian $R$ 上で通るため、主張した底変更後の性質も従う。
\end{proof}

\subsection{正則基底上の実 window と有限 Hom}

\begin{lemma}[相対三対象 window]
\label{mbrn:lem:window}
$S=\Spec R$ が滑らかなアフィン基底であるとする。
有限射影 twisted graded MF の Morita 完備化を
$\mathscr M_i=\MF^{\mathbb Z}_R(S_i,w_i;\mathscr L)$ と書き、
$K_i$ を augmentation $R$ の stabilization とする。
このとき
\begin{equation}\label{mbrn:eq:factor-window}
 \mathscr M_i\simeq\Perf(E_i),\qquad
 K_i\longleftrightarrow\OO_{E_i}[1]
\end{equation}
である。$\mathscr M_i$ は $K_i,K_i(1),K_i(2)$ によって $R$-線型に厚生成される。

総和については $K_w=\operatorname{stab}(R)$ と置くと、
実際の $R$-線型 dg quasi-functor
\begin{equation}\label{mbrn:eq:total-window}
 \Psi_R=\operatorname{stab}\circ q_{\mathrm{sg}}
                          \circ(R\omega(-))_{\geq1-r}
 :\mathscr K_A\xrightarrow{\sim}
       \operatorname{thick}_R\{K_w,K_w(1),K_w(2)\}
\end{equation}
がある。右辺の Hom と全ての一因子 Hom は $R$-perfect である。
これらの比較は関係直線を含む指定周期を保持する。
\end{lemma}
\begin{proof}
$d=\operatorname{gldim}R<\infty$ とする。
定理\ref{mbsd:thm:ambient-duality}の長さ $N$ の相対双加群解消を
任意の $B$-加群へ tensor すると、各項は $R$-加群からの誘導加群である。
それらの $B$-射影次元は高々 $d$ なので、
$\operatorname{gldim}B\leq N+d$ である。
中心関係の二項解消による change of rings は、左右で
$\operatorname{id}_A A\leq N+d-1$ と
\[
 R\underline{\operatorname{Hom}}_A(R,A)
     \simeq R(r)\otimes\mathscr D^\vee[-N+1],
 \qquad \mathscr D=\det(\bigoplus_iW_i)\otimes\ell^\vee
\]
を与える。
\cite[Definition~5.1, Remark~5.2]{LP-BS}は、ここに現れる cohomological shift と
injective dimension の一致を要求しない。
従って正の次元を持つ $R$ でも適用できる。
係数直線 $\mathscr D$ は、その frame を取った開で同定理の仮定を確認すればよい。
section と window の関手自体はこの frame に依存しない。
一因子にも $N=3,r=0$ として同じ議論を適用する。

補題\ref{mbrn:lem:noetherian}、$R$-flatness と
\cite[Lemmas~6.1--6.2, Proposition~4.18, Theorem~6.4]{LP-BS}により、
導来 section、有限 cohomological dimension、$\chi$ および Orlov 分解を得る。
有限 MF と特異圏の比較には \cite[Theorem~7.2]{AMChen}を用い、
補題\ref{tw:lem:stabilization}と\ref{mbnp:lem:window}の
exact 圏・有限 totalization による dg 化を行う。
右加群には opposite 環を使い、射影表示と持上げの次数零成分を取る。
係数直線付きの場合は局所 frame で同じ余核構成を行い、
その二拡張を関係直線とともに貼り合わせる。

一因子の tails は $\operatorname{coh}E_i$ であり、
$E_i$ は正則なので $D^b(\operatorname{coh}E_i)=\Perf(E_i)$ である。
parameter zero の window1 が式\eqref{mbrn:eq:factor-window}を与える。
三段 augmentation Koszul 解消を tails へ写すと、
三つの連続する線対象が全 twist を $\Perf(R)$-線型に生成する。
ample pair の全 twist は perfect 圏を生成するので、三対象の生成性も従う。
$R\Gamma_{\geq1}\OO_{E_i}=(T_i)_{\geq1}$ により
$K_i\leftrightarrow\OO_{E_i}[1]$ となる。

総和の window の完全忠実性は、上の Burke--Stevenson 定理を
$i=1-r$ で直接適用して得る。
その核は $\OO_A,\ldots,\OO_A(r-1)$ の $R$-線型生成部分である。
定理\ref{mbsd:thm:serre}の compact generator
$\bigoplus_{j=0}^{N-1}\OO_A(j)$ 上では
$(R\omega\OO_A(j))_{\geq1-r}=A(j)_{\geq1-r}$ である。
$U_j=A(j)_{<1-r}$ の有限 filtration は
\[
 \operatorname{gr}U_j=
       \bigoplus_{u=0}^{j-r}A_u\otimes_RR(j-u)
       \quad(j=r,r+1,r+2)
\]
であり、$j<r$ では $U_j=0$ である。
係数 $A_u$ は有限射影であり、
$\Psi_R\OO_A(j)=\operatorname{stab}(U_j)[-1]$ だから、
補題\ref{mbnp:lem:window}と同じ三角的な計算が通る。
$K_w(r),K_w(r+1),K_w(r+2)$ を $K_w,K_w(1),K_w(2)$ に直す際には、
式\eqref{mbrn:eq:mf-period}の $\mathscr L$ の tensor も保持する。

一因子 Hom の $R$-perfect 性は $E_i/R$ の smooth proper 性による。
総和の augmentation 部分の Hom は、
式\eqref{mbrn:eq:total-window}と定理\ref{mbsd:thm:serre}から $R$-perfect となる。
全 coherent 圏の生成性やファイバーからの smoothness の持上げを、
この段階の前提にはしていない。
period は式\eqref{tw:eq:period-extension}の余核二拡張を
関係直線付きで取ったものなので、選んだ曲率の scalar 倍を途中で挿入しない。
\end{proof}

\begin{corollary}[位数二の平行移動を含む一因子 window]
\label{mbrn:cor:factor-boundary}
$R$ を有限大域次元を持つ正則 noetherian $\mathbb C$-代数とする。
この系では、滑らかな有限型基底の局所化、特にその DVR も許す。
$E/R$ は滑らかな平面 genus-one 曲線族、$L=\OO_E(1)$、
$\sigma=t_p$ とし、幾何的な二次関係から構成した実 Sklyanin 族 $D$ を考える。
全ての幾何ファイバーが三次元 quadratic AS-regular であると仮定する。
直線束 $\ell\to D_3$ は中心に入り、全ての幾何ファイバーで
非零の正則中心三次を与えるものとする。
各幾何ファイバーで、この中心商は当該点曲線の THCR に一致するとする。
後の条件は通常の滑らかな楕円 Sklyanin 中心商に対する条件であり、
相対 window の存在を仮定するものではない。

このとき $2p\ne0$ を仮定せずに
\[
 \MF^{\mathbb Z}_R(D,w;\ell^\vee)\simeq\Perf(E),\qquad
 \operatorname{stab}(R)\longleftrightarrow\OO_E[1],\qquad
 (1)\longleftrightarrow
       R_p=T_{\OO_E/R}\circ\sigma_*\circ(-\otimes L)
\]
が成り立つ。指定周期は
$R_p^3\otimes\ell^\vee\simeq[2]$ であり、
実際の $\ell\to D_3$ の余核二拡張から定まる。
これらの compact/perfect 模型と period は任意の noetherian 底変更に両立する。
\end{corollary}
\begin{proof}
この系の仮定では、補題\ref{mbsd:lem:potential-line}の
最上 syzygy と $\det W$ の同定を使用しない。
代わりに、二次関係を $Q$ として、そのままの直線束
\[
 P=(W\otimes_RQ)\cap(Q\otimes_RW)\subset W^{\otimes3}
\]
を保持する。
各次数の $D_n$ は有限行列の余核として表され、
幾何ファイバーの Hilbert 級数は $(1-t)^{-3}$ で一定である。
$R$ は被約なので、一定階数の小行列式の議論により
各 $D_n$ は有限射影となり、その形成は任意底変更と両立する。
$P$ も同じ議論で直和因子となる直線束である。

Sklyanin の cyclic potential から得る長さ三の
Koszul 双加群複体と片側 augmentation 複体を使う。
その各内部次数は有限射影 $R$-加群の有界複体であり、
全幾何ファイバーで完全である。
右端から kernel を局所自由にする有限射影複体の Nakayama の議論により、
これらは次数別に $R$ 上分裂完全である。
従って底変更後も完全で、相対 Koszul 解消を与える。
cyclic potential の自己双対性では top line を $P$ のまま残し、
\[
 R\underline{\operatorname{Hom}}_D(R,D)
       \simeq R(3)\otimes P^\vee[-3],\qquad
 R\underline{\operatorname{Hom}}_{D/(\ell)}(R,D/(\ell))
       \simeq R\otimes P^\vee\otimes\ell[-2]
\]
を得る。
$P\to\det W$ が零になるファイバーがあっても、この自己双対性には影響しない。

幾何的な二次関係からの評価は $D$ から相対 THCR への準同型を与える。
その $\ell$ 上の像は全幾何ファイバーで零であり、
$R$ の被約性と THCR の次数別局所自由性により零である。
また中心元の乗法は各ファイバーで単射なので、
$D/(\ell)$ も次数別に有限射影である。
商から相対 THCR への写像は各幾何ファイバーで同型であるから、
次数ごとの束の同型となる。
この THCR は補題\ref{mbrn:lem:noetherian}の一因子の議論で両側 noetherian である。
その次数別射影な商の切断と中心イデアルの濾過から $D$ も両側 noetherian となる。
有限双加群解消により $\operatorname{gldim}D\leq3+\operatorname{gldim}R$、
中心二項解消により商の左右 injective dimension は有限である。
係数直線 $P^\vee\otimes\ell$ の frame を取って
\cite[Definition~5.1, Remark~5.2, Theorem~6.4(3)]{LP-BS}を適用できる。
window1 自体はその frame に依存せず、既に構成した
余核、section、内部切断と相対評価 cone から定まる。
補題\ref{mbrn:lem:finite-models}の finite folding と Hom の底変更の証明も、
この三段解消と指定中心元だけを使うので適用できる。
曲線側の三つの線対象の生成と perfect 底変更により、
非正則な変更後にも、引き戻した compact/perfect 模型の比較を得る。

具体的に、滑らかな Hesse 曲線の
$a=b$、$c\ne0$ で、対応する平行移動が非零の位数二である点はこの系に含まれる。
その Sklyanin regularity と中心商は
\cite[Theorems~6.6, 6.8]{GRSklyaninATV}
\cite[Section~1]{GRSklyaninST}の楕円データの条件に含まれる。
このとき raw potential は非零であり、
中心元 $\Omega_H=a(a^3-c^3)(U-V_0)$ も該当する滑らかな領域で非零である
（式\eqref{mbpl:eq:hesse-center}の記号と式を用いる）。
実際、式\eqref{mbpl:eq:hesse-invariant-relation}の自由 tensor の計算は $a=b$ でも有効であり、
$U-V_0$ は $cC+aU+aV_0$ の倍数ではない。
$a^3=c^3$ の場合は点三次が特異となるので、残る係数も非零である。
したがって、この一因子の比較のために全O認識の基本開を拡張する必要はない。
\end{proof}

\begin{lemma}[有限 MF 模型、相対外積、底変更]
\label{mbrn:lem:finite-models}
補題\ref{mbrn:lem:window}の augmentation 対象には、有限射影な Koszul--Shamash 模型を取れる。
その外積は総和の augmentation stabilization と同定される。
有限生成対象間の MF Hom は、有界な有限射影 $R$-複体と homotopy 同値であり、
任意の係数変更後の raw MF Hom を導来底変更によって計算する。
\end{lemma}
\begin{proof}
各因子の相対 augmentation 分解 $P_{i,\bullet}\to R$ は、
定理\ref{mbsd:thm:ambient-duality}により有限 graded projective 分解である。
その End 複体の負 Ext は零であり、曲率は $R$ を殺すので、
補題\ref{mbnp:lem:shamash}の higher homotopies を $R$ 上で直接選べる。
関係直線を含めた次数は
\[
 s_a:P_j\longrightarrow P_{j+2a-1}(3a)\otimes\mathscr L^a
\]
である。一因子では $s_1,s_2$ だけで足りる。
ここではファイバーごとに選んだ homotopy が貼り合うとは仮定していない。
folding は
\[
 F^0=\bigoplus_{j\ {\rm even}}P_j(3j/2)\otimes\mathscr L^{j/2},
 \qquad
 F^{-1}=\bigoplus_{j\ {\rm odd}}
             P_j(3(j-1)/2)\otimes\mathscr L^{(j-1)/2}
\]
であり、対応する Shamash 分解は
\[
 C_n=\bigoplus_{q\geq0}
       \overline P_{n-2q}(-3q)\otimes\ell^q
\]
である。
その第一頁は
$\operatorname{Tor}^{B}_j(R,A)=R$ $(j=0)$、
$R(-3)\otimes\ell$ $(j=1)$ から成る。
$s_1$ の連結射は隣接する二項を $\pm\operatorname{id}$ で結ぶので、
第二頁には augmentation の $R$ だけが残る。
濾過は各全次数で有限であり、この証明は体を必要としない。
全式を任意の係数環へ引き戻しても同じ恒等式が成り立つ。

外積では共通周期 $\Theta=(\chi)\otimes\mathscr L$ を使い、
奇成分どうしの tensor に $\Theta$ を一つ付ける。
式\eqref{mbrp:eq:exterior-components}の有限行列がこの外積を与える。
各 higher homotopy は奇 homological degree を持つので、
異なる因子の cross term は Koszul 符号により消える。
従って
\[
 K_1\boxtimes\cdots\boxtimes K_k\simeq K_w
 \quad\text{in }\MF^{\Lambda_k}_R(B,w;\mathscr L).
\]
floor の恒等式を用いた folding の同定は
補題\ref{mbnp:lem:shamash}と同じであり、
奇成分の補正には $\mathscr L$ の冪も含まれている。

次に Hom を考える。
有限射影次数付き因子について、式\eqref{mbnp:eq:hom}の Hom 分解は
体の K\"unneth 公式を使わず、複体そのものの同型である。
係数環を $R$ に替えても、有限射影加群の Hom と
$\mathbb Z^k\to\Lambda_k$ の指標の配分による同じ証明が通る。
正次数性により固定した全 cohomological degree の和は有限である。
共通直線の冪は両辺で同じなので、frame の変更にも両立する。

各 raw Hom 複体の項は有限射影 $R$-加群であり、
その導来同型類は補題\ref{mbrn:lem:window}により $R$-perfect である。
有界な有限射影模型からの準同型を取り、その cone を考える。
$R$ の大域次元が有限なので、射影加群からなる acyclic 複体の
全 cycle は dimension shifting により射影である。
従ってこの cone は可縮であり、raw Hom は有界射影模型と homotopy 同値である。
特に $K$-flat で、任意の $R\to R'$ に対し
\begin{equation}\label{mbrn:eq:hom-basechange}
 \operatorname{Hom}_{\MF_R}^\bullet(P,Q)\otimes_R^{\mathbf L}R'
  \simeq
 \operatorname{Hom}_{\MF_{R'}}^\bullet(P_{R'},Q_{R'})
\end{equation}
となる。右辺は選んだ有限 MF をそのまま引き戻した対象の Hom を意味する。
$R'$ 上の全 coherent MF 圏との一致をこの式から結論してはいない。
\end{proof}

\subsection{相対積同値と実 Morita 核}

\begin{theorem}[実表示上の相対非可換積比較]
\label{mbrn:thm:relative-product}
式\eqref{mbrn:eq:relation-line}の順序付き実表示について、
$S$-線型 dg Morita 同値
\begin{equation}\label{mbrn:eq:relative-product}
 \Phi_S:\Perf(Y/S)\xrightarrow{\sim}\mathscr K_A^{Z_0}
\end{equation}
がある。これは一因子 window、有限 MF の外積、総和 window の逆の合成であり、
実際の perfect Morita 核によって表される。
基底の縮小、有限被覆、関係直線の三乗根の選択は必要ない。
$\mathscr K_A^{Z_0}$、$\mathscr K_A$、$\mathscr C_A$ は $S$ 上 smooth proper である。
\end{theorem}
\begin{proof}
まず $S=\Spec R$ とする。
補題\ref{mbrn:lem:finite-models}の Hom 分解と $K$-flat 性により、外積は
\[
 \mathscr M_1\otimes_R^{\mathbf L}\cdots\otimes_R^{\mathbf L}\mathscr M_k
       \longrightarrow\MF^{\Lambda_k}_R(B,w;\mathscr L)
\]
からの完全忠実な dg 関手を与える。
各因子は三つの augmentation shift で生成され、
その外積は $K_w$ である。
全内部 shift と $\Perf(R)$ の作用を取れば、像は
$\operatorname{thick}_R\{K_w(\lambda):\lambda\in\Lambda_k\}$ となる。

window は内部 grading の作用と可換する。
$\Gamma_k$ の分裂を一つ選んで通常の $Z_0$ 同変化として書けば、
有限群の誘導と平均化により、augmentation 部分の同変化は
全 $\Lambda_k$ shift で生成される。
従って総和 window はこの像を $\mathscr K_A^{Z_0}$ と同定する。
分裂を変えたときの比較は命題\ref{mbza:prop:gamma-action}の
scalar 自明化によって与えられる。
一因子 window と smooth proper scheme の相対外積
\cite[Section~8.3, Lemma~8.11, Theorem~8.15]{MBRPToenMorita}を合成すると、
式\eqref{mbrn:eq:relative-product}を得る。

$Y/S$ が smooth proper なので $\mathscr K_A^{Z_0}$ は $R$ 上 smooth proper である。
双対群の指標作用と二重同変化は $\mathbb C$ の全三乗根および有限 Fourier 分解で定まり、
$R$-線型にも成立する\cite[Theorem~4.2, Section~8]{MEDElagin}。
有限軌道の直和を generator に取って作用を厳密な dg 代数作用へ移し、
補題\ref{med:lem:crossed-perfect}を適用すれば $\mathscr K_A$ も smooth proper である。
さらに有限例外部分は $\Perf(R)$ の有限個のコピーである。
proper な gluing 加群は smooth な両成分上 perfect なので、
半直交分解の対角の有限 cell 表示から $\mathscr C_A$ も smooth となる。
properness は定理\ref{mbsd:thm:serre}で既に得られている。

この関手は実際の Morita 核を与える。
両辺の有限 generator の End 代数をそれぞれ $D_R,C_R$ と書けば、
構成した dg 関手は $D_R^{\mathrm{op}}\otimes_RC_R$ の双加群で表される。
両圏が smooth proper なので、この双加群と Morita 逆は perfect であり、
単位・余単位も実際の双加群の射として存在する。
非可換側を scheme と見なした積上の層を仮定しているのではなく、
ここでの核はこれらの dg 代数間の perfect Morita 双加群である。

一般の $S$ では、上の構成をアフィン開ごとに行う。
関手の構成は higher homotopy の選択によって定義したものではない。
section、内部切断、特異 dg 商、有限 MF の余核比較、
および twisted MF の外積によって定義しており、
higher homotopies は像の生成性を確認するための有限模型である。
これらの構成は開制限に自然である。
余核同値と window 同値の逆は Morita $\infty$-圏で取り、
同値の逆の選択に関する可縮な空間を用いて合成の coherence を保持する。

関係直線の frame を $e'=u e$ と変えると、局所曲率は全因子で同時に
$w_i'=u w_i$ となる。
有限 factorization の一方の矢印と、その反復外積の
$\mathscr L^{\lfloor|\epsilon|/2\rfloor}$ の遷移がこれを補正する。
これは式\eqref{mbrn:eq:mf-period}を持つ twisted MF の同じ対象の二つの表示であり、
三重重なりでは通常の直線束の cocycle 条件を満たす。
従って関係直線を自明化する追加の選択を大域的に残さない。
perfect 加群と有限係数の加群作用の降下により、
関手・核・逆・単位・余単位が $S$ 上へ貼り合う。
smoothness と properness もこの局所表示によって従う。
\end{proof}

\begin{proposition}[相対 rotation、指定周期、双対群]
\label{mbrn:prop:rotations}
一因子の相対評価 cone を
\[
 T_{\OO_{E_i}/S}(F)=
 \operatorname{Cone}\bigl(\pi_i^*R(\pi_i)_*F\longrightarrow F\bigr)
\]
とし、
\[
 R_{i,p_i}=T_{\OO_{E_i}/S}\circ(\sigma_i)_*\circ(-\otimes L_i)
\]
を $Y$ の当該因子に作用させる。
一因子 window は内部 $(1)$ を $R_{i,p_i}$ へ移し、
\begin{equation}\label{mbrn:eq:factor-period}
 R_{i,p_i}^3\otimes\mathscr L\xrightarrow{\sim}[2]
\end{equation}
は指定した関係直線の余核二拡張から移された周期である。
これらの関手は実際の相対 perfect 核を持つ。

基準因子 $b$ による線型化付き residual rotation を $\rho_{A,b}$ とすると、
\begin{equation}\label{mbrn:eq:rotation-comparison}
 \Phi_SR_{b,p_b}\simeq\rho_{A,b}\Phi_S,\qquad
 \mathsf T_{e_i/e_b}\Phi_S\simeq
                    \Phi_S(R_{i,p_i}^{-1}R_{b,p_b}).
\end{equation}
後式は双対群 $\widehat Z_0$ の coherent な作用を表す。
第一式を三回合成して $\mathscr L$ を tensor した図式は、
式\eqref{mbrn:eq:factor-period}と
$\rho_{A,b}^3\otimes\mathscr L\simeq[2]$ を保持する。
これらの作用、二重同変化および period 図式は底変更に自然である。
\end{proposition}
\begin{proof}
一因子について、補題\ref{lem:nfm-twisted-elliptic-window}の切断三角を
$R$ 上でそのまま取る。
$K_i\leftrightarrow\OO_{E_i}[1]$ なので、内部 shift は相対評価 cone と
右 THCR shift の合成になる。
$E_i/S$ は smooth proper であり、評価 cone の核は
$\operatorname{Cone}(\OO_{E_i\times_SE_i}\to\OO_{\Delta_i})$ である。
graph の核および $L_i$ との tensor を合成すれば、$R_{i,p_i}$ の perfect 核を得る。

総和の切断・mutation の比較も式\eqref{tw:eq:rotation-cut}によって与えられる。
余核から得る二拡張は
\[
 0\longrightarrow M(-3)\otimes\ell\longrightarrow
       \overline P^{-1}\longrightarrow\overline P^0
                       \longrightarrow M\longrightarrow0
\]
である。その ambient への持上げの合成は、
指定した $\ell\to B_3$ の乗法である。
従って period 比較の正方形\eqref{tw:eq:rotation-period-square}は
$\mathscr L=\ell^\vee$ を含めて成立する。
この構成で得る period を独立な scalar 倍へ取り替えてはいない。

内部 $(1)$ は表現の指標 $e_i^{-1}$ の tensor なので、
命題\ref{mbnp:prop:rotations}の符号と同じ
$\mathsf T_{e_i/e_b}\leftrightarrow R_{i,p_i}^{-1}R_{b,p_b}$ を得る。
三乗・交換の全関係は共通の $\Gamma_k$ 次数付き MF の関係から移す。
二重同変化に用いる指標射影は
$|Z_0|^{-1}\sum_g\eta(g)^{-1}g$ という係数 $\mathbb C$ の平均化であり、
任意の係数変更と可換する。
これにより作用の関手だけでなく、その coherence と二重同変化も底変更に両立する。
\end{proof}

\begin{corollary}[任意の noetherian 底変更後の compact 比較]
\label{mbrn:cor:base-change}
任意の準コンパクト・分離的 noetherian $S$-scheme $T$ に対し、
定理\ref{mbrn:thm:relative-product}の Morita 核の導来引戻しは
\begin{equation}\label{mbrn:eq:basechange}
 \Perf(Y_T/T)\xrightarrow{\sim}\mathscr K_{A_T}^{Z_0}
\end{equation}
を与える。$T\to S$ の flatness や $T$ の被約性・正則性を仮定しない。
作用、rotation、関係直線付き指定周期、および同値の単位・余単位は
合成した底変更とも coherent に両立する。
$T$ が滑らかな準射影基底である場合、この同値は引戻した実表示からの直接構成に一致する。
\end{corollary}
\begin{proof}
アフィン局所で証明する。
定理\ref{mbsd:thm:serre}と残余 projector の底変更により、
有限 generator $G$ の引戻し $G_T$ は $\mathscr K_{A_T}$ を生成し、
\[
 R\operatorname{End}(G)\otimes_R^{\mathbf L}R'
      \simeq R\operatorname{End}(G_T)
\]
である。有限群への誘導と平均化を施した generator にも同じことが成り立つ。
曲線積側の有限外積 generator とその Hom も、
smooth proper な $Y/S$ の perfect 底変更により同じ性質を持つ。
従って両辺の実際の compact/perfect 圏は、
元の dg 代数模型の導来底変更で計算される。

perfect Morita 核、逆、単位・余単位を tensor すれば、
同値を表す二つの双加群恒等式はそのまま保たれる。
式\eqref{mbrn:eq:hom-basechange}は有限 MF 模型でも同じ係数変更を計算する。
rotation は有限 mutation と係数 perfect な評価で定まり、
period の二拡張も有限 MF の曲率恒等式で定まるから、
命題\ref{mbrn:prop:rotations}の図式も引き戻せる。
導来 tensor の結合同型と平均化の自然性が合成の coherence を与える。

滑らかな $T$ での直接構成との比較については、
window の生成線対象の section は $A(j)_{\geq1-r}$、
その有限切断は $U_j$ という次数別射影な加群で表される。
stabilization には補題\ref{mbrn:lem:finite-models}の有限解消を使える。
従ってこれらの対象・射の自然な係数変更は直接構成のものと一致する。
一因子では smooth proper な曲線の導来 section の perfect 底変更を使う。
有限 cone、直和因子、相対外積と mutation へ延長して関手の比較を得る。
局所のこれらの同定は核と作用写像の降下により $T$ 上へ貼り合う。

一般の非正則な $T$ でここに現れるのは、
表示 $A_T$ の compact 残余圏と $Y_T$ の perfect 圏である。
この結論を $D^b(\operatorname{coh}Y_T)$ 全体や、
独立に定義した非正則基底上の全 coherent MF 圏との同値へ拡大しない。
\end{proof}

\begin{corollary}[元の正則アフィン基底上の coherent 飽和]
\label{mbrn:cor:regular-coherent}
$S=\Spec R$ が上の滑らかなアフィン基底なら、
$\mathscr C_A=D^b_{\mathrm{dg}}(\operatorname{qgr}A)$、
$\mathscr K_A=\mathcal K_A^{\mathrm{coh}}$ である。
\end{corollary}
\begin{proof}
Burke--Stevenson の導来 section により、
任意の bounded coherent 対象 $M$ と固定整数 $j$ について
$R\operatorname{Hom}(\OO_A(j),M)$ は有界有限生成 $R$-cohomology を持つ。
$R$ は正則で有限大域次元を持つので、これは $R$-perfect である。
全ての線対象は同じ section によって coherent 対象を検出する。
$G=\bigoplus_{j=0}^{N-1}\OO_A(j)$、$C=R\operatorname{End}(G)$ と置くと、
定理\ref{mbrn:thm:relative-product}によって $C$ の対角は $C^e$ 上 perfect である。
その有限 cell/retract 表示へ
$R\operatorname{Hom}(G,M)\otimes_C^{\mathbf L}(-)$ を適用すれば、
$R\operatorname{Hom}(G,M)$ は $C$-perfect となる。
評価射の cone は全線対象に右直交するので零であり、$M\in\operatorname{thick}(G)$ である。
これは定理\ref{crs:thm:abstract}の証明を係数 perfect 性を保って実施したものである。
非正則な変更後の係数環には、この $R$-perfect 性の推論を適用していない。
\end{proof}

\begin{remark}[表示の上での核と、その後の降下]
\label{mbrn:rem:scope}
上の定理は、順序付き実表示を固定した $S$ 全体で構成する核の定理であり、
非空開への縮小による存在だけの主張ではない。
一方、次数付き表示を共通スカラーで割った底や全O stack 上への
核の降下は、比較関手の scalar coherence と局所核型の実現をさらに検討する問題である。
また通常残余圏を自然な可換三次族へ比較するためには、
積側の $R_{i,p_i}^{-1}R_{b,p_b}$ を可換因子の作用へ移す Picard 補正と、
指定された関係直線同期に関する period 比較が必要である。
本節はその入力となる実相対積核と、source の period を保持する因子作用を与える。

因子置換は外積の Koszul braiding と $\Gamma_k$ の指標置換によって扱える。
window は共通の整数 charge と全 $Z_0$ 指標を保持するので、
この操作と両立する。
特定の基準因子で線型化した rotation を同時に固定する場合には、
命題\ref{mbza:prop:shift-character}の指標補正を残す。
この置換の整合性を共通表示 gerbe の中立化と読み替えることはしない。
\end{remark}

% !TEX root = ../nc-cubic.tex
\section{点の Picard 類による多ブロック比較と有限定義体}
\label{mbrc:sec:rotation-comparison}

本節では、定理\ref{mbnp:thm:product}の非可換積表示を、
自然な可換三次への比較に接続する。
積圏上での因子 rotation は点の平行移動を含む。
これを実際の次数零線束によって共役し、指定周期の差を三乗根で補正する。
その結果は体上の実際の比較核の存在であり、
モジュライ全体における核の大域切断とは区別する。

\subsection{一因子の Picard 補正と指定周期}

$K\supset\C$ を体、$E/K$ を滑らかな幾何学的連結 genus-one 曲線、
$L$ を次数三の実際の線束とする。
$p\in J_E(K)$ の平行移動を $t_p$ と書き、
右加群の規約で
\[
 R_p=T_{\OO_E}\circ(t_p)_*\circ(-\otimes L),\qquad
 R_0=T_{\OO_E}\circ(-\otimes L)
\]
と置く。
実際の三次方程式と window から移した、coherent な中心周期
\[
 \pi_p:R_p^3\xrightarrow{\sim}[2],\qquad
 \pi_0:R_0^3\xrightarrow{\sim}[2]
\]
も固定する。以下でこれらを任意の scalar 倍に取り替えることはしない。

\begin{lemma}[実 Picard 線束による周期付き補正]
\label{mbrc:lem:factor-correction}
$[M]=p$ を満たす $E$ 上の実際の次数零線束 $M$ があるとする。
$F_M=-\otimes M$ に対し、$K$ 上で同型
\[
 \eta:F_MR_p\xrightarrow{\sim}R_0F_M
\]
が存在する。その全体は $K^\times$-torsor であり、指定周期と両立する
$\eta$ の fpqc 層は $K$ 上の $\mu_3$-torsor である。
従って高々三次の有限拡大後には周期付きの $\eta$ が存在する。
\end{lemma}
\begin{proof}
定理\ref{thm:nfm-elliptic-normalized-intertwiners}において
$C=E$、$\phi=\mathrm{id}$ を取れば、$K$ 上の $\eta$ の存在を得る。
これらの functor は perfect Fourier--Mukai 核を持つ。
その自己準同型環は
$\mathrm{HH}^0(E)=H^0(E,\OO_E)=K$ なので、同型の集合は
$K^\times$-torsor である。

$\eta$ を三回合成し、両側の $\pi_p,\pi_0$ を適用して得る
$F_M[2]$ の自己同型を $u\,\mathrm{id}$ と書く。
ここで $u\in K^\times$ であり、$\eta$ を $v\eta$ に取り替えると
その scalar は $v^3u$ となる。
従って周期条件は $v^3=u^{-1}$ に等しい。
この方程式は有限 \'etale な $\mu_3$-torsor を定め、
$K$ は三乗根を含むので、その分裂体を高々三次に取れる。
中心周期と rotation の coherence は window のものを用いる。
これは定理\ref{mbks:thm:kernel-criterion}の比較直線と
三乗写像を、体上の一因子へ適用した計算でもある。
\end{proof}

一般の $K$ では $p\in\Pic^0(E)(K)$ が実線束で表されるとは限らない。
その障害は
$\delta_E(p)\in\operatorname{Br}(K)$ であり、
系\ref{cor:nfm-elliptic-any-intertwiner-criterion}は、この障害が
全数値型を許した一因子 intertwiner にも必要であることを示す。
ただし、この一因子の必要条件を、そのまま多因子の
\emph{全ての}比較核の必要条件とすることはしない。

\begin{lemma}[次数三の偏極を持つ曲線での有限定義体]
\label{mbrc:lem:degree-three-point}
上の $(E,L)$ に対し、次数一または三の有限分離拡大 $K_1/K$ 上で
$E(K_1)\ne\varnothing$ となる。
従って任意の $p\in J_E(K)$ は $E_{K_1}$ 上の実際の次数零線束で表される。
\end{lemma}
\begin{proof}
$L$ による平面三次表示を用いる。
$K$ は無限体で標数零なので、$K$-有理な直線を、三つの相異なる幾何点で
$E$ と交わるように選べる。
その交わりは $K$ 上次数三の有限 \'etale スキームである。
次数三の分解型は $3$、$2+1$、$1+1+1$ なので、
次数一または三の閉点が存在する。その剰余体を $K_1$ とする。
選んだ点 $o$ により $E_{K_1}$ を楕円曲線とみなし、
$\OO_E(t_p(o)-o)$ を取ればよい。
\end{proof}

\subsection{通常の残余圏への比較}

定理\ref{mbs:thm:presentation-stack}の基本楕円領域に属する
\emph{順序付き実表示}を $K\supset\C$ 上に取り、
\[
 B=\bigotimes_{i=1}^k S_i,\quad
 w=\sum_{i=1}^k a_i\Omega_i,\quad
 A=B/(w),\qquad a_i\in K^\times
\]
とする。各 $S_i$ の点データを $(E_i,L_i,t_{p_i})$ と書く。
定理\ref{mbpl:thm:global-line-comparison}の自然な線同型を用いて
\[
 f_i=a_i\vartheta_i(\Omega_i),\qquad
 X_{\mathrm{pt}}=V\!\left(\sum_i f_i\right)
       \subset\mathbb P\!\left(\bigoplus_i W_i^\vee\right)
\]
と置く。各 $f_i$ は非零で $V(f_i)=E_i$ が滑らかなので、
$X_{\mathrm{pt}}$ も滑らかである。
source と target の周期は、関係直線の
自然な同期 $\ell\simeq\ell_{\mathrm{pt}}$ とこれらの実方程式から固定する。

\begin{theorem}[固定体値点の有限比較と次数上界]
\label{mbrc:thm:finite-point-comparison}
上の仮定の下で、ある有限分離拡大 $K'/K$ 上に
\begin{equation}\label{mbrc:eq:ordinary-comparison}
 \Ku_{\mathrm{coh}}(A_{K'})
       \xrightarrow{\sim}\Ku(X_{\mathrm{pt},K'})
\end{equation}
を与える実際の dg Morita 核が存在する。
この比較は満支持表示の中心群
$Z_0=(\mu_3^k)/\Delta\mu_3$ の作用、rotation、
および自然な関係直線同期に関する指定周期を保つ。
拡大は
\[
 [K':K]=3^e,\qquad 0\le e\le2k
\]
を満たすように選べる。この上界の最小性は主張しない。
特に代数閉体上の全ての上記表示に対して、拡大なしで
式\eqref{mbrc:eq:ordinary-comparison}が成立する。
一変数変形の一般点という仮定は要らない。
\end{theorem}
\begin{proof}
まず補題\ref{mbrc:lem:degree-three-point}を各曲線へ順に適用し、
高々 $k$ 回の三次拡大後に $[M_i]=p_i$ を満たす実線束を選ぶ。
次いで補題\ref{mbrc:lem:factor-correction}を各因子へ適用し、
高々 $k$ 回の三次拡大後に
\[
 F_iR_{i,p_i}\xrightarrow{\sim}R_{i,0}F_i,\qquad
 F_i=-\otimes M_i
\]
を指定周期と両立するように選ぶ。これで次数上界を得る。

非可換側では定理\ref{mbnp:thm:product}と
系\ref{mbnp:cor:coherent}、可換側では
系\ref{mbrp:cor:comparison-field}と
系\ref{mbrp:cor:base-change}を用い、共通の積
$Y=\prod_iE_{i,K'}$ 上に
\[
 \Ku_{\mathrm{coh}}(A_{K'})^{Z_0}
       \simeq\Perf(Y)
       \simeq\Ku(X_{\mathrm{pt},K'})^{Z_0}
\]
を得る。ここで可換側も実方程式 $f_i$ を用い、
各因子の周期を、その方程式の MF/window から移したものとする。
可換側の比較には追加の体拡大を要しない。
実際、$f_i$ の有限個の係数を含む $\C$ 上有限生成な部分体 $K_0\subset K'$ を選ぶ。
$K_0$ 上の曲線も忠実平坦性により滑らかなので、
系\ref{mbrp:cor:comparison-field}の実 perfect 核を適用できる。
その証明に現れる滑らかな有限型模型から
系\ref{mbrp:cor:base-change}で $\Spec K'$ へ引き戻せばよい。
ここで field $K'$ は noetherian であり、$K'/\C$ の有限生成性は不要である。
積上の実際の Fourier--Mukai 自己同値
$F=F_1\boxtimes\cdots\boxtimes F_k$ によって、この二つの積表示を結ぶ。

作用の coherence を確認する。
命題\ref{mbnp:prop:rotations}と命題\ref{mbrd:prop:factor-rotations}により、
双対指標作用の生成元は、基準因子 $b$ に対して
$Q_{i,p}=R_{i,p_i}^{-1}R_{b,p_b}$ と
$Q_{i,0}=R_{i,0}^{-1}R_{b,0}$ $(i\ne b)$ である。
可換側の公式は、上で用いた相対外積と window の構成にもそのまま両立する。
各因子の上の比較射は、三乗周期を保つように正規化済みである。
異なる因子の関手・比較射は、外積の標準交換同型の下で可換する。
比較射の次数は零、共通周期の shift は偶数であるため、
交換の符号も両辺で一致する。
これらの比較射とその逆から
$FQ_{i,p}\simeq Q_{i,0}F$ を作ると、
$Q_{i,p}^3=R_{i,p_i}^{-3}R_{b,p_b}^3\simeq[-2][2]\simeq\mathrm{id}$
を定める正負の周期は、可換側の同じ正負の周期へ移り、
異なる $Q_i$ の交換図式も保たれる。

より具体的には、$\widehat Z_0\simeq(\mathbb Z/3)^{k-1}$ の各元を
$0\le d_i<3$ の正規形 $\prod_{i\ne b}Q_i^{d_i}$ で表す。
その積の coherence は、異なる因子を交換して同じ順序へ直し、
各 $Q_i^3$ を上の周期によって消す操作からなる。
周期の中心 coherence により、交換と三乗の消去の順序は結果を変えない。
各比較射はこの二種類の操作と可換するため、全ての積の coherence 図式が可換する。
これは
\[
 \mathbb Z^k/\langle3e_i-3e_j\rangle
\]
の graded MF 模型から来る指定された双対作用を、coherence ごと intertwine する構成である。
これは生成元の functor 同型類だけを比較しているのではない。
元の $\Gamma_k$-MF 模型上では指標作用は厳密な dg 作用である。
そこへ移した比較双加群には上の図式による $\widehat Z_0$-同変構造があり、
有限群の平均化と\cite[Lemma~8.6]{MEDElagin}により
実際の同変 perfect 双加群として取れる。
従って、同変化後の比較は単なる三角圏上の抽象同値ではなく dg Morita 比較である。

有限群の双対同変化をこの実際の比較へ施す。
系\ref{mbrd:cor:dual-action}で用いた Elagin の二重同変化により
式\eqref{mbrc:eq:ordinary-comparison}を得る。
原文\cite[Theorem~4.2]{MEDElagin}は代数閉体について述べられているが、
その証明で用いるのは $|Z_0|$ の可逆性と Fourier 分解
\[
 K'[Z_0]\simeq\bigoplus_{\chi\in\widehat Z_0}K'\chi
\]
である。$K'\supset\C$ では両方が成立するので、同じ証明が $K'$ 上で適用できる。
induction、上の Fourier 射影および冪等完備化によるこの構成は、
同変 dg 関手とその自然同型に自然である
\cite[Proposition~4.1, Lemma~8.6, Theorems~8.7, 8.9]{MEDElagin}。
双対群の bidual を両辺で同じ規約で $Z_0$ と同定すると、
元の $Z_0$ 作用が回復するため、得られた比較は $Z_0$-同変である。

rotation についても、この自然性を適用するための同変性を確認する。
基準因子の $R_{b,p_b}$ と $R_{b,0}$ は、元の rotation の
$Z_0$-線型化を積表示へ移したものである。
それらと $Q_i$ の交換同型は、因子の交換と同じ $R_b$ の冪の結合から作られる。
従って $FR_{b,p_b}\simeq R_{b,0}F$ の比較図式も、
同じ因子比較射の反復と外積の交換則により各 $Q_i$ と可換する。
三乗を消す際の図式は既に保った指定周期の図式である。
よってこの rotation 比較自体が $\widehat Z_0$-同変であり、
二重同変化後に元の残余圏の rotation 比較を与える。
元の指定周期も同じ自然性によって保持される。
これらは全て dg functor または perfect 双加群の構成であり、
smooth proper な圏間の同値は実際の Morita 核で表される。
代数閉体上では実線束と三乗根が既に存在するので拡大は不要である。
\end{proof}

\begin{corollary}[偏極点曲線の組を固定した幾何ファイバー]
\label{mbrc:cor:point-fiber-constancy}
$\kappa\supset\C$ を代数閉体とし、基本楕円領域に属する二つの満支持実表示
$A,A'$ を取る。
それらの偏極点曲線の順序を忘れた組
\[
 \{(E_i,L_i)\}_{i=1}^k,\qquad
 \{(E'_i,L'_i)\}_{i=1}^k
\]
が同型なら、標識を付けない残余 dg 圏は
\[
 \Ku_{\mathrm{coh}}(A)\simeq\Ku_{\mathrm{coh}}(A')
\]
となる。従って命題\ref{mbfm:prop:point-map}の点曲線組への射の
幾何ファイバー上で、通常の残余圏の Morita 同値類は一定である。
これは圏の族や選択した同値関手の降下を意味しない。
\end{corollary}
\begin{proof}
ブロックを並べ替えて、偏極同型
$\phi_i:(E_i,L_i)\xrightarrow{\sim}(E'_i,L'_i)$ を選ぶ。
その次数一切断への作用を用いて射影平面を同定すると、
二つの点三次の方程式はそれぞれ同じ直線を生成する。
従って $X_{\mathrm{pt}}$ と $X'_{\mathrm{pt}}$ の実方程式は、
共通のブロック変数の下で
\[
 \sum_i f_i,\qquad \sum_i u_i f_i,
 \qquad u_i\in\kappa^\times
\]
と書ける。
$\lambda_i^3=u_i$ を満たす $\lambda_i\in\kappa^\times$ を選び、
第 $i$ ブロックの座標を $\lambda_i$ 倍すれば、二つの三次超曲面は射影的に同型となる。
この同型は $\OO(1)$ を保つので、その Kuznetsov 圏を同定する。
定理\ref{mbrc:thm:finite-point-comparison}を両表示へ適用し、これらの同値を合成すればよい。

ここでは中心群作用、rotation 比較および関係直線同期の標識を忘れている。
同値関手の一意性、大域的な比較核の存在、
また残余圏から偏極点曲線の組を回復する逆向きの主張は含まない。
\end{proof}

\begin{remark}[体値点の比較と大域障害]
\label{mbrc:rem:scope}
定理\ref{mbrc:thm:finite-point-comparison}は、
基本楕円領域内の固定された全ての順序付き実表示に適用できる。
系\ref{mbrr:cor:stack-isomorphism}の全 $\OO$-認識開では、
ここでの中心群作用は全 $\OO$-局所型から内在的に回復される作用に一致する。
従って、全モジュライを支配する表示チャートの一般点にも適用できる。
この主張は一変数部分族の一般点からの推測ではない。
順序や実表示が元の体上にない場合、その持上げを得る拡大は
上の $3^{2k}$ という上界に含めていない。

また、点ごとに有限拡大上の核を構成しても、
モジュライ上の核 gerbe の類 $\alpha_P$ や固定同期の
周期 torsor の類 $\gamma(P,\vartheta)$ が消えるとは限らない。
本証明は実線束の選択と周期の三乗根を用いるので、
それらの選択を全 $\OO$-開上で降ろす問題は残る。
順序付き実表示上の相対模型と任意底変更後の compact 比較は、
定理\ref{mbrn:thm:relative-product}と系\ref{mbrn:cor:base-change}で別に構成する。
その上の Picard 補正型については、定理\ref{mbpd:thm:coherent-corrections}が
核 gerbe と位相 torsor を分類し、
命題\ref{mbpd:prop:boundary-phase}は自然な同期でも位相が非零となることを示す。
rotation を標識しない場合の改良上界は系\ref{mbpd:cor:field-degree}による。
全 $\OO$-stack 上の全ての核の型や、非正則な底での全 coherent 圏の比較は、
これらの実表示上の結果とは別の問題である。
\end{remark}

% !TEX root = ../nc-cubic.tex
\section{中心変数を持つ分離和の直接 Ku 比較}
\label{csc:sec:comparison}

本節は、楕円因子の中心三次和へ中心多項式を加えた場合の可換実現を与える。
中心多項式を同じ行列分解因子として両側に残し、楕円因子だけを
周期付き Picard 比較で取り替える。
全 coherent 圏の smoothness は有限 MF の比較から導き、
その比較の前提にはしない。
形式変形の可換復元、一般点への移行、閉点への spread は用いない。
以下は体上の実表示についての結果であり、
その表示を持たない全 intrinsic 対象の認識や大域核の降下を含まない。
この適用の包括的な新規性調査は未完了である。

\subsection{因子の仮定と環論的条件}

基礎体を $\C$ とする。$k,m\geq1$、$N=3k+m\geq5$、$r=N-3$ と置く。
各 $1\leq i\leq k$ について、滑らかな平面楕円曲線 $E_i$、
実際の線束 $L_i=\OO_{E_i}(1)$、平行移動 $\sigma_i=t_{p_i}$ と、
連結局所有限 quadratic Koszul AS-regular 代数 $S_i$ を取る。
$S_i$ の global dimension と Gorenstein parameter はともに三とし、
実際の正則中心元 $\Omega_i\in Z(S_i)_3$ が
\begin{equation}\label{csc:eq:factor}
 S_i/(\Omega_i)\simeq B(E_i,L_i,\sigma_i)
\end{equation}
を満たすと仮定する。
これは補題\ref{mbnp:lem:noetherian}で用いた実因子の仮定である。
$a_i\in\C^\times$、$w_i=a_i\Omega_i$ とし、
\begin{equation}\label{csc:eq:data}
 T=\C[u_1,\ldots,u_m],\qquad
 B=\left(\bigotimes_{i=1}^kS_i\right)\otimes T,\qquad
 w=\sum_{i=1}^kw_i+g(U),\qquad A=B/(w)
\end{equation}
と置く。$g$ は斉次三次で、
\begin{equation}\label{csc:eq:isolated}
 T/(\partial_{u_1}g,\ldots,\partial_{u_m}g)
 \quad\text{は有限次元の $\C$-ベクトル空間}
\end{equation}
とする。$m\geq2$ ではこれは $V(g)\subset\mathbb P^{m-1}$ の滑らかさに等しく、
$m=1$ では $g=cu^3$、$c\ne0$ に等しい。
係数 $a_i$ と方程式 $g$ は指定した frame の一部である。

\begin{lemma}[中心因子を持つ分離型の環論と三対象 window]
\label{csc:lem:algebra-window}
$B,A$ は両側 noetherian である。
$B$ は global dimension と parameter が $N$ の Koszul AS-regular 代数、
$w$ は正則中心三次、$A$ は parameter $r$ の AS-Gorenstein 代数である。
特に導来 section、$\chi$ 条件、有限 cohomological dimension を持つ。
有限線束生成部分の残余圏を $\mathcal K_A^{\mathrm{line}}$ とすると、
実際の dg window は
\begin{equation}\label{csc:eq:window}
 \mathcal K_A^{\mathrm{line}}
 \simeq\operatorname{thick}\{K_w,K_w(1),K_w(2)\}
       \subset\MF^{\mathbb Z}(B,w)
\end{equation}
を与える。ここで $K_w$ は augmentation の stabilization である。
この比較は residual rotation と内部 shift $(1)$ を同定し、
指定した $w$ による周期 $(3)\simeq[2]$ を保つ。
\end{lemma}
\begin{proof}
$k\geq2$ での $\bigotimes_iS_i$ の noetherian 性は
補題\ref{mbnp:lem:noetherian}による。
$k=1$ でも、ample な THCR $S_1/(\Omega_1)$ の noetherian 性と、
同補題の正次数中心元による初項の持上げがそのまま適用される。
中心多項式拡大の Hilbert 基底定理から $B$、従ってその商 $A$ も両側 noetherian である。
各 $S_i$ は $\C[w_i]$ 上自由である。
従って $B$ は中心多項式環
$\C[w_1,\ldots,w_k,u_1,\ldots,u_m]$ 上自由であり、
$w=\sum_iw_i+g$ はこの多項式環の非零元なので $B$ 上正則である。
各因子の有限 Koszul 分解と双加群分解を tensor すれば
ambient の dimension と parameter はともに $3k+m$ となる。
中心元の二項分解による change of rings は
$R\underline{\operatorname{Hom}}_A(\C,A)\simeq\C(r)[-N+1]$
と左右有限 injective dimension を与える。
従って\cite[Proposition~4.18, Lemmas~6.1--6.2, Theorem~6.4]{LP-BS}の
全仮定を満たす。全 coherent 対象に対する有限 Hom と線束による検出も
補題\ref{crs:lem:sections}の結論である。

有限 MF と graded 特異圏の比較は、正則中心元と有限大域次元だけを使う
補題\ref{mbnp:lem:window}の余核比較による。
同じ補題の window を $i=1-r$ で取り、
$\OO_A(0),\ldots,\OO_A(N-1)$ を用いる。
その像を計算する三つの有限 filtration は
式\eqref{mbnp:eq:triangular}と同じで、
像は $K_w(r),K_w(r+1),K_w(r+2)$ によって厚生成される。
ここで $r$ が三の倍数である必要はない。
任意の連続する三整数は三つの剰余類を一度ずつ含み、
$K_w(j+3)\simeq K_w(j)[2]$ であるから
その厚生成部分は式\eqref{csc:eq:window}に等しい。
rotation と周期についても、同じ切断三角
\eqref{tw:eq:rotation-cut}と比較正方形
\eqref{tw:eq:rotation-period-square}を用いる。
\end{proof}

\subsection{中心多項式の有限 MF 因子}

\begin{lemma}[孤立三次の中心因子]
\label{csc:lem:central-factor}
式\eqref{csc:eq:isolated}の下で
$\mathcal C_g=\MF^{\mathbb Z}(T,g)$ は smooth proper であり、
augmentation の stabilization $K_g$ の三つの shift
$K_g,K_g(1),K_g(2)$ によって厚生成される。
内部 shift $R_g=(1)$ は、実方程式 $g$ による指定周期
$\pi_g:R_g^3\simeq[2]$ を持つ。
\end{lemma}
\begin{proof}
正則 ambient $T$ により、有限 MF の余核は
$D^{\mathrm{gr}}_{\mathrm{sg}}(T/(g))$ と比較できる。
有限自由行列分解の微分 $d$ に対して
\[
 d\,\partial_{u_j}d+(\partial_{u_j}d)d
      =(\partial_{u_j}g)\operatorname{id}
\]
である。従って各 $\partial_{u_j}g$ の乗法は特異圏で零である。
ここで偏微分は選んだ斉次自由基底の行列成分に施し、
内部次数二の乗法に対応する斉次 homotopy を取る。

任意の有限生成 maximal Cohen--Macaulay 加群 $M$ に、
$\partial_{u_1}g,\ldots,\partial_{u_m}g$ の Koszul 複体を tensor する。
特異圏でこの複体は零となる乗法の cone を順に取ったものであり、
$M$ はその直和因子になる。
一方、その有界 cohomology は Jacobi イデアルで消される有限生成加群なので、
式\eqref{csc:eq:isolated}により有限長である。
有限長次数付き加群は $\C(j)$ の有限 filtration を持つ。
従って全有限 MF は $K_g(j)$、$j\in\mathbb Z$ から生成され、
$(3)\simeq[2]$ により三つの shift で十分である。
この周期は余核の二拡張を $T$ へ持ち上げた積が $g$ となることから固定する。

properness も同じ homotopy で分かる。
次数を忘れた二周期 Hom 複体の cohomology は、有限生成 $T$-加群で、
Jacobi イデアルにより消される。従って有限次元である。
その内部次数を三の剰余類で分解すると、各 graded Hom の全 cohomology の
直和が得られるので、graded Hom も $\C$ 上 perfect である。

$Y_g=V(g)\subset\mathbb P^{m-1}$ は smooth proper である。
$T/(g)$ の Gorenstein parameter は $m-3$ であり、
\cite[Theorem~6.4]{LP-BS}を適用する。
$m>3$ なら $\mathcal C_g$ は $D^b(Y_g)$ の admissible 成分、
$m=3$ なら $D^b(Y_g)$ 自体である。
$m=1,2$ なら $\mathcal C_g$ は $D^b(Y_g)$ と $3-m$ 個の
$\Perf(\C)$ による半直交分解を持つ。
この最後の場合も、既に示した properness と各成分の smoothness により
貼合せ双加群は perfect であり、全体は smooth である。
これは命題\ref{crs:prop:residual-smooth}の対角の有限貼合せと同じ議論である。
$m=1$ の空な $Y_g$ もこの記述に含める。
\end{proof}

\begin{lemma}[一変数の $A_2$ 模型と内部 shift]
\label{csc:lem:a2}
$g=cu^3$、$c\in\C^\times$ とする。
$Q=\C[u]/(u^3)$、$K=Q/(u)$ と置くと
\[
 \mathcal C_{cu^3}\simeq D^b(\operatorname{mod}\C\vec A_2).
\]
具体的には $E_1=K$、$E_2=K(-1)[1]$ が例外生成対で、
$\operatorname{RHom}(E_1,E_2)=\C$、
$\operatorname{RHom}(E_2,E_1)=0$ である。
内部 shift $R=(1)$ は
\begin{equation}\label{csc:eq:a2-rotation}
 R(E_1)=\operatorname{Cone}(E_1\longrightarrow E_2),\qquad
 R(E_2)=E_1[1],\qquad R^3\simeq[2]
\end{equation}
を満たす。最後の同型には実方程式 $cu^3$ の MF 周期を用いる。
\end{lemma}
\begin{proof}
$K$ の完全次数付き分解の第 $2j$ 項と第 $2j+1$ 項は、
それぞれ $Q(-3j)$ と $Q(-3j-1)$ である。
$K(l)$ への Hom の微分は零なので、全整数 $j$ に対して
\[
 \operatorname{Hom}_{D_{\mathrm{sg}}^{\mathrm{gr}}(Q)}
    (K,K(l)[q])=
 \begin{cases}
 \C,& q=2j,\ -l=3j,\\
 \C,& q=2j+1,\ -l=3j+1,\\
 0,&\text{それ以外}
 \end{cases}
\]
となる。$l=0,-1,1$ を代入すれば例外対の Hom を得る。
$M_2=Q/(u^2)$ とすると $\operatorname{Syz}(K)=M_2(-1)$ であり、
$0\to K(-1)\to M_2\to K\to0$ の三角を回転して
\[
 E_1\longrightarrow E_2\longrightarrow K(1)
       \longrightarrow E_1[1]
\]
を得る。これは式\eqref{csc:eq:a2-rotation}の第一式を与え、第二式は定義から従う。
また $K(2)=E_2[1]$ なので、$E_1,E_2$ は三つの augmentation shift を生成する。
$Q$ の全有限生成次数付き加群は有限長であるから、これは全特異圏を生成する。
従って endomorphism dg 代数の cohomology は次数零の $A_2$ path algebra である。
次数零以外の cohomology が零なので、その標準的な切断を介して
この path algebra と準同型で結ばれる。
その有限大域次元から smooth proper 性も直接従う。
$c$ を変えると指定周期の正規化は変わり得るため、
商環 $Q$ が同じであることだけから周期も同じとはしない。
\end{proof}

\subsection{楕円因子と中心因子の積・coherent 飽和}

中心多項式全体を第零因子と数え、
\[
 \Gamma=\{(\lambda_0,\ldots,\lambda_k)\in\mathbb G_m^{k+1}:
                  \lambda_0^3=\cdots=\lambda_k^3\},\quad
 \Lambda=\mathbb Z^{k+1}/\langle3e_i-3e_j\rangle,
 \quad Z=(\mu_3^{k+1})/\Delta\mu_3
\]
と置く。$T_1$ の全変数の重みを $e_0$、$(S_i)_1$ の重みを $e_i$ とする。
$Z$ の作用は、ブロックごとの scalar 作用と、対角 scalar の
次数付き加群上の自明化による coherent 作用である。
基準因子を第零因子に取れば、中心変数を固定し各 Sklyanin ブロックを
独立な三乗根で倍する通常の有限群作用として表せる。
この構成は全慣性群の分類を必要としない。

\begin{proposition}[中心因子を残す直接積比較]
\label{csc:prop:mixed-product}
式\eqref{csc:eq:factor}--\eqref{csc:eq:isolated}の下で
\begin{equation}\label{csc:eq:mixed-product}
 (\mathcal K_A^{\mathrm{line}})^Z
       \simeq \Perf(E_1\times\cdots\times E_k)
                         \otimes^{\mathbf L}\mathcal C_g
\end{equation}
という実際の dg Morita 同値がある。
両辺と $\mathcal K_A^{\mathrm{line}}$ は smooth proper であり、
\begin{equation}\label{csc:eq:coherent}
 \operatorname{thick}\{\OO_A(j)\}_{j\in\mathbb Z}
       =D^b_{\mathrm{dg}}(\operatorname{qgr}A),\qquad
 \mathcal K_A^{\mathrm{line}}=\Ku_{\mathrm{coh}}(A)
\end{equation}
が成り立つ。圏は全て冪等完備化して読む。
\end{proposition}
\begin{proof}
各 $\MF^{\mathbb Z}(S_i,w_i)$ は、式\eqref{mbnp:eq:factor-window}により
$\Perf(E_i)$ と同値で、三つの augmentation shift により生成される。
第零因子には補題\ref{csc:lem:central-factor}を用いる。
有限射影 MF の外積を取ると、定理\ref{mbnp:thm:product}の
有限 Hom の分解\eqref{mbnp:eq:hom}は第零因子を加えても同じである。
実際、独立な環の有限自由加群に対する Hom の tensor 分解と、
$\mathbb Z^{k+1}\to\Lambda$ の核に沿う偶数 cohomological shift の配分だけを使う。
各 Hom 複体は正次数性により下に有界なので、全次数を固定した和は有限である。
この議論に因子の global dimension が三であるという条件はない。

augmentation の外積も補題\ref{mbnp:lem:shamash}の有限模型で計算できる。
$T$ の augmentation の長さ $m$ の Koszul 分解に、$g$ の乗法の
higher homotopy を選ぶ。
同補題の障害は負次数の Ext に属して消え、分解の長さが有限なので
必要な higher homotopy は有限個である。
各因子で選んだものを tensor 分解へ Koszul 符号付きで足せば、
異なる因子の奇写像の cross term が消え、曲率は $\sum_iw_i+g$ になる。
従って
\[
 K_{w_1}\boxtimes\cdots\boxtimes K_{w_k}\boxtimes K_g\simeq K_w
\]
であり、外積の像は全 $K_w(\lambda)$、$\lambda\in\Lambda$ の厚生成部分である。

window は上のブロック作用と可換する。
$K_w$ とその整数 shift は線型化を持つので、有限群の誘導と平均化により
式\eqref{csc:eq:window}の同変化は、その全指標 twist で生成される。
これはまさに上の $\Lambda$-shift 全体である。
従って外積は式\eqref{csc:eq:mixed-product}を与える。
ここまで全 MF や全 coherent 圏の生成性を仮定していない。

積側は smooth proper である。
Elagin の二重同変化\cite[Theorem~4.2, Section~8]{MEDElagin}により
\[
 \mathcal K_A^{\mathrm{line}}
       \simeq\bigl((\mathcal K_A^{\mathrm{line}})^Z\bigr)^{\widehat Z}.
\]
有限群の位数は可逆なので、補題\ref{med:lem:crossed-perfect}の
有限 Fourier 分解と separability により左辺も smooth proper である。
次いで補題\ref{csc:lem:algebra-window}の環論的条件、
補題\ref{crs:lem:sections}の有限 Hom と検出、
補題\ref{crs:lem:finite-lines}の有限線束生成対象を用いる。
例外部分との貼合せと命題\ref{crs:prop:residual-smooth}、
定理\ref{crs:thm:coherent}から式\eqref{csc:eq:coherent}を得る。
従ってこの段階で初めて、外積の像は全 $\Lambda$ 次数付き MF とも一致する。
\end{proof}

\subsection{同じ中心多項式を持つ可換三次への比較}

\begin{theorem}[中心変数を持つ分離和の可換実現]
\label{csc:thm:comparison}
式\eqref{csc:eq:factor}--\eqref{csc:eq:isolated}を仮定する。
各 $E_i\subset\mathbb P(H^0(E_i,L_i)^\vee)$ の任意の非零方程式 $f_i$ を選び、
\begin{equation}\label{csc:eq:target}
 X=V\!\left(\sum_{i=1}^kf_i+g(U)\right)\subset\mathbb P^{N-1}
\end{equation}
と置く。この $X$ は滑らかな三次 $(N-2)$-fold であり、
\begin{equation}\label{csc:eq:comparison}
 \Ku_{\mathrm{coh}}(A)\simeq\Ku(X)
\end{equation}
という $\C$ 上の実際の dg Morita 同値が存在する。
比較は中心ブロック作用、rotation と、選んだ実方程式から定まる指定周期を保つように取れる。
因子の $j$ の相異性や特殊 $j$ の除外は不要である。
point graph の固定点なしという条件も、この Ku 比較自体には不要であり、
必要なのは式\eqref{csc:eq:factor}に記した実因子の条件である。
\end{theorem}
\begin{proof}
各 $f_i$ の偏微分の共通零点はその affine 三次元空間の原点だけである。
式\eqref{csc:eq:isolated}と変数の独立性から、
$\sum_i f_i+g$ の affine critical locus も原点だけである。
従って $X$ は滑らかである。

命題\ref{csc:prop:mixed-product}を非可換側に適用する。
可換側でも、各 $S_i,w_i$ を三変数多項式環と $f_i$ に替え、
同じ有限外積、window、生成性の証明を適用する。
三変数因子の parameter は零であり、その中心商は $B(E_i,L_i,\mathrm{id})$ なので、
両側の $Z$ 同変圏は同じ
\[
 \mathcal P=\Perf\left(\prod_iE_i\right)\otimes^{\mathbf L}\mathcal C_g
\]
と比較される。可換側の全 coherent 圏は、通常の smooth projective $X$ の圏である。

非可換楕円因子の内部 shift は
$R_{i,p_i}=T_{\OO_{E_i}}(t_{p_i})_*(-\otimes L_i)$、
可換因子では $R_{i,0}=T_{\OO_{E_i}}(-\otimes L_i)$ である。
$\C$ 上では $p_i$ を表す実際の次数零線束 $M_i$ が存在する。
補題\ref{mbrc:lem:factor-correction}により
$F_i=-\otimes M_i$ と比較射
\[
 \eta_i:F_iR_{i,p_i}\xrightarrow{\sim}R_{i,0}F_i
\]
を取れる。三回の合成と実方程式 $w_i,f_i$ の指定周期を比べた差は
非零 scalar である。$\eta_i$ をその逆の三乗根で倍すれば、指定周期を保つ。
ここでは中心元の自然な点方程式への写像やその非消滅を使わない。
第零因子では $g$ は両側で同じ実方程式なので、
$\mathrm{id}_{\mathcal C_g}$ が内部 shift と指定周期をそのまま保つ。

$F=(\boxtimes_iF_i)\otimes\mathrm{id}_{\mathcal C_g}$ を取る。
双対群 $\widehat Z$ の生成元は第零因子を基準として
\[
 Q_{i,p}=R_{i,p_i}^{-1}R_g,\qquad
 Q_{i,0}=R_{i,0}^{-1}R_g\qquad(1\leq i\leq k)
\]
である。各 $\eta_i$ とその逆は $FQ_{i,p}\simeq Q_{i,0}F$ を与える。
異なる因子の交換図式は外積の交換則で保たれ、
$Q_i^3\simeq[-2][2]\simeq\mathrm{id}$ の図式は上の周期の正規化で保たれる。
従って、各群元を $0\leq d_i<3$ の正規形 $\prod_iQ_i^{d_i}$ で書き、
交換と三乗の消去で積を定めれば、全ての群法則の coherence が保たれる。
これは指標作用を持つ実 $\Gamma$-MF 模型から移した作用についての主張である。
有限群の平均化により実際の同変 perfect 双加群として表せるので、
二重同変化は式\eqref{csc:eq:comparison}を dg Morita 同値として与える。
基準因子の $R_g$ は線型化付きの全体の rotation であり、
その比較も同じ図式と可換する。従って rotation と指定周期も保たれる。
\end{proof}

\begin{corollary}[独立な中心立方項と最小の三様体例]
\label{csc:cor:fermat}
$g=\sum_{j=1}^mc_ju_j^3$、$c_j\ne0$ とすれば、
定理\ref{csc:thm:comparison}は中心因子の smoothness に関する
一般の Orlov 分解を使わずにも証明できる。
特に $k=1,m=2$ では、式\eqref{csc:eq:factor}の任意の実楕円因子について
\[
 \Ku_{\mathrm{coh}}\!\left(
    S[u,v]/(a\Omega+cu^3+dv^3)\right)
   \simeq
 \Ku\!\left(V(f+cu^3+dv^3)\subset\mathbb P^4\right),
 \qquad a,c,d\in\C^\times,
\]
が $\C$ 上で直接成り立つ。ここで $V(f)=E$ である。
\end{corollary}
\begin{proof}
各中心変数を独立な因子として数え、因子数を $k+m$ とする。
各因子 $\MF^{\mathbb Z}(\C[u_j],c_ju_j^3)$ に補題\ref{csc:lem:a2}を適用し、
群を $(\mu_3^{k+m})/\Delta\mu_3$ に替える。
命題\ref{csc:prop:mixed-product}の有限 Hom、Shamash 分解、平均化の証明は
この因子分解にそのまま適用できる。
積側は $\Perf(\prod_iE_i)$ と $m$ 個の $A_2$ 圏の Morita tensor なので smooth proper である。
中心因子上の恒等比較はそれぞれの実方程式の周期を保ち、
楕円因子の Picard 補正と二重同変化によって結論を得る。
$k=1,m=2$ では $N=5,r=2$ であり、
補題\ref{csc:lem:algebra-window}の三つの剰余類の議論が適用される。
\end{proof}

\begin{remark}[存在結果の範囲]
\label{csc:rem:scope}
定理\ref{csc:thm:comparison}は $m>0$ の分離表示について主予想を与える。
全 $a_i\ne0$、実楕円因子の商表示、中心三次の孤立 critical point が仮定である。
$g=0$ や非孤立な $g$、active 変数と中心変数の混合項を含む中心三次、
楕円因子の退化、これらの表示を持たない対象へは無条件に延長しない。
全 point scheme と Poisson 零スキームの比較もこの定理の結論ではない。

比較には実際の次数零線束と三乗根を選んでいる。
従って固定された $\C$ 値表示での実核の存在を与えるが、
全族の核の自然な選択、全 $\OO$-モジュライへの大域降下、
一般体上の障害の消滅を主張しない。
各 $f_i$ の非零 scalar を変えた標的は、$\C$ 上では
その scalar の三乗根によるブロック座標変換で射影的に同型である。
\end{remark}

% !TEX root = ../nc-cubic.tex
\section{二ブロック型の完備復元と一般点比較}
\label{tbr:sec:recovery}

本節では周囲代数を$B$、中心三次商を$A$と書く。
第\ref{gdr:sec:main}節の局所的な記号$A_R,B_R$とは順序が逆である。
二ブロックの量子化の存在は外部入力であり、ここでは与えられた族に対して
既存の完備移送定理と一般点復元定理の仮定を確認する。
新たに確認するのは、ブロックの相対三乗根作用を保つ完備移送を選べることと、
その復元potentialに混合三次項が現れないことである。
この形式的な主張と、代数的な$\mathcal O$-markedモジュライへの降下は区別する。

\subsection{中心三次和の滑らかな特殊ファイバー}

$(R,\mathfrak m)$を剰余体$\mathbb C$のnoetherian局所$\mathbb C$-代数とする。
$D_1,D_2$は非負次数付き中央$R$-代数で、
$(D_i)_0=R$、各次数成分は有限自由、
\[
 D_1/\mathfrak mD_1=\mathbb C[x_1,x_2,x_3],\qquad
 D_2/\mathfrak mD_2=\mathbb C[y_1,y_2,y_3]
\]
と仮定する。中心元$\Omega_i\in Z(D_i)_3$の剰余$f_i$は
それぞれ滑らかな平面三次曲線を定めるものとし、$c_1,c_2\in R^\times$を取る。
通常のテンソル積に全次数を入れ、
\begin{equation}\label{tbr:eq:input}
 B=D_1\otimes_RD_2,\qquad
 w=c_1(\Omega_1\otimes1)+c_2(1\otimes\Omega_2),\qquad A=B/(w)
\end{equation}
と置く。テンソル積の交換式にsuper符号は付けない。

\begin{lemma}[二ブロック入力の次数別自由性と正則中心商]
\label{tbr:lem:input}
上の入力では$B_d$は階数$\binom{d+5}{5}$の有限自由$R$-加群である。
$w$は左右正則で、$A_d$は有限自由、形成は任意の通常の底変更と両立する。
そのHilbert級数は$(1-t^3)/(1-t)^6$である。
特殊ファイバーの方程式
\[
 F(x,y)=\overline c_1f_1(x)+\overline c_2f_2(y)
\]
は滑らかな三次四様体を定める。
\end{lemma}
\begin{proof}
次数$d$のテンソル積は
$\bigoplus_{p+q=d}(D_1)_p\otimes_R(D_2)_q$である。
各因子の階数は剰余の多項式環のHilbert関数なので、畳込みにより表示した階数を得る。
$F$の全偏微分が消える点では、$f_1$の三偏微分と$f_2$の三偏微分が
それぞれ消える。滑らかな平面三次の偏微分の共通アフィン零点は原点だけであるから、
六変数すべてが零となる。従って射影的特異点は存在しない。

固定次数の乗法$w:B_{d-3}\to B_d$を有限自由加群間の行列で表す。
その剰余は非零多項式$F$の乗法であり、単射である。
従って列数と同じ大きさの小行列式の一つは$R$の単元となる。
基底変換でこの行列を$\binom{I}{0}$にできるので、分裂単射で余核は自由である。
中心性により左右の乗法は一致する。各次数の分裂完全列を底変更しても完全であり、
Hilbert級数の公式もこの完全列から従う。
\end{proof}

\subsection{相対三乗根作用を保つ完備移送}

以下では$R$をさらに完備とする。原始三乗根$\zeta\in\mathbb C$を固定し、
$G=\mu_3$の生成元を$D_1$の次数$d$成分で$\zeta^d$倍、
$D_2$上で恒等として作用させる。この作用は$w$を保つ。
可換座標側では$(x,y)\mapsto(\zeta x,y)$である。

\begin{proposition}[完備復元potentialの二ブロック分解]
\label{tbr:prop:complete-split}
定理\ref{cft:thm:complete}の構成を式\eqref{tbr:eq:input}へ適用するとき、
復元potentialと有限rank graded MFの同値を$G$同変に選べる。
そのpotentialは
\begin{equation}\label{tbr:eq:split-potential}
 W=W_1(x_1,x_2,x_3)+W_2(y_1,y_2,y_3),\qquad
 W_i\bmod\mathfrak m=\overline c_i f_i
\end{equation}
となる。各$V(W_i)\subset\mathbb P^2_R$および$V(W)\subset\mathbb P^5_R$は滑らかである。
同値は内部shift、剰余の恒等同値、指定された三次周期を保ち、
$R/\mathfrak m^{q+1}$への切断と両立する。
\end{proposition}
\begin{proof}
まず次数別の加群同定を$G$同変に選ぶ。
剰余の単項式基底の任意の持上げが定める線形写像を
$|G|^{-1}\sum_{g\in G}g(-)g^{-1}$で平均する。
その剰余は元の$G$同変な同定であり、有限自由加群間のNakayamaにより同型である。
次数零では単位を固定する。この同定で戻した結合積は$G$不変となる。

補題\ref{cft:lem:model}のnormalizedかつ内部weight零のHochschild複体と
polyvector複体を考える。有限群の不変部分を取る操作は標数零で完全であり、
HKRおよび線形共変なKontsevich写像は$G$同変である。
従って同補題で使う準同型は不変部分へ制限しても準同型となる。
各Artin切断の冪零Maurer--Cartan比較をこの不変部分で行えば、
formal Poisson tensor $\pi$と厳密代数同型$\phi$を$G$同変に選べる。
具体的には、既存の低次jetを固定した各平方零の持上げ段階の線形な補正方程式に
Reynolds平均を適用する。
完成した非線形のPoisson解や代数同型を平均する操作ではない。
これは命題\ref{gre:prop:hesse-normalform}で用いた有限対称群の議論と同じである。

補題\ref{cft:lem:center}の中心写像$C_\pi$は$G$同変であり、三次成分上で可逆である。
$w$は不変なので、$W=C_\pi^{-1}(\phi(w))$も不変となる。
三次単項式の$x$次数は$0,1,2,3$のいずれかである。
$G$不変性はこの次数が三の倍数であることを要求するため、
不変三次空間は
\[
 R[x_1,x_2,x_3]_3\oplus R[y_1,y_2,y_3]_3
\]
であり、式\eqref{tbr:eq:split-potential}を得る。
剰余は$F$なので各$W_i$の剰余も表示通りである。
各平面三次の判別式は剰余で非零、従って$R$の単元である。
六変数の$W$も同じ理由、または補題\ref{tbr:lem:input}の勾配計算で滑らかとなる。

最後に実際の同値の同変性を確認する。
補題\ref{cft:lem:primitive}の$\Xi$は$\pi=[\Xi,W]$を満たす唯一の定数trivectorであり、
$\pi,W$の不変性から$\Xi$も不変である。
周期変数$t$を$G$で固定すると、$\eta=\Xi/t$、曲率付き経路と
補題\ref{cft:lem:arity}のcoderivationはすべて$G$同変である。
従って補題\ref{cft:lem:flow}の順序付き積分と逆flowも$G$同変となる。
各Taylor成分が同変なので、MF行列を代入する構成は係数への$G$作用と可換する。
次数別収束、有限rank、内部shift、周期とArtin切断の両立は
定理\ref{cft:thm:complete}の同じ評価から従う。
\end{proof}

この命題は既存の完備MF比較を入力とする帰結であり、その基礎定理の独立な再証明ではない。
また、$W_i$の係数はここでは計算していない。
とくに$V(W_i)$を各Sklyanin因子の点曲線と同一視する結果を含まない。

\subsection{一変数の一般点と代数的降下の境界}

\begin{corollary}[二ブロック一変数族への一般点復元の適用]
\label{tbr:cor:generic}
式\eqref{tbr:eq:input}の入力を$R_0=\mathbb C[h]_{(h)}$上で与え、
各$(D_i)_d$が有限自由であるとする。$K=\mathbb C(h)$と置く。
このとき$K$上に滑らかな三次四様体$Y\subset\mathbb P^5_K$が存在し、
ある有限拡大$K'/K$上で
\[
 \bigl(\mathcal K^{\rm coh}_{A_{K'}},\rho_A\bigr)
 \simeq\bigl(\operatorname{Ku}(Y_{K'}),\rho_Y\bigr)
\]
となる。ここで
$\mathcal K^{\rm coh}_{A_K}=D^b(\operatorname{qgr}A_K)
\cap\langle\mathcal O,\mathcal O(1),\mathcal O(2)\rangle^\perp$
であり、同値は指定した三次周期も保つ。
\end{corollary}
\begin{proof}
補題\ref{tbr:lem:input}は$B$の各次数の有限自由性、
六変数多項式環への特殊化、正則な中心三次元と滑らかな特殊ファイバーを与える。
これらは$N=6,n=4$の定理\ref{gdr:thm:generic}の仮定そのものである。
同定理の周囲代数・商の記号をそれぞれ$B,A$へ読み替えて適用する。
\end{proof}

系\ref{tbr:cor:generic}だけでは、群作用とそのcoherenceを保つ降下も、
$K$上の方程式の二ブロック分裂も結論しない。
これらは次節の定理\ref{med:thm:multiblock-generic}において、
不変な積分空間とcrossed product上のperfect核を用いて別に証明する。
そこで得るのは有限拡大上の同変比較の存在であり、
命題\ref{tbr:prop:complete-split}で選んだ形式核の全jetを再現する代数化ではない。
任意の固定非零パラメータ、全二変数底上の相対同値、全$\mathcal O$-保存同値への降下、
元の古典三次そのもののKuznetsov成分との一致は別の課題である。

\subsection{既知の可換二ブロック模型との比較}

複素数体上の滑らかな三次四様体
$Y=V(g_1(x)+g_2(y))$と、$G$の作用$(x,y)\mapsto(\zeta x,y)$を考える。
$C_i=V(g_i)\subset\mathbb P^2$とすると、既知の定理により
\begin{equation}\label{tbr:eq:known-equivariant-product}
 \operatorname{Ku}(Y)^G\simeq D^b(C_1\times C_2)
\end{equation}
である。Hu\cite[Theorems~0.1, 5.8]{TBHu}はこの形を明示的に扱い、
同論文のIntroductionはBallard--Favero--Katzarkov、Lim、
Beckmann--Oberdieck、Hirano--Ouchiによる先行する一般形を挙げている。
この既知の同値自体を本研究の成果とは数えない。

式\eqref{tbr:eq:known-equivariant-product}の左辺は$G$同変対象の圏である。
通常の$\operatorname{Ku}(Y)$を楕円曲線の積の導来圏と同定する式ではない。
命題\ref{tbr:prop:complete-split}は、この既知模型を二ブロック非可換族と比較するための
形式的な出発点を与える。
有限拡大上の同変比較とこの既知模型への接続は系\ref{med:cor:elliptic-product}で行う。
点曲線との同定と全$\mathcal O$-降下はなお別問題である。

% !TEX root = ../nc-cubic.tex
\section{満支持ブロック型の点対応と次数一復元の障害}
\label{tbp:sec:points}

本節では $k$ を標数零の代数閉体とし、
$S_i=T(W_i)/(R_i)$ $(i=1,2)$ を三次元楕円型の二次 AS 正則代数とする。
$\dim W_i=\dim R_i=3$ とし、二次関係の双線型化の零スキームが
\[
 \Gamma_i=\Gamma_{\sigma_i}\subset
 \PP(W_i^*)\times\PP(W_i^*)
\]
であると仮定する。ここで $E_i\subset\PP(W_i^*)$ は滑らかな平面三次曲線、
$\sigma_i:E_i\to E_i$ は固定点を持たない自己同型であり、
$\Gamma_{\sigma_i}$ はそのグラフである。
さらに $\Omega_i\in Z(S_i)_3$ は正則元で、
$S_i/(\Omega_i)=B(E_i,\mathcal L_i,\sigma_i)$ とする。
この幾何学的表示と中心三次商は Sklyanin 代数の既知の入力である
\cite[\S1, pp.~20--21]{GRSklyaninST}。
本節の主張はこれらを満たす表示について述べ、
固定点を持つ退化や一般の非楕円型へは拡張しない。

全次数で次数付けした tensor 積とその中心三次商を
\begin{equation}\label{tbp:eq:algebra}
 V=W_1\oplus W_2,\qquad
 B=S_1\otimes_k S_2,\qquad
 A_{a,b}=B/(a\Omega_1+b\Omega_2)
 \quad ((a,b)\ne(0,0))
\end{equation}
と書く。点加群は定理\paperref{moduli}{thm:point-nfold}と同じく、
次数零から生成され、各次数成分が可逆層で、次数零の自明化を備えるものとする。

\begin{theorem}[二ブロックの点対応]\label{tbp:thm:point-graph}
$B$ の双線型関係の零スキーム、および $B,A_{a,b}$ の点加群関手は
\begin{equation}\label{tbp:eq:point-graph}
 \Gamma_B=\Gamma_1\amalg\Gamma_2,\qquad
 \operatorname{Pt}(B)=\operatorname{Pt}(A_{a,b})=E_1\amalg E_2
\end{equation}
である。shift は各 $E_i$ 上で $\sigma_i$ に等しい。
これらはスキームとしての等式であり、非被約な試験基底も含めた
任意の基底変更に両立する。特にブロックをまたぐ横断的な冪零構造はない。
\end{theorem}

\begin{proof}
$p=(x,y)$, $q=(x',y')$ を二つの射影座標とする。
$x,x'$ は第一ブロック、$y,y'$ は第二ブロックの三座標である。
ブロック間の交換関係は
\begin{equation}\label{tbp:eq:cross-bilinear}
 x_\alpha y'_\beta-y_\beta x'_\alpha=0
 \qquad(1\leq\alpha,\beta\leq3)
\end{equation}
を与える。

まず幾何点を調べる。$x,y$ がともに非零ならば、
\eqref{tbp:eq:cross-bilinear}と $q\ne0$ により $x',y'$ もともに非零で、
$x'=c x$, $y'=c y$ となる。
純ブロックの関係は $[x]=\sigma_1([x])$ および
$[y]=\sigma_2([y])$ を要求するので、固定点なしの仮定に矛盾する。
$y=0$ ならば $x\ne0$ であり、同じ交換関係から $y'=0$ となる。
従ってこの点は $\Gamma_1$ に属する。$x=0$ の場合は $\Gamma_2$ に属する。

この集合論的な議論だけでは冪零構造を排除できないので、局所環を調べる。
$\Gamma_1$ 上では固定点がないため、各点の近傍である小行列式
\[
 \Delta_{\alpha\gamma}
   =x_\alpha x'_\gamma-x_\gamma x'_\alpha
\]
が可逆となる。各 $\beta$ について
\eqref{tbp:eq:cross-bilinear}の $\alpha,\gamma$ の二式を連立すれば、
この可逆行列式により $y_\beta=y'_\beta=0$ が環の中で従う。
残る関係はちょうど $\Gamma_1$ の関係である。
$\Gamma_2$ 上でも同じ議論で $x=x'=0$ を得る。
幾何点についての前段の議論により、これらの小行列式の開集合は
$\Gamma_B$ 全体を覆う。従ってスキームとして最初の等式を得る。

二次対応が自己同型のグラフなので、それを満たす点列は初期点から
一意に決まり、$B$ の点関手は $E_1\amalg E_2$ で表される。
第一成分上の普遍点加群では、第二ブロックの生成元は全て零に作用し、
第一ブロックの $\Omega_1$ も twisted homogeneous coordinate ring への
商表示により零に作用する。第二成分についても同様である。
従って $a\Omega_1+b\Omega_2$ の三重線型化は全ての普遍点列上で零であり、
中心三次関係を課しても点関手は変わらない。
以上は可逆小行列式による環上の消去と普遍点加群の恒等式なので、
任意の基底変更にも両立する。
\end{proof}

\begin{corollary}[埋込み点対応による二つのブロックの認識]
\label{tbp:cor:block-spans}
埋込み点スキームから、順序を除いて二つの平面
$\PP(W_1^*),\PP(W_2^*)$、従って分解 $V=W_1\oplus W_2$ を復元できる。
同じ仮定を満たす二つの中心三次商の次数付き同型は、
その二つの生成元ブロックを順序の交換を除いて保つ。
\end{corollary}

\begin{proof}
$E_1,E_2$ は点スキームの二つの連結成分であり、
各々がその埋込み平面を張る。
二つの平面を張るベクトル部分空間は $V^*$ の直和分解を与えるので、
その双対として $V$ の分解を得る。
次数付き同型は点関手と次数一の埋込みを同定するので、連結成分とその線型包を保つ。
\end{proof}

\begin{proposition}[点対応の双線型イデアルは二次関係より大きい]
\label{tbp:prop:quadratic-saturation}
$\Gamma=\Gamma_1\amalg\Gamma_2$ とすると、$V^{\otimes2}$ の部分空間として
\begin{equation}\label{tbp:eq:saturation}
 H^0(\mathcal I_\Gamma(1,1))
 =R_1\oplus R_2\oplus(W_1\otimes W_2)\oplus(W_2\otimes W_1).
\end{equation}
従って左辺の次元は $24$ である。一方、実際の二次関係空間は
\begin{equation}\label{tbp:eq:actual-relations}
 R_B=R_1\oplus R_2\oplus
 \langle x\otimes y-y\otimes x:
         x\in W_1,\ y\in W_2\rangle
\end{equation}
であり、次元は $15$ である。
特に一ブロック型の定理\paperref{moduli}{thm:point-envelope}の復元式は、
この二ブロック型では成立しない。
\end{proposition}

\begin{proof}
$V^{\otimes2}$ を四つのブロックに分解すると、
$\Gamma_i$ への制限で寄与するのは $W_i^{\otimes2}$ だけである。
そこでの核が $R_i$ であることを確認する。
$L=\mathcal L_i$, $M=\sigma_i^*\mathcal L_i$ と置けば、
その核は乗法
$H^0(L)\otimes H^0(M)\to H^0(LM)$ の核である。
$L\not\simeq M$ なら、$L$ の基点のない二次元線型系の評価完全列と
$H^1(ML^{-1})=0$ によりこの乗法は全射である。
$L\simeq M$ なら、平面三次曲線には二次の方程式がないことにより全射である。
$h^0(LM)=6$ だから核は三次元で、含まれている $R_i$ に等しい。
混合した二つのブロックは全て $\Gamma$ 上で零なので
\eqref{tbp:eq:saturation}を得る。
tensor 積表示から \eqref{tbp:eq:actual-relations}が従い、
二つの次元はそれぞれ $3+3+9+9$ および $3+3+9$ である。
\end{proof}

\begin{remark}[中心直線の不可視性の意味]\label{tbp:rem:pencil}
点スキーム・埋込み・shift は、固定した $B$ の座標の下で
中心直線 $k(a\Omega_1+b\Omega_2)$ に依存しない。
さらに前命題により、点対応だけではブロック間の九本の交換関係も選べない。
これは点対応から中心直線を直接回復する方法の障害である。
ただし $ab\ne0$ の全ての中心直線は、代数閉体上では
ブロック別のスカラー変換による一つの軌道をなす。
実際、$W_1,W_2$ 上のスカラー $c_1,c_2$ は係数を
$(a,b)\mapsto(c_1^3a,c_2^3b)$ と変える。
従ってこの不可視性だけから、同じ点対応を持つ非同型な
中心三次商の一次元族があるとは結論しない。
\end{remark}

\begin{proposition}[同じ三次 Veronese を持つ異なる中心三次表示]
\label{tbp:prop:zhang-ambiguity}
$\zeta\in k$ を原始三乗根とし、$B$ の自己同型 $\tau$ を
\[
 \tau|_{W_1}=\operatorname{id},\qquad
 \tau|_{W_2}=\zeta\operatorname{id}
\]
で定める。斉次元の積を
$u*v=u\tau^{\deg u}(v)$ とする Zhang twist を $B^\tau$ と書く。
$w=a\Omega_1+b\Omega_2$ とすると、次が成り立つ。
\begin{enumerate}
 \item $w$ は $B^\tau$ でも中心であり、
 $A_{a,b}^\tau=B^\tau/(w)$ である。
 $w$ が $B$ 上で正則なら $B^\tau$ 上でも正則である。
 \item 同じ次数付きベクトル空間を用いた文字通りの等式
 \[
   (A_{a,b}^\tau)^{(3)}=A_{a,b}^{(3)}
 \]
 が成り立ち、構造対象を保つ同値
 $\operatorname{QGr}A_{a,b}\simeq\operatorname{QGr}A_{a,b}^\tau$ がある。
 \item $A_{a,b}$ と $A_{a,b}^\tau$ の埋込み点スキームと shift は等しいが、
 この二つの代数の間に次数付き同型は存在しない。
\end{enumerate}
\end{proposition}

\begin{proof}
$\tau^3=1$ かつ $\tau(w)=w$ である。
任意の斉次元 $u$ について
$w*u=wu$ および $u*w=uw$ だから、中心性・正則性および商との可換性を得る。
$\deg u$ が三の倍数なら $u*v=uv$ なので、三次 Veronese の等式も従う。

Zhang の twist による同値は既知である
\cite[Theorem~1.1]{TBPZhangTwist}。
ここでは右加群 $M$ に新しい作用
$m*u=m\tau^{\deg m}(u)$ を入れれば直接構成できる。
この関手は次数付きベクトル空間を変えず、逆 twist が逆関手となり、
正則右加群 $A_{a,b}$ を正則右加群 $A_{a,b}^\tau$ に送る。
また torsion を保つため $\operatorname{QGr}$ の同値を誘導し、構造対象を保つ。
noetherian な場合は有限生成対象に制限して $\operatorname{qgr}$ の同値も得る。

twist の純ブロックの二次関係は元と同じであり、混合関係は
\begin{equation}\label{tbp:eq:twisted-cross}
 x*y=\zeta\,y*x\qquad(x\in W_1,\ y\in W_2)
\end{equation}
となる。第二ブロックの三次語には係数 $\zeta^{1+2}=1$ が掛かるため、
各 $\Omega_i$ の三次 tensor 代表も同じに取れる。
双線型混合関係は
$x_\alpha y'_\beta-\zeta y_\beta x'_\alpha=0$ である。
両ブロックが非零の幾何点では、各ブロックの二つの射影座標がそれぞれ比例し、
再び固定点を生じる。純ブロック上では、先の二行二列の消去の
行列式が可逆スカラー $\zeta$ 倍されるだけである。
従って定理\ref{tbp:thm:point-graph}と同じスキーム論的証明が通り、
点対応は全く変わらない。

次数付き同型 $A_{a,b}\to A_{a,b}^\tau$ があると仮定する。
その点対応への作用と系\ref{tbp:cor:block-spans}の線型包の議論により、
次数一の同型は $W_1,W_2$ を保つか交換する。
元の代数ではこの二つの空間の元は互いに可換なので、
像の二つの空間の元も全て可換でなければならない。
しかし \eqref{tbp:eq:twisted-cross}により、非零の
$x\in W_1$, $y\in W_2$ に対して
$x*y-y*x=(\zeta-1)yx\ne0$ である。
最後の非零性は、三次商が次数二を変えず、混合次数二部分が
$W_1\otimes W_2$ に同型であることから従う。
ブロックを交換しても同じ矛盾を得る。
\end{proof}

\begin{corollary}[任意個の満支持ブロック]
\label{tbp:cor:all-blocks}
上と同じ仮定を満たす $\ell\ge2$ 個の因子を取り、
\[
 B=\bigotimes_{i=1}^{\ell}S_i,\qquad
 V=\bigoplus_{i=1}^{\ell}W_i,\qquad
 A=B/\left(\sum_{i=1}^{\ell}a_i\Omega_i\right),\qquad a_i\ne0
\]
とする。このとき
\begin{equation}\label{tbp:eq:all-block-points}
 \Gamma_B=\coprod_{i=1}^{\ell}\Gamma_{\sigma_i},\qquad
 \operatorname{Pt}(A)=\coprod_{i=1}^{\ell}E_i
\end{equation}
であり、任意の基底変更に両立する。
実際の二次関係空間 $R_B$ と点対応の双線型イデアルの次元は
\begin{equation}\label{tbp:eq:all-block-defect}
 \dim R_B=3\ell+9\binom\ell2,\qquad
 \dim H^0(\mathcal I_{\Gamma_B}(1,1))
       =3\ell+18\binom\ell2.
\end{equation}
従って復元式の次元差は $9\binom\ell2$ である。

さらに $\varepsilon_i\in\mu_3(k)$ を選び、
$\tau|_{W_i}=\varepsilon_i\operatorname{id}$ とする。
$n=3\ell-2$, $d=n-1=3\ell-3$ と置くと、
$w=\sum_i a_i\Omega_i$ は $B^\tau$ でも中心であり、
\[
 (A^\tau)^{(d)}=A^{(d)},\qquad
 (\operatorname{QGr}A^\tau,\OO)\simeq(\operatorname{QGr}A,\OO).
\]
$\varepsilon_i$ が全て等しくなければ、$A^\tau$ と $A$ は次数付き同型でない。
これは異なる twist 同士の同型類の個数を決定する主張ではない。
\end{corollary}

\begin{proof}
幾何点の第一座標で複数のブロックが非零なら、二つずつの混合関係から
各非零ブロックで第一・第二の射影座標が比例し、固定点を生じる。
従って幾何点は一つの $\Gamma_{\sigma_i}$ に属する。
そのブロックの固定点なしの条件による可逆な二行二列小行列式を使うと、
各 $j\ne i$ との混合関係から第 $j$ ブロックの全座標が零となる。
この環上の消去は横断的冪零構造も除く。
全ての $\Omega_i$ が各普遍点列上で零に作用するので、中心商を課しても変わらない。

各対角ブロックは三次元の関係空間を持つ。
異なるブロックの各非順序対について、実際の関係は九次元の交換関係であるが、
点対応上では二つの向きの全十八次元が零になる。
これで \eqref{tbp:eq:all-block-defect} を得る。

$\tau^3=1$ かつ $\tau(w)=w$ なので、
命題\ref{tbp:prop:zhang-ambiguity}と同じ式が中心性と同値を与える。
$3\mid d$ より $\tau^d=1$ であり、$d$-Veronese の積も元の積と一致する。
twist 後の混合関係は
\[
 x_i*x_j=(\varepsilon_j/\varepsilon_i)x_j*x_i
 \qquad(x_i\in W_i,\ x_j\in W_j)
\]
である。可逆小行列式による証明には非零スカラー倍しか生じないので、
twist 後も点対応は同じである。
次数付き同型はその連結成分と線型包を置換する。
元の代数では異なるブロックが全て可換だから、同型が存在すれば
twist 後の全てのブロック対も可換でなければならない。
$\varepsilon_i\ne\varepsilon_j$ となる対では混合次数二の積が非零であり、
この必要条件に反する。
\end{proof}

\begin{remark}[O-marked 復元への正確な帰結]\label{tbp:rem:scope}
命題\ref{tbp:prop:zhang-ambiguity}と系\ref{tbp:cor:all-blocks}は、
一般の中心三次表示を許した範囲では、
O-marked 圏も Serre 指数 $d$ の Veronese も次数一表示を一意に決めないことを示す。
従って中心一次変数を用いる
定理\paperref{moduli}{thm:nfm-central-veronese}を $U=0$ の場合へ無条件には延長できない。
ここで $A_{a,b}^\tau$ はブロック間が $\zeta$ 可換な表示であり、
もとの可換 tensor 積表示のチャートには属さない。
この例だけで、二つの可換 tensor 積表示の間の全 O-同値が
次数付き同型から来ないとは結論しない。
また同じ O-marked 対象に異なる表示を与えているのであり、
新たな intrinsic 成分を構成したという主張でもない。
二ブロック型の内在的モジュライには、可換 tensor 積表示の認識と、
このような次数一の根の曖昧さを区別した議論が必要である。
\end{remark}

% !TEX root = ../nc-cubic.tex
\section{満支持型での可換比較と Poisson 零スキーム}
\label{pps:sec:split}

本節では、主予想\ref{conj:nc-main}の可換 Ku 実現と
点スキーム・Poisson 零スキームの一致を、同じ可換三次超曲面について示す。
中心一次変数を持たない満支持型では、原代数の点スキームは
既に系\ref{tbp:cor:all-blocks}でスキームとして求められ、
自然な可換標的への残余圏比較も
定理\ref{mbrc:thm:finite-point-comparison}で構成されている。
追加するのは、その同じ標的上の Poisson tensor の零スキームの計算である。
第三 Veronese の中心還元を介した点加群の分類を、
原代数の点スキームの代用には用いない。

\subsection{分離三次和の零スキーム}

$S$ を $\C$ 上のスキームとし、各 $i=1,\ldots,k$ に三変数
$x_i=(x_{i1},x_{i2},x_{i3})$ を取る。$k\ge2$ とする。
三次斉次式 $f_i\in\Gamma(S,\OO_S)[x_i]_3$ が
滑らかな平面三次族 $E_i=V(f_i)\subset\PP^2_S$ を定めると仮定する。
各 $E_i$ をその純ブロック平面へ埋め込み、
\begin{equation}\label{pps:eq:split-data}
 F=\sum_{i=1}^k f_i,\qquad
 X=V(F)\subset\PP^{3k-1}_S,\qquad
 P_i=\partial_{x_{i1}}\wedge\partial_{x_{i2}}\wedge\partial_{x_{i3}}
\end{equation}
と置く。以下の微分と tangent 束は全て $S$ に相対的に取る。
$X/S$ は滑らかである。
実際、射影点では少なくとも一つのブロックが非零であり、
その $f_i$ の偏微分のいずれかが非零となる。
この Jacobian 判定は相対的にも適用できる。

任意の単元 $c_i\in\Gamma(S,\OO_S)^\times$ に対して
\begin{equation}\label{pps:eq:poisson-tensor}
 \beta=\sum_{i=1}^k c_i\,\iota_{df_i}P_i
\end{equation}
を affine 空間上の二次 Poisson tensor とする。
各項は三変数の Jacobian 括弧であり、異なる項の係数は互いに別の
変数にしか依存しないため、その和も Poisson である。
$F$ は Casimir なので、この tensor の射影化は
$\pi\in H^0(X,\bigwedge^2T_{X/S})$ を定める。
零スキーム $Z_X(\pi)$ は切断 $\pi$ の係数イデアルで定義する。

\begin{proposition}[満支持分離和の Poisson 零スキーム]
\label{pps:prop:zero-scheme}
上の仮定の下で、同じ射影空間内の閉部分スキームとして
\begin{equation}\label{pps:eq:zero-scheme}
 Z_X(\pi)=\coprod_{i=1}^k E_i
\end{equation}
が成立する。この等式は任意の基底変更に両立し、
非被約な基底上でも横断的な冪零構造を加えない。
\end{proposition}

\begin{proof}
Euler ベクトル場を
$\xi=\sum_{i,a}x_{ia}\partial_{x_{ia}}$ とする。
affine 錐の原点を除いた部分では、射影 tangent 束は
自明 tangent 束を直線 $\OO\xi$ で割って得られる。
従って射影化した $\beta$ の零イデアルは
$\xi\wedge\beta$ の係数イデアルで与えられる。
これは点ごとの判定だけではない。
ある $x_{ia}$ が可逆な chart で $\xi$ を自由加群の基底の一つに取れば、
$\bigwedge^2(T/\OO\xi)$ から $\bigwedge^3T$ への
$\xi\wedge(-)$ は直和因子への同型となり、係数イデアルが一致する。
また $X/S$ の滑らかさにより
$\bigwedge^2T_{X/S}\hookrightarrow
 \bigwedge^2T_{\PP^{3k-1}_S/S}|_X$
は局所的な直和因子の包含なので、$X$ 上の零イデアルも同じ式を制限して得られる。

まず幾何ファイバー上で台を調べる。
ある点で非零となるブロックの数を $s$ とする。
滑らかさと斉次 Euler 恒等式により、非零の $x_i$ では
$df_i(x_i)\ne0$ である。
従ってそのブロックの $c_i\iota_{df_i}P_i$ は階数二を持ち、
ブロック直和である $\beta$ の階数は $2s$ となる。
射影化が零ならば $\beta=\xi\wedge v$ と書けるので階数は高々二であり、
$s\ge2$ は不可能である。
純ブロック $i$ では
\[
 \xi\wedge\beta
   =c_i\left(\sum_{a=1}^3x_{ia}\frac{\partial f_i}{\partial x_{ia}}\right)P_i
   =3c_if_iP_i.
\]
従って零スキームの幾何学的な台は $\coprod_iE_i$ である。

次に $E_i$ の近傍で構造層を調べる。
滑らかさから、各点である
$g_{ia}=\partial f_i/\partial x_{ia}$ が、射影 chart 上の自明化の下で
可逆となる近傍を取れる。
$\xi\wedge\beta$ において、第 $i$ ブロックから二つ、
第 $j$ ブロックから一つの tangent 基底を取った係数は、
符号を除いて
\begin{equation}\label{pps:eq:transverse-coefficient}
 c_i x_{jb}g_{ia}\qquad(j\ne i,\ 1\le b\le3)
\end{equation}
である。他の $\beta$ の項は、このブロック次数の係数に寄与しない。
従って零イデアルは、可逆な $c_ig_{ia}$ を用いて
全ての $x_{jb}$ $(j\ne i)$ を含む。
純ブロックの係数 $3c_if_i$ から $f_i$ も含む。
逆に他ブロックの座標を零とし $f_i=0$ を課すと、
$\xi\wedge\beta$ の全係数は零になる。
よってこの近傍で零イデアルはちょうど
\[
 (f_i,\ x_{jb}\mid j\ne i,\ 1\le b\le3)
\]
である。幾何学的な台の議論により、これらの開集合は零スキーム全体を覆う。
純ブロック平面同士は射影空間で互いに交わらないので、
式\eqref{pps:eq:zero-scheme}を得る。

この証明は Euler 恒等式、係数イデアルおよび単元を用いる局所消去による。
基底の被約性を使わず、全ての等式は任意の基底変更に保たれる。
\end{proof}

\subsection{同じ可換標的での Ku 比較と零スキームの一致}

\begin{theorem}[楕円 tensor 表示に対する存在版]
\label{pps:thm:elliptic-existence}
$\kappa\supset\C$ を代数閉体とする。
各 $i=1,\ldots,k$ について次のデータを仮定する。
\begin{enumerate}
 \item $S_i=T(W_i)/(R_i)$ は三次元の連結二次 Koszul AS 正則な
 楕円 Sklyanin 代数で、$\dim W_i=\dim R_i=3$、
 Gorenstein parameter も三である。
 正則中心三次 $\Omega_i$ があり、
 $S_i/(\Omega_i)\simeq B(E_i,L_i,t_{p_i})$ である。
 ここで $E_i$ は滑らかな平面三次、$L_i=\OO_{E_i}(1)$ である。
 \item $S_i$ の双線型関係の零スキームは
 $\Gamma_{t_{p_i}}\subset\PP(W_i^\vee)\times\PP(W_i^\vee)$ であり、
 $t_{p_i}$ は固定点を持たない。
\end{enumerate}
$k\ge2$、$a_i\in\kappa^\times$ とし、
\[
 A=\left(\bigotimes_iS_i\right)\Big/\left(\sum_i a_i\Omega_i\right)
\]
と置く。各 $E_i$ の非零三次方程式 $f_i$ を任意に選ぶと、
$X=V(\sum_i f_i)$ は滑らかであり、
\begin{equation}\label{pps:eq:elliptic-existence}
 \Ku_{\mathrm{coh}}(A)\simeq\Ku(X),\qquad
 \operatorname{Pt}(A)=Z_X(\pi)=\coprod_iE_i
\end{equation}
を満たす実際の dg Morita 同値と Poisson tensor $\pi$ が存在する。
従って、この範囲では可換 Ku 実現と点スキームの一致を同じ $X$ 上で実現できる。
$j$ 不変量の相異性、$j\ne0,1728$、$2p_i\ne0$ は仮定しない。
\end{theorem}

\begin{proof}
最初の仮定は式\eqref{mbnp:eq:factor-data}--\eqref{mbnp:eq:sum-data}の
直接積比較の仮定を満たす。
ここで用いる $Z_0$ はブロック別の三乗根スカラーによる指定作用であり、
全自己同値群の中心を分類する結果は使わない。
実際、$\Gamma_k=\{(\lambda_i):\lambda_i^3\text{ が共通}\}$ の
次数付き作用と共通 grading scalar の自明化から、
命題\ref{mbza:prop:gamma-action}と同じ coherent 作用を直接得られる。
定理\ref{mbnp:thm:product}と系\ref{mbnp:cor:coherent}により、
中心群 $Z_0=(\mu_3^k)/\Delta\mu_3$ について
\[
 \Ku_{\mathrm{coh}}(A)^{Z_0}\simeq\Perf\left(\prod_iE_i\right)
\]
を得る。これらの結果の仮定に $j$ の相異性や $2p_i\ne0$ は含まれない。
可換側にも系\ref{mbrp:cor:comparison-field}と
系\ref{mbrp:cor:base-change}を適用し、同じ積への比較を得る。
係数体が $\C$ 上有限生成でない場合は、各 $f_i$ の有限個の係数を含む
有限生成部分体上で核を作って $\kappa$ へ引き戻す。

基準因子 $b$ を一つ選ぶ。
命題\ref{mbnp:prop:rotations}および命題\ref{mbrd:prop:factor-rotations}により、
二つの積表示の双対群作用は、それぞれ因子 rotation
\[
 R_{i,p_i}=T_{\OO_{E_i}}(t_{p_i})_*(-\otimes L_i),\qquad
 R_{i,0}=T_{\OO_{E_i}}(-\otimes L_i)
\]
の比 $R_{i,p_i}^{-1}R_{b,p_b}$ と $R_{i,0}^{-1}R_{b,0}$ で与えられる。
ここでは非可換側の $a_i\Omega_i$ と可換側の $f_i$ を実方程式として用い、
両側の三乗周期をそれぞれの MF/window から取る。
因子の Koszul 分解と実際の正則中心元による余核二拡張を使うので、
中心元の可換化や自然な線比較の値で周期を割る操作はない。
二つの方程式直線の自然な同型は、この存在証明では仮定しない。

$\kappa$ は代数閉体なので、各 $p_i$ は実際の次数零線束 $M_i$ で表せる。
補題\ref{mbrc:lem:factor-correction}を、今固定した二つの周期に適用すると、
$F_i=-\otimes M_i$ と同型
$F_iR_{i,p_i}\simeq R_{i,0}F_i$ を得る。
三回の合成の差は $\kappa^\times$ の元であり、その三乗根で補正すれば
指定した両周期と両立させられる。
この補題は任意の滑らかな genus-one 曲線と $p_i$ に対する結果で、
基本表示開の有限位数条件を使わない。

外積 $F=\boxtimes_iF_i$ は二つの双対群作用を coherent に intertwine する。
実際、異なる因子の比較射は外積の交換同型と可換し、
各生成元の三乗を消す図式は、今合わせた周期によって可換する。
群元を $0\le d_i<3$ の rotation 比の積として書けば、
これらの交換と三乗の二種類の図式が全ての群法則の coherence を与える。
従って定理\ref{mbrc:thm:finite-point-comparison}の証明と同じ
実同変 perfect 双加群と二重同変化の構成が適用でき、
式\eqref{pps:eq:elliptic-existence}の残余圏同値を得る。

第二の仮定と系\ref{tbp:cor:all-blocks}から、原代数の点スキームは
$\coprod_iE_i$ である。
命題\ref{pps:prop:zero-scheme}を選んだ $f_i$ に適用すると、
この同じ $X$ 上で残るスキームの等式を得る。
\end{proof}

\begin{corollary}[基本満支持型に対する同時比較]
\label{pps:cor:simultaneous-comparison}
$\kappa\supset\C$ を代数閉体とする。
定理\paperref{moduli}{mbs:thm:presentation-stack}の基本楕円領域に属する実表示
\[
 A=\left(\bigotimes_{i=1}^kS_i\right)
       \Big/\left(\sum_{i=1}^ka_i\Omega_i\right),
 \qquad k\ge2,\quad a_i\in\kappa^\times
\]
を取る。すなわち、各因子の点データ $(E_i,L_i,t_{p_i})$ は
$p_i,2p_i,3p_i\ne0$、$j(E_i)\ne0,1728$ を満たし、
$j(E_i)$ は相異なるものとする。
自然な線比較を用いて
\begin{equation}\label{pps:eq:point-target}
 f_i=a_i\vartheta_i(\Omega_i),\qquad
 X_{\mathrm{pt}}=V\left(\sum_i f_i\right)
\end{equation}
と置く。この同じ滑らかな可換三次 $X_{\mathrm{pt}}$ に対し、
実際の dg Morita 同値と Poisson tensor が存在し、
\begin{equation}\label{pps:eq:simultaneous}
 \Ku_{\mathrm{coh}}(A)\simeq\Ku(X_{\mathrm{pt}}),\qquad
 \operatorname{Pt}(A)=Z_{X_{\mathrm{pt}}}(\pi)
          =\coprod_iE_i
\end{equation}
となる。第二の等式は、元の次数一表示と同じ射影空間での
閉部分スキームの等式である。
\end{corollary}

\begin{proof}
定理\paperref{moduli}{mbpl:thm:global-line-comparison}により
$\vartheta_i$ は中心三次直線を $E_i$ の方程式直線へ同型に移す。
従って $f_i$ は $E_i$ の非零方程式である。
残余圏の同値は定理\ref{mbrc:thm:finite-point-comparison}の
代数閉体の場合である。原代数の点関手は
系\ref{tbp:cor:all-blocks}により $\coprod_iE_i$ で表される。
各ブロックの座標を取り、式\eqref{pps:eq:poisson-tensor}を
任意の $c_i\ne0$ で定めれば、
命題\ref{pps:prop:zero-scheme}が残る等式を与える。
\end{proof}

\begin{remark}[表示・体・族に関する範囲]
\label{pps:rem:scope}
点スキームの等式自体に必要なのは、
二次関係が滑らかな平面三次上の固定点のない自己同型のグラフを定め、
各 $\Omega_i$ がその twisted homogeneous coordinate ring を切り出すことだけである。
系\ref{pps:cor:simultaneous-comparison}の追加の有限位数・$j$ の条件は、
自然な $\vartheta_i$ と $X_{\mathrm{pt}}$ を持つ基本表示領域を示す。
存在版の定理\ref{pps:thm:elliptic-existence}には必要ない。
「無限位数」や very general という条件へ置き換える必要はない。

任意の体 $K\supset\C$ 上の基本領域の順序付き実表示でも、
$X_{\mathrm{pt}}$ と上の $\pi$ は $K$ 上で選べ、
点スキームと零スキームの等式は同じ体上で成立する。
ただし引用した残余圏比較は一般には有限拡大後の結果であり、
元の体での比較核の存在とは区別する。

固定した順序付き実表示の族で、ブロック座標と可逆な係数 $c_i$ を
選んだ場合には、点・零スキームの等式は非被約基底を含む任意底変更に両立する。
この等式から、比較核の族全体への延長や、全 $\OO$-marked stack 上の
単一の自然な Poisson tensor の選択・降下は結論しない。
$X_{\mathrm{pt}}$ 自体の自然な降下は
定理\paperref{moduli}{mbpl:thm:commutative-family}による別の結果である。

注意\paperref{moduli}{mbpl:rem:two-torsion-boundary}の二捩れ境界では、
自然な線比較 $\vartheta_i$ が零となる。
しかし、定理\ref{pps:thm:elliptic-existence}の代数的仮定と
固定点なしの点グラフを満たす非零二捩れ因子については、
$f_i$ を別に選ぶ存在証明がそのまま適用できる。
自然な方程式直線の同型を境界へ延ばせないことと、
Ku 比較と零スキーム比較を同時に満たす可換三次の不存在は異なる。
一方、$p_i=0$ や三捩れへの退化では点グラフの仮定を改めて確認する必要があり、
そのような境界全体を上の定理で扱ったとはしない。
\end{remark}

\begin{remark}[中心係数と未解決部分]
\label{pps:rem:coefficients}
固定した $S_i$ に対する点スキームは、全ての $a_i\ne0$ の中心係数に依存しない。
しかし、代数閉体上ではブロック別スカラー $x_i\mapsto s_ix_i$ が
$\Omega_i\mapsto s_i^3\Omega_i$ と作用するので、
これらの満支持中心直線は一つの軌道をなす。
従って係数の不可視性だけから、同じ点スキームを持つ非同型な
中心商の族が得られるとは言えない
（注意\ref{tbp:rem:pencil}）。

本節は $m=0$ の満支持表示を扱う。
中心一次変数を持つ $m>0$ の混合族では定常成分も比較する必要があり、
式\eqref{pps:eq:simultaneous}をそのまま流用しない。
また、任意の内在的な非可換三次がここで扱った表示を持つとは証明していない。
本節の局所計算と既存比較定理の結合について、
先行研究上の新規性の網羅的確認は未了である。
\end{remark}

% !TEX root = ../nc-cubic.tex
\section{中心変数を持つ分離テンソル型の点関手と零スキーム}
\label{csp:sec:central-split}

中心一次変数を持つ場合、満支持型の点スキームに何が加わるかを
原代数の点関手から計算する。
本節の対象は、各楕円因子の関係式が中心変数に依存しない
\emph{分離テンソル表示}である。
中心変数を係数として関係式自体が変形される混合型の表示は含めない。
中心変数が一つなら非零の中心三次項が追加成分を消す一方、
二つ以上では非被約成分の局所 Artin 環が残る。
この区別は閉点集合の計算だけでは検出できない。

\subsection{設定、中心三次式、可換化}

$\kappa\supset\C$ を代数閉体、$k\ge1$、$m\ge1$ とする。
各 $i=1,\ldots,k$ について
\[
 S_i=T(W_i)/(R_i),\qquad \dim W_i=\dim R_i=3
\]
を三次元の二次楕円 Sklyanin 代数とする。
その双線型関係スキームが、滑らかな平面三次
$E_i\subset\PP(W_i^\vee)$ 上の固定点のない自己同型 $\sigma_i$ の
グラフ $\Gamma_{\sigma_i}$ であると仮定する。
また正則中心三次 $\Omega_i$ により
\[
 S_i/(\Omega_i)=B(E_i,\OO_{E_i}(1),\sigma_i)
\]
となるものとする。
$\dim U=m$ とし、全ての生成元に次数一を与えて
\begin{equation}\label{csp:eq:pair}
 \begin{aligned}
 B&=\left(\bigotimes_{i=1}^kS_i\right)\otimes_\kappa\operatorname{Sym}U,
 &V&=\bigoplus_iW_i\oplus U,\\
 w&=\sum_{i=1}^k a_i\Omega_i+g,
 &A&=B/(w),\qquad g\in\operatorname{Sym}^3U
 \end{aligned}
\end{equation}
と置く。$a_i\in\kappa^\times$ とする。
以下の点関手の計算自体は、一部の $a_i$ が零でも成立する。
この表示の三次超曲面の次元は $n=3k+m-2$ である。

\begin{lemma}[分離表示の中心三次式]
\label{csp:lem:center}
$Z(S_i)_1=Z(S_i)_2=0$ を仮定すると
\begin{equation}\label{csp:eq:centers}
 Z(B)_1=U,\qquad Z(B)_2=\operatorname{Sym}^2U,\qquad
 Z(B)_3=\left(\bigoplus_iZ(S_i)_3\right)
                  \oplus\operatorname{Sym}^3U.
\end{equation}
さらに $Z(S_i)_3=\kappa\Omega_i$ ならば、全ての中心三次式は
式\eqref{csp:eq:pair}の分離和の形である。
後者の仮定は、例えば $\sigma_i^j\ne\operatorname{id}$
$(1\le j\le3)$ により保証される。
\end{lemma}

\begin{proof}
補題\ref{tba:lem:tensor-center}のテンソル積の中心の等式を反復し、
全次数三以下の成分を取る。
異なる二因子に正次数を配分した三次項には、いずれかの因子の
一次または二次中心が必要なので、そのような項は消える。
最後の十分条件も同補題による。
\end{proof}

一次・二次中心の消滅\emph{だけ}から三次中心の一次元性を結論しない。
また、式\eqref{csp:eq:centers}は関係式を固定したテンソル表示の結果であり、
一般の中心変数付き二次関係をこの形に直す主張ではない。

各因子の可換化と中心三次式の像を
\begin{equation}\label{csp:eq:artin-data}
 J_i=(S_i)_{\mathrm{ab}}
       =\operatorname{Sym}W_i/(\operatorname{ab}R_i),\qquad
 J=\bigotimes_iJ_i,\qquad
 \omega_i=[\Omega_i]\in(J_i)_3,\qquad
 \theta=\sum_i a_i\omega_i\in J_3
\end{equation}
と書く。$J_i$ を Jacobian 環と呼ぶのは、別途
$\operatorname{ab}R_i=\langle\partial q_i\rangle$ という
三次式 $q_i$ の表示がある場合に限る。
その $q_i$ と点曲線 $E_i$ の三次方程式は、一般には異なる。

\begin{lemma}[可換化の長さ八の構造]
\label{csp:lem:artin}
各 $J_i$ は三つの二次式の Artin 完全交叉であり、
\begin{equation}\label{csp:eq:hilbert}
 H_{J_i}(t)=(1+t)^3,\qquad H_J(t)=(1+t)^{3k},\qquad
 \dim_\kappa J=8^k=:L
\end{equation}
となる。$(J_i)_3$ は socle であり、$\omega_i^2=0$ である。
\end{lemma}

\begin{proof}
$\operatorname{Proj}J_i$ は $\Gamma_{\sigma_i}$ と対角の交差である。
固定点がないので空であり、三つの可換化した二次関係の共通零点は
affine 原点だけである。
従ってこれらは三変数多項式環の正則列をなし、
Hilbert 級数は $(1-t^2)^3/(1-t)^3=(1+t)^3$ である。
完全交叉の Gorenstein 性と最高次数三から socle の主張を得る。
テンソル積を取れば $J$ の式が従う。
\end{proof}

\subsection{任意の試験環上での点関手}

\begin{theorem}[分離テンソル型の点グラフ]
\label{csp:thm:point-functor}
式\eqref{csp:eq:pair}の下で、
\begin{equation}\label{csp:eq:central-schemes}
 C=\operatorname{Proj}(J\otimes\operatorname{Sym}U),\qquad
 D=\operatorname{Proj}\bigl((J\otimes\operatorname{Sym}U)/(g+\theta)\bigr)
\end{equation}
と置く。
$B$ の双線型関係スキームと $A$ の点関手はそれぞれ
\begin{align}
 \Gamma_B&=\left(\coprod_{i=1}^k\Gamma_{\sigma_i}\right)\amalg\Delta_C,
 \label{csp:eq:graph}\\
 \operatorname{Pt}(A)&=\left(\coprod_{i=1}^kE_i\right)\amalg D
 \label{csp:eq:points}
\end{align}
で与えられる。点加群の shift は $E_i$ 上では $\sigma_i$、
$D$ 上では恒等写像である。
これらは同じ射影空間内のスキームの等式であり、
任意の可換 $\kappa$-試験環への基底変更と両立する。
\end{theorem}

\begin{proof}
点対の二つの座標を $p,q$ とする。
ある中心座標 $u$ の $p$ における値が可逆な chart では、
中心生成元との可換関係
\[
 u(p)v(q)-v(p)u(q)=0\qquad(v\in V)
\]
により $q$ は $p$ に比例する。
$q$ が射影点なので、必要なら chart を細分すれば $u(q)$ も可逆であり、
両方を一に正規化して $p=q$ を得る。
各ブロックの残る関係はちょうど $\operatorname{ab}R_i$ である。
従ってこの開集合上の関係スキームは $\Delta_C$ となる。

中心座標が全て零である幾何点については、
定理\ref{tbp:thm:point-graph}および系\ref{tbp:cor:all-blocks}の
ブロック間可換関係の議論を適用できる。
二つ以上のブロックが非零なら、ある非零ブロック内の $p_i,q_i$ が
比例し、$\sigma_i$ の固定点を与えてしまう。
従って残る幾何点は純ブロックの $\Gamma_{\sigma_i}$ に属する。

この台の議論に加えて、純ブロック $i$ の近傍でイデアルを計算する。
$p_i,q_i$ は比例しないので、その二列から作るある二次小行列式が
可逆な近傍を取れる。
他の任意の座標 $v$、すなわち別ブロックの座標または中心座標に対し、
ブロック間の二つの可換関係は
\[
 p_{ia}v(q)-q_{ia}v(p)=0,\qquad
 p_{ib}v(q)-q_{ib}v(p)=0
\]
である。係数行列の行列式が可逆なので $v(p)=v(q)=0$ となる。
よって横断方向の冪零成分は残らず、この近傍は正確に
$\Gamma_{\sigma_i}$ である。
以上の開集合は関係スキーム全体を覆い、
$C_{\mathrm{red}}=\PP(U^\vee)$ は各純ブロック平面と交わらない。
これで式\eqref{csp:eq:graph}を得る。

各射影が同型であるこのグラフの反復は、任意の試験環上で
$B$ の無限点加群を与える。
$E_i$ 上の普遍点加群は $S_i/(\Omega_i)$ を経由するので、
$w$ はその上で零となる。
一方、定常な点加群では任意の非可換語の値は可換化した語の値と一致する。
従って追加の三次関係は正確に $g+\theta=0$ であり、
shift が恒等であるため全ての次数で同じ条件となる。
これが式\eqref{csp:eq:points}を与える。

局所計算は全て多項式関係と単元による消去であり、
試験環の被約性を用いない。
従って点の集合だけでなく、任意の試験環上の関手として等式が成立する。
\end{proof}

\subsection{中心超曲面の厚みと中心変数一つの場合}

$P=\PP(U^\vee)$ と書く。
中心座標による射影 $f:C\to P$ は有限平坦で、
\begin{equation}\label{csp:eq:central-algebra}
 f_*\OO_C=\bigoplus_{d=0}^{3k}J_d\otimes\OO_P(-d)
\end{equation}
である。ここで積は $J$ の次数付き積による。
特に $C$ は階数 $L$ の、中心射影空間の非被約な厚みである。

\begin{proposition}[中心成分の消滅と完全交叉]
\label{csp:prop:central-component}
次が成立する。
\begin{enumerate}
 \item $m=1$、$g=cu^3$、$c\ne0$ ならば $D=\varnothing$ である。
 従って $k=2,m=1,n=5$ の場合を含め、
 \[
 \operatorname{Pt}(A)=\coprod_i E_i
 \]
 がスキームとして成立する。
 \item $m\ge2$、$g\ne0$ ならば $D$ は
 $\PP^{3k+m-1}$ 内の型 $(2^{3k},3)$ の完全交叉であり、
 \begin{equation}\label{csp:eq:complete-intersection}
 \dim D=m-2,\qquad \deg D=3\cdot8^k,\qquad
 \omega_D=\OO_D(3k+3-m).
 \end{equation}
 $D_{\mathrm{red}}=V_P(g)_{\mathrm{red}}$ であり、
 中心射影の第零 Fitting イデアルは $(g^L)$ である。
 \item $s=\#\{i:a_i\omega_i\ne0\}$ とする。
 $D\to P$ のスキーム論的像は $V_P(g^{s+1})$ である。
 特に $\theta\ne0$、$m\ge2$、$g\ne0$ なら、この自然な射影は
 $V_P(g)$ を経由しない。
\end{enumerate}
最後の主張では $m=1$ の空な像も同じ式で解釈する。
\end{proposition}

\begin{proof}
中心 chart $u_j=1$ 上で $C$ は $\operatorname{Spec}(J\otimes\OO_P)$
となるので、式\eqref{csp:eq:central-algebra}を得る。
$m=1$ ならこの chart は $C$ 全体を覆い、
$D=\operatorname{Spec}J/(c+\theta)$ である。
$\theta$ は冪零なので $c+\theta$ は可逆となり、第一の主張が従う。

次に $g\ne0$ とする。
$J\otimes\operatorname{Sym}U$ を $J$ の次数で濾過すると、
$g+\theta$ の乗法の次数零部分は $g$ の乗法である。
各次数成分は $\operatorname{Sym}U$ 上自由であり、$g$ は非零因子なので、
最小の非零次数を比較して $g+\theta$ も非零因子であることが分かる。
補題\ref{csp:lem:artin}の二次正則列にこの三次式を加えると、
完全交叉の次元・次数・双対化層の公式から
式\eqref{csp:eq:complete-intersection}を得る。
$\theta$ が冪零なので被約化の式も従う。
また有限射 $f$ について
\begin{equation}\label{csp:eq:pushforward-sequence}
 0\longrightarrow f_*\OO_C(-3)
 \xrightarrow{\ g+\theta\ } f_*\OO_C
 \longrightarrow f_*\OO_D\longrightarrow0
\end{equation}
である。
局所的な行列式は $\det(gI+m_\theta)=g^L$ なので、
第零 Fitting イデアルは $(g^L)$ となる。

非零な $a_i\omega_i$ は互いに異なるテンソル因子の socle に属する。
従って
\[
 \theta^s=s!\prod_{a_i\omega_i\ne0}a_i\omega_i\ne0,
 \qquad \theta^{s+1}=0
\]
である。$s=0$ のときは $\theta^0=1$ と読む。
乗法 $m_\theta$ の最大 Jordan block の大きさは $s+1$ である。
中心 chart でこの作用素を $\kappa$ 上の Jordan 標準形に取ると、
$gI+m_\theta$ の余核は、各 block の大きさ $\ell$ に対する
$\OO_P/(g^\ell)$ の直和となる。
その零化イデアルは $(g^{s+1})$ であり、
単位を持つ $\OO_P$-代数 $f_*\OO_D$ の零化イデアルは
$\OO_P\to f_*\OO_D$ の核と等しい。
従ってスキーム論的像は $V_P(g^{s+1})$ である。
$D$ 上では $g=-\theta$ であるから、$\theta\ne0$ の場合に
$V_P(g)$ を経由するという結論にはならない。
\end{proof}

Fitting イデアル $(g^L)$ と、像を定めるイデアル $(g^{s+1})$ は
異なる情報である。
特に全ての $\omega_i$ が非零なら像は $(k+1)$ 重の中心超曲面だが、
その横断的な長さは $8^k$ である。
$\omega_i$ の非零性は固定点がないことだけからは仮定しない。
一方 $g=0$ の境界では
$D=\operatorname{Proj}((J/(\theta))\otimes\operatorname{Sym}U)$ となり、
上の三次非零因子の議論は適用しない。

\begin{proposition}[滑らかな中心超曲面に沿う局所 Artin 環]
\label{csp:prop:local-artin}
$m\ge2$、$g\ne0$、$H=V_P(g)$ が滑らかであると仮定する。
任意の閉点 $q\in H$ において
\begin{equation}\label{csp:eq:completed-local}
 \widehat{\OO}_{D,q}\simeq
 J[[t_1,\ldots,t_{m-2}]]
\end{equation}
である。この同型は局所座標の選択に依存する。
$m=2$ ならば
\begin{equation}\label{csp:eq:binary-central}
 D\simeq\coprod_{q\in H}\operatorname{Spec}J
       =\coprod_{a=1}^3\operatorname{Spec}J.
\end{equation}
\end{proposition}

\begin{proof}
中心 chart で $g$ を非斉次化する。
$dg(q)\ne0$ なので、完備局所環において $g$ 自身を正則局所座標の
一つとし、残りを $t_1,\ldots,t_{m-2}$ とできる。
従って $D$ の完備局所環は
\[
 J[[g,t_1,\ldots,t_{m-2}]]/(g+\theta)
 \simeq J[[t_1,\ldots,t_{m-2}]]
\]
となる。$\theta$ が冪零であるため $g=-\theta$ の代入は正当である。
滑らかな二変数三次は射影直線上の三つの相異なる点を切り出す。
この場合の局所環は既に Artin 環で完備なので、最後の式を得る。
\end{proof}

式\eqref{csp:eq:completed-local}は $D$ が大域的に
$H\times\operatorname{Spec}J$ であることを意味しない。
また、自然な中心射影 $D\to P$ の $q$ 上のファイバーは
$\operatorname{Spec}(J/(\theta))$ であり、
局所環全体を記述する $J$ と区別する必要がある。

\subsection{同じ分離三次上の Poisson 零スキーム}

各 $E_i$ の非零方程式を $f_i\in\operatorname{Sym}^3W_i$ とし、
\begin{equation}\label{csp:eq:commutative-target}
 F=\sum_i f_i+g,\qquad
 X=V(F)\subset\PP(V^\vee),\qquad
 J_i^{\mathrm{cl}}=\operatorname{Jac}(f_i),\qquad
 J^{\mathrm{cl}}=\bigotimes_iJ_i^{\mathrm{cl}}
\end{equation}
と置く。
$m=1$ なら $g=cu^3$、$c\ne0$、
$m\ge2$ なら $g\ne0$ かつ $H=V_P(g)$ は滑らかと仮定する。
これらの仮定により $X$ は滑らかである。
各 $W_i$ の座標を $x_{i1},x_{i2},x_{i3}$ とし、
\begin{equation}\label{csp:eq:block-poisson}
 \beta=\sum_i b_i\,\iota_{df_i}
  (\partial_{x_{i1}}\wedge\partial_{x_{i2}}\wedge\partial_{x_{i3}}),
 \qquad b_i\in\kappa^\times
\end{equation}
を取る。$\beta$ は Poisson で $F$ は Casimir だから、
射影化して $X$ 上の Poisson tensor $\pi$ を得る。

\begin{theorem}[中心変数を持つ分離和の零スキーム]
\label{csp:thm:poisson-zero}
式\eqref{csp:eq:commutative-target}--\eqref{csp:eq:block-poisson}の下で
\begin{equation}\label{csp:eq:poisson-zero}
 Z_X(\pi)=\left(\coprod_i E_i\right)\amalg D_{\mathrm{cl}},\qquad
 D_{\mathrm{cl}}=\operatorname{Proj}
       \bigl((J^{\mathrm{cl}}\otimes\operatorname{Sym}U)/(g)\bigr).
\end{equation}
この式は任意の $\kappa$-基底変更に両立する。
$m=1$ なら $D_{\mathrm{cl}}=\varnothing$ であり、
\begin{equation}\label{csp:eq:one-central-equality}
 \operatorname{Pt}(A)=Z_X(\pi)=\coprod_iE_i
\end{equation}
が同じ射影空間内のスキームの等式として成立する。
$m\ge2$ なら $D_{\mathrm{cl}}\to H$ は有限平坦で階数 $L=8^k$、
その代数は
\begin{equation}\label{csp:eq:commutative-thickening}
 \bigoplus_{d=0}^{3k}(J^{\mathrm{cl}})_d\otimes\OO_H(-d)
\end{equation}
である。特に被約スキームは常に一致し、
\begin{equation}\label{csp:eq:reduced-equality}
 \operatorname{Pt}(A)_{\mathrm{red}}
   =Z_X(\pi)_{\mathrm{red}}
   =\left(\coprod_iE_i\right)\amalg H
\end{equation}
となる。
\end{theorem}

\begin{proof}
Euler ベクトル場を $\xi$ とすると、
命題\ref{pps:prop:zero-scheme}の局所直和因子の議論により、
$Z_X(\pi)$ のイデアルは $\xi\wedge\beta$ の係数で与えられる。
中心座標 $u_j$ が可逆な chart では、
第 $i$ ブロックの二つの tangent 基底と $\partial_{u_j}$ からなる
係数が、符号を除いて $u_jb_i\partial_{x_{ia}}f_i$ となる。
従って零イデアルは全ての $\partial_{x_{ia}}f_i$ を含み、
逆の包含も $\beta$ の式から成り立つ。
各ブロックの Euler 恒等式により
$f_i=0$ in $J_i^{\mathrm{cl}}$ なので、$F$ の残る条件は $g=0$ である。
これが $D_{\mathrm{cl}}$ を与える。

中心座標が全て零である幾何点では、
非零 active ブロックの数を $r$ とすると $\beta$ の階数は $2r$ である。
射影化が零なら $\beta=\xi\wedge v$ なので階数は高々二であり、
純ブロック以外の点は存在しない。
純ブロック $E_i$ の近傍では、ある偏微分 $\partial_{x_{ia}}f_i$ が可逆である。
$\xi\wedge\beta$ の混合係数は、この偏微分に他ブロックの座標または
中心座標を掛けたものとなる。
従って全ての横断座標が零になり、残るイデアルは $(f_i)$ である。
この局所消去は非被約な基底でも有効なので、
式\eqref{csp:eq:poisson-zero}をスキームとして得る。

$m=1$ の中心 chart では $g=c$ が可逆なので $D_{\mathrm{cl}}$ は空である。
命題\ref{csp:prop:central-component}と合わせると
式\eqref{csp:eq:one-central-equality}が従う。
$m\ge2$ では $g$ が中心変数だけの式であるため、
中心射影は $H$ を経由し、式\eqref{csp:eq:commutative-thickening}を得る。
被約化の等式は両側の明示式から従う。
\end{proof}

\begin{corollary}[中心変数一つの分離型における同時存在]
\label{csp:cor:simultaneous-one-center}
基礎体を $\C$ とし、$m=1$、$N=3k+1\ge5$ とする。
本節の点グラフの仮定に加えて、各 $S_i,\Omega_i$ が
式\eqref{csc:eq:factor}の実楕円因子の仮定を満たすとする。
すなわち $S_i$ は連結局所有限な二次 Koszul AS 正則代数で、
大域次元と Gorenstein parameter は三、$\Omega_i$ は次数三の実際の正則中心元であり、
$S_i/(\Omega_i)\simeq B(E_i,\OO_{E_i}(1),t_{p_i})$ とする。
さらに、本節で仮定した通り二次関係スキームは固定点のない
$t_{p_i}$ のグラフであるとする。
$a_i,c\in\C^\times$ とし、
\[
 A=\left(\bigotimes_iS_i\right)[u]
       \Big/\left(\sum_i a_i\Omega_i+cu^3\right),\qquad
 X=V\!\left(\sum_i f_i+cu^3\right)\subset\PP^{3k}_\C
\]
と置く。ここで $f_i$ は同じ埋込み平面三次 $E_i$ の任意の非零方程式である。
この $X$ は滑らかで、実際の dg Morita 同値と Poisson tensor が存在し、
\begin{equation}\label{csp:eq:simultaneous-one-center}
 \Ku_{\mathrm{coh}}(A)\simeq\Ku(X),\qquad
 \operatorname{Pt}(A)=Z_X(\pi)=\coprod_iE_i
\end{equation}
を同時に満たす。第二の等式は元の次数一射影空間でのスキームの等式である。
特に $k=2,m=1,n=5$ の分離型で、同じ可換標的が二つの条件を満たす。
\end{corollary}

\begin{proof}
$\C[u]/(\partial_u(cu^3))=\C[u]/(u^2)$ は有限次元だから、
定理\ref{csc:thm:comparison}の中心因子の孤立 critical point 条件を満たす。
指定した実楕円因子の仮定により同定理を適用でき、第一の同値を得る。
式\eqref{csp:eq:block-poisson}で任意の $b_i\ne0$ を選ぶと、
定理\ref{csp:thm:poisson-zero}の式\eqref{csp:eq:one-central-equality}が
残るスキームの等式を与える。
\end{proof}

\subsection{抽象スキーム同型にも残る Artin 環の障害}

\begin{proposition}[局所 Artin 環による必要条件]
\label{csp:prop:abstract-obstruction}
$m\ge2$、$g\ne0$、$H$ が滑らかであるという上の仮定を保つ。
抽象 $\kappa$-スキームとして
$\operatorname{Pt}(A)\simeq Z_X(\pi)$ ならば
\begin{equation}\label{csp:eq:necessary-artin}
 J\simeq J^{\mathrm{cl}}
 \quad\text{as graded }\kappa\text{-algebras}
\end{equation}
である。
$m=2$ の場合、この条件は抽象スキーム同型の存在の必要十分条件である。
また任意の $m\ge2$ で $\theta=0$ なら、
式\eqref{csp:eq:necessary-artin}は同じく十分である。
\end{proposition}

\begin{proof}
各 $E_i$ は被約だが、$D$ と $D_{\mathrm{cl}}$ の全ての点で
局所環は非被約である。
従って全体のスキーム同型は、$D$ の点を $D_{\mathrm{cl}}$ の点へ移す。
命題\ref{csp:prop:local-artin}と可換側の同じ計算により、
対応する完備局所環は
\[
 J[[t_1,\ldots,t_{m-2}]],\qquad
 J^{\mathrm{cl}}[[t_1,\ldots,t_{m-2}]]
\]
である。その極大イデアルに関する随伴次数付き環はそれぞれ
$J[t_1,\ldots,t_{m-2}]$ と
$J^{\mathrm{cl}}[t_1,\ldots,t_{m-2}]$ で、全ての変数の次数は一である。
これらの次数付き環の冪零根基の次数一部分は、それぞれ $J_1$ と
$(J^{\mathrm{cl}})_1$ である。
その次数一部分が生成する単位的部分代数は、正確に $J$ と
$J^{\mathrm{cl}}$ なので、環同型から式\eqref{csp:eq:necessary-artin}を回復できる。

$m=2$ なら式\eqref{csp:eq:binary-central}により、
中心成分はそれぞれ三つの $\operatorname{Spec}J$ と
$\operatorname{Spec}J^{\mathrm{cl}}$ の非交和である。
従って次数付き環同型は中心成分の同型を与え、
$E_i$ 上の恒等写像と合わせて全体の同型を得る。
$\theta=0$ の場合は一般の $m$ でも、式\eqref{csp:eq:central-algebra}を
$H$ に制限した次数付き代数が両側で同型となるので十分性が従う。
\end{proof}

\begin{corollary}[ブロック別の $j$ 不変量による不一致の判定]
\label{csp:cor:block-j-obstruction}
命題\ref{csp:prop:abstract-obstruction}の仮定の下で、さらに
$J_i=\operatorname{Jac}(q_i)$ という滑らかな平面三次式 $q_i$ の表示を持ち、
各 $q_i,f_i$ が Sebastiani--Thom 分解不能であるとする。
例えば、各平面三次の $j$ 不変量が非零なら分解不能である。
多重集合として
\begin{equation}\label{csp:eq:j-obstruction}
 \{j(V(q_i))\}_{i=1}^k\ne\{j(E_i)\}_{i=1}^k
\end{equation}
ならば、$\operatorname{Pt}(A)$ と式\eqref{csp:eq:block-poisson}の
$Z_X(\pi)$ は抽象スキームとしても同型でない。
\end{corollary}

\begin{proof}
両スキームが同型なら、命題により
\[
 \operatorname{Jac}\left(\sum_iq_i\right)
       \simeq\operatorname{Jac}\left(\sum_if_i\right)
\]
という次数付き環同型を得る。
次数一・二の積から勾配の二次式空間が線形同型で対応するので、
命題\ref{prop:ej-projective-uniqueness}により二つの滑らかな三次和は
線形座標変換で移り合う。
補題\ref{ej:lem:symmetrizer-algebra}の最細分解の一意性と
各因子の分解不能性により、$q_i$ と $f_i$ は順序を除いて
因子ごとに射影同型となる。これは式\eqref{csp:eq:j-obstruction}に矛盾する。
なお、滑らかな三変数三次式が非自明に直和分解すると、
二変数の滑らかな三次式と一変数の三乗との和に直せる。
二変数三次の三つの根は射影変換で互いに移せるので、これは Fermat 型であり、
その $j$ 不変量は零である。
\end{proof}

\begin{remark}[二つの予想への適用範囲]
\label{csp:rem:scope}
本節は点関手と、明示した可換分離三次上の block Poisson tensor の
零スキームを計算するものである。
$\C$ 上で定理\ref{csc:thm:comparison}の実楕円因子と中心多項式の仮定を
満たす分離型では、この同じ $X$ への残余圏同値も構成されている。
$m=1$ での同時存在は系\ref{csp:cor:simultaneous-one-center}である。
$m=0$ の比較は既に定理\ref{pps:thm:elliptic-existence}で構成した。
これらの正の結果は分離テンソル表示に対するものである。
関係式自体が中心変数に依存する混合型については、
第\ref{cmo:sec:points}節の有限 norm による点スキームの判定を用い、
上の分離型の残余圏比較をそのまま適用しない。

$m\ge2$ では、被約部分の一致を full scheme の一致と取り違えない。
系\ref{csp:cor:block-j-obstruction}は、この $X$ とこの block tensor
という自然な処方の障害である。
同じ $X$ の別の Poisson tensor や別の残余圏標的の存在まで排除しない。
その全称的な排除は、三次三様体の場合の
定理\ref{ppz:thm:all-poisson-obstruction}で別途扱った。
\end{remark}

% !TEX root = ../nc-cubic.tex
\section{一つの複素 Sklyanin 因子の厳密な入力確認}
\label{esi:sec:input}

本節では、形式パラメータの特殊化に関する主張を使わず、
複素数上の一つの楕円因子とその指定中心三次式を直接確認する。
以下で非可換な語と可換多項式を区別し、
\[
 S=x^3+y^3+z^3,\qquad
 C_+=xyz+yzx+zxy,\qquad C_-=yxz+xzy+zyx
\]
は自由結合代数の元とする。

\begin{proposition}[指定中心元を持つ具体的楕円因子]
\label{esi:prop:factor}
次数一の生成元を持つ複素代数と中心元の候補を
\begin{equation}\label{esi:eq:factor}
 \begin{split}
 D&=\mathbb C\langle x,y,z\rangle/
 (yz-3zy-2x^2,\ zx-3xz-2y^2,\ xy-3yx-2z^2),\\
 \Omega&=\frac{79}{144}S+\frac5{96}(C_++C_-)
 \end{split}
\end{equation}
と定める。また
\begin{equation}\label{esi:eq:elliptic-curve}
 \begin{gathered}
  \lambda=\frac{17}{9},\qquad
  g_\lambda=\frac{x^3+y^3+z^3}{3}+\lambda xyz,\\
  E=V(g_\lambda)\subset\mathbb P^2,\qquad
  o=(1:-1:0),\quad p=(1:-3:-2)
 \end{gathered}
\end{equation}
と置き、$o$を$E$の原点とする。次が成り立つ。
\begin{enumerate}
 \item $E$は滑らかな楕円曲線で、$p\ne o$は変曲点ではない。
 二次関係の多重線型化は、スキームとして平行移動$t_p:E\to E$のグラフである。
 特にこのグラフと対角線の交わりは空である。
 \item $D$は Noetherian domain で、Koszul、AS 正則である。
 大域次元と Gorenstein parameter はともに$3$、Hilbert 級数は$(1-t)^{-3}$であり、
 Nakayama 自己同型は恒等である。
 \item $\Omega$は非零の中心正則元である。
 標準的な Sklyanin 中心生成元$c_3$を
 \cite[Sections~1--2, pp.~1, 3]{GRSklyaninDLL}の規約で取ると
 \begin{equation}\label{esi:eq:center-normalization}
  \Omega=-\frac1{48}c_3.
 \end{equation}
 \item 次数付き代数の同型
 \begin{equation}\label{esi:eq:thcr}
  D/(\Omega)\simeq B(E,\OO_E(1),t_p)
 \end{equation}
 が成り立つ。
\end{enumerate}
\end{proposition}
\begin{proof}
まず$g_\lambda$の偏微分は$x^2+\lambda yz$とその巡回置換である。
三つの偏微分が射影点で同時に零なら、いずれかの座標が零の場合は
残る座標も零となる。全座標が非零の場合は三式の積から$1+\lambda^3=0$を得る。
しかし
\[
 1+\lambda^3=\frac{5642}{729}\ne0
\]
なので$E$は滑らかである。$o,p\in E$は直接代入で分かる。
$o$は変曲点なので楕円曲線の原点として選べる。
Hessian の計算は
\[
 \det\operatorname{Hess}(g_\lambda)
  =(8+2\lambda^3)xyz-2\lambda^2(x^3+y^3+z^3)
\]
を与え、$E$上では$8(1+\lambda^3)xyz$となる。
特に$p$での値は$90272/243\ne0$であり、$p$は変曲点でない。
従って$3p\ne o$であり、$p\ne o$でもある。

左側の射影座標を$(x:y:z)$、右側を$(X:Y:Z)$とすると、
多重線型化は
\begin{equation}\label{esi:eq:graph-matrix}
 \begin{gathered}
  M(x,y,z)\begin{pmatrix}X\\Y\\Z\end{pmatrix}=0,
  \qquad
  M(x,y,z)=
  \begin{pmatrix}
   -2x&-3z&y\\ z&-2y&-3x\\ -3y&x&-2z
  \end{pmatrix},\\
  \det M=-18g_\lambda
 \end{gathered}
\end{equation}
である。$E$上で$M$の階数はちょうど$2$となる。
実際、階数が$1$以下なら余因子がすべて零となり、
$\det M$の全偏微分も零となって$E$の滑らかさに反する。
従って、可逆な二次小行列式の開集合上で二つの右座標を消去すると、
残る方程式は$\det M=0$となる。
これは多重線型化で定まる閉部分スキームが
$E$上の核直線束の射影化であり、第一射影が$E$との同型であることを示す。
左右を交換した行列の行列式も$-18g_\lambda$なので第二射影も同型である。
よってこの閉部分スキームは自己同型$\tau:E\to E$のグラフである。

対角線上の三関係は、非零定数倍を除いて
\[
 x^2+yz=0,\qquad y^2+zx=0,\qquad z^2+xy=0
\]
となる。座標が一つでも零ならすべて零となり、全座標が非零なら三式の積が
$(xyz)^2=-(xyz)^2$を与えるので、射影共通零点はない。
従って対角線とのスキーム交叉も空である。
楕円曲線の自己同型を$t_q\alpha$と書き、$\alpha$が原点を固定するとする。
$\alpha\ne1$なら$1-\alpha$は非零 isogeny で全射なので、
$t_q\alpha$は固定点を持つ。よって$\tau$は平行移動である。
さらに$M(o)(1,-3,-2)^{\mathsf t}=0$より$\tau(o)=p$、従って$\tau=t_p$である。
以上のグラフの確認は閉点の対応だけでなく、スキームとしての確認である。

式\eqref{esi:eq:factor}は Sklyanin parameters
$(a:b:c)=(1:-3:-2)$の表示である。
その標準的な曲線の方程式
\[
 abc(x^3+y^3+z^3)-(a^3+b^3+c^3)xyz=0
\]
は$6(x^3+y^3+z^3)+34xyz=18g_\lambda=0$に一致する。
この表示と$(E,t_p,\OO_E(1))$の対応、および中心商の記述は
\cite[Section~1, p.~1]{GRSklyaninDLL}による。
滑らかな楕円曲線、平行移動、次数$3$の可逆層という入力に対し、
Tate--Van den Bergh の Theorem~1.1 と Corollary~1.3
\cite[著者版 pp.~2--3; 構成は Section~4.1]{GRSklyaninTV}は
Hilbert 級数、Koszul 性、Noetherian 性、domain 性と AS 正則性を与える。
これらの主張には$p$が有限位数であるという仮定はない。
Koszul 性と Hilbert 級数から単純加群の極小分解は
\[
 0\longrightarrow D(-3)\longrightarrow D(-2)^{\oplus3}
 \longrightarrow D(-1)^{\oplus3}\longrightarrow D
 \longrightarrow\mathbb C\longrightarrow0
\]
となり、Gorenstein parameter は$3$である。
また$D$は cyclic potential
\[
 \Psi=xyz-3yxz-\frac23(x^3+y^3+z^3)
\]
の三つの cyclic derivative を関係式とする Jacobian 代数である。
三生成元の Jacobian 代数が AS 正則なら untwisted $3$-Calabi--Yau であること
\cite[Corollary~2.9]{MoriSmith}から、Nakayama 自己同型は恒等となる。

指定した中心元の規格化を確認する。
\cite[Section~2, p.~3]{GRSklyaninDLL}の公式は
\[
 \begin{aligned}
 3c_3&=a(b^3-c^3)C_++b(c^3-a^3)C_-+c(a^3-b^3)S\\
     &=-19C_++27C_--56S
 \end{aligned}
\]
を与える。一方、三関係をそれぞれ右から$x,y,z$で掛けて加えると
$C_+-3C_--2S=0$である。従って
\[
 \begin{aligned}
 3c_3&=-94S-30C_-,\\
 -144\Omega&=-79S-\frac{15}{2}(C_++C_-)
              =-94S-30C_-
 \end{aligned}
\]
となり、式\eqref{esi:eq:center-normalization}と中心性を得る。

非零性は有限次元の明示的な証明でも確認できる。
自由結合代数の次数三部分に線型汎関数$\ell$を
\[
 \ell(x^3)=\ell(y^3)=\ell(z^3)=1,\qquad
 \ell(xyz)=\ell(yzx)=\ell(zxy)=2
\]
と定め、他の語では零とする。
各二次関係を左または右から対応する二乗項の変数で掛けたとき、
その像は$2-3\cdot0-2\cdot1=0$となる。
別の変数で掛けたときは、すべての語で$\ell$が零となる。
従って$\ell$は二次関係イデアルの次数三部分を消し、$D_3$上に降下する。
そこで
\[
 \ell(\Omega)=\frac{79}{144}\cdot3+\frac5{96}\cdot6
             =\frac{47}{24}\ne0
\]
なので$\Omega\ne0$である。$D$は domain だから$\Omega$は左右とも正則である。
最後に、上記の標準中心元は
$D/(c_3)=B(E,\OO_E(1),t_p)$を満たす
\cite[Section~1, p.~1]{GRSklyaninDLL}。
式\eqref{esi:eq:center-normalization}の係数は非零なので、
$(\Omega)=(c_3)$となり式\eqref{esi:eq:thcr}を得る。
\end{proof}

\begin{corollary}[二つの中心変数を加えた実三次対]
\label{esi:cor:ambient}
命題\ref{esi:prop:factor}の$D,\Omega$を用いて
\begin{equation}\label{esi:eq:ambient}
 B=D[u,v],\qquad w=\Omega+\frac{u^3+v^3}{3},\qquad A=B/(w)
\end{equation}
と置く。$u,v$は中心の次数一変数である。
$B$は Noetherian、Koszul、AS 正則な domain で、
大域次元と Gorenstein parameter はともに$5$、Nakayama 自己同型は恒等である。
$w$は中心正則な次数三元で、
\[
 H_A(t)=\frac{1-t^3}{(1-t)^5}.
\]
$A$は非可換であり、対応する可換式
\[
 g_{17/9}(x,y,z)+\frac{u^3+v^3}{3}
\]
は滑らかな複素三次三様体を定める。
\end{corollary}
\begin{proof}
中心多項式環とのテンソル積により、命題\ref{esi:prop:factor}の
Noetherian 性、domain 性と Koszul 性は保たれる。
左右の Koszul 分解とその双対をテンソル積すると、
AS 正則性、次元$5$、parameter $5$を得る。
同様に Calabi--Yau 双加群分解のテンソル積から Nakayama 自己同型は恒等である。
$w$は$v$に関して最高次係数$1/3$の多項式なので非零であり、
domain $B$上で中心正則である。正則元の短完全列から Hilbert 級数が従う。

$yz-zy$は$D$の次数二の関係空間に属さない。
実際、その空間で各関係の二乗項の係数を見ると、
$yz-zy$を表す線型結合の三つの係数はすべて零でなければならない。
次数三の元で割っても次数二は変わらないので、$A$も非可換である。
最後の可換式の偏微分が同時に零なら$u=v=0$となり、
残る三式は$g_{17/9}$の偏微分である。
命題\ref{esi:prop:factor}で確認した滑らかさにより、
その共通零点はアフィン原点だけである。
\end{proof}

系\ref{esi:cor:ambient}は、実際の複素代数と中心商を与えるための入力である。
$\operatorname{qgr}A$の滑らかさと全 coherent Kuznetsov 成分の比較は、
楕円因子の比較、中心 Fermat 因子との外積、および coherent 飽和を適用して別に得る。
本節の正則性の証明は、有限個の Hilbert 関数の検算や、
形式族のある開集合が$h=1$を含むという推論には依存しない。

% !TEX root = ../nc-cubic.tex
\section{点スキームと Poisson 零スキームを比較する際の障害}
\label{ppz:sec:comparison}

点スキームと Poisson tensor の零スキームとの一致を検証する際には、
抽象スキームとしての同型と、既存の次数一偏極に関する埋込み同型を区別する。
まず一ブロック型の零スキームを計算し、埋込み同型に対する障害を示す。
次に長さ八の局所 Artin 環と Hessian 階数二の幾何を用い、
複素数体上の具体的な非可換三次三様体について、
同じ Kuznetsov 成分を持つ全ての可換三次と全ての Poisson tensor を調べる。
その例では可換 Kuznetsov 実現が直接構成できる一方、
全点スキームと零スキームの抽象同型は存在しない。
圏比較には系\ref{csc:cor:fermat}の分離テンソル比較を用いる。

\subsection{射影 Poisson tensor の零イデアル}

$k$ を標数零の代数閉体、$n\geq3$、$r=n-1$ とし、
$V=\langle x,y,z,u_1,\ldots,u_r\rangle$、
$K=\operatorname{ann}\langle u_1,\ldots,u_r\rangle\subset V^*$ と置く。
$H\in\Sym^3V$ に対し、$X=V(H)\subset\PP(V^*)$ と
$g_H=H|_{u=0}$ の定める平面三次がともに滑らかとする。
$0\ne P\in\det K$ を選び、
\[
 \Pi_H=\iota_{dH}P,\qquad \pi_H=\Pi_H|_X
\]
とする。$Z_X(\pi_H)$ は、切断 $\pi_H\in H^0(X,\bigwedge^2T_X)$
の零スキーム、すなわち
$\bigwedge^2\Omega_X\xrightarrow{\pi_H}\OO_X$ の像イデアルで定義する。

\begin{proposition}[一ブロック Poisson 零スキームの分解]
\label{ppz:prop:zero-scheme}
上の仮定の下で、スキームとして
\begin{equation}\label{ppz:eq:zero-decomposition}
 Z_X(\pi_H)=E_H\amalg D_H,\qquad
 E_H=V(u_1,\ldots,u_r,H),\quad
 D_H=V(H_x,H_y,H_z,H)
\end{equation}
となる。
$C_H=V(H_x,H_y,H_z)$ 上で $H$ が非零因子ならば、
$D_H$ は型 $(2,2,2,3)$、次元 $n-3$、次数 $24$ の完全交叉である。
相対的な滑らかさと、この相対 Cartier 条件を仮定した族では、
表示した零イデアルの同定は任意の基底変更と両立する。
\end{proposition}

\begin{proof}
$P$ の非零定数倍は零イデアルを変えないので、
$P=\partial_x\wedge\partial_y\wedge\partial_z$ としてよい。
射影空間の Euler 列は局所的に分裂し、その外積により、
斉次 bivector の射影像が零となるイデアルは
Euler vector field $\operatorname{Eu}$ との積
$\operatorname{Eu}\wedge\Pi_H$ の係数で計算できる。
また $X$ が滑らかなので $T_X\subset T_{\PP(V^*)}|_X$ は部分束であり、
$\bigwedge^2T_X$ の切断を周囲の外積束で計算しても零イデアルは変わらない。
従って $X$ 上の零イデアルは
\begin{equation}\label{ppz:eq:zero-ideal}
 \bigl(u_aH_x,u_aH_y,u_aH_z\ (1\leq a\leq r),
       xH_x+yH_y+zH_z\bigr)
\end{equation}
で与えられる。

ある $u_a$ が可逆な chart では、このイデアルは
$(H_x,H_y,H_z)$ となり、$D_H$ を得る。
$u=0$ の点では $g_H$ が滑らかなので、
$H_x,H_y,H_z$ の少なくとも一つが可逆な近傍を取れる。
その近傍では最初の $r$ 組の方程式が全ての $u_a$ を消す。
最後の方程式は
\[
 xH_x+yH_y+zH_z=3H-\sum_a u_aH_{u_a}
\]
により $X$ 上で冗長であり、零スキームは $E_H$ となる。
$C_H\cap\PP(K)=\varnothing$ は $g_H$ の滑らかさから従うため、
この二つの閉スキームは交わらない。
二種類の chart は零スキーム全体を覆い、式
\eqref{ppz:eq:zero-decomposition}を得る。

三つの active 偏微分は、定理\ref{thm:relative-point-nfold}の
濾過 Koszul 計算と同様に正則列である。
その商は中心多項式環上階数八で有限自由となる。
非零因子 $H$ を加えると完全交叉型・次元・次数が従う。
相対版の局所計算では、幾何ファイバー上で非零な偏微分を可逆にしている。
各 chart のイデアル等式は係数環上の等式なので、
非平坦・非被約のものを含む基底変更後にも成立する。
\end{proof}

\subsection{同じ埋込みを保つ一致の不存在}

量子側には第\ref{sec:points}節と第\ref{sec:poisson-morphism}節の
良い一ブロック中心商 $A=B/(w)$ を取り、
\[
 \operatorname{Pt}(A)=E\amalg D,\qquad
 Q=I(D)_2=\langle F_x,F_y,F_z\rangle,
 \qquad C=V(Q)
\]
とする。
特に $w\notin\Sym^3U$ とし、境界 Jacobi 線への中心評価
\[
 \epsilon:Z(B)_3/\Sym^3U\longrightarrow
       \bigl(k[x,y,z]/(g_x,g_y,g_z)\bigr)_3
\]
が同型であると仮定する。
この写像は中心三次空間全体の同型ではなく、商線束間の同型である。

\begin{corollary}[元の Poisson tensor を保つ切断の障害]
\label{ppz:cor:fixed-tensor}
元の $\Pi=h\Pi_F$ とその support を固定する。
滑らかな古典 Casimir 三次 $H\in kF\oplus\Sym^3U$ に対し、
$D_H=D$ をスキームとして実現することはできない。
従って trace--Hessian の中心同定だけでは、
定常点成分を Poisson 零点の定常成分と同定できない。
\end{corollary}
\begin{proof}
滑らかな $H$ では $H=cF+G(u)$ の係数 $c$ は非零である。
従って $C_H=C$ であり、$D_H=C\cap V(H)$ となる。
命題\ref{prop:point-casimir-restriction-obstruction}により、
この Cartier 切断は $D=C\cap V(w_{\mathrm{ab}})$ と一致しない。
\end{proof}

\begin{proposition}[良い一ブロック型の埋込み一致に対する障害]
\label{ppz:prop:embedded-obstruction}
上の仮定に加え、$g=F|_{u=0}$ が
式\eqref{eq:point-hessian-scalar}の Hessian スカラー条件を満たすとする。
滑らかな一ブロック Poisson 対 $(X_H,\pi_H)$ について、
active 平面三次が滑らかであり、$D_H$ が上の型 $(2,2,2,3)$ の
完全交叉であると仮定する。
このとき、射影線型同型によって
\[
 \operatorname{Pt}(A)\simeq Z_{X_H}(\pi_H)
\]
を埋込みスキームとして実現することはできない。
ここでは標的の Poisson tensor を元の $\Pi$ に固定していない。
\end{proposition}

\begin{proof}
仮にこの射影同型があれば、標的のデータを引き戻して
同じ $\PP(V^*)$ の中で比較する。
$n=3$ では楕円成分は唯一の正次元成分、$n>4$ では唯一の一次元成分である。
$n=4$ では $D,D_H$ は連結な次数 $24$ の完全交叉曲線であり、
次数 $3$ の連結成分 $E,E_H$ と区別できる。
従って射影同型は $E$ を $E_H$ に、$D$ を $D_H$ に移す。
$E$ と $E_H$ の線型包はそれぞれの support 平面なので、
引戻し後の support は元の $K$ と一致する。

完全交叉イデアルの二次部分から
\[
 \langle H_x,H_y,H_z\rangle=I(D_H)_2=I(D)_2=Q
\]
を得る。
命題\ref{prop:point-casimir}の Hessian スカラー条件による積分の一意性から
$H=cF+G(u)$ である。
よって前系の切断障害に反する。
\end{proof}

\begin{remark}[第二予想との論理的な関係]
この命題が排除するのは、良い一ブロック標的について、
既存の次数一埋込みを射影同型まで保つ強い一致である。
点スキームと零スキームの抽象同型、
別の次数一表示や偏極を用いた比較、
複数ブロックまたは退化した Poisson 標的を含む存在問題を
この命題だけで否定したことにはならない。
特に楕円成分の一致と、定常成分を含む全零スキームの一致は別である。
命題\ref{prop:nfm-correction-point-cubic}が与える Hessian 補正の整合は前者であり、
定理\ref{thm:nfm-residual-formal-poisson}の残余比較も
全定常点スキームとの一致を結論していない。
\end{remark}

\subsection{自然な残余 Poisson 対では抽象スキームも異なり得る}

次の結果は、既に構成した特定の残余 Poisson 対が
第二予想\ref{conj:nc-second}を自動的に実現するという推論を排除する。
別の Poisson tensor の存在を排除する結果ではない。
被約な台では消える差を、定常点の局所 Artin 環で検出する。

\begin{proposition}[具体的な分離型 Hesse 対の定常局所環の不一致]
\label{ppz:prop:split-artin-obstruction}
命題\ref{esi:prop:factor}の複素 Sklyanin 因子と指定中心三次式を
$S_\star,\Omega_\star$ と書く。すなわち
\begin{equation}\label{ppz:eq:explicit-pair}
 \begin{split}
 S_\star&=\C\langle x,y,z\rangle/
 (yz-3zy-2x^2,\ zx-3xz-2y^2,\ xy-3yx-2z^2),\\
 \Omega_\star&=\frac{79}{144}(x^3+y^3+z^3)
       +\frac5{96}(xyz+yzx+zxy+yxz+xzy+zyx),\\
 A_\star&=S_\star[u,v]/(\Omega_\star+G),\qquad
 G=(u^3+v^3)/3.
 \end{split}
\end{equation}
$u,v$ は中心変数である。
\[
 g_t=(x^3+y^3+z^3)/3+txyz,\qquad
 \lambda=17/9,\qquad X_\star=V(g_\lambda+G)\subset\PP^4_\C
\]
と置く。
この滑らかな可換三次の自然な一ブロック tensor を
\[
 \pi_\star
 =\left.\iota_{d(g_\lambda+G)}
       (\partial_x\wedge\partial_y\wedge\partial_z)
   \right|_{X_\star}
\]
とする。その非零スカラー倍でも以下の結論は同じである。
このとき
\begin{equation}\label{ppz:eq:split-artin-nonisom}
 \operatorname{Pt}(A_\star)\not\simeq Z_{X_\star}(\pi_\star)
\end{equation}
が抽象スキームとして成立する。
一方、両者の被約部分は、元の $\PP^4_\C$ の中でも一致する。
さらに実際の dg Morita 同値
$\Ku_{\mathrm{coh}}(A_\star)\simeq\Ku(X_\star)$ が存在する。
\end{proposition}

\begin{proof}
$J_t=\C[x,y,z]/(x^2+tyz,y^2+tzx,z^2+txy)$ と置く。
$1+t^3\ne0$ ならば $J_t$ は Hilbert 級数 $(1+s)^3$ の
長さ八の局所 Artin 環であり、極大イデアルの四乗は零である。
三つの非可換二次関係の可換化は、非零定数倍を除いて
$x^2+yz,y^2+zx,z^2+xy$ である。
$J_1$ では $x^3=y^3=z^3=-xyz$ なので、指定した中心三次式の像は
\[
 \eta=[\Omega_\star]
   =\left(-\frac{79}{48}+\frac5{16}\right)xyz
   =-\frac43xyz\ne0.
\]
従って境界中心評価の非消滅条件も満たす。
定理\ref{csp:thm:point-functor}により、点スキームは
$E_\lambda\amalg D_\star$ であり、その定常成分は
$D_\star=\operatorname{Proj}(J_1[u,v]/(G+\eta))$ である。
ここで $E_\lambda=V(g_\lambda)$ は命題\ref{esi:prop:factor}で確認した点曲線である。
定常成分の被約部分は
$x=y=z=0$ と $u^3+v^3=0$ で与えられる三点である。
全ての点は $u\ne0$ に属するので、$u=1$ とする。
各 $a\in\C$、$a^3=-1$ に対して
\[
 v=a-\eta/a^2
\]
は $J_1$ の中で $(1+v^3)/3+\eta=0$ を厳密に満たす。
これは $\eta^2=0$ による。
三つの根の差は単元なので、可換環上の中国剰余定理から
\begin{equation}\label{ppz:eq:source-local-artin}
 D_\star\simeq\coprod_{a^3=-1}\operatorname{Spec}J_1
\end{equation}
を得る。従って中心三次の非零 socle 項を捨てたわけではなく、
その項を中心座標の冪零補正に吸収して局所環を計算した。

命題\ref{ppz:prop:zero-scheme}を $(X_\star,\pi_\star)$ へ適用すると、
その定常成分は
\[
 D_{\mathrm{cl}}=V((g_\lambda)_x,(g_\lambda)_y,(g_\lambda)_z,g_\lambda+G)
 \simeq\coprod_{a^3=-1}\operatorname{Spec}J_\lambda
\]
である。最後の式には Euler 恒等式による $[g_\lambda]=0$ in $J_\lambda$
を用いた。

仮に $J_1\simeq J_\lambda$ が局所 $\C$-代数として成立すれば、
極大イデアルの冪に関する associated graded 環も同型になる。
両環は標準次数付きで斉次二次イデアルによる商なので、
この associated graded 環はそれぞれ $J_1,J_\lambda$ 自身である。
次数一と二の乗法から二つの勾配空間の線型同型を得て、
命題\ref{prop:ej-projective-uniqueness}により
$V(g_1)\simeq V(g_\lambda)$ を得るはずである。
しかし Hesse の不変量の公式は厳密に
\begin{equation}\label{ppz:eq:explicit-j}
 \begin{split}
 j(t)&=\frac{27t^3(8-t^3)^3}{(1+t^3)^3},\qquad j(1)=\frac{9261}{8},\\
 j(17/9)&=\frac{3813232609367}{4849120870776},\qquad
 j(17/9)-j(1)=-\frac{1402412578855675}{1212280217694}\ne0
 \end{split}
\end{equation}
を与える。従ってこの局所代数同型は存在しない。
点スキームの一次元成分は楕円曲線一つだけなので、
全スキームの抽象同型があれば零次元成分同士が対応する。
三つ全ての局所環が上の非同型な二種類であるため、
式\eqref{ppz:eq:split-artin-nonisom}が従う。

両側の楕円成分は同じ埋込み平面三次 $E_\lambda$ であり、
定常成分の被約部分も上の同じ三点である。
これで被約部分の一致も得る。
最後に、命題\ref{esi:prop:factor}は実楕円因子の正則性と中心商
$S_\star/(\Omega_\star)=B(E_\lambda,\OO_{E_\lambda}(1),t_p)$ を与える。
系\ref{csc:cor:fermat}を $a=1$、$c=d=1/3$、$f=g_\lambda$ に適用すると、
複素数体上で直接 $\Ku_{\mathrm{coh}}(A_\star)\simeq\Ku(X_\star)$ を得る。
\end{proof}

\begin{remark}[自然な処方への反例と存在問題の区別]
この例は、自然な残余 support と tensor をそのまま使う処方に対する反例である。
$X_\star$ 上で別の support を持つ Poisson tensor を選ぶ可能性や、
同じ Kuznetsov 成分を持つ別の可換三次を選ぶ可能性は扱っていない。
従ってこの命題だけでは「ある可換三次とある Poisson tensor が同時に存在する」という
第二予想\ref{conj:nc-second}を否定したことにはならない。
次小節で任意の support を調べ、この例についてその可能性も排除する。
一方、被約な台や楕円成分だけの比較では失われる情報が、
この例では長さ八の局所環に残っている。
一般パラメータにおける Hesse の有理恒等式は
第\ref{sec:nfm-residual-th-obstruction}節に記載した。
ここで用いた有理数の厳密検算は
\texttt{computations/verify\_point\_poisson\_obstruction.py}に含まれる。
局所環の分解と非同型の証明は上の議論であり、
その有限多項式検算と同一視しない。
\end{remark}

\subsection{三次三様体における任意の Poisson tensor の排除}

三次三様体では、全ての非零 Poisson tensor が分解可能な定数三ベクトルから来る。
これは Mayanskiy\cite[Theorem~1, \S2.2]{Mayanskiy}の既知結果である。
同論文の Schouten 括弧の式では、Pl\"ucker 関係式が五つの偏微分の係数となる。
滑らかな三次の偏微分は一次独立なので、Poisson 条件は全 Pl\"ucker 関係に等しい。
従ってこの結論は係数の特殊な三次にも適用でき、
同じ多項式計算によって標数零の代数閉体上でも成立する。
以下ではこの分類と、孤立零点の局所環を組み合わせる。

\begin{lemma}[埋込み次元三の孤立 Poisson 零点]
\label{ppz:lem:isolated-artin-eckardt}
$k$ を標数零の代数閉体、$H\in k[W]_3$、$\dim W=5$ とし、
$X=V(H)\subset\PP(W)$ を滑らかな三次三様体とする。
$0\ne\pi\in H^0(X,\bigwedge^2T_X)$ を Poisson tensor とし、
$p\in Z_X(\pi)$ を、局所環が長さ八・埋込み次元三を持つ孤立点とする。
$p$ の非零ベクトル代表を同じ文字で表す。
このとき $B_p=D^2H(p)$ の階数は二である。
さらに
\[
 R_p=\ker B_p,\qquad G_p=H|_{R_p}
\]
と置くと、$\dim R_p=3$、$V(G_p)\subset\PP(R_p)$ は滑らかな平面三次であり、
局所環の極大イデアル $\mathfrak m$ に対して
\begin{equation}\label{ppz:eq:eckardt-graded-local}
 \operatorname{gr}_{\mathfrak m}\OO_{Z_X(\pi),p}
 \simeq\operatorname{Jac}(G_p)
\end{equation}
となる。右辺は定数倍および線型座標変換を除いて $X,p$ だけから定まる。
\end{lemma}

\begin{proof}
上の分類により、$\pi=(\iota_{dH}P)|_X$、
$0\ne P\in\det K$、$\dim K=3$ と書ける。
$X\cap\PP(K)$ は Poisson 零集合に含まれる。
実際、この平面上では Euler vector と bivector の外積は
$3H P$ であり、$X$ 上で零になる。
この交叉は平面曲線、または平面全体なので、
孤立点 $p$ は $\PP(K)$ に属さない。

$p$ で非零な中心一次式を一にした affine chart を取る。
零点条件は $H_x=H_y=H_z=0$ であるから、
三つの $K$ 方向は $p$ で $X$ の接方向である。
滑らかさにより残りの affine 座標 $s$ を $H_s(p)\ne0$ となるよう選べ、
$x,y,z$ は $X$ の局所座標となる。
零イデアルはその正則局所環の中で $(H_x,H_y,H_z)$ である。
商の埋込み次元が三であることは、これらの元に一次項がないことを意味する。
従って $B_p|_{K\times K}=0$ である。

affine cone の接超平面を $T=\ker dH(p)\subset W$ とすると、
$T=K\oplus kp$ である。
Euler の式から $B_p(p,v)=2dH(p)(v)$、$B_p(p,p)=6H(p)=0$ なので、
$B_p|_{T\times T}=0$ となる。
余次元一の部分空間上で零となる対称形式の階数は高々二であり、
実際には二である。階数一なら $B_p=c\ell^2$ と書け、
$B_p(p,p)=0$ から $\ell(p)=0$、従って
$B_p(p,-)=0$ となって $dH(p)\ne0$ に反する。
従って $R_p$ は三次元で、$R_p\subset T$、$p\notin R_p$、
\[
 T=R_p\oplus kp
\]
を得る。

$B_p|_{T\times T}=0$ と $H(p)=0$ により、$H|_T$ は $p$ 方向に依らず
$H(v+tp)=H(v)$ for $v\in R_p$ である。
$G_p$ が特異と仮定し、$0\ne a\in R_p$ をその特異点の代表とする。
全ての $t$ について $a+tp\in V(H)$ であり、
その点での $T$ 方向の微分は全て零である。
$n\notin T$ を取ると、残る微分
$dH(a+tp)(n)$ は $t$ の二次式で、その最高次係数は
$dH(p)(n)\ne0$ である。
代数閉体上でその零点を選ぶと $X$ の特異点となり、矛盾する。
よって $G_p$ は滑らかである。

$K$ は $T=R_p\oplus kp$ の $kp$ に対する補空間だから、
$R_p$ の $p$ 方向への graph であり、$H|_K$ は $G_p$ と線型同型である。
$p$ を原点とする上の affine chart では、方程式を
\[
 H(x,y,z,s)=G_p(x,y,z)+sA(x,y,z,s),\qquad A(0)\ne0
\]
と書ける。形式的に $s$ を解くと $s\in(x,y,z)^3$ であり、
三つの零イデアルの生成元の二次先頭項は
$(G_p)_x,(G_p)_y,(G_p)_z$ である。
従って $\operatorname{Jac}(G_p)$ から
$\operatorname{gr}_{\mathfrak m}\OO_{Z_X(\pi),p}$ への次数付き全射がある。
前者は滑らかな平面三次の勾配正則列による長さ八の環であり、
後者も仮定によって長さ八である。
従ってこの全射は同型で、式\eqref{ppz:eq:eckardt-graded-local}を得る。
\end{proof}

\begin{lemma}[分離型 Hesse 三様体の階数二 Hessian 点]
\label{ppz:lem:hesse-eckardt}
$\lambda\ne0$、$\lambda^3\ne8,-1$ とし、
\[
 H_\lambda=g_\lambda+(u^3+v^3)/3
\]
と置く。
$V(H_\lambda)$ 上で Hessian 行列の階数が高々二となる点は、
中心直線上の三点と active 平面上の九つの flex、合計十二点だけである。
前補題の平面三次 $G_p$ の $j$ 不変量は、
中心三点では $j(g_\lambda)$、九つの flex では $0$ である。
\end{lemma}

\begin{proof}
Hessian 行列は
\[
 \operatorname{diag}\bigl(M_\lambda(x,y,z),2u,2v\bigr),\qquad
 M_\lambda=\begin{pmatrix}
 2x&\lambda z&\lambda y\\
 \lambda z&2y&\lambda x\\
 \lambda y&\lambda x&2z
 \end{pmatrix}
\]
である。$(x,y,z)\ne0$ では $M_\lambda$ の階数は二以上である。
実際、階数が一以下で一座標が零なら、主小行列式と $\lambda\ne0$ により
残りの二座標も零となる。
全て非零なら、非主小行列式から $2x^2=\lambda yz$、
主小行列式から $4yz=\lambda^2x^2$ となり、$\lambda^3=8$ に反する。

従って active 座標が非零な階数二点では $u=v=0$ である。
$g_\lambda=0$ 上では
\[
 \det M_\lambda=8(1+\lambda^3)xyz
\]
なので、該当点は一座標が零となる九点、すなわち九つの flex である。
active 座標が零なら $u^3+v^3=0$ であり、
$u,v$ はともに非零で、中心直線上の三点を得る。

中心三点では Hessian の核は active 三平面であり、$G_p=g_\lambda$ である。
flex は座標の置換により
$p=(1,-\epsilon,0,0,0)$、$\epsilon^3=1$ としてよい。
このとき Hessian の核は $e_u,e_v$ と
\[
 r=(\lambda\epsilon/2,\lambda\epsilon^2/2,1,0,0)
\]
で張られ、その制限は
\[
 H_\lambda(sr+u e_u+v e_v)
 =\bigl((1+\lambda^3)s^3+u^3+v^3\bigr)/3
\]
となる。これは滑らかな Fermat 平面三次と同型で、$j=0$ である。
\end{proof}

\begin{theorem}[分離型三次三様体における全 Poisson 零スキームの障害]
\label{ppz:thm:all-poisson-obstruction}
$k$ を標数零の代数閉体とし、
\[
 \lambda\ne0,\qquad \lambda^3\ne8,-1,\qquad
 1+t^3\ne0,\qquad j(t)\notin\{0,j(\lambda)\}
\]
とする。
$X_\lambda=V(g_\lambda+(u^3+v^3)/3)$ 上のどの Poisson tensor の
零スキームも、局所環が
$J_t=k[x,y,z]/(x^2+tyz,y^2+tzx,z^2+txy)$ と同型な孤立点を持たない。
従って、そのような孤立点を持つ任意のスキームは、
どの $Z_{X_\lambda}(\pi)$ とも同型でない。

特に命題\ref{ppz:prop:split-artin-obstruction}の具体的な複素代数について、
$X_\star$ 上のどの Poisson tensor $\pi$ に対しても
\[
 \operatorname{Pt}(A_\star)\not\simeq Z_{X_\star}(\pi)
\]
である。ここでは support の指定、active 制限の滑らかさ、
零スキームの完全交叉性を標的への追加仮定として課さない。
\end{theorem}

\begin{proof}
$\pi=0$ の零スキームは滑らかな三次元の $X_\lambda$ 全体であり、孤立点を持たない。
$\pi\ne0$ の場合、仮に $J_t$ と同型な局所環を持つ孤立零点があれば、
その長さは八、埋込み次元は三である。
補題\ref{ppz:lem:isolated-artin-eckardt}と
補題\ref{ppz:lem:hesse-eckardt}により、その associated graded 環は
$J_\lambda$ または Fermat 三次の Jacobian 環になる。
$J_t$ は標準次数付きなので、自身の極大イデアル filtration の
associated graded 環に等しい。
よって命題\ref{prop:ej-projective-uniqueness}から
$j(t)=j(\lambda)$ または $j(t)=0$ が必要になり、仮定に反する。

具体例では $t=1$、$\lambda=17/9$ である。
式\eqref{ppz:eq:explicit-j}に加えて
\[
 \lambda^3-8=-\frac{919}{729}\ne0,\qquad
 \lambda^3+1=\frac{5642}{729}\ne0
\]
だから全ての仮定を満たす。
式\eqref{ppz:eq:source-local-artin}の三つの孤立点により、最後の結論を得る。
\end{proof}

この判定は固定された有理パラメータの例だけに限らない。
第\ref{sec:nfm-residual-th-obstruction}節の Hesse 有理恒等式を用いれば、
同節の幾何学的一般点でも同じ $j$ の非一致条件が成立する。
この一般パラメータでの計算と、以下の複素数体上の直接構成を区別する。

\subsection{具体的な複素非可換三次三様体による反例}

\begin{corollary}[全零スキーム版の第二予想に対する複素数体上の反例]
\label{ppz:cor:complex-counterexample}
式\eqref{ppz:eq:explicit-pair}の $A_\star$ は、複素数体上の
点スキームの良い一ブロック条件を満たす滑らかな楕円一ブロック型の
非可換三次三様体であり、次の二つを同時に満たす。
\begin{enumerate}
 \item 滑らかな可換三次三様体
 \[
 X_\star=V\!\left(
 \frac{x^3+y^3+z^3+u^3+v^3}{3}+\frac{17}{9}xyz\right)
 \subset\PP^4_\C
 \]
 と実際の dg Morita 同値
 $\Ku_{\mathrm{coh}}(A_\star)\simeq\Ku(X_\star)$ が存在する。
 \item $\Ku_{\mathrm{coh}}(A_\star)\simeq\Ku(Y)$ を満たす
 どの滑らかな可換三次三様体 $Y$ と、
 その上のどの Poisson tensor $\pi$ に対しても、
 $\operatorname{Pt}(A_\star)\not\simeq Z_Y(\pi)$ が抽象スキームとして成立する。
\end{enumerate}
従って、非被約構造を含む全 point scheme と全 Poisson 零スキームの同型を要求する
第二予想\ref{conj:nc-second}は成立しない。
主予想の可換 Kuznetsov 実現はこの例では成立している。
\end{corollary}

\begin{proof}
系\ref{esi:cor:ambient}は、指定した複素代数の AS 正則な周囲環と
中心正則三次商を与える。
命題\ref{ppz:prop:split-artin-obstruction}の証明で、
系\ref{csc:cor:fermat}による全 coherent 残余圏の比較を構成した。
従って第一項が成立し、同じ比較の smoothness と coherent 飽和により
$A_\star$ は滑らかな非可換三次三様体である。
active 可換化 $g_1$、点曲線 $g_{17/9}$ はともに滑らかであり、
中心三次の socle 評価 $-4xyz/3$ は非零である。
定常三次は $J_1[u,v]$ 上の非零因子であるから、
点スキームの良い一ブロック条件も満たす。
ここでいう良い条件は点関手の条件であり、
全 $\OO$-表示 stack の全ての相対入力への所属を同時に主張するものではない。
それらの入力は本系の反例の証明には用いない。

$X_\star$ 上の全ての Poisson tensor は
定理\ref{ppz:thm:all-poisson-obstruction}により排除される。
別の標的 $Y$ に対しては、第一項の同値を合成して
$\Ku(Y)\simeq\Ku(X_\star)$ を得た後、
三次三様体の categorical Torelli
\cite[Theorem~1.1]{NFMThreefoldCategoricalTorelli}から $Y\simeq X_\star$ を得る。
$Y$ 上の Poisson tensor をこの同型で $X_\star$ へ移せば、
その零スキームも移送されるため、第二項が従う。
\end{proof}

\begin{remark}[反例が排除する範囲と残る問題]
この反例は $n=3$ の非被約な定常点を持つ分離型であり、
被約 point scheme を持つ一般混合族の一致や、被約化後の一致を否定しない。
実際、命題\ref{ppz:prop:split-artin-obstruction}では
自然な残余 Poisson 対と被約部分が一致する。
従って被約化を取った命題、滑らかな定常成分を仮定した命題、
または別の点関手による命題は、別の検証問題として残る。
これらのいずれかを元の第二予想の代わりに自動採用するものではない。
Hessian と制限三次の恒等式は
\texttt{computations/verify\_point\_poisson\_obstruction.py}で検算する。
全 support の排除、局所環、Ku 比較、Torelli の適用は本文で証明し、
この計算で代用しない。
\end{remark}


\section{残余圏の形式変形と次元に依存する算術}
\label{sec:nf-residual}

本節では $k=\mathbb C$、$n\geq3$、$N=n+2$ とし、
$X=V(F)\subset\mathbb P^{N-1}$ を滑らかな三次超曲面とする。
残余圏と rotation の規約は
\[
 \mathcal K_X=\operatorname{Ku}(X)
   =\langle\mathcal O_X,\ldots,\mathcal O_X(n-2)\rangle^\perp,
 \qquad
 \mathsf R_X=L_{\mathcal O_X}\circ(-\otimes\mathcal O_X(1))
\]
である。dg enhancement と冪等完備化を伴う圏を用い、
圏の変形を考える際には
$\mathcal D_X=\operatorname{Ind}(\mathcal K_X)$ を用いる。
ここで扱うのは、通常の可換局所 Artin $k$-代数を基底とする変形である。
全ての変形には特殊ファイバーと $X$、または $\mathcal D_X$ との同定を付ける。

\subsection{Hochschild 計算と幾何的 Kodaira--Spencer 写像}

次数付き Jacobian 環を
\[
 J_F=k[x_1,\ldots,x_N]/(F_{x_1},\ldots,F_{x_N})
\]
と置く。負の次数成分は零と解釈する。
次の計算は matrix factorization の既知の sector 公式の三次の場合であり、
本稿の新しい Hochschild 計算とは数えない。

\begin{proposition}[三次超曲面の Hochschild 次数]
\label{prop:nf-hh}
ベクトル空間として
\begin{equation}
\label{eq:nf-hh}
 HH^q(\mathcal K_X)\simeq
 \begin{cases}
   (J_F)_{3q/2},&q\text{ が偶数},\\
   0,&q\text{ が奇数}
 \end{cases}
 \ \oplus\
 \begin{cases}
   k^2,&3\mid N,\ q=N/3,\\
   0,&\text{それ以外}
 \end{cases}
\end{equation}
である。特に $HH^q=0$ for $q<0$、$HH^0=k$、$HH^1=0$ であり、
\[
 \dim HH^2(\mathcal K_X)=
 \begin{cases}
  22,&n=4,\\
  \binom{n+2}{3},&n\neq4.
 \end{cases}
\]
幾何的変形から誘導される実際の写像
\begin{equation}
\label{eq:nf-geometric-ks}
 \eta_X:H^1(X,T_X)\longrightarrow HH^2(\mathcal K_X)
\end{equation}
は単射である。$n\neq4$ では同型、$n=4$ では余核は二次元である。
\end{proposition}

\begin{proof}
滑らかさと Euler の公式により、$F$ の affine cone の特異点は原点だけである。
graded matrix factorization による dg 比較と
\cite[Theorem~5.39 の一般式とその証明]{NFBFKKernels}、
\cite[Theorem~3.5]{NFLZSerre} を使う。
後者の表示式の parity は、cohomological degree $q$ から conormal rank を引いたものと読む。
前者の一般式では Koszul コホモロジーの偶奇が区別され、
孤立特異点の場合には次数零の Jacobian 成分だけが残る。
恒等 sector の $q$ 次成分は、$q$ が偶数なら $(J_F)_{3q/2}$、
奇数なら零である。
非恒等な $\mu_3$ の二つの元は、それぞれ原点だけを固定する。
その conormal rank は $N$、determinant の重みは $-N$ なので、
不変部分が現れる条件は
\[
 q\equiv N\pmod2,\qquad
 N+\frac32(q-N)=0
\]
となる。これは $3\mid N$ かつ $q=N/3$ と同値である。
この場合には各元から一次元が寄与し、
\eqref{eq:nf-hh} を得る。

偏微分は $N$ 個の二次式からなる正則列であるから
\[
 \operatorname{Hilb}(J_F,t)=(1+t)^N,\qquad
 \dim(J_F)_3=\binom N3.
\]
Euler 列と正規列により
$H^1(X,T_X)\simeq(J_F)_3$ である。
写像 \eqref{eq:nf-geometric-ks} の単射性には、
単なるベクトル空間の同型でなく、
\cite[Lemma~5.2]{NFPirozhkov} の幾何的写像を用いる。
同補題の入力は、直交部分が例外列で生成されることと
$X$ の Hodge diamond が対角的でないことである。
前者は上の半直交分解、後者は同論文 Lemma~5.1 により $n\geq3$ の三次の場合に成立する。
次元を比較すれば結論を得る。
四次元の場合の余核二次元は
\cite[Remark~3.8 and the proof of Proposition~3.5]{NFHuybrechtsRennemo}
にも明記されている。
\end{proof}

\begin{lemma}[三次超曲面の通常の変形空間]
\label{lem:nf-geometric-smooth}
特殊ファイバーを固定した $X$ の通常の変形群oidは、
次元 $\binom N3$ の滑らかな形式円板に同値である。
各 Artin 基底上の自己同型で特殊ファイバー上恒等なものは恒等だけである。
\end{lemma}

\begin{proof}
超曲面の制限列から、$0<i<n$ と任意の整数 $a$ について
$H^i(X,\mathcal O_X(a))=0$ である。
また $X$ は Fano なので $H^n(X,\mathcal O_X)=0$ である。
Euler 列と
\[
 0\longrightarrow T_X\longrightarrow
 T_{\mathbb P^{N-1}}|_X\longrightarrow\mathcal O_X(3)
 \longrightarrow0
\]
から $H^2(X,T_X)=0$ を得る。
線型ベクトル場が $X$ に接していれば
$\sum_i\ell_iF_{x_i}=cF$ であり、Euler の公式により
\[
 \sum_i\left(\ell_i-\frac c3x_i\right)F_{x_i}=0
\]
となる。正則な二次式列の Koszul 複体には線型係数の一次 syzygy がない。
従ってこのベクトル場は scalar であり、射影空間上では零である。
よって $H^0(X,T_X)=0$ である。
標準的な小拡大の変形・障害理論から、変形関手は滑らかで、
形式自己同型群は自明となる。
接空間は前命題の $(J_F)_3$ であり、次元は $\binom N3$ である。

この変形を埋込み付き三次として扱うこともできる。
$H^1(X,\mathcal O_X)=H^2(X,\mathcal O_X)=0$ によって
$\mathcal O_X(1)$ は同型類として小拡大へ一意に延長し、
$H^1(X,\mathcal O_X(1))=0$ によってその切断を持ち上げられる。
従って局所 Artin 基底上で射影空間への埋込みが延長し、
方程式は再び三次の一式となる。
\end{proof}

\subsection{滑らかな幾何的族からの形式認識}

次の補題は、接空間の次元一致だけから全形式変形を結論するものではない。
全ての幾何的変形を送る実際の自然変換と、
その出発点の形式的滑らかさを同時に使う。

\begin{lemma}[滑らかな源と接写像同型]
\label{lem:nf-smooth-source}
$k$ を標数零の体、$\mathfrak g$ を dg Lie 代数とし、
$H^0(\mathfrak g)=0$、$V=H^1(\mathfrak g)$ と
$H^2(\mathfrak g)$ は有限次元とする。
その形式 moduli problem を通常の局所 Artin $k$-代数へ制限し、
一次切断した群oidを $\tau_{\leq1}\mathcal F_{\mathfrak g}$ と書く。
$d=\dim V$ とし、自然変換
\[
 \phi:\operatorname{Spf}k[[t_1,\ldots,t_d]]
       \longrightarrow\tau_{\leq1}\mathcal F_{\mathfrak g}
\]
の接写像が同型なら、$\phi$ は同値である。
$H^2(\mathfrak g)=0$ は仮定しない。
\end{lemma}

\begin{proof}
標数零での形式 moduli problem と dg Lie 代数の比較
\cite[Section~2]{NFDAGX} を用い、
$H=H^\bullet(\mathfrak g)$ 上へ極小 $L_\infty$ 構造
$\ell_j$ を移す。通常の Artin 環 $A$ に対する
Maurer--Cartan 元は $x\in V\otimes\mathfrak m_A$ であり、
方程式は
\begin{equation}
\label{eq:nf-mc}
 Q(x)=\sum_{j\geq2}\frac1{j!}\ell_j(x,\ldots,x)=0
       \quad\text{in }H^2(\mathfrak g)\otimes\mathfrak m_A
\end{equation}
である。$A$ ごとにこの和は有限である。
$H^0(\mathfrak g)=0$ なので gauge による異なる $x$ の同一視はなく、
一次切断の自己同型群も自明である。
従ってこの群oidは、
$k[[V^\vee]]$ を $Q$ の成分で割った完備局所環による離散群oidで表される。

$V$ の座標を選ぶと、$\phi$ は
$x=f(t)\in V\otimes(t_1,\ldots,t_d)k[[t_1,\ldots,t_d]]$
で表される。これは全 Artin 切断で両立する自然変換から得られる。
接写像の仮定により $f$ の線型項は可逆である。
形式逆写像定理から $f$ は形式座標変換である。
また $\phi$ が全ての次数で定義されるので、
$Q(f(t))=0$ は形式冪級数の恒等式である。
$f^{-1}$ を代入すれば $Q(x)=0$ が恒等的に成り立つ。
従って上の商環は $k[[V^\vee]]$ 自身であり、
$\phi$ は形式座標変換による同値である。
この議論は、高次の $L_\infty$ 演算全部が零であるとは主張しない。
\end{proof}

以下の $\operatorname{CatDef}^{c}_{\mathcal D_X}(A)$ は、
特殊ファイバーを $\mathcal D_X$ と同定した
compactly generated な $A$-線型圏の変形空間である。
$\tau_{\leq1}$ を取ると、射は圏同値を自然同型で割ったものとなる。
自然変換の scalar 自己同型まで保持する二群oidとの同一視ではない。

\begin{theorem}[四次元を除く残余圏の通常の Artin 形式認識]
\label{thm:nf-residual-formal}
$n\geq3$、$n\neq4$ とする。
滑らかな三次超曲面 $X$ について、自然な残余操作は
通常の局所 Artin $k$-代数上の群oidの同値
\begin{equation}
\label{eq:nf-formal-recognition}
 \operatorname{Def}_X
 \ \xrightarrow{\ \sim\ }\
 \tau_{\leq1}\operatorname{CatDef}^{c}_{\mathcal D_X}
\end{equation}
を与える。
両辺の特殊ファイバーは固定している。
同型類関手は次元 $\binom{n+2}{3}$ の形式円板である。
これは全 Abel 圏の変形を中心三次対で表示する定理ではない。
\end{theorem}

\begin{proof}
$X_A$ の相対半直交分解から得られる
$\operatorname{Ind}\operatorname{Ku}(X_A)$ は、
実際の compactly generated な圏の変形を与える。
相対射影・mutation の核は有限個の線束から作られ、
通常の Artin 係数変更と両立する。
線束の延長の自明化に含まれる基底の単元は、
射影核を作る $E^\vee\boxtimes E$ では相殺される。
従って上の自然変換はその選択に依存しない。
これで \eqref{eq:nf-formal-recognition} の自然変換を得る。

その微分は \eqref{eq:nf-geometric-ks} の写像である。
実際、幾何的な一次変形の Hochschild 類は
$H^1(T_X)\to HH^2(X)$ で与えられ、
相対射影核による変形の制限は $HH^2(X)\to HH^2(\mathcal K_X)$ に対応する。
両者の合成が、上で構成した相対半直交成分の一次変形を記述する。
これは幾何的変形・Hochschild 類・admissible 射影の自然性であり、
接空間の抽象的な次元一致を代用するものではない。

$\mathcal D_X$ は proper な compact 部分圏で生成されるので
tamely compactly generated である。
\cite[Theorems~5.1.9, 5.3.16, Proposition~5.3.21]{NFDAGX}
および \cite[Sections~3.3 and~4.3]{NFBKP} により、
形式 hull への自然変換
\[
 \theta:\operatorname{CatDef}^{c}_{\mathcal D_X}
       \longrightarrow\operatorname{CatDef}^{\wedge}_{\mathcal D_X}
\]
は $(-1)$-truncated である。
従って $\pi_0$ 上単射であり、
既に存在する対象における全ての高次 homotopy group 上同型である。
通常の可換 Artin 環へ制限した hull を支配する dg Lie 代数は
\[
 H^i(\mathfrak g)=HH^{i+1}(\mathcal K_X)
\]
を満たす。
命題 \ref{prop:nf-hh} により $H^0(\mathfrak g)=0$、
幾何的接写像は同型である。
補題 \ref{lem:nf-geometric-smooth} と
補題 \ref{lem:nf-smooth-source} により
\[
 \operatorname{Def}_X
 \xrightarrow{\sim}
 \tau_{\leq1}\operatorname{CatDef}^{\wedge}_{\mathcal D_X}.
\]
この全ての対象は幾何的変形から来るので、
$\theta$ の $\pi_0$ 上の像は hull の全体である。
$(-1)$-truncated 性による単射性と
$\pi_1$ の比較を合わせると主張を得る。
\end{proof}

\begin{remark}[七次元と高次自己同型]
\label{rem:nf-seven}
$n=7$ では $N=9$ であり、
$HH^3(\mathcal K_X)=k^2$ が実際に残る。
従って $HH^3=0$ を前提にする証明はこの次元に適用できない。
上の証明は滑らかな幾何的変形からの全次数の自然変換を使うため、
この場合も成り立つ。
また、一般に $HH^0=k$ に対応する高次 scalar は消えない。
例えば $A=k[\epsilon]/(\epsilon^2)$ では
$\pi_2$ に $k\epsilon$ が現れる。
定理は $\tau_{\leq1}$ と通常の基底上の主張であり、
未切断の高次群oidや導来 Artin 基底について幾何変形との同値を主張しない。
\end{remark}

\begin{proposition}[四次元の形式的障害]
\label{prop:nf-fourfold-obstruction}
$n=4$ のとき、残余圏の形式 hull の同型類関手は
二十二次元の滑らかな形式円板である。
幾何的変形の二十次元形式円板はその中の滑らかな余次元二の形式部分空間となる。
従って定理 \ref{thm:nf-residual-formal} の hull 版は四次元では成立しない。
\end{proposition}

\begin{proof}
\eqref{eq:nf-hh} より $HH^\bullet(\mathcal K_X)$ は次数 $0,2,4$ にのみ存在し、
その次元は $1,22,1$ である。
従って $\mathfrak g$ の極小模型は次数 $-1,1,3$ にのみ存在する。
$j$ 項演算の次数は $2-j$ であり、
奇数次数の $j$ 個の入力に対する出力は偶数次数になる。
極小模型には偶数次数成分がないので全ての演算は零である。
通常の Artin 基底上の $\pi_0$ は二十二次元形式円板となる。
幾何的写像の接写像は命題 \ref{prop:nf-hh} により単射である。
二十個の標的座標を選んで接写像を可逆にし、形式逆写像を使うと、
写像は残り二座標の形式グラフで表される。
標的の形式座標を変更すれば、このグラフは最後の二座標が零の部分空間となる。
\end{proof}

この余次元二は、指定した非可換中心三次族の残余圏が
非幾何的になることを示すものではない。
その族の形式写像が上の幾何的部分空間に入ることは、
中心性や matrix factorization の追加構造によって起こり得る。
また、ここで二十二次元と計算したのは明示した形式 hull である。

\subsection{rotation、Serre 安定性と三を法とする場合分け}

\begin{proposition}[rotation を回復する算術]
\label{prop:nf-rotation}
$S_X$ を $\mathcal K_X$ の Serre 自己同値とすると、
\begin{equation}
\label{eq:nf-serre-rotation}
 \mathsf R_X^3\simeq[2],\qquad
 S_X\simeq\mathsf R_X^{-(n-1)}[n]
          \simeq\mathsf R_X^{-N}[N]
\end{equation}
である。従って
\begin{equation}
\label{eq:nf-rotation-cases}
 \begin{array}{c|c|c}
 n&S_X&\mathsf R_X\\ \hline
 3q&\mathsf R_X[q]&S_X[-q]\\
 3q+1&[q+1]&\text{Serre と整数 shift の積では表せない}\\
 3q+2&\mathsf R_X^{-1}[q+2]&S_X^{-1}[q+2]
 \end{array}
\end{equation}
となる。中段の否定は、他の圏論的構造を使った rotation の復元まで否定するものではない。
\end{proposition}

\begin{proof}
\eqref{eq:nf-serre-rotation} は graded matrix factorization の周期性と
Serre 公式である
\cite[Section~3.3, Corollary~3.18, Lemma~3.19]{NFKuznetsovFCY},
\cite[proof of Proposition~3.7]{NFLZSerre}。
$n$ を三で割った余りに従って指数を整理すると表を得る。
中段では Serre と整数 shift が生成する群は整数 shift だけからなる。
もし $\mathsf R_X\simeq[a]$ なら $\mathsf R_X^3\simeq[2]$ から $[3a]\simeq[2]$ となる。
非零 proper 圏において shift は無限位数である。
実際、非零対象が非零 shift と同型なら、その自己 Ext は無限に多くの次数で非零となり、
proper 性に反する。従って $3a=2$ が必要となり、矛盾する。
\end{proof}

rotation の算術的回復と、それを用いる categorical Torelli は
\cite[Proposition~3.7, Theorems~4.1 and~5.1]{NFLZSerre}
の既知結果である。本稿で必要なのは、これを残余射の構成に使う際、
三次元で自動的だったどの条件を一般次元で記憶すべきかを明示することである。

\begin{proposition}[rotation に安定な像の生成]
\label{prop:nf-rotation-generation}
$\mathcal T\subset\mathcal K_X$ を非零の thick 部分圏とする。
$\mathsf R_X$ とその逆で安定なら $\mathcal T=\mathcal K_X$ である。
従って $n\not\equiv1\pmod3$ なら、
$S_X$ とその逆に安定な非零 thick 部分圏も全体に等しい。
特に、冪等完備な非零三角圏から $\mathcal K_X$ への完全忠実関手の像が
該当する自己同値で安定なら、その関手は同値である。
\end{proposition}

\begin{proof}
graded Orlov 比較で $\mathcal K_X$ を
$k[x_1,\ldots,x_N]/(F)$ の graded 特異圏と同定する。
任意の非零対象 $M$ に対し
$M\oplus\mathsf R_XM\oplus\mathsf R_X^2M$ は全体を生成する
\cite[Theorem~5.15]{NFBFKOrlov}。
これは次数三・孤立特異点の場合への既知定理の直接の適用である。
従って $\mathcal T$ は生成対象を含み、全体に等しい。
Serre に関する主張は \eqref{eq:nf-rotation-cases} に従う。
完全忠実関手については、冪等完備性により像が thick となるので同じ議論を使う。
像の Serre 安定性そのものはこの命題の仮定であり、
完全忠実性だけからは結論しない。
\end{proof}

\begin{remark}[四次元での Serre 安定性の不足]
\label{rem:nf-serre-insufficient}
滑らかな Pfaffian 三次四様体 $X$ には、
K3 曲面 $Y$ と同値 $\mathcal K_X\simeq D^b(Y)$ がある
\cite[Theorem~3.1]{NFKuznetsovCubicFourfolds}。
$y\in Y$ に支持された複体全体は、非零で真の thick 部分圏である。
これは $S=[2]$ に安定である。
従って四次元では「Serre に安定な非零 thick 部分圏は全体」という命題自体が偽である。
rotation 安定性を仮定する前命題は、この次元でも成立する。
\end{remark}

\subsection{既存の一般点復元と新しいモジュライ問題の境界}

研究稿 \cite{NFNCCubicB} の主定理は既に任意の $n\geq3$ を対象としている。
その入力は、$R=\mathbb C[h]_{(h)}$ 上の次数別有限自由な中心三次対
$(B_R,w_R)$ で、
特殊ファイバーが多項式環と滑らかな三次式になることである。
同稿の定理を入力すれば、一般点の全 coherent 残余圏に、
rotation と中心元からの指定三次周期を保つ可換模型が存在する。
方程式は $\mathbb C(h)$ 上、同値核はある有限拡大上に得られ、
幾何学的な射影同型類は一意である。
これは本研究で初めて得る一般次元化ではなく、
同稿の未公表の一般点復元定理を用いる結論として扱う。

以上から、一般次元のモジュライ研究では次を分ける必要がある。
第一に、線型偏極を記憶した中心三次対の族では rotation が与えられる。
第二に、構造対象 $\mathcal O$ だけを保つ全 Abel 同値へ移るには、
例外列と残余構成の独立性・相対降下を別に証明する必要がある。
第三に、可換模型の存在、模型の一意性、全 Artin 形式認識は異なる入力である。
$n\neq4$ では定理 \ref{thm:nf-residual-formal} が形式認識を与えるが、
$n\equiv1\pmod3$ では rotation の Serre による回復は依然として使えない。
$n=4$ では一般の圏変形の形式認識も成立せず、
指定された中心三次族の追加構造を使う必要がある。

定理\ref{thm:point-trace-hessian}は、局所一ブロック型の良い開上で
別の構成による Poisson 三次 $X_w$ を与える。
残余復元の三次 $X_{\mathrm{res}}$ と $X_w$ の一致、および
この同じ三次に対する残余圏同値の存在と降下は、なお比較すべき問題である。
形式変形の次元一致や零ファイバーでの一致だけでは、この比較は従わない。
本節から全 $\mathcal O$-marked Abel モジュライの成分性、
残余圏による大域射の相対次元・滑らかさ・全射性を結論しない。
そのためには、残余の相対構成、rotation と三次周期、同値核の有限表示と降下が必要である。

% !TEX root = ../nc-cubic.tex
\section{可換 Kuznetsov 実現を持つ幾何点の開性}
\label{grl:sec:geometric-locus}

主予想\ref{conj:nc-main}の「任意の対象」へ進むには、既知の一般点を増やすだけでなく、
可換模型の存在が族の中でどこまで保たれるかを調べる必要がある。
本節では、定理\ref{thm:nf-residual-formal}の通常の Artin 形式認識を、
有限表示の比較核のモジュライを通して代数的な開性へ移す。
四次元を除くという仮定と、実際の smooth proper な圏族という仮定を保つ。
この議論から境界での閉性や全内在的成分の被覆を結論しない。

\subsection{圏族と比較核の有限表示}

以下 $S$ は有限型 $\C$-スキームとする。
$S$ 上の圏族 $\mathcal C_S$ は、小さい冪等完備な $S$-線型 dg 圏の族で、
\'etale 局所的に
\[
 \mathcal C_S\simeq\operatorname{Perf}(\Gamma_S),\qquad
 \Gamma_S\in\operatorname{Perf}(S),\qquad
 \Gamma_S\in\operatorname{Perf}(\Gamma_S^{\mathrm{op}}
                      \otimes_S^{\mathbf L}\Gamma_S)
\]
と表示されるものとする。従って族は相対的に smooth proper である。
係数変更はこの表示の導来 tensor で定め、compact 対象とその Ind 完備化を
相対 Morita 理論で扱う。以下で核というときは perfect 双加群を意味する。
一般の Abel 圏族に、このような相対 dg 模型が自動的に存在するとは仮定しない。

$M_n\subset\PP(\operatorname{Sym}^3(\C^{n+2})^*)$ を
滑らかな埋込み三次超曲面の係数スキームとし、$\mathcal X\to M_n$ を普遍族とする。
相対例外列の射影核から $\operatorname{Ku}(\mathcal X/M_n)$ を構成する。
この圏族は smooth proper であり、\'etale 局所的な有限 dg 生成対象を持つ。

\begin{lemma}[可換比較を表示する核スタック]
\label{grl:lem:comparison-stack}
$S\times M_n$ 上に、対 $(X_T,P_T)$、すなわち
$M_n$ が与える三次族 $X_T$ と実際の Morita 同値
\[
 P_T:\mathcal C_T\xrightarrow{\sim}\operatorname{Ku}(X_T/T)
\]
をパラメータ付ける局所有限表示の Artin スタック $\mathcal E$ が存在する。
射は同じ埋込み三次族上の核の同型であり、核の scalar 自己同型は保持する。
幾何点の核の自己同型群は $\mathbf G_m$ である。
\end{lemma}

\begin{proof}
アファインな係数 chart 上で両側の smooth proper な有限 dg 模型を選ぶ。
二つの模型の導来 tensor 積も smooth proper であり、
To\"en--Vaqui\'e の対象モジュライ
\cite[Corollaries~3.15, 3.21, 3.22]{NFMToenVaquieModuli}
により、その perfect 双加群のうち負の自己 Ext がない部分は
局所有限表示の Artin スタックをなす。

perfect 双加群の左右随伴は相対双対で与えられる。
同値性は単位・余単位の二つの cone が零となる条件である。
これらの cone は底上 perfect なので、その零化 locus は開である
\cite[Lemma~3.25 and its proof]{NFMToenVaquieModuli}。
同補題は自己関手について記されているが、二つの dg 代数を用いる場合も、
各側の自由生成対象で単位・余単位を検査する同じ証明が適用される。
従って同値核の部分は開部分スタックである。

幾何点で同値核 $P$ を固定すれば、その自己 Hom は標的圏の
Hochschild cohomology である。命題\ref{prop:nf-hh}より
\[
 \operatorname{Ext}^{q}(P,P)=0\quad(q<0),\qquad
 \operatorname{End}(P)=\kappa .
\]
従って同値核は上の Artin 部分に属し、自己同型は $\kappa^\times$ となる。
核・随伴・単位・余単位と有限表示は相対 Morita 変更に自然であるため、
係数 chart 間で降下して $\mathcal E$ を与える。
ここでは核の同型類を先に選んで gerbe を中立化する操作は行わない。
\end{proof}

\subsection{形式認識から代数的な近傍への移行}

\begin{theorem}[四次元を除く可換実現 locus の開性]
\label{grl:thm:open-locus}
$n\ge3$、$n\ne4$ とし、上の smooth proper な族 $\mathcal C_S$ を取る。
幾何的なファイバーが滑らかな $n$ 次元三次超曲面の Kuznetsov 成分に
dg Morita 同値となる点の集合は、Zariski 開部分 $S_{\mathrm{cub}}\subset S$ である。
この開上では \emph{\'etale 局所的に}滑らかな三次族と実際の perfect 比較核が存在する。

より具体的に、$s\in S(\C)$ における
\[
 P_s:\mathcal C_s\xrightarrow{\sim}\operatorname{Ku}(X_s)
\]
を固定すると、$s$ の \'etale 近傍 $S'\to S$ 上で
$X_s$ と $P_s$ を、滑らかな三次族と Morita 同値へ延長できる。
この同値は全ての導来底変更に保たれる。
rotation、指定三次周期、構造対象および point scheme の標識は結論に含めない。
\end{theorem}

\begin{proof}
補題\ref{grl:lem:comparison-stack}の射 $q:\mathcal E\to S$ を考える。
まず $e=(X_s,P_s)$ が $\C$-点である場合に $q$ の形式的滑らかさを示す。
剰余体 $\C$ を持つ通常の局所 Artin 環の小拡大 $A'\twoheadrightarrow A$ と、
$\operatorname{Spec}A'\to S$、および $e$ を特殊ファイバーとする
$(X_A,P_A)\in\mathcal E(A)$ が与えられたとする。

特殊ファイバーの $P_s$ によって
$\operatorname{Ind}\mathcal C_{A'}$ を
$\operatorname{Ind}\operatorname{Ku}(X_s)$ の標識付き変形とみなす。
定理\ref{thm:nf-residual-formal}の本質的全射性から、
滑らかな三次の変形 $X_{A'}$ と、実際の圏同値の同型類
\[
 Q_{A'}:\mathcal C_{A'}\xrightarrow{\sim}\operatorname{Ku}(X_{A'})
\]
が得られる。compactly generated な圏の同値は compact 対象を保ち、
相対 Morita 理論によって perfect 双加群で表される。

ここでは存在だけでなく、既に指定された $P_A$ との整合が必要である。
$A$ 上では $P_AQ_{A'}^{-1}|_A$ が二つの幾何的変形の残余圏間の射を与える。
定理\ref{thm:nf-residual-formal}の\emph{群oidとしての完全忠実性}により、
この射は、特殊ファイバー上恒等な幾何的同型
$X_{A'}\times_{A'}A\simeq X_A$ によって実現される。
同定後には $Q_{A'}|_A\simeq P_A$ という核の実際の同型を選べる。
$\tau_{\le1}$ における射の等しさは、まさにこの自然同型の存在を意味する。
核の scalar 自己同型を自明と主張する必要はなく、
持上げに選んだ同型の非一意性も許す。

補題\ref{lem:nf-geometric-smooth}の最後の議論により、
$\OO_{X_A}(1)$ とその切断の基底を $A'$ へ持ち上げられる。
従って $X_{A'}$ を、指定された $M_n$ の $A$-点を持ち上げる
埋込み三次として選べる。これで $(X_A,P_A)$ の持上げを得た。

有限型の smooth atlas $U\to\mathcal E$ と $e$ の上の $\C$-点 $u$ を取る。
上の持上げと atlas の形式的滑らかさを合成すると、$U\to S$ は
$u$ において Artin 小拡大の持上げ条件を満たす。
局所有限表示と局所的な滑らかさの判定
\cite[Tag~02HX]{NFMStacksSmoothLocal}
より、$U\to S$ は $u$ で滑らかである。
この議論は $\mathcal E$ の全ての $\C$-点に適用できる。
有限型 $\C$-スキームは Jacobson であり、有限型 atlas の非滑らか locus が
非空なら $\C$-点を含むので、$q$ は全体で滑らかとなる。

従って $q$ の像は開であり、定義からその像が $S_{\mathrm{cub}}$ である。
smooth atlas を通し、滑らかな射の \'etale 局所切断
\cite[Tag~055U]{NFMStacksSmoothLocal}
を取れば、この開の \'etale 被覆上に実際の三次族と核が得られる。
固定した $\C$-点 $u$ を通る切断も取れる。
実際、$u$ の近傍の標準的な滑らか表示を
$U\to\mathbb A^d_S$ という \'etale 射にし、
$u$ の相対座標を持つ $S$ の切断に沿って引き戻せばよい。
これにより $X_s,P_s$ を保つ延長を得る。
最後に、核と逆核の単位・余単位の同型は導来 tensor に保たれるので、
構成した同値は全ての導来底変更後も同値である。
\end{proof}

\begin{remark}[主予想への寄与と残る障害]
\label{grl:rem:scope}
この定理は、既知の幾何的実現点を含む smooth proper な圏族の近傍についての結果である。
連結な $S$ の非空開で可換実現があることから、$S$ の全点での存在は従わない。
主予想へ進むには、境界の DVR に沿って三次模型と比較核を延長する条件、
および全内在的対象をこのような族で扱えることを別に示す必要がある。
系\paperref{moduli}{cor:nfm-residual-good-reduction}の選択した形式比較に依存する
良い還元も、比較核の境界延長を自動的には与えない。

四次元では命題\ref{prop:nf-fourfold-obstruction}の二つの余分な
Hochschild 変形方向があるため、この証明を適用できない。
rotation が実際の自己同値として持ち上がる場合は、
定理\ref{ifk:thm:rotation-formal}が固定形式空間を幾何的変形と同定する。
その結果、定理\ref{ifk:thm:rotation-neighborhood}では
相対 rotation と原点の rotation 付き核を入力として、四次元の局所比較を得る。
これは本定理の $n\ne4$ を無条件に除くものではない。
また本節では相対 smooth proper 性を入力とし、
任意の Abel 圏変形の飽和性や全 $D^b(\mathrm{coh})$ との一致を証明したとはしない。

形式認識と核モジュライの既知の代数性を組み合わせた帰結であり、
この開性の先行研究上の新規性の網羅的確認は未了である。
\end{remark}

\subsection{trait 上の比較核に残る二つの有限検査}

可換模型の滑らかな延長と、比較核の延長は別の条件である。
さらに、比較核を perfect 双加群として延長できても、
その特殊ファイバーが同値を与えるとは限らない。
以下は実際の延長核を得た後に残る検査を、有限複体として記述する。
核の存在をこの検査自体から導くものではない。

\begin{proposition}[DVR 上の同値核の有限判定]
\label{grl:prop:trait-kernel-test}
$R$ を DVR、$t$ をその一様化元、$K$ を商体、$k$ を剰余体とする。
$A,B$ を $R$ 上 smooth proper な dg 代数とし、
$P$ を $A$--$B$ perfect 双加群とする。
右加群の圏を用い、
\[
 F_P=-\otimes_A^{\mathbf L}P:
 \operatorname{Perf}(A)\longrightarrow\operatorname{Perf}(B)
\]
と置く。$F_{P_K}$ は同値であると仮定する。
$Q=\operatorname{RHom}_B(P,B)$ とし、単位・余単位の cone を
\[
 C_A=\operatorname{Cone}\left(
 A\longrightarrow\operatorname{RHom}_B(P,P)\right),\qquad
 C_B=\operatorname{Cone}\left(
 Q\otimes_A^{\mathbf L}P\longrightarrow B\right)
\]
とする。このとき以下が成立する。
\begin{enumerate}
\item $C_A,C_B$ は $R$ 上 perfect であり、その各 cohomology は有限長である。
\item 次の三条件は同値である：
  $F_P$ が同値であること、$F_{P_k}$ が同値であること、
  $(C_A\oplus C_B)\otimes_R^{\mathbf L}k=0$ であること。
\item 固定した $A,B,P$ に対して
\[
 d_R(P;A,B)=\sum_j\left(
 \operatorname{length}_R H^j(C_A)
 +\operatorname{length}_R H^j(C_B)\right)
\]
は有限な非負整数であり、上の条件は $d_R(P;A,B)=0$ と同値である。
\item $R\to R'$ を DVR の有限平坦な局所拡大、$e$ を分岐指数とすると、
\[
 d_{R'}(P_{R'};A_{R'},B_{R'})=e\,d_R(P;A,B).
\]
従って、固定した延長核の不良な特殊化は、有限分岐底変更だけでは解消しない。
\end{enumerate}
\end{proposition}

\begin{proof}
$P$ の $B$-perfect 性から、$F_P$ の右随伴は
$G_P=\operatorname{RHom}_B(P,-)\simeq-\otimes_B^{\mathbf L}Q$ で与えられる。
smooth proper 性と有限 cell--retract 表示により、
$P,Q$、二つの合成核、二つの cone は底上 perfect である。
随伴、単位・余単位、これらの cone の形成は導来底変更と両立する。

$A$ は $\operatorname{Perf}(A)$ の生成対象である。
従って $F_P$ の完全忠実性は $A$ 上で単位を検査すること、
すなわち $C_A=0$ と同値である。
同様に $B$ 上で余単位を検査することで、$C_A=C_B=0$ は
同値性と同値となる。
これは To\"en--Vaqui\'e
\cite[Lemma~3.25 and its proof]{NFMToenVaquieModuli}
の二つの代数に対する同じ生成対象の議論である。

一般ファイバーで同値であるため $(C_A)_K=(C_B)_K=0$ である。
従って有限生成な各 cohomology は $R$-torsion、従って有限長である。
底上 perfect な複体 $C$ について、導来 Nakayama により
$C\otimes_R^{\mathbf L}k=0$ なら $C=0$ である。
これを二つの cone に適用して第2項を得る。
第3項は有限長の非負性による。

最後に $R'$ の平坦性と底変更から
$H^j(C\otimes_R R')=H^j(C)\otimes_R R'$ である。
有限長 $R$-加群の組成列を用いると、
$\operatorname{length}_{R'}(M\otimes_R R')
=e\operatorname{length}_R M$ となる。
これを各 cohomology へ適用すれば第4項を得る。
\end{proof}

\begin{remark}[行列での検査とその依存性]
\label{grl:rem:trait-smith}
上の二つの cone を有界有限自由 $R$-複体で表せば、
特殊ファイバーでの判定は有限個の行列の $t=0$ での exactness の検査である。
DVR 上の Smith 標準形を逐次使うと、一般ファイバーで acyclic な
有限自由複体は、shift を除いて
$[R\xrightarrow{t^{m_\nu}}R]$、$m_\nu\ge0$、の有限直和に分解できる。
$m_\nu=0$ の項は contractible であり、残る指数の和が
上で定めた長さとなる。

この数値は固定した dg 生成模型と延長核の数値であり、
圏だけの新しい内在的不変量とはしない。
異なる延長核を選んだときの値の不変性も主張しない。
判定が有限であることと、具体的な族でその行列を既に計算したことは別である。
\end{remark}

\begin{example}[perfect 延長だけでは同値性は残らない]
\label{grl:ex:perfect-kernel-not-equivalence}
$A=B=R$ とし、$P=R\oplus R/(t)$ と置く。
$R/(t)$ は二項有限自由分解を持つので、$P$ は perfect 双加群である。
一般ファイバーでは $P_K=K$ で同値を与えるが、
特殊ファイバーでは
\[
 P\otimes_R^{\mathbf L}k\simeq k^{\oplus2}\oplus k[1]
\]
であり、tensor 関手は完全忠実でない。
この例では別の延長核 $R$ を選べば同値になる。
従って問題は有限底変更の有無だけでなく、適切な核の整格子の構成にある。
\end{example}

\begin{remark}[指定混合族の比較と選択済み一般点核の特殊化]
\label{grl:rem:mixed-trait-inputs}
指定された中心係数 $c$ の族については、
定理\paperref{mixed}{cmo:thm:all-central-commutative-kernel}が
滑らかな範囲全体で可換三次への実比較核を与える。
$c=0$ での直接構成は
定理\paperref{mixed}{cmo:thm:fixed-zero-commutative-kernel}である。
従って、各滑らかな $c$ における Ku 同値の存在を、
一般点核の特殊化に関する未解決問題として扱わない。

別に、非 PI 変形の $h$-族について選択済みの一般点核を固定した場合は、
定理\paperref{mixed}{mix:thm:h-exact-candidate}が一般点での実同値を、
定理\paperref{mixed}{mix:thm:h-integral-kernel}が perfect な整延長を与える。
この延長について残るのは、一般点での指定を保ちながら、
$h=1$ における unit と counit の cone がともに零となるように
整核を選べるかという点である。
これは特殊ファイバーにおける同値そのものの存在とは異なる条件であり、
smooth proper 性、判別式の一致、Hochschild 次元の一致だけからは従わない。
\end{remark}

% !TEX root = ../nc-cubic.tex
\section{残余圏の形式認識と到達範囲}
\label{real:sec:scope}
\subsection{通常 Artin 変形の認識}
\begin{maintheorem}[残余圏の通常の Artin 形式認識]
滑らかな三次 $X$ の dg 残余圏を $\mathcal K_X=\operatorname{Ku}(X)$ とし、
$\mathcal D_X=\operatorname{Ind}(\mathcal K_X)$ とする。
$n\ge3$、$n\ne4$ なら、特殊ファイバーを同定した通常の可換局所 Artin 基底上で
\[
 \operatorname{Def}_X\simeq
 \tau_{\le1}\operatorname{CatDef}^{c}_{\mathcal D_X}.
\]
両辺の同型類は $\binom{n+2}{3}$ 次元の滑らかな形式円板となる。
$n=4$ では圏変形の形式 hull の通常 Artin 同型類関手が
二十二次元形式円板となり、幾何的変形の円板は二十次元である。
\end{maintheorem}

証明は定理\ref{thm:nf-residual-formal}に与える。
既知の sector 公式と幾何的 Kodaira--Spencer 単射性から得る接写像同型を、
実際の幾何的変形の全次数の自然変換と合わせる。
$HH^3$ の消滅に依存しないため、$HH^3=\C^2$ となる $n=7$ も含む。
圏変形の形式 hull と実際の compactly generated 変形の比較には
Lurie \cite{NFDAGX} の tame 圏に対する結果を用いる。
未切断の高次群oidや導来 Artin 基底上での同値は主張しない。

rotation $\mathsf R$ と Serre 関手 $S$ の既知の関係は
$\mathsf R^3=[2]$、$S=\mathsf R^{-(n-1)}[n]$ である。
$n\not\equiv1\pmod3$ では rotation を Serre と shift から回復できるが、
$n\equiv1\pmod3$ ではこの方法は使えない。
これは四次元の形式認識の例外とは異なる区分である。


\subsection{相対比較と降下の位置付け}

構造対象付きモジュライ、trace--Hessian と残余三次の差、
点曲線による尺度の規格化、比較核の Picard 障害と周期位相は、
モジュライ・降下稿\cite{NCCubicModuliDescent}で扱う。
自然な標的を指定する比較や族上の比較では、本稿からもそこでの
中心直線・相対尺度・scalar 自明化の結果を参照する。
固定された複素係数で任意の平面方程式を選ぶ直接 Ku 比較は、
これらの精密な降下条件を前提としない。
幾何学的な同値の存在、核の同型類の定義体、実際の核の scalar 障害、
rotation と指定周期の位相は、それぞれ異なる問題である。

\subsection{可換実現の境界と残る問題}

中心課題は主予想\ref{conj:nc-main}であり、
既知の可換実現 locus の境界延長と各 intrinsic 成分を覆う表示を優先する。
点スキームと Poisson 零スキームの一致を一般に期待できないことは、
注意\ref{ppz:rem:naive-expectation}とそこで参照する三様体・五様体の例で説明した。
既知の分離型における一致とその限界は、可換実現の幾何を理解するために保持する。
中心変数を持つ分離型の可換実現は定理\ref{csc:thm:comparison}で扱い、
定数型混合族では、有限中心性と相対対角の証明から
系\paperref{mixed}{cmo:cor:cofinite}の余有限可換実現を得た。
指定二因子では、定理\paperref{mixed}{cmo:thm:morse-singular-sixty-four}により
norm と交わらない六十四係数で全 coherent Ku の非 proper 性が確定した。
定理\paperref{mixed}{cmo:thm:fourteen-node-orders}で十四節点の局所代数と全単純加群を決定し、
中央切断・相対対角・射影境界の証明を合わせて
定理\paperref{mixed}{cmo:thm:exact-sixty-four}の正確な非 smooth 集合 $V(P_{64})$ を得た。
例\paperref{mixed}{cmo:ex:sixty-four}の norm の全六十四係数と $c=0$ は smooth proper である。
定理\paperref{mixed}{cmo:thm:all-central-commutative-kernel}は、
この指定族の全 smooth locus 上に、明示可換三次族への相対 perfect 核を構成する。
norm の全六十四根にも同じ核が適用される。
具体的な比較先として、定理\paperref{mixed}{cmo:thm:commutative-discriminant-candidate}の
可換三次族と、零係数での命題\paperref{mixed}{cmo:prop:zero-eckardt-resolution}の模型を使う。
零係数では定理\paperref{mixed}{cmo:thm:eckardt-root-residual}の根スタック上の十六直交条件と
指定された指標作用を比較核の標的として使える。
さらに定理\paperref{mixed}{cmo:thm:root-order-residual}の非可換側の実二成分分解と、
境界上の twist 射および高次構造の比較から、
定理\paperref{mixed}{cmo:thm:fixed-zero-commutative-kernel}の
$\Ku_{\mathrm{coh}}(A_0)\simeq\Ku(V(G-u^2L))$ の実 perfect 核を得る。
その構成は $\C$ 上の通常の dg Morita 同値であり、
全標識の同時保存や有理数体への核降下までは主張しない。
全中心係数の構成では、次数一の中心超平面制限を使う。
その comonad の二項・三項の構造を混合係数と中心係数へ同定し、
四項以降の比較の一意性から実際の Morita 同値族を得る。
特異係数を含む残余線束圏の比較が先にあり、
既知の判別集合はその後の尺度決定に用いる。
系\paperref{mixed}{cmo:cor:all-alpha-central-kernel}は
固定した二つの Sklyanin 因子の任意の定数補正も扱う。
交換子パラメータ $h$ については、定理\paperref{mixed}{mix:thm:h-exact-candidate}により
指定した可換候補への全形式次数の MF 同値と一般点の実 perfect 核を得た。
$h=0,1$ を含む相対 smooth proper 模型と、
$h=1$ の DVR 上の perfect 整延長は
定理\paperref{mixed}{mix:thm:h-semilocal-smooth}と
\paperref{mixed}{mix:thm:h-integral-kernel}で構成できる。
選択済み $h$ 一般点核を特殊点でも可逆に延長し、
固定 $h=1$ の核と相対的に整合させることは残る。
通常の Hochschild 群と有理 Euler 形式は全 smooth 係数で一致するが、
それだけで比較核は得られない。
非次数付き affine MF 比較は同領域で零となるため、
本稿の実比較は、内部次数と実随伴を持つ射影残余圏から構成している。
他の族や選択済み一般点核に対しては、
命題\ref{grl:prop:trait-kernel-test}の単位・余単位の有限長判定が引き続き必要となる。
通常の周囲代数の derived Morita 比較と、形式的な $h$ 進核の
$h=1$ への無条件な代入は、この手順の代わりにならない。
可換実現 locus の開性だけでは全 smooth locus への拡張は従わない。
次の中心課題は、他の Sklyanin パラメータや線形補正、
複数の中心変数を含む混合表示への実比較の拡張と、全 intrinsic 成分の被覆である。

精密な降下の側には、周期位相の他の境界での分岐、正準位相体上の scalar index、
代数的核の coherent な中心線型性、高次元の任意の full class と非分裂軌道が残る。
これらの記法と条件はモジュライ・降下稿に置く。
特に同稿の補題\paperref{moduli}{lem:nfm-period-two-boundary-criterion}は、
一つの完備 trait 上の位相値だけから大域式を結論できないことを明記する。
同稿の第\paperref{moduli}{subsubsec:nfm-boundary-products-period}項では、
全 orbit Hom と指定周期を保つ整模型で低次積・飽和関係・三次積の消失次数を調べる。
この積の計算と局所尺度の比較は、滑らかな残余模型や full 比較核自体の
境界延長とは別である。
形式座標の代数性と内在的なモジュライの代数性も区別し、
境界や全 Abel 形式近傍は、一般次元の可換実現を妨げる障害の解消に必要な範囲で扱う。


\begin{thebibliography}{99}

\bibitem{NCCubicA}
\emph{滑らかな三次超曲面のPoisson構造：正規形・退化軌跡・代数的量子化},
\texttt{nc-cubic} 論文A、未公表研究稿 A-v61、2026年11月14日。
\href{https://github.com/uedakazushi/nc-cubic/blob/eabe601c1341777742964584dde5314c0bcf64a9/papers/paper_a.tex}{固定 commit の主原稿}。
支持節は \texttt{papers/support/papers/sections/poisson\_geometry.tex}
の命題 \texttt{prop:pg-represent}, \texttt{prop:pg-jacobi-extension}、
定理 \texttt{thm:pg-normal}、系 \texttt{cor:pg-generic}、および
\texttt{general\_dimension\_universal\_quantization.tex}
の定理 \texttt{thm:pg-all-dimensions}、命題 \texttt{prop:pg-ad-relative-chart}。

\bibitem{NFNCCubicB}
\emph{中心三次非可換変形の一般点における可換復元},
\texttt{nc-cubic} 論文B、未公表研究稿 B-v59、2026年11月14日。
\href{https://github.com/uedakazushi/nc-cubic/blob/eabe601c1341777742964584dde5314c0bcf64a9/papers/paper_b.tex}{固定 commit の主原稿}。
\texttt{papers/generic\_recovery/} の
\texttt{main\_result.tex}, \texttt{generic\_descent.tex},
\texttt{coherent\_saturation.tex}, \texttt{field\_model.tex} を参照。

\bibitem{Threefold}
\emph{非可換三次三様体のAbel圏モジュライ：二次包絡とKuznetsov残余射},
未公表研究稿 v22、2026年11月20日。
\href{https://github.com/uedakazushi/nc-cubic-threefold-moduli/blob/6c35d2c84671c2538e61836a6bf3a511392eacce/papers/nc-cubic-threefold-moduli.tex}{固定 commit の主原稿}。
点スキームは \texttt{papers/sections/point\_scheme\_poisson.tex}、
全 Hom の比較は \texttt{jacobian\_abel\_comparison.tex} を参照。

\bibitem{NFMEG}
P.~Etingof and V.~Ginzburg,
\emph{Noncommutative del Pezzo surfaces and Calabi--Yau algebras},
J. Eur. Math. Soc. \textbf{12} (2010), 1371--1416.
\url{https://doi.org/10.4171/JEMS/235}.
\url{https://ems.press/content/serial-article-files/31746},
Theorem~3.3.2, Section~9.

\bibitem{MoriSmith}
I.~Mori and S.~P.~Smith,
\emph{The classification of 3-Calabi--Yau algebras with 3 generators and 3 quadratic relations},
Math. Z. \textbf{287} (2017), 215--241.
\url{https://doi.org/10.1007/s00209-016-1824-5}.
\url{https://arxiv.org/abs/1502.07403v2}, Theorem~4.2.

\bibitem{LBSV}
L.~Le Bruyn, S.~P.~Smith and M.~Van den Bergh,
\emph{Central extensions of three dimensional Artin--Schelter regular algebras},
Math. Z. \textbf{222} (1996), 171--212.
\url{https://doi.org/10.1007/PL00004532}.
\url{https://sites.math.washington.edu/~smith/Research/Central-extns.pdf},
Theorem~4.2.2, Proposition~4.2.5, Section~4.3.

\bibitem{Mayanskiy}
E.~Mayanskiy,
\emph{Poisson cohomology of two Fano threefolds},
J. Algebra \textbf{424} (2015), 21--45.
\url{https://doi.org/10.1016/j.jalgebra.2014.08.049}.
\url{https://arxiv.org/abs/1303.5916v1}, Theorem~1, Section~2.2.
Theorem~1 は Loray--Pereira--Touzet の分類に帰属されている。
スキーム構造には Section~2.2 の Schouten 計算を用いる。

\bibitem{BCMStandardMonomial}
N.~C.~Bonala, O.~Clarke and F.~Mohammadi,
\emph{Standard monomial theory and toric degenerations of Richardson varieties in the Grassmannian},
\url{https://arxiv.org/abs/2103.16208v1}, 2021,
Section~2.4, Theorem~2.3.

\bibitem{NFMStacksProjective}
The Stacks Project Authors,
\emph{The Stacks Project}, Section~27.21, Projective bundles, Tag~01OA.
\url{https://stacks.math.columbia.edu/tag/01OA}.

\bibitem{NFMStacksAffineDescent}
The Stacks Project Authors,
\emph{The Stacks Project}, Lemma~35.37.1, Tag~0245.
\url{https://stacks.math.columbia.edu/tag/0245}.

\bibitem{NFLZSerre}
X.~Lin and S.~Zhang,
\emph{Serre algebra, matrix factorization and categorical Torelli theorem for hypersurfaces},
Math. Ann. \textbf{391} (2025), 163--177.
\url{https://doi.org/10.1007/s00208-024-02915-8}.
定理番号は \url{https://arxiv.org/abs/2310.09927v2}
の Theorem~3.5, Proposition~3.7, Theorems~4.1 and~5.1 による。

\bibitem{NFBFKKernels}
M.~Ballard, D.~Favero and L.~Katzarkov,
\emph{A category of kernels for equivariant factorizations and its implications for Hodge theory},
Publ. Math. Inst. Hautes \'Etudes Sci. \textbf{120} (2014), 1--111.
\url{https://doi.org/10.1007/s10240-013-0059-9}.
\url{https://numdam.org/item/10.1007/s10240-013-0059-9.pdf},
Theorem~5.39 の一般式とその証明、pp.~82--84。

\bibitem{NFPirozhkov}
D.~Pirozhkov,
\emph{A categorical Torelli theorem for hypersurfaces},
Bull. London Math. Soc. \textbf{56} (2024), no.~10, 3075--3089.
\url{https://doi.org/10.1112/blms.13117}.
本文の Lemmas~5.1 and~5.2、Theorem~5.3 は
\url{https://arxiv.org/abs/2208.13604v1} による。

\bibitem{NFHuybrechtsRennemo}
D.~Huybrechts and J.~V.~Rennemo,
\emph{Hochschild cohomology versus the Jacobian ring and the Torelli theorem for cubic fourfolds},
Algebraic Geometry \textbf{6} (2019), no.~1, 76--99.
\url{https://doi.org/10.14231/AG-2019-005}.
\url{https://content.algebraicgeometry.nl/2019-1/2019-1-005.pdf},
Proposition~3.5 and Remark~3.8.

\bibitem{NFDAGX}
J.~Lurie,
\emph{Derived Algebraic Geometry X: Formal Moduli Problems}, November 5, 2011.
\url{https://www.math.ias.edu/~lurie/papers/DAG-X.pdf},
Section~2, Theorems~5.1.9 and~5.3.16, Proposition~5.3.21.

\bibitem{NFBKP}
A.~Blanc, L.~Katzarkov and P.~Pandit,
\emph{Generators in formal deformations of categories},
Compos. Math. \textbf{154} (2018), 2055--2089.
\url{https://doi.org/10.1112/S0010437X18007303}.
\url{https://arxiv.org/abs/1705.00655v1}, Sections~3.3 and~4.3.

\bibitem{NFKuznetsovFCY}
A.~Kuznetsov,
\emph{Calabi--Yau and fractional Calabi--Yau categories},
J. reine angew. Math. \textbf{753} (2019), 239--267.
\url{https://doi.org/10.1515/crelle-2017-0004}.
\url{https://arxiv.org/abs/1509.07657v2},
Section~3.3, Corollaries~3.18 and~4.1, Lemma~3.19.

\bibitem{NFBFKOrlov}
M.~Ballard, D.~Favero and L.~Katzarkov,
\emph{Orlov spectra: bounds and gaps},
Invent. Math. \textbf{189} (2012), 359--430.
\url{https://doi.org/10.1007/s00222-011-0367-y}.
定理番号は \url{https://arxiv.org/abs/1012.0864v3} の Theorem~5.15 による。

\bibitem{NFKuznetsovCubicFourfolds}
A.~Kuznetsov,
\emph{Derived categories of cubic fourfolds},
in \emph{Cohomological and Geometric Approaches to Rationality Problems},
Progress in Mathematics \textbf{282}, Birkh\"auser, 2010, 219--243.
\url{https://arxiv.org/abs/0808.3351}, Theorem~3.1.
著者稿 \url{https://cims.nyu.edu/~tschinke/books/rationality/kuznetsov.pdf}.

\bibitem{NFMGabrielQuotient}
P.~Gabriel, \emph{Des cat\'egories ab\'eliennes},
Bull. Soc. Math. France \textbf{90} (1962), 323--448,
Chapitre~III, \S1, p.~365.
\url{https://doi.org/10.24033/bsmf.1583}.
\url{https://www.numdam.org/article/BSMF_1962__90__323_0.pdf}.

\bibitem{Fraenkel}
A.~McMillan Fraenkel,
\emph{Extensions of Poisson structures on singular hypersurfaces},
\url{https://arxiv.org/abs/1310.6083v2}, 2013,
Proposition~4.1, Theorem~4.3.

\bibitem{Carles}
R.~Carles,
\emph{D\'eformations dans les sch\'emas d\'efinis par les identit\'es de Jacobi},
Ann. Math. Blaise Pascal \textbf{3} (1996), no.~2, 33--62.
\url{https://doi.org/10.5802/ambp.67}.
\url{https://www.numdam.org/item/AMBP_1996__3_2_33_0.pdf},
Section~8, Lemma~8.35.

\bibitem{CarlesMarquez}
R.~Carles and M.~C.~M\'arquez,
\emph{Different methods for the study of obstructions in the schemes of Jacobi},
Ann. Inst. Fourier \textbf{61} (2011), no.~2, 453--490.
\url{https://doi.org/10.5802/aif.2620}.
\url{https://www.numdam.org/item/AIF_2011__61_2_453_0.pdf},
Section~3, Theorem~3.2, Propositions~3.9 and~3.10.

\end{thebibliography}

\end{document}
